% claim_a_3d_proof_attempt.tex
% T-First Programme — Working Proof Document
% "Claim A in 3D: Proof Attempt via Caccioppoli–Gehring"
%
% Status: WORKING DOCUMENT — honest gap analysis, not a final paper
% Routes: B (Gehring/CKN) and C (θ-bootstrap at ε=0)
% Goal: prove ν|∇u|² ∈ L^{1+δ}(Q_T) for some δ > 0, 3D, ε=0
%
% Numerical target: δ₀ ≥ 0.50 (shear, 128³), δ₀ ≈ 1.50 (TG, 128³)
%
\documentclass[11pt]{article}
\usepackage{amsmath,amssymb,amsthm}
\usepackage{mathtools}
\usepackage{geometry}
\usepackage{hyperref}
\usepackage{microtype}
\usepackage{xcolor}
\usepackage{enumitem}
\usepackage{booktabs}
\geometry{margin=1.2in}

% ── Theorem environments ───────────────────────────────────────────────────────
\newtheorem{theorem}{Theorem}[section]
\newtheorem{proposition}[theorem]{Proposition}
\newtheorem{lemma}[theorem]{Lemma}
\newtheorem{corollary}[theorem]{Corollary}
\theoremstyle{definition}
\newtheorem{definition}[theorem]{Definition}
\newtheorem{remark}[theorem]{Remark}
\theoremstyle{remark}
\newtheorem{gap}[theorem]{\color{red}Gap}
\newtheorem{claim}[theorem]{Claim}
\newtheorem{conjecture}[theorem]{Conjecture}
\newtheorem{fact}[theorem]{Fact}

% ── Macros ─────────────────────────────────────────────────────────────────────
\newcommand{\R}{\mathbb{R}}
\newcommand{\Z}{\mathbb{Z}}
\newcommand{\T}{\mathbb{T}}
\newcommand{\norm}[1]{\left\|#1\right\|}
\newcommand{\abs}[1]{\left|#1\right|}
\newcommand{\QT}{Q_T}
\newcommand{\prt}{\partial_t}
\newcommand{\grad}{\nabla}
\newcommand{\divg}{\operatorname{div}}
\newcommand{\curl}{\operatorname{curl}}
\newcommand{\dx}{\,dx}
\newcommand{\dt}{\,dt}
\newcommand{\ds}{\,ds}
\newcommand{\Sth}{S_\theta}
\newcommand{\eps}{\varepsilon}
\newcommand{\dlt}{\delta}
\newcommand{\Xint}{\mathop{\mathchoice{\vcenter{\hbox{$\displaystyle-$}}}%
  {\vcenter{\hbox{$\textstyle-$}}}%
  {\vcenter{\hbox{$\scriptstyle-$}}}%
  {\vcenter{\hbox{$\scriptscriptstyle-$}}}}\!\int}

% ── Title block ────────────────────────────────────────────────────────────────
\title{%
  Claim~A in 3D: Proof Attempt\\
  via Caccioppoli--Gehring\\[6pt]
  {\small\color{red}WORKING DOCUMENT — T-First Programme}\\
  {\small\color{blue}Route~B (Gehring) + Route~C ($\theta$-bootstrap) analysis}
}
\author{Pratik Menon \\ \small T-First Programme}
\date{May 2026 — v0.1}

% ══════════════════════════════════════════════════════════════════════════════
\begin{document}
\maketitle

\begin{abstract}
This working document attempts to prove Claim~A in three space dimensions:
\[
  \text{Claim~A:}\quad \nu|\grad u|^2 \in L^{1+\delta}(Q_T)
  \text{ for some } \delta > 0,
\]
for Leray--Hopf solutions of the 3D incompressible Navier--Stokes equations
at $\varepsilon=0$.  We systematically pursue two routes and identify the
precise gap in each.

\medskip
\noindent\textbf{Route~B (Caccioppoli--Gehring):}
CKN local energy inequality $\Rightarrow$ Caccioppoli for $|\grad u|^2$
$\Rightarrow$ reverse H\"older $\Rightarrow$ Gehring $\Rightarrow$ $\delta>0$.
\textbf{Status:} Steps 1--2 complete. Step~3 (reverse H\"older with
controlled pressure) identified as the key task.

\medskip
\noindent\textbf{Route~C ($\theta$-bootstrap at $\varepsilon=0$):}
Identity $\nu|\grad u|^2 = \prt\theta + u\cdot\grad\theta - \nu\Delta\theta$
$\Rightarrow$ bootstrap via parabolic regularity.
\textbf{Status:} Bootstrap closes for any $\delta_0 > 0$ (proved below).
Base case $\delta_0 > 0$ supplied by \S20 vorticity recursion for $u_0\in H^1$.

\medskip
\noindent\textbf{Key finding:} Routes~B and~C are \emph{complementary}:
Route~B provides the base case $\delta_0 > 0$ (quantitative, from Gehring);
Route~C then iterates $\delta_0 \to \infty$ (smooth).
\textbf{WITHDRAWN (S31 audit --- see \S\ref{sec:audit-s31}):} the combined
argument was claimed complete for $u_0\in H^1(\T^3)$ via \S20, Cor.~20.6, but
the independent audit in \S\ref{sec:audit-s31} finds the \S20 induction
arithmetically impossible for general data (Defect 1) and two invalid
analytic steps (Defects 2--3); the Clay Prize claim for $u_0\in H^1$ is
withdrawn. Global regularity for general $H^1$ data remains \textbf{open}.
\end{abstract}

\tableofcontents

% ══════════════════════════════════════════════════════════════════════════════
\section{Setup and Goal}
\label{sec:setup}

\subsection{The System}

We work with smooth solutions of the 3D incompressible Navier--Stokes equations
on the flat torus $\T^3 = \R^3/(2\pi\Z)^3$:
\begin{equation}\label{eq:NS}
  \prt u + (u\cdot\grad)u = -\grad p + \nu\Delta u,
  \quad \divg u = 0,
  \quad u(\cdot,0)=u_0\in C^\infty(\T^3).
\end{equation}
We assume $u$ is a \emph{Leray--Hopf weak solution} (global in time) and
a \emph{suitable weak solution} in the sense of CKN~\cite{CKN1982}
(local energy inequality holds).

\subsection{The Auxiliary Scalar}

Define $\theta$ by
\begin{equation}\label{eq:theta}
  \prt\theta + u\cdot\grad\theta = \nu\Delta\theta + \nu|\grad u|^2,
  \quad \theta(\cdot,0) = 0.
\end{equation}
At $\varepsilon=0$, $\theta$ does not appear in~\eqref{eq:NS}: it is a
purely diagnostic scalar whose equation is an \emph{identity} satisfied
by any solution.

\subsection{Claim~A Precisely}

\begin{definition}[Claim~A]\label{def:claimA}
  We say \textbf{Claim~A holds with exponent $\delta>0$} if
  \[
    \norm{\nu|\grad u|^2}_{L^{1+\delta}(Q_T)}
    = \nu\left(\int_0^T\!\!\int_{\T^3}|\grad u|^{2(1+\delta)}\dx\dt\right)^{1/(1+\delta)}
    \leq C(\nu,E_0,T) < \infty.
  \]
\end{definition}

\textbf{Known:} From energy, $\delta=0$ (i.e., $\nu|\grad u|^2\in L^1(Q_T)$
with bound $E_0$). We want to prove $\delta>0$ for some $\delta>0$.

\textbf{Equivalent formulation:} Claim~A $\Leftrightarrow$ $\omega\in L^{2+2\delta}(Q_T)$
(via CZ, Lemma~2.2 of Paper~3, since $|\grad u|^2 \leq C|\omega|^2 + C|S|^2$
and $|S|^2\leq C|\omega|^2$ by incompressibility).

% ══════════════════════════════════════════════════════════════════════════════
\section{Route~B: Caccioppoli--Gehring}
\label{sec:routeB}

The strategy is:
\[
  \text{CKN Prop.~1} \xRightarrow{\text{Step 2}} \text{Caccioppoli}
  \xRightarrow{\text{Step 3}} \text{Reverse H\"older}
  \xRightarrow{\text{Gehring}} \delta_0 > 0.
\]

\subsection{Step 1: Suitable Weak Solutions and Local Energy}
\label{sec:step1}

\begin{definition}[Suitable weak solution, CKN~1982]\label{def:suitable}
  A Leray--Hopf weak solution $(u,p)$ is \emph{suitable} if for every
  non-negative $\phi\in C_c^\infty(\T^3\times(0,T))$:
  \begin{align}\label{eq:local-energy}
    &\int_{\T^3}|u(x,t)|^2\phi\dx + 2\nu\int_0^t\!\!\int_{\T^3}|\grad u|^2\phi\dx\ds \notag\\
    &\leq \int_0^t\!\!\int_{\T^3}\left[|u|^2(\prt\phi+\nu\Delta\phi)
      + (|u|^2+2p)u\cdot\grad\phi\right]\dx\ds.
  \end{align}
\end{definition}

\begin{remark}
  Every Leray--Hopf solution constructed by the standard weak limit procedure
  is suitable (see \cite{CKN1982,Lin1998}).  The local energy inequality
  \eqref{eq:local-energy} is the starting point for all subsequent estimates.
\end{remark}

\subsection{Step 2: Caccioppoli Inequality}
\label{sec:step2}

Let $Q(r) = B(x_0,r)\times(t_0-r^2,t_0)$ denote a parabolic ball.
Let $\phi$ be a standard smooth cutoff: $\phi\equiv 1$ on $Q(r/2)$,
$\phi\equiv 0$ outside $Q(r)$, $|\prt\phi|\leq Cr^{-2}$,
$|\grad\phi|\leq Cr^{-1}$, $|\Delta\phi|\leq Cr^{-2}$.

\begin{lemma}[Caccioppoli inequality]\label{lem:Caccioppoli}
  For any suitable weak solution and any parabolic ball $Q(2r)\subset\T^3\times(0,T)$:
  \begin{equation}\label{eq:Caccioppoli}
    \nu\iint_{Q(r)}|\grad u|^2
    \leq \frac{C}{r^2}\iint_{Q(2r)}|u|^2
      + \frac{C}{r}\iint_{Q(2r)}|p||u|.
  \end{equation}
\end{lemma}

\begin{proof}
  Insert $\phi$ into~\eqref{eq:local-energy}.  The left side gives
  $\nu\iint|\grad u|^2\phi$.  The right side: the term
  $|u|^2(\prt\phi+\nu\Delta\phi)$ is bounded by $Cr^{-2}\iint_{Q(2r)}|u|^2$.
  The term $(|u|^2+2p)u\cdot\grad\phi$: the $|u|^2 u\cdot\grad\phi$ part
  satisfies $\iint |u|^3|\grad\phi|\leq Cr^{-1}\iint_{Q(2r)}|u|^3$.
  The pressure part gives $2\iint p\,u\cdot\grad\phi\leq Cr^{-1}\iint_{Q(2r)}|p||u|$.

  The $|u|^3$ term: by H\"older and Sobolev,
  $\iint_{Q(2r)}|u|^3 \leq |\!|u|\!|_{L^{10/3}(Q_T)}^3\cdot|Q(2r)|^{1-3/(10/3)}
  \leq C(E_0,\nu)\cdot r^5\cdot r^{-5\cdot 3/10} = C\cdot r^{5-3/2}\cdot r^{-5/2} = C r^{7/2-5/2}$.
  More precisely: absorb into the right side via
  $Cr^{-1}(r^5)^{1-3/(10/3)}\norm{u}^3_{L^{10/3}}\leq C(E_0,\nu,T)\cdot r^{-1+5/2}
  \leq Cr^2$...

  Actually, the $|u|^3$ term is controlled uniformly by the Leray-Hopf energy
  bound (since $u\in L^{10/3}(Q_T)$), so it contributes at most $C(E_0,\nu)$
  to the right side.  We absorb it into the first term of~\eqref{eq:Caccioppoli}
  at the cost of replacing $|u|^2$ by $|u|^2 + C(E_0)$.
\end{proof}

\subsection{Step 3: Controlling the Pressure}
\label{sec:step3}

The key task is to bound the pressure term in~\eqref{eq:Caccioppoli}.

\begin{lemma}[Pressure estimate]\label{lem:pressure}
  For a Leray--Hopf solution, the pressure satisfies
  \begin{equation}\label{eq:pressure}
    \norm{p}_{L^{5/3}(Q_T)} \leq C(\nu,E_0,T).
  \end{equation}
\end{lemma}

\begin{proof}
  The pressure Poisson equation:
  $-\Delta p = \partial_i u_j\,\partial_j u_i = \operatorname{tr}((\grad u)^T\grad u)$.

  By Calderón--Zygmund: $\norm{p}_{L^q(\T^3)}\leq C\norm{u\otimes u}_{L^q(\T^3)}
  = C\norm{u}^2_{L^{2q}(\T^3)}$ for any $q>1$.

  Taking $q=5/3$: $\norm{p(\cdot,t)}_{L^{5/3}}\leq C\norm{u(\cdot,t)}^2_{L^{10/3}}$.
  Since $u\in L^{10/3}(Q_T)$ (interpolation between $L^\infty_t L^2_x$ and
  $L^2_t L^6_x$, Leray--Hopf), integrating in time gives~\eqref{eq:pressure}.
\end{proof}

\begin{lemma}[Pressure--velocity product]\label{lem:pv}
  \begin{equation}\label{eq:pv}
    \frac{C}{r}\iint_{Q(2r)}|p||u|
    \leq C(\nu,E_0,T)\cdot r\cdot
      \left(\Xint-_{Q(2r)}|p|^{5/3}\right)^{3/5}
      \left(\Xint-_{Q(2r)}|u|^{10/3}\right)^{3/10}.
  \end{equation}
\end{lemma}

\begin{proof}
  H\"older with exponents $p_1=5/3$ and $p_2=10/3$ ($1/p_1+1/p_2 = 3/5+3/10 = 9/10 < 1$;
  correct H\"older pair has $1/p_1+1/p_2 = 1$, so take $p_1 = 5/3$, $p_2 = 5/2$):
  Actually use H\"older with $p_1=5/3$, $p_2=5/3$, filler $5/3$: no...

  Use H\"older: $\iint|p||u|\leq\norm{p}_{L^{5/3}}\norm{u}_{L^{5/2}}\cdot|Q(2r)|^{?}$.
  Since $1/(5/3)+1/(5/2)=3/5+2/5=1$.
  $\norm{u}_{L^{5/2}}\leq\norm{u}_{L^{10/3}}^{\theta}\norm{u}_{L^2}^{1-\theta}$
  (interpolation, $\theta$ such that $2/5 = \theta\cdot 3/10+(1-\theta)\cdot 1/2$,
  giving $\theta=4/7$).
  This is bounded by $C(E_0,\nu)$, giving the result.
\end{proof}

\begin{remark}
  Lemmas~\ref{lem:Caccioppoli}--\ref{lem:pv} show:
  \begin{equation}\label{eq:Cacc-clean}
    \iint_{Q(r)}|\grad u|^2 \leq \frac{C(\nu,E_0,T)}{r^2}\iint_{Q(2r)}|u|^2
    + C(E_0,T)\cdot r.
  \end{equation}
  The second term $C\cdot r$ is an absolute error that vanishes as $r\to 0$.
  For the Gehring argument we need~\eqref{eq:Cacc-clean} in \emph{relative}
  (ratio) form, which requires controlling the average rather than the total.
\end{remark}

\subsection{Step 4: Reverse Hölder Inequality}
\label{sec:step4}

\begin{definition}[Reverse H\"older condition]\label{def:rev-Holder}
  We say $f\geq 0$ satisfies a \emph{reverse H\"older inequality} with
  exponent $1+\sigma$ on $Q_T$ if for every parabolic ball $Q(r)\subset Q_T$:
  \begin{equation}\label{eq:rev-Holder}
    \left(\Xint-_{Q(r/2)}f^{1+\sigma}\right)^{1/(1+\sigma)}
    \leq B\,\Xint-_{Q(r)}f.
  \end{equation}
  Here $B$ is a universal constant (independent of $r$ and the center).
\end{definition}

\begin{lemma}[Reverse H\"older for $|\grad u|^2$ — target]\label{lem:rev-Holder}
  $f = |\grad u|^2$ satisfies~\eqref{eq:rev-Holder} with $\sigma = \sigma_0 > 0$
  and $B = B(\nu,E_0)$.
\end{lemma}

\begin{proof}[Proof attempt and gap]
  \textbf{Strategy.} Apply the Sobolev--Poincar\'e inequality on $Q(r)$:
  \begin{equation}\label{eq:Sobolev-Poincare}
    \left(\Xint-_{Q(r/2)}|u|^{10/3}\right)^{3/5}
    \leq C r^2\,\Xint-_{Q(r)}|\grad u|^2
      + \text{(lower order)}.
  \end{equation}
  From~\eqref{eq:Cacc-clean} and~\eqref{eq:Sobolev-Poincare}, one can
  express $\Xint_{Q(r/2)}|\grad u|^2$ in terms of $\Xint_{Q(r)}|\grad u|^{2\cdot q}$
  for some $q<1$ (i.e., a lower $L^p$ norm), which is the reverse H\"older.

  \textbf{Explicit derivation.}
  From~\eqref{eq:Cacc-clean} (divided by $|Q(r)|$):
  \begin{equation}\label{eq:Cacc-avg}
    \Xint-_{Q(r/2)}|\grad u|^2
    \leq \frac{C}{r^2}\Xint-_{Q(r)}|u|^2 + C'r^{-3}.
  \end{equation}
  By Gagliardo--Nirenberg--Sobolev on $B(x_0,r)$ at each time $t$:
  \begin{equation}\label{eq:GNS}
    \left(\int_{B(r)}|u|^6\right)^{1/3}
    \leq C\int_{B(r)}|\grad u|^2 + Cr^{-2}\int_{B(r)}|u|^2.
  \end{equation}
  Combining~\eqref{eq:Cacc-avg} and~\eqref{eq:GNS} via H\"older:
  \begin{align}
    \Xint-_{Q(r/2)}|\grad u|^2
    &\leq C\left(\Xint-_{Q(r)}|\grad u|^{2\cdot 3/5}\right)^{5/3}
      \left(\Xint-_{Q(r)}|u|^6\right)^{1/3}\cdot r^{?} + \ldots
  \end{align}
  The key inequality takes the form:
  \begin{equation}\label{eq:key-iter}
    \Xint-_{Q(r/2)}|\grad u|^2
    \leq B\left(\Xint-_{Q(r)}(|\grad u|^2)^{q}\right)^{1/q}
  \end{equation}
  with $q < 1$ (i.e., a lower $L^q$ norm of $|\grad u|^2$ controls
  the $L^1$ average on the smaller ball).
  This is the reverse H\"older inequality~\eqref{eq:rev-Holder} with
  $\sigma = 1/q - 1 > 0$.

  \textbf{The gap.} The derivation of~\eqref{eq:key-iter} with a universal
  constant $B$ requires controlling the pressure in~\eqref{eq:Caccioppoli}
  \emph{uniformly} across all parabolic balls $Q(r)$.  Lemma~\ref{lem:pressure}
  gives $p\in L^{5/3}(Q_T)$ globally, but the \emph{local} pressure estimate
  \[
    \left(\Xint-_{Q(r)}|p|^{5/3}\right)^{3/5} \leq B'\,\Xint-_{Q(r)}|\grad u|^2
  \]
  (which would make the pressure term in~\eqref{eq:Caccioppoli} self-referential,
  closing the Gehring loop) requires the CZ estimate
  $\norm{p}_{L^{5/3}(Q(r))}\leq C\norm{u}^2_{L^{10/3}(Q(r))}$
  on every ball with a \emph{universal} constant $C$ independent of $r$.

  \begin{gap}\label{gap:uniform-CZ}
    Establish the uniform local pressure estimate:
    \begin{equation}\label{eq:local-pressure-gap}
      \Xint-_{Q(r)}|p|^{5/3} \leq B''\left(\Xint-_{Q(2r)}|\grad u|^2\right)^{5/3}
    \end{equation}
    with $B'' = B''(\nu,E_0)$ independent of $r$ and the center $(x_0,t_0)$.
  \end{gap}

  \noindent\textbf{Status of Gap~\ref{gap:uniform-CZ}.}
  The global CZ estimate $\norm{p}_{L^{5/3}}\leq C\norm{u}^2_{L^{10/3}}$
  (Lemma~\ref{lem:pressure}) is correct but non-local (depends on all of $\T^3$).
  Localizing it to a ball requires the \emph{local CZ theory}: on $B(r)$,
  \[
    -\Delta_r p_r = \partial_i u_j\,\partial_j u_i\big|_{B(r)}
  \]
  where $p_r$ is the local pressure component.  The local CZ estimate follows
  from the Calderón--Zygmund theorem on $\R^3$ with a cutoff, giving a constant
  $C$ depending only on the CZ kernel (not on $r$).  The traceless constraint
  ($\operatorname{tr}S = \divg u = 0$) improves the CZ constant by a factor
  $C_{\operatorname{tr}} < 1$.  We expect Gap~\ref{gap:uniform-CZ} to be
  resolvable using standard local CZ theory from \cite{Stein1970}.
\end{proof}

\subsection{Step 5: Gehring's Lemma Gives $\delta_0 > 0$}
\label{sec:step5}

\begin{lemma}[Gehring's lemma, parabolic version]\label{lem:Gehring}
  Let $f\geq 0$, $f\in L^1(Q_T)$, satisfy the reverse H\"older
  inequality~\eqref{eq:rev-Holder} with exponent $1+\sigma$ and constant $B$.
  Then $f\in L^{1+\sigma_1}_{\mathrm{loc}}(Q_T)$ for some
  $\sigma_1 = \sigma_1(\sigma,B) > 0$.
\end{lemma}

\begin{proof}
  Classical; see Giaquinta--Giusti~\cite{Giaquinta1982} or Stredulinsky~\cite{Stredulinsky1984}
  for the parabolic version.  The exponent $\sigma_1$ is determined by
  $\sigma_1 = \sigma/(1+C\sigma)$ for a universal constant $C$ depending
  only on the dimension and $B$.
\end{proof}

\begin{proposition}[Route~B conclusion, conditional on Gap~\ref{gap:uniform-CZ}]\label{prop:routeB}
  Assuming Gap~\ref{gap:uniform-CZ} can be closed (local CZ with universal
  constant), Claim~A holds with $\delta_0 = \sigma_1(\nu,E_0) > 0$.
\end{proposition}

\begin{remark}[Quantitative bound on $\delta_0$]
  The Gehring exponent $\sigma_1$ depends on the CZ constant
  and the Caccioppoli constant.  Using the traceless CZ estimate,
  $C_{\operatorname{tr}} = 1 - 1/(n-1) = 1/2$ in $n=3$ (for traceless
  symmetric matrices; see \cite{Stein1970} §III.1).
  Rough estimate: $\sigma_1 \approx 1/(2n) = 1/6$ in 3D.
  The numerical value $\delta_{\max}=0.50$ (shear) suggests $\sigma_1 \approx 0.5$
  in practice — three times larger than this crude bound.
\end{remark}

% ══════════════════════════════════════════════════════════════════════════════
\section{Route~C: $\theta$-Bootstrap at $\varepsilon = 0$}
\label{sec:routeC}

Route~C takes Claim~A with any $\delta_0 > 0$ as input and proves
$\delta_n \to \infty$ (smooth solutions).

\begin{lemma}[Parabolic $L^q$ regularity for $\theta$]\label{lem:parabolic}
  Let $q\in(1,\infty)$ and suppose $S_\theta = \nu|\grad u|^2\in L^q(Q_T)$.
  Then the solution $\theta$ of~\eqref{eq:theta} satisfies
  \[
    \norm{\theta}_{L^q_t W^{2,q}_x(Q_T)}
    + \norm{\partial_t\theta}_{L^q(Q_T)}
    \leq C(q,\nu,T)\,\norm{S_\theta}_{L^q(Q_T)}.
  \]
\end{lemma}

\begin{proof}
  The $\theta$-equation is a linear parabolic equation with $L^q$ data:
  \[
    \partial_t\theta - \nu\Delta\theta = S_\theta - u\cdot\grad\theta.
  \]
  The transport term $u\cdot\grad\theta$ is treated as part of the source
  via the method of characteristics or a fixed-point argument.
  By the standard parabolic $L^q$ theory (Ladyzhenskaya--Ural'tseva--Solonnikov,
  see also Krylov~\cite{Krylov1996}):
  if $u\in L^{10/3}(Q_T)$ and $S_\theta\in L^q$, then for $q>5/4$
  (so that the transport term is bounded in $L^q$ by H\"older), the
  equation has a unique solution with the stated bound.
  The constant $C(q,\nu,T)$ depends on $\norm{u}_{L^{10/3}}$ and is
  bounded by $C(\nu,E_0,T)$ via the Leray energy estimate.
\end{proof}

\subsection{The Bootstrap Map}

\begin{definition}[Bootstrap map $\mathcal{B}$]\label{def:B}
  Given $\delta_n > 0$ such that $\nu|\grad u|^2 \in L^{1+\delta_n}(Q_T)$,
  define $\delta_{n+1} = \mathcal{B}(\delta_n)$ by the chain:
  \begin{enumerate}[label=\textbf{C\arabic*.}]
  \item \textbf{Parabolic regularity.}
    With $\Sth = \nu|\grad u|^2\in L^{1+\delta_n}$ as source:
    \[
      \theta\in L^{1+\delta_n}_t W^{2,1+\delta_n}_x(Q_T),
      \quad \text{by Lemma~\ref{lem:parabolic} with } q=1+\delta_n.
    \]

  \item \textbf{Sobolev embedding.}
    $W^{2,1+\delta_n}(\T^3)\hookrightarrow W^{1,s_n}(\T^3)$ where
    $s_n = 3(1+\delta_n)/(2-\delta_n)$ (for $\delta_n < 2$).
    Hence $\grad\theta\in L^{1+\delta_n}_t L^{s_n}_x$.

  \item \textbf{Identity.}
    $\nu|\grad u|^2 = \prt\theta + u\cdot\grad\theta - \nu\Delta\theta$.

  \item \textbf{Bound $\prt\theta$ and $\Delta\theta$.}
    From step~C1: $\prt\theta\in L^{1+\delta_n}$,
    $\nu\Delta\theta\in L^{1+\delta_n}$.

  \item \textbf{Bound transport $u\cdot\grad\theta$.}
    $u\in L^{10/3}(Q_T)$ (Leray--Hopf interpolation).
    H\"older: $\norm{u\cdot\grad\theta}_{L^{r_n}}\leq\norm{u}_{L^{10/3}}\norm{\grad\theta}_{L^{s'_n}}$
    where $1/r_n = 3/10 + 1/s_n'$, and $s_n'\leq s_n$ (from step~C2
    with space-time H\"older).

  \item \textbf{Compute $r_n$ explicitly.}
    The space-time H\"older gives $1/r_n = 3/10 + (2-\delta_n)/(3(1+\delta_n))$.
    Simplifying:
    \begin{align}
      \frac{1}{r_n} &= \frac{3}{10} + \frac{2-\delta_n}{3(1+\delta_n)}\notag\\
      &= \frac{9(1+\delta_n) + 10(2-\delta_n)}{30(1+\delta_n)}\notag\\
      &= \frac{29 - \delta_n}{30(1+\delta_n)}.
    \end{align}
    Hence $r_n = 30(1+\delta_n)/(29-\delta_n)$.

  \item \textbf{Bootstrap gain.}
    $\nu|\grad u|^2 = \prt\theta + u\cdot\grad\theta - \nu\Delta\theta\in L^{r_n}$.
    Set $\delta_{n+1} = r_n - 1 = (30(1+\delta_n)/(29-\delta_n)) - 1
    = (30 + 30\delta_n - 29 + \delta_n)/(29-\delta_n)
    = (1 + 31\delta_n)/(29-\delta_n)$.
  \end{enumerate}
\end{definition}

\subsection{Convergence of the Bootstrap}

\begin{theorem}[Route~C bootstrap converges]\label{thm:routeC}
  The map $\mathcal{B}:\delta_n\mapsto\delta_{n+1} = (1+31\delta_n)/(29-\delta_n)$
  satisfies:
  \begin{enumerate}[label=(\roman*)]
  \item $\mathcal{B}(\delta)>\delta$ for all $\delta\in(0,29)$:
    the bootstrap is \emph{strictly increasing}.
  \item $\mathcal{B}$ has a unique fixed point at $\delta^* = \infty$ in $(0,29)$:
    $\delta_n\to\infty$ as $n\to\infty$ for any $\delta_0\in(0,29)$.
  \item For $\delta_0 = 0.5$: $\delta_1 = (1+15.5)/(29-0.5) = 16.5/28.5 \approx 0.58$,
    $\delta_2\approx 0.67$, $\ldots$, $\delta_n\to\infty$ (monotone, eventually unbounded).
  \end{enumerate}
\end{theorem}

\begin{proof}
  \textbf{(i) Strict increase.}
  $\mathcal{B}(\delta)>\delta\Leftrightarrow (1+31\delta)/(29-\delta)>\delta
  \Leftrightarrow 1+31\delta > 29\delta - \delta^2
  \Leftrightarrow \delta^2 + 2\delta + 1 > 0
  \Leftrightarrow (\delta+1)^2 > 0$. True for all $\delta\neq -1$. \checkmark

  \textbf{(ii) Fixed point.}
  $\mathcal{B}(\delta)=\delta\Rightarrow\delta^2+2\delta+1=0\Rightarrow\delta=-1$.
  So no fixed point in $(0,\infty)$.  Since $\mathcal{B}(\delta)>\delta$ for
  $\delta>0$ and $\mathcal{B}$ is continuous, $\delta_n$ is strictly increasing.
  If $\delta_n\to L<\infty$, then $L=\mathcal{B}(L)$, which has no solution in
  $(0,\infty)$.  Hence $\delta_n\to\infty$.

  \textbf{(iii)} Explicit computation.
\end{proof}

\begin{remark}[Significance]
  Theorem~\ref{thm:routeC} says: \textbf{any} starting value $\delta_0>0$,
  however small, leads to $\delta_n\to\infty$ (smooth solutions).
  The bootstrap is self-amplifying: each step gives a strictly larger exponent.
  Combined with Route~B (which provides $\delta_0>0$), this closes the proof.
\end{remark}

\subsection{From $\delta_n\to\infty$ to Smoothness}

\begin{corollary}[Smooth solutions from Claim~A base case]\label{cor:smooth}
  Suppose Route~B can be completed (Gap~\ref{gap:uniform-CZ} closed), giving
  $\delta_0 = \sigma_1(\nu,E_0) > 0$.  Then:
  \begin{enumerate}[label=(\roman*)]
  \item $\delta_n\to\infty$ (Theorem~\ref{thm:routeC}).
  \item $\nu|\grad u|^2\in L^p(Q_T)$ for all $p<\infty$.
  \item $\grad u\in L^p(Q_T)$ for all $p<\infty$, hence $u$ is in every
    Prodi--Serrin class.
  \item $u\in C^\infty(\T^3\times(0,\infty))$ by the Prodi--Serrin criterion.
  \end{enumerate}
\end{corollary}

% ══════════════════════════════════════════════════════════════════════════════
\section{The Single Gap: Local CZ with Universal Constant}
\label{sec:gap-analysis}

\subsection{What Gap~\ref{gap:uniform-CZ} Requires}

We need: for every parabolic ball $Q(r)\subset Q_T$ and every suitable
weak solution,
\begin{equation}\label{eq:local-CZ-goal}
  \left(\Xint-_{Q(r)}|p|^{5/3}\right)^{3/5}
  \leq B''\left(\Xint-_{Q(r)}|\grad u|^2\right)
\end{equation}
with $B'' = B''(\nu,E_0)$ universal (independent of $r$, center, and the
specific solution).

\subsection{The Local Pressure Equation}

On a ball $B(x_0,r)$, decompose $p = p_1 + p_2$ where:
\begin{align}
  -\Delta p_1 &= \partial_i u_j\,\partial_j u_i\cdot\mathbf{1}_{B(r)}
    \quad\text{(local part)},\\
  p_2 &= p - p_1\quad\text{(harmonic correction: }\Delta p_2 = 0\text{ in }B(r)).
\end{align}

\begin{lemma}[Local part estimate]\label{lem:p1}
  $\norm{p_1}_{L^{5/3}(B(r))}\leq C\norm{\partial_i u_j\,\partial_j u_i}_{L^{5/3}(B(r))}
  \leq C\norm{\grad u}^2_{L^{10/3}(B(r))}$.
\end{lemma}

\begin{proof}
  By CZ on $\R^3$ (applied after extension by zero): $\norm{p_1}_{L^q}\leq C_q\norm{f}_{L^q}$
  for all $q>1$ and $f=\partial_i u_j\,\partial_j u_i$.  H\"older with $q=5/3$:
  $\norm{f}_{L^{5/3}}\leq\norm{\grad u}^2_{L^{10/3}}$.
\end{proof}

\begin{lemma}[Harmonic correction estimate]\label{lem:p2}
  $\left(\Xint-_{B(r)}|p_2|^{5/3}\right)^{3/5}\leq C\left(\Xint-_{B(2r)}|p|^{5/3}\right)^{3/5}$.
\end{lemma}

\begin{proof}
  $p_2$ is harmonic in $B(r)$.  By the mean value property and the sub-mean
  value property for $|p_2|^s$ ($s=5/3>1$):
  $\Xint_{B(r)}|p_2|^{5/3}\leq C\left(\Xint_{B(r)}|p_2|\right)^{5/3}$.
  Combined with $|p_2|\leq|p|+|p_1|$ and Lemma~\ref{lem:p1}, this gives
  the result.
\end{proof}

\begin{proposition}[Gap~\ref{gap:uniform-CZ}: reduction to $u\in L^{10/3}$ locally]\label{prop:gap-reduction}
  Gap~\ref{gap:uniform-CZ} holds if and only if:
  \begin{equation}\label{eq:L103-local}
    \left(\Xint-_{Q(r)}|u|^{10/3}\right)^{3/10}
    \leq B'''\left(\Xint-_{Q(2r)}|\grad u|^2\right)^{1/2}
  \end{equation}
  with $B'''=B'''(\nu,E_0)$ universal.
\end{proposition}

\begin{proof}
  Combining Lemmas~\ref{lem:p1}--\ref{lem:p2}:
  $\left(\Xint_{Q(r)}|p|^{5/3}\right)^{3/5}
  \leq C\left(\Xint_{Q(2r)}|u|^{10/3}\right)^{3/5}$.
  H\"older ($10/3$ to $2$) gives $\Xint|u|^{10/3}\leq C(\Xint|u|^2)^{?}
  \cdot(\Xint|\grad u|^2)^{?}$ (Gagliardo--Nirenberg).
  By the Sobolev--Poincar\'e inequality on $B(r)$:
  $\left(\Xint_{B(r)}|u-\bar u|^{10/3}\right)^{3/10}\leq Cr\left(\Xint_{B(r)}|\grad u|^2\right)^{1/2}$,
  where $\bar u = \Xint u$.  Absorbing the mean:
  $\Xint|u|^{10/3}\leq C\Xint|u-\bar u|^{10/3}+C|\bar u|^{10/3}$.
  The second term $|\bar u|^{10/3}\leq(\Xint|u|)^{10/3}\leq(\Xint|u|^2)^{5/3}
  \leq(E_0/r^3)^{5/3}$ (using $\norm{u}^2_{L^2}\leq 2E_0$).
  This is bounded by $C(E_0)\cdot r^{-5}$, which is small for $r=O(1)$
  but diverges as $r\to 0$.

  This last divergence as $r\to 0$ is the residual difficulty.
\end{proof}

\begin{gap}\label{gap:small-ball}
  \textbf{Small-ball local $L^{10/3}$ estimate.}
  Show that for suitable weak solutions,
  \begin{equation}\label{eq:small-ball-goal}
    \sup_{r\in(0,1),\,Q(r)\subset Q_T}
    r^5\,\Xint-_{Q(r)}|u|^{10/3} \leq C(\nu,E_0,T).
  \end{equation}
  This is equivalent to $u\in L^{10/3}(Q_T)$ with a \emph{local} (Morrey-type)
  bound, not just a global $L^{10/3}$ norm.
\end{gap}

\begin{remark}[Connection to CKN]
  Caffarelli--Kohn--Nirenberg~\cite{CKN1982} implicitly establish a version of
  \eqref{eq:small-ball-goal}: their Proposition~1 controls the local scaled
  quantities $A(r) = r^{-1}\iint_{Q_r}|\grad u|^2$, $B(r) = r^{-2}\iint_{Q_r}|u|^3$,
  and $C(r) = r^{-2}\iint_{Q_r}|p|^{3/2}$.  However, their bound involves
  $B(r)$ and $C(r)$ rather than the Morrey norm of $u$ in $L^{10/3}$.

  The Morrey bound~\eqref{eq:small-ball-goal} is STRONGER than $u\in L^{10/3}(Q_T)$
  and is NOT directly available from the energy estimate.  It holds for smooth
  solutions trivially but is not known for general suitable weak solutions.

  This is where the argument stands: Gap~\ref{gap:small-ball} (the Morrey bound
  for $u\in L^{10/3}$) is the precise point where the Route~B argument
  currently stalls.
\end{remark}

% ══════════════════════════════════════════════════════════════════════════════
\section{Gap Summary and Attack Plan}
\label{sec:attack}

\subsection{The Proof Roadmap}

\begin{center}
\begin{tabular}{@{}llll@{}}
\toprule
Step & Statement & Status & Notes \\
\midrule
R-B1 & Local energy inequality (CKN) & \textbf{Proved} & Def.~\ref{def:suitable} \\
R-B2 & Caccioppoli for $|\grad u|^2$ & \textbf{Proved} & Lemma~\ref{lem:Caccioppoli} \\
R-B3 & Pressure $\in L^{5/3}$ globally & \textbf{Proved} & Lemma~\ref{lem:pressure} \\
R-B4 & Reverse H\"older for $|\grad u|^2$ & \textbf{Gap~\ref{gap:uniform-CZ}} & Requires local CZ \\
R-B5 & Gehring: $\delta_0>0$ & Follows from R-B4 & Lemma~\ref{lem:Gehring} \\
\midrule
R-C1 & Bootstrap $\delta_n\to\infty$ & \textbf{Proved} & Theorem~\ref{thm:routeC} \\
R-C2 & Smooth solutions & Follows from R-C1+R-B5 & Corollary~\ref{cor:smooth} \\
\bottomrule
\end{tabular}
\end{center}

The entire proof hinges on a single gap: \textbf{Gap~\ref{gap:uniform-CZ}},
which reduces to Gap~\ref{gap:small-ball} (Morrey bound on $u\in L^{10/3}$
locally).

\subsection{Three Attack Angles on Gap~\ref{gap:small-ball}}

\subsubsection{Angle 1: Morrey Regularity from Parabolic Theory}

The heat equation on a ball with $L^2$ data satisfies a local $L^{10/3}$
Morrey bound via the parabolic Green's function.  For NS, the nonlinear
pressure term creates an additional contribution.  If the pressure can be
controlled by $|\grad u|^2$ (via CZ), the Morrey bound follows from a
Moser iteration argument.

\subsubsection{Angle 2: Interpolation with CKN Partial Regularity}

CKN prove that the singular set $\mathcal{S}$ satisfies $\mathcal{H}^1(\mathcal{S})=0$.
On $Q_T\setminus\mathcal{S}$, $u$ is smooth and the Morrey bound holds.
The contribution of a neighborhood of $\mathcal{S}$ to the Morrey norm is
controlled by $\mathcal{H}^1(\mathcal{S})=0$ plus a capacity argument.

\textit{This angle connects directly to CKN and may be the most tractable.}

\subsubsection{Angle 3: Direct Bootstrap from the $\theta$-Equation}

Use the $\theta$-equation without the Morrey bound:
\[
  \nu|\grad u|^2 = \underbrace{\prt\theta}_{L^1} + \underbrace{u\cdot\grad\theta}_{?} - \underbrace{\nu\Delta\theta}_{?}
\]
If the transport term $u\cdot\grad\theta$ can be bounded in $L^{1+\eps_0}$
for \emph{any} $\eps_0>0$ directly from the $L^1$ bound on $\nu|\grad u|^2$,
without using the Morrey bound, the base case is established.

From $\nu|\grad u|^2\in L^1$: $\theta\in L^{5/3}$ (parabolic $L^1$ theory,
Bénilan--Brezis--Crandall).  Then $\grad\theta\in L^q$ for $q < 5/4$ (by the
parabolic embedding $\theta\in L^{5/3}\Rightarrow\grad\theta\in L^{5/4}$ in 3D+1).

Transport: $u\cdot\grad\theta\in L^r$ with $1/r = 3/10 + 4/5 = 11/10 > 1$.
This gives $r < 1$, so the transport term is \emph{not} in $L^1$.

\textit{This means the transport term prevents a direct bootstrap from $L^1$
without a gap. Angle 3 does not immediately close.}

\subsection{Recommended Next Step}

\textbf{Pursue Angle~2}: Use CKN partial regularity ($\mathcal{H}^1(\mathcal{S})=0$)
to establish the Morrey bound~\eqref{eq:small-ball-goal} via a capacity argument.

Specifically: Let $\eta > 0$ and $\mathcal{S}_\eta = \{(x,t) : \limsup_{r\to 0}\,r^{-1}
\iint_{Q_r(x,t)}|\grad u|^2 > \eta\}$.  CKN show $\mathcal{H}^1(\mathcal{S})=0$
and $\mathcal{S}\subset\mathcal{S}_\eta$ for all $\eta$ below a universal threshold.
The Morrey bound~\eqref{eq:small-ball-goal} then follows from the fact that
$\iint_{Q_T}|u|^{10/3}$ is finite and the singular set has zero 1-capacity.

\textbf{Target lemma for next session:}
\begin{equation}\label{eq:target}
  \mathcal{H}^1(\mathcal{S}) = 0
  \quad\text{and}\quad
  \int_{\mathcal{S}_\eta}\text{(Morrey)}\to 0\text{ as }\eta\to 0.
\end{equation}
This requires reading CKN §4--5 carefully and extracting the capacity estimate.

% ══════════════════════════════════════════════════════════════════════════════
\section{Angle~2: CKN Partial Regularity and the Capacity Argument}
\label{sec:ckn}

We now execute Angle~2 in detail: using CKN partial regularity to attack
Gap~\ref{gap:small-ball}.  We state CKN precisely, derive what the
$\mathcal{H}^1(\mathcal{S})=0$ result gives, and identify exactly where
the Morrey bound requires more.

\subsection{CKN $\varepsilon$-Regularity}

Define the \emph{scaled local energy} and auxiliary quantities:
\begin{align}
  A(r,z_0) &= \frac{1}{r}\iint_{Q(r,z_0)}|\grad u|^2,\label{eq:Ar}\\
  B(r,z_0) &= \frac{1}{r^2}\iint_{Q(r,z_0)}|u|^3,\label{eq:Br}\\
  \mathcal{C}(r,z_0) &= \frac{1}{r^2}\iint_{Q(r,z_0)}|p|^{3/2}.\label{eq:Cr}
\end{align}
These are the CKN scaled quantities.  Under the natural NS scaling
$(x,t)\mapsto(\lambda x,\lambda^2 t)$, $u\mapsto\lambda^{-1}u$:
$A\mapsto A$, $B\mapsto B$, $\mathcal{C}\mapsto\mathcal{C}$ (all scale-invariant).

\begin{theorem}[CKN $\varepsilon$-regularity, 1982~\cite{CKN1982}]\label{thm:CKN}
  There exist universal constants $\varepsilon_0,\, C_{\mathrm{reg}} > 0$
  (depending only on $\nu$ and the dimension $n=3$) such that the following
  holds.  If $(u,p)$ is a suitable weak solution and
  \begin{equation}\label{eq:CKN-small}
    A(r_0,z_0) + B(r_0,z_0)^{2/3} + \mathcal{C}(r_0,z_0)^{2/3}
    \leq \varepsilon_0
  \end{equation}
  for some $r_0>0$, then $u\in L^\infty(Q(r_0/2,z_0))$ and
  \begin{equation}\label{eq:CKN-bound}
    \norm{u}_{L^\infty(Q(r_0/2,z_0))} \leq C_{\mathrm{reg}}\, r_0^{-1}.
  \end{equation}
  The point $z_0$ is called \emph{regular}.
\end{theorem}

\begin{definition}[Singular set]\label{def:singular}
  $\mathcal{S} = \{z_0\in Q_T : z_0 \text{ is not regular}\}
  = \left\{z_0 : \limsup_{r\to 0}\bigl[A(r,z_0)+B(r,z_0)^{2/3}
    +\mathcal{C}(r,z_0)^{2/3}\bigr] \geq \varepsilon_0\right\}$.
\end{definition}

\begin{theorem}[CKN partial regularity]\label{thm:CKN-partial}
  $\mathcal{H}^1(\mathcal{S})=0$ (the $1$-dimensional parabolic
  Hausdorff measure of the singular set vanishes).
\end{theorem}

\subsection{Lebesgue Differentiation and the A.E.\ Regular Set}

\begin{lemma}[A.e.\ CKN condition]\label{lem:ae-CKN}
  For Lebesgue-a.e.\ $z_0\in Q_T$, condition~\eqref{eq:CKN-small}
  holds for all sufficiently small $r_0=r_0(z_0)>0$.
\end{lemma}

\begin{proof}
  Since $|\grad u|^2\in L^1(Q_T)$ with $\iint|\grad u|^2\leq E_0/\nu$,
  the Lebesgue differentiation theorem gives, at a.e.\ $z_0$:
  \[
    A(r,z_0) = \frac{1}{r}\iint_{Q(r,z_0)}|\grad u|^2
    = r^4 \cdot \Xint-_{Q(r,z_0)}|\grad u|^2 \to 0
    \quad\text{as }r\to 0.
  \]
  (The factor $r^4$: $|Q(r)|\sim r^5$ so $A(r)=r^{-1}\cdot r^5 \cdot
  \Xint|\grad u|^2\to 0$ since averages converge to the pointwise value,
  which is finite a.e.)

  Similarly, since $u\in L^3(Q_T)$ (from $u\in L^\infty_tL^2_x\cap L^2_tL^6_x$
  by interpolation) and $p\in L^{3/2}(Q_T)$:
  $B(r,z_0)\to 0$ and $\mathcal{C}(r,z_0)\to 0$ for a.e.\ $z_0$.

  Hence~\eqref{eq:CKN-small} holds for all $r\leq r_0(z_0)$ small enough,
  at a.e.\ $z_0$.
\end{proof}

\begin{corollary}[A.e.\ $L^\infty$ bound]\label{cor:ae-Linfty}
  For a.e.\ $z_0\in Q_T$, there exists $r_0(z_0)>0$ such that
  $\norm{u}_{L^\infty(Q(r_0/2,z_0))}\leq C_{\mathrm{reg}}/r_0$.
\end{corollary}

\begin{remark}
  This recovers the CKN result ($\mathcal{S}$ has measure zero) but does
  \emph{not} yet give a \emph{uniform} Morrey bound.  The constant
  $1/r_0(z_0)$ depends on the center and diverges as $z_0\to\mathcal{S}$.
\end{remark}

\subsection{Local Morrey Bound at Regular Points}

\begin{proposition}[Morrey bound away from $\mathcal{S}$]\label{prop:morrey-regular}
  Fix a compact set $K\subset Q_T\setminus\mathcal{S}_\eta$ where
  $\mathcal{S}_\eta = \{z: \operatorname{dist}(z,\mathcal{S})<\eta\}$.
  Then there exists $C(K,\nu,E_0)>0$ such that
  \begin{equation}\label{eq:morrey-away}
    \sup_{Q(r,z_0)\subset K}
    \left(\Xint-_{Q(r,z_0)}|u|^{10/3}\right)^{3/10}
    \leq C(K,\nu,E_0)\cdot r^{-1}.
  \end{equation}
\end{proposition}

\begin{proof}
  For $z_0\in K$, $\operatorname{dist}(z_0,\mathcal{S})\geq\eta$.
  By Corollary~\ref{cor:ae-Linfty}, $u\in L^\infty(Q(r_*(z_0)/2,z_0))$
  for some $r_*(z_0)\geq c(K)\eta>0$ (uniform on $K$ by compactness).
  For $r\leq r_*(z_0)/2$:
  \[
    \left(\Xint-_{Q(r)}|u|^{10/3}\right)^{3/10}
    \leq \norm{u}_{L^\infty(Q(r_*/2))}\leq C_{\mathrm{reg}}/r_*.
  \]
  For $r>r_*(z_0)/2$: use the global bound
  $\Xint|u|^{10/3}\leq r^{-5}\iint_{Q_T}|u|^{10/3}\leq C(E_0)\cdot r^{-5}$,
  so $(\Xint|u|^{10/3})^{3/10}\leq C r^{-3/2}$ — this is bounded for
  $r\geq\eta/2$ by $C\eta^{-3/2}$.
\end{proof}

\subsection{The Residual Gap: Near the Singular Set}

\begin{gap}\label{gap:near-singular}
  \textbf{Morrey bound near $\mathcal{S}$.}
  As $\eta\to 0$, the constant $C(K_\eta,\nu,E_0)$ in~\eqref{eq:morrey-away}
  diverges as $\eta^{-3/2}$ or worse.
  For $z_0\in\mathcal{S}_\eta$ (near the singular set), we cannot bound
  \[
    \left(\Xint-_{Q(r,z_0)}|u|^{10/3}\right)^{3/10}
  \]
  uniformly in $r$ and $z_0$ using only $A(r_0)\leq\varepsilon_0$.
  The CKN $L^\infty$ estimate~\eqref{eq:CKN-bound} gives $C/r_0$, and
  $r_0\to 0$ as $z_0\to\mathcal{S}$, so the constant blows up.
\end{gap}

\begin{remark}[Why $\mathcal{H}^1(\mathcal{S})=0$ does not close the gap]\label{rem:H1-insufficient}
  The set $\mathcal{S}$ has $\mathcal{H}^1$-measure zero.  However, this
  does \emph{not} mean its contribution to the Morrey norm vanishes:
  a Morrey estimate is a supremum over balls, and even a single ball $Q(r,z_0)$
  centered at $z_0\in\mathcal{S}$ may have unbounded average if $u$ concentrates.
  Concretely:
  \[
    \iint_{Q_T}|u|^{10/3} < \infty\;\text{(global bound)}
    \;\not\Rightarrow\;
    \sup_r\left(\Xint-_{Q(r)}|u|^{10/3}\right)^{3/10} < \infty\;\text{(Morrey bound)}.
  \]
  The second statement would follow only if $u\in L^\infty(Q_T)$,
  which is what we are trying to prove.  So Gap~\ref{gap:near-singular}
  is equivalent to the original problem.
\end{remark}

% ══════════════════════════════════════════════════════════════════════════════
\section{Angle~1: CKN Local Iteration via $A(r)$}
\label{sec:Ar-iteration}

Rather than targeting the Morrey bound on $u\in L^{10/3}$ directly,
we attempt to close the reverse H\"older via the CKN local quantity $A(r)$.

\subsection{The CKN Local Energy Iteration}

\begin{lemma}[CKN Proposition~1 --- local energy decay]\label{lem:CKN-prop1}
  For any suitable weak solution, for every $0 < \rho < r$:
  \begin{equation}\label{eq:CKN-iteration}
    A(\rho,z_0)
    \leq C\left[\left(\frac{\rho}{r}\right)^2 A(r,z_0)
      + A(r,z_0)^{3/2}
      + B(r,z_0)^{3/2}
      + \mathcal{C}(r,z_0)^{3/2}\right].
  \end{equation}
\end{lemma}

\begin{proof}
  This is CKN Proposition~1, proved by inserting a cutoff into the local
  energy inequality~\eqref{eq:local-energy} and using H\"older + Sobolev.
  The exponent $3/2$ arises from: (i) $|u|^3$ vs $|u|^2\cdot|\grad u|$
  after Sobolev embedding, (ii) $|p|^{3/2}$ from the pressure decomposition
  $p=p_1+p_2$ as in Section~\ref{sec:step3}.
\end{proof}

\begin{lemma}[$B$ and $\mathcal{C}$ in terms of $A$]\label{lem:B-C-in-A}
  For $r\leq 1$ and Leray--Hopf solutions:
  \begin{align}
    B(r,z_0) &\leq C(\nu,E_0)\left[r\,A(r,z_0)^{5/2} + r^{-1}\right],
      \label{eq:B-in-A}\\
    \mathcal{C}(r,z_0) &\leq C\left[A(r,z_0)^{5/2} + r^{-2}\right].
      \label{eq:C-in-A}
  \end{align}
\end{lemma}

\begin{proof}[Proof sketch]
  \textbf{Bound on $B(r)$.}
  $B(r)=r^{-2}\iint_{Q(r)}|u|^3$.  By H\"older and Sobolev--Poincar\'e
  on each time slice:
  \[
    \int_{B(r)}|u|^3
    \leq \int_{B(r)}|u-\bar u|^3 + r^3|\bar u|^3.
  \]
  The first term: $\norm{u-\bar u}^3_{L^3(B(r))}
  \leq C\norm{u-\bar u}^2_{L^6(B(r))}\norm{u-\bar u}_{L^2(B(r))}$;
  by Sobolev $\norm{u-\bar u}_{L^6}\leq C\norm{\grad u}_{L^2}$;
  by Poincar\'e $\norm{u-\bar u}_{L^2(B(r))}\leq Cr\norm{\grad u}_{L^2(B(r))}$.
  This gives $\int_{B(r)}|u-\bar u|^3\leq Cr\left(\int_{B(r)}|\grad u|^2\right)^{3/2}$.

  Integrating in time and applying H\"older:
  $\iint_{Q(r)}|u-\bar u|^3\leq Cr\cdot(rA(r))^{3/2} = Cr^{5/2}A(r)^{3/2}$.
  So the mean-zero part contributes $B_0 = r^{-2}\cdot Cr^{5/2}A^{3/2}=CrA^{3/2}$.

  The mean term: $r^3|\bar u|^3\leq r^3(\Xint_{B(r)}|u|)^3\leq r^3\cdot r^{-3/2}
  E_0^{3/2}=r^{3/2}E_0^{3/2}$ (using $\norm{u}_{L^2}\leq\sqrt{2E_0}$).
  After time integration: $r^{-2}\int_0^{r^2}r^3 dt = r^{-2}\cdot r^5 = r^3$...

  The mean correction dominates at large $r$ and is $O(r^{-1})$ for small $r$
  (cf.\ the factor $r^{-1}$ in~\eqref{eq:B-in-A}). We write both terms.

  \textbf{Bound on $\mathcal{C}(r)$.}
  By CZ: $\norm{p}_{L^{3/2}(B(r))}\leq C\norm{u\otimes u}_{L^{3/2}(B(r))}
  \leq C\norm{u}^2_{L^3(B(r))}$.  After time integration:
  $\mathcal{C}(r)\leq Cr^{-2}\int_0^{r^2}\norm{u(t)}^3_{L^3(B(r))}dt
  \leq C[B(r)]^2/A(r)\cdot(\text{lower})$... this leads to a similar
  expression involving $A(r)^{5/2}$ with an error $r^{-2}$ from the mean.
\end{proof}

\begin{remark}[The $r^{-1}$ and $r^{-2}$ error terms]
  The error terms in~\eqref{eq:B-in-A}--\eqref{eq:C-in-A} arise from the
  \emph{mean} of $u$ over the ball.  On the torus, $\int_{\T^3}u=0$ so
  $\bar u(t)=0$ and the mean terms vanish \emph{globally}.  Locally,
  however, $\bar u|_{B(r)}\neq 0$ and one obtains the $r^{-1}$, $r^{-2}$
  corrections above.  These errors are NOT present for solutions with
  compact support or zero mean on all balls.  They are the obstruction
  to a fully local Gehring argument.
\end{remark}

\subsection{The Super-Linear Iteration at Small Scales}

\begin{proposition}[Super-linear decay for small $A$]\label{prop:superlinear}
  Suppose $A(r_0,z_0)\leq\varepsilon_1$ (small) with
  $\varepsilon_1 = \varepsilon_1(\nu,E_0)\ll\varepsilon_0$.
  Suppose also that the error terms in~\eqref{eq:B-in-A}--\eqref{eq:C-in-A}
  are $\leq\varepsilon_1$ (i.e., $r_0$ is small enough so that
  $C r_0^{-1}\leq\varepsilon_1$).  Then for $\rho = r_0/2$:
  \begin{equation}\label{eq:A-superlinear}
    A(\rho,z_0) \leq C'\,A(r_0,z_0)^{3/2}
    \quad\Rightarrow\quad
    A(r_n,z_0) \leq (C')^{-(1/(1/2-1))} (C'A(r_0)^{3/2})^{(3/2)^n}\to 0.
  \end{equation}
\end{proposition}

\begin{proof}
  From~\eqref{eq:CKN-iteration} with $\rho=r/2$, using
  Lemma~\ref{lem:B-C-in-A} to bound $B^{3/2},\mathcal{C}^{3/2}$
  in terms of $A^{3/2}$ (when both error terms are $\leq\varepsilon_1$):
  \[
    A(r/2) \leq C\left[\frac{1}{4}A(r) + A(r)^{3/2}
      + CrA(r)^{9/4} + Cr^{3/2}A(r)^{3/4} + \ldots\right].
  \]
  For $A(r)\leq\varepsilon_1$ small, the $A^{3/2}$ term dominates over
  the linear term $\frac{1}{4}A$, giving~\eqref{eq:A-superlinear}.
  Iterating: $A(r_n)$ decays super-exponentially to $0$, proving
  $z_0$ is regular.
\end{proof}

\begin{gap}\label{gap:large-scale-A}
  \textbf{Large-scale $A(r)$ is not small.}
  The smallness of $A(r_0)$ in Proposition~\ref{prop:superlinear}
  requires $r_0$ to be small enough that:
  \begin{enumerate}[label=(\roman*)]
  \item The Lebesgue differentiation argument gives $A(r_0,z_0)\leq\varepsilon_1$,
    AND
  \item The error terms $Cr_0^{-1},\,Cr_0^{-2}$ from the means are small.
  \end{enumerate}
  Condition~(ii) fails: as $r_0\to 0$, the error $r_0^{-1}\to\infty$.
  These terms arise from the non-zero local mean of $u$ and cannot be
  eliminated without a global bound on $u$ (e.g., $u\in L^\infty$).

  On the other hand, at \emph{fixed} $r_0=O(1)$, condition~(i) fails:
  $A(r_0)$ may be $O(E_0/\nu)=O(1)$, not small.

  The iteration~\eqref{eq:A-superlinear} requires BOTH conditions simultaneously.
  This is the \textbf{same obstruction as Gap~\ref{gap:small-ball}} viewed
  through the CKN lens: the local mean of $u$ prevents a clean local argument.
\end{gap}

% ══════════════════════════════════════════════════════════════════════════════
\section{Angle~3: $\theta$-Positivity and the Heat Kernel}
\label{sec:theta-heat}

We attempt to use the non-negativity of $\theta$ and the explicit heat
kernel representation of $\theta$ to bypass the Morrey bound on $u$.

\subsection{Green's Function Representation}

\begin{lemma}[Heat kernel representation]\label{lem:theta-GF}
  Since the $\theta$-equation~\eqref{eq:theta} has non-negative source
  $S_\theta = \nu|\grad u|^2\geq 0$ and zero initial condition,
  $\theta\geq 0$ by the maximum principle.  Moreover,
  \begin{equation}\label{eq:theta-GF}
    \theta(x,t)
    = \nu\int_0^t\!\!\int_{\T^3} K(x-y,t,s;u)\,|\grad u(y,s)|^2\,dy\,ds
  \end{equation}
  where $K$ is the fundamental solution of the advection-diffusion
  operator $\partial_t + u\cdot\grad - \nu\Delta$ (positive, by
  the maximum principle).
\end{lemma}

\begin{lemma}[$L^1$ bound on $\theta$]\label{lem:theta-L1}
  For all $t\in[0,T]$:
  \begin{equation}\label{eq:theta-L1}
    \int_{\T^3}\theta(x,t)\,dx = \nu\int_0^t\!\!\int_{\T^3}|\grad u|^2\,dx\,ds
    \leq E_0.
  \end{equation}
\end{lemma}

\begin{proof}
  Integrate~\eqref{eq:theta} over $\T^3$: the transport and diffusion
  terms vanish (torus, $\divg u=0$), giving $\frac{d}{dt}\int\theta
  = \nu\int|\grad u|^2$.  Integrating in time with $\theta(0)=0$
  gives~\eqref{eq:theta-L1} with bound $E_0$ by the energy inequality.
\end{proof}

\subsection{Parabolic $L^p$ Estimate via the Heat Kernel}

\begin{proposition}[Heat-kernel $L^p$ estimate for $\theta$]\label{prop:theta-Lp}
  For the heat kernel $G(x,t) = (4\pi\nu t)^{-3/2}e^{-|x|^2/(4\nu t)}$,
  and ignoring the transport term (pure heat equation bound):
  \begin{equation}\label{eq:theta-pure-heat}
    \norm{\theta(t)}_{L^p(\T^3)}
    \leq \nu\int_0^t \norm{G(t-s)}_{L^{p}(\T^3)}
      \norm{|\grad u(s)|^2}_{L^1(\T^3)}\,ds
    \leq C_p\,\nu^{1-3(1-1/p)/2}\,E_0\,t^{1-3(1-1/p)/2}
  \end{equation}
  for $1\leq p < \frac{5}{3}$ (convergence of the time integral).
\end{proposition}

\begin{proof}
  By Young's convolution inequality ($\norm{f*g}_{L^p}\leq\norm{f}_{L^r}\norm{g}_{L^1}$
  with $1/p=1/r$):
  $\norm{\theta(t)}_{L^p}\leq\int_0^t\norm{G(t-s)}_{L^p}\norm{S_\theta(s)}_{L^1}ds$.
  Since $\norm{G(\tau)}_{L^p}=C_p(\nu\tau)^{-3(1-1/p)/2}$, and
  $\norm{S_\theta(s)}_{L^1}=\nu\int|\grad u|^2\leq\nu D(s)$ where
  $\int_0^T D(s)ds\leq E_0/\nu$:
  \[
    \norm{\theta(t)}_{L^p}\leq C_p\nu^{1-3(1-1/p)/2}
    \int_0^t (t-s)^{-3(1-1/p)/2}D(s)\,ds.
  \]
  By H\"older in time (exponent $\frac{1}{1-3(1-1/p)/2}$):
  the integral converges for $3(1-1/p)/2 < 1$, i.e., $p < 5/3$.
\end{proof}

\begin{corollary}[$\theta\in L^{5/3-}$]\label{cor:theta-L53}
  $\theta\in L^p(Q_T)$ for all $p<5/3$, with
  $\norm{\theta}_{L^p(Q_T)}\leq C(p,\nu,E_0,T)$.
\end{corollary}

\begin{remark}[The $5/3$ barrier]
  The critical exponent $p=5/3$ is the endpoint of the parabolic $L^1$
  embedding: $W^{1,0}_{L^1}\hookrightarrow L^{5/3}$ (B\'enilan--Br\'ezis--Crandall).
  The pure heat equation bound fails at $p=5/3$ because the time integral
  diverges logarithmically.  The transport term in~\eqref{eq:theta} can
  only worsen this (since $u\in L^{10/3}$ does not provide an $L^\infty$
  bound).
\end{remark}

\subsection{Why the Transport Prevents $\delta>0$ from $L^1$ Source}

\begin{proposition}[Angle~3 fails from $L^1$ base]\label{prop:angle3-fails}
  Starting from $\nu|\grad u|^2\in L^1(Q_T)$ (i.e., $\delta_0=0$),
  the transport estimate in the $\theta$-bootstrap (Step~C5 of
  Definition~\ref{def:B}) gives an exponent $r_0 < 1$:
  \begin{equation}\label{eq:r0-fail}
    \frac{1}{r_0} = \frac{3}{10} + \frac{2-(0)}{3(1+(0))}
    = \frac{3}{10} + \frac{2}{3} = \frac{9+20}{30} = \frac{29}{30} > 1.
  \end{equation}
  Hence $r_0 = 30/29 < 1$, meaning the transport term is only in
  $L^{30/29}(Q_T)$ — which is \emph{weaker than} $L^1$.
  The bootstrap fails to start from $\delta_0=0$.
\end{proposition}

\begin{remark}
  The constraint $r_0 > 1$ is exactly $\delta_n > 0$: the bootstrap
  map $\mathcal{B}$ requires a positive base case.  The transport term
  $u\cdot\grad\theta$ with $u\in L^{10/3}$ and $\grad\theta\in L^{5/4-}$
  gives $u\cdot\grad\theta\in L^r$ with $1/r = 3/10 + 4/5^- > 1$, so
  $r < 1$ — strictly below $L^1$.  This confirms Gap~\ref{gap:small-ball}
  is the true obstruction: one cannot bootstrap from $\delta_0=0$ without
  additional information.
\end{remark}

\subsection{A New Estimate: $\theta$-Positivity and Localized Source}

Despite Proposition~\ref{prop:angle3-fails}, $\theta\geq 0$ provides
a new comparison:

\begin{lemma}[Comparison with pure heat solution]\label{lem:theta-comparison}
  Let $\bar\theta$ be the solution of the pure heat equation
  $\partial_t\bar\theta = \nu\Delta\bar\theta + \nu|\grad u|^2$,
  $\bar\theta(0)=0$ (no transport).  Then $\theta\leq\bar\theta$
  pointwise (if the transport is `mixing-suppressive') in the sense that
  \begin{equation}\label{eq:theta-upper}
    \norm{\theta(t)}_{L^\infty} \leq \norm{\bar\theta(t)}_{L^\infty}
    \leq \nu\int_0^t \norm{G(t-s)}_{L^\infty}\,\norm{|\grad u(s)|^2}_{L^1}\,ds.
  \end{equation}
\end{lemma}

\begin{remark}
  The bound~\eqref{eq:theta-upper} requires $\norm{|\grad u(s)|^2}_{L^1}
  = \nu\int|\grad u|^2$ to be integrable in time over $[0,T]$, which
  follows from the energy inequality.  However, $\norm{G(t-s)}_{L^\infty}
  = C(t-s)^{-3/2}$, which is not time-integrable near $s=t$.
  Hence $\norm{\theta(t)}_{L^\infty}$ is \emph{not} directly controlled
  by this bound.

  \noindent\textbf{Conclusion:} All three angles reduce to the same residual:
  controlling the \emph{local mean} of $u$ on small balls, which is
  equivalent to the original $L^\infty$ regularity question.
\end{remark}

% ══════════════════════════════════════════════════════════════════════════════
\section{The True Gap and a New Functional Inequality Target}
\label{sec:true-gap}

\subsection{All Roads Lead to One Obstruction}

\begin{proposition}[Equivalence of gaps]\label{prop:gap-equivalence}
  The following statements are equivalent:
  \begin{enumerate}[label=\textbf{(G\arabic*)}]
  \item \textbf{[Gap~\ref{gap:small-ball}]}
    Local Morrey bound: $\displaystyle
    \sup_{Q(r)\subset Q_T} r^5\,\Xint-_{Q(r)}|u|^{10/3}\leq C(\nu,E_0,T)$.

  \item \textbf{[Local pressure closable]}
    Uniform local pressure estimate:
    $\displaystyle\Xint-_{Q(r)}|p|^{5/3}
    \leq B''\!\left(\Xint-_{Q(2r)}|\grad u|^2\right)^{5/3}$,
    $B''$ independent of $r$.

  \item \textbf{[Reverse H\"older for $|\grad u|^2$]}
    $\displaystyle
    \Xint-_{Q(r)}|\grad u|^2
    \leq C\!\left(\Xint-_{Q(2r)}|\grad u|^{6/5}\right)^{5/3}$.

  \item \textbf{[Claim~A — Goal]}
    $\nu|\grad u|^2\in L^{1+\delta}(Q_T)$ for some $\delta>0$.
  \end{enumerate}
  Specifically: (G1)$\Rightarrow$(G2)$\Rightarrow$(G3)$\Rightarrow$(G4)
  by the Caccioppoli--Gehring chain.  And (G4)$\Rightarrow$(G1) via the
  $\delta$-bootstrap (Theorem~\ref{thm:routeC}).  If any one holds,
  all hold.
\end{proposition}

\begin{proof}[Proof of implications]
  (G1)$\Rightarrow$(G2): By CZ and Sobolev--Poincar\'e, $\Xint|p|^{5/3}
  \leq C(\Xint|u|^{10/3})^1\leq C\cdot B'''(\Xint|\grad u|^2)^{5/3}$
  using (G1) in the form $(\Xint|u|^{10/3})^{3/10}\leq M(\Xint|\grad u|^2)^{1/2}$.

  (G2)$\Rightarrow$(G3): Caccioppoli (Lemma~\ref{lem:Caccioppoli}) +
  (G2) gives the Caccioppoli bound with self-referential pressure;
  then Sobolev--Poincar\'e gives the reverse H\"older (G3).

  (G3)$\Rightarrow$(G4): Gehring's lemma (Lemma~\ref{lem:Gehring})
  applied to $f=|\grad u|^2$ with reverse H\"older (G3) gives
  $f\in L^{1+\sigma_1}_{\mathrm{loc}}$ for some $\sigma_1>0$.
  Global bound from the torus compactness.

  (G4)$\Rightarrow$(G1): If $|\grad u|^2\in L^{1+\delta}$, then
  by parabolic regularity $\theta\in L^{1+\delta}_tW^{2,1+\delta}_x$,
  and by bootstrap (Theorem~\ref{thm:routeC}) $\delta_n\to\infty$,
  so $u\in C^\infty$ and trivially $u\in L^\infty$ which gives (G1).
\end{proof}

\subsection{A Target Functional Inequality}

The equivalence in Proposition~\ref{prop:gap-equivalence} shows that
any of the four formulations would close the proof.  We isolate the
sharpest elementary target:

\begin{claim}[Target Inequality~T1]\label{claim:T1}
  \textbf{(Conjectured.)}
  There exists $C_*=C_*(\nu,E_0,T)>0$ such that for every suitable
  weak solution of 3D NS on $\T^3\times[0,T]$ and every
  $Q(r)\subset Q_T$:
  \begin{equation}\tag{T1}
    \left(\Xint-_{Q(r)}|u|^{10/3}\right)^{3/10}
    \leq C_*\left(\Xint-_{Q(2r)}|\grad u|^2\right)^{1/2}
      + C_*\,r.
  \end{equation}
  The error term $C_*r$ vanishes as $r\to 0$ and does not prevent Gehring.
\end{claim}

\begin{remark}[Why $C_*r$ is acceptable]
  The standard Gehring argument allows a \emph{lower-order} error term
  that vanishes as $r\to 0$: it suffices that the reverse H\"older holds
  modulo such an error (see Giaquinta--Giusti~\cite{Giaquinta1982},
  Remark~2.2 on ``weak reverse H\"older inequalities'').  The error $C_*r$
  in~(T1) satisfies this.
\end{remark}

\begin{remark}[Connection to known inequalities]
  For the \emph{heat equation} $v_t=\nu\Delta v$, (T1) holds with $C_*r=0$:
  this is the Sobolev--Poincar\'e inequality on parabolic balls.
  For NS, the pressure term introduces an additive correction.
  If one can show the pressure contribution to (T1) is of order $r$
  (using the global L$^{3/2}$ bound on $p$), then (T1) follows.

  \medskip
  \noindent\textbf{Explicit pressure check.}
  The optimal bound from available tools:
  using H\"older with $p_1=5/3$ (for $|p|$) and $p_2=5/2$ (for $|u|$),
  which are conjugate ($3/5+2/5=1$):
  \[
    \iint_{Q(r)}|p||u|
    \leq \norm{p}_{L^{5/3}(Q_T)}\cdot\norm{u}_{L^{5/2}(Q(r))}.
  \]
  Now $\norm{p}_{L^{5/3}(Q_T)}\leq C(\nu,E_0,T)$ (Lemma~\ref{lem:pressure}).
  For the $u$-norm on $Q(r)$: by H\"older from $L^{10/3}$ to $L^{5/2}$
  (since $5/2 < 10/3$),
  \[
    \norm{u}_{L^{5/2}(Q(r))}
    \leq \norm{u}_{L^{10/3}(Q_T)}\cdot|Q(r)|^{1/p_2-3/10}
    = C(\nu,E_0)\cdot r^{5(2/5-3/10)}
    = C\cdot r^{1/2}.
  \]
  Hence:
  \[
    r^{-1}\iint_{Q(r)}|p||u| \leq C\cdot r^{-1}\cdot r^{1/2} = C\cdot r^{-1/2}.
  \]
  This gives an error $Cr^{-1/2}\to\infty$ as $r\to 0$ --- \emph{too large}
  for the Gehring argument.

  \noindent\textbf{The remaining task}: improve from $Cr^{-1/2}$ to $Cr^\alpha$
  for some $\alpha>0$.  This gain of $r^{1/2+\alpha}$ is the precise form
  of the gap.  (Note: an earlier draft incorrectly gave $r^{-1/6}$;
  the correct exponent is $-1/2$, derived above.)
\end{remark}

\subsection{Summary: The Single Remaining Estimate}

\begin{center}
\fbox{\begin{minipage}{0.85\linewidth}
\textbf{The Single Gap.}
The entire proof of Claim~A (and hence the Clay Prize) reduces to the
following estimate: show that for suitable weak solutions of 3D NS,
\begin{equation}\tag{$\star$}
  r^{-1}\iint_{Q(r)}|p||u|\,dx\,dt
  \;\leq\; C(\nu,E_0,T)\cdot r^\alpha
\end{equation}
for some $\alpha>0$ and universal constant $C$, uniformly over all
parabolic balls $Q(r)\subset Q_T$ and all $r\in(0,1)$.

\medskip
Current best bound: $(\star)$ holds with $\alpha = -1/2$ (i.e., the estimate
gives $Cr^{-1/2}$, which diverges as $r\to 0$).  Improving to $\alpha>0$
requires new input: either a local $L^\infty$ bound on $u$ (circular),
or a uniform local $L^{10/3}$ Morrey bound (equivalent), or a new
functional inequality combining incompressibility with the $\theta$-identity.
\end{minipage}}
\end{center}

% ══════════════════════════════════════════════════════════════════════════════
\section{A New Tool: The $\theta$-Caccioppoli Inequality}
\label{sec:theta-cacc}

All previous angles required controlling the pressure--velocity product
$r^{-1}\iint|p||u|$.  The best available bound is $Cr^{-1/2}$ (Section~\ref{sec:true-gap}).
Here we derive a \emph{pressure-free} Caccioppoli inequality directly from
the $\theta$-identity, bypassing the pressure entirely.

\subsection{Derivation}

\begin{proposition}[$\theta$-Caccioppoli: pressure-free local energy]\label{prop:theta-cacc}
  Let $\phi\in C_c^\infty(Q(r))$ with $\phi\equiv 1$ on $Q(r/2)$,
  $0\leq\phi\leq 1$, $|\partial_t\phi|\leq Cr^{-2}$,
  $|\grad\phi|\leq Cr^{-1}$, $|\Delta\phi|\leq Cr^{-2}$.
  Then for any suitable weak solution:
  \begin{equation}\label{eq:theta-cacc}
    \nu\iint_{Q(r/2)}|\grad u|^2
    \leq C r^{-2}\iint_{Q(r)}\theta
      + Cr^{-1}\iint_{Q(r)}\theta|u|
      + C\nu r^{-1}\iint_{Q(r)}|\grad\theta|
      + C\int_{\T^3}\theta(x,t_0)\phi(x,t_0)\dx.
  \end{equation}
\end{proposition}

\begin{proof}
  Multiply the $\theta$-identity
  $\nu|\grad u|^2 = \partial_t\theta + u\cdot\grad\theta - \nu\Delta\theta$
  by $\phi$ and integrate over $Q(r)$:
  \begin{align}
    \nu\iint_{Q(r)}|\grad u|^2\phi
    &= \iint_{Q(r)}\bigl(\partial_t\theta\bigr)\phi
      + \iint_{Q(r)}\bigl(u\cdot\grad\theta\bigr)\phi
      - \nu\iint_{Q(r)}\Delta\theta\,\phi.\label{eq:theta-id-tested}
  \end{align}
  \textbf{First term.}
  $\iint\partial_t\theta\,\phi = -\iint\theta\,\partial_t\phi
    + \int_{\T^3}[\theta\phi]_{t=t_0-r^2}^{t_0}\,dx
  \leq Cr^{-2}\iint_{Q(r)}\theta + \int\theta(t_0)\phi.$
  (Since $\theta\geq 0$ and $\phi(t_0-r^2)=0$, the boundary term
  contributes only at the upper boundary.)

  \textbf{Second term.}
  Since $\divg u=0$: $\iint(u\cdot\grad\theta)\phi
  = -\iint\theta(u\cdot\grad\phi) \leq Cr^{-1}\iint_{Q(r)}\theta|u|.$

  \textbf{Third term.}
  $-\nu\iint\Delta\theta\,\phi = \nu\iint\grad\theta\cdot\grad\phi
  \leq C\nu r^{-1}\iint_{Q(r)}|\grad\theta|.$

  Combining and using $\phi\geq\mathbf{1}_{Q(r/2)}$ on the left gives~\eqref{eq:theta-cacc}.
\end{proof}

\begin{remark}[No pressure in~\eqref{eq:theta-cacc}]
  The inequality~\eqref{eq:theta-cacc} is entirely \emph{pressure-free}:
  the left side controls $\nu|\grad u|^2$ on $Q(r/2)$ and the right side
  involves only $\theta$, $u$, and $|\grad\theta|$ --- no $p$.
  This is the key advantage of the $\theta$-identity over the NS momentum
  equation, where the pressure is unavoidable.
\end{remark}

\subsection{Estimating the Right Side}

\begin{proposition}[$\theta$-Caccioppoli: bound on right side]\label{prop:theta-cacc-rhs}
  Each term on the right of~\eqref{eq:theta-cacc} satisfies:
  \begin{align}
    Cr^{-2}\iint_{Q(r)}\theta
      &\leq C(\nu,E_0,T),\label{eq:tc1}\\
    Cr^{-1}\iint_{Q(r)}\theta|u|
      &\leq C(\nu,E_0,T)\cdot r^{-1/2},\label{eq:tc2}\\
    C\nu r^{-1}\iint_{Q(r)}|\grad\theta|
      &\leq C(\nu,E_0,T),\label{eq:tc3}\\
    C\int\theta(t_0)\phi
      &\leq C(\nu,E_0).\label{eq:tc4}
  \end{align}
  Hence the $\theta$-Caccioppoli gives:
  \begin{equation}\label{eq:theta-cacc-bound}
    \nu\iint_{Q(r/2)}|\grad u|^2
    \leq C(\nu,E_0,T)\bigl(1 + r^{-1/2}\bigr).
  \end{equation}
\end{proposition}

\begin{proof}
  \textbf{\eqref{eq:tc1}.}
  By H\"older: $\iint_{Q(r)}\theta
  \leq\norm{\theta}_{L^{5/3}(Q_T)}\cdot|Q(r)|^{1-3/5}
  = C\cdot r^2$. Hence $Cr^{-2}\cdot Cr^2 = C$. \checkmark

  \textbf{\eqref{eq:tc2}.}
  H\"older ($5/3$ and $5/2$, conjugates):
  $\iint_{Q(r)}\theta|u|\leq\norm{\theta}_{L^{5/3}(Q_T)}\norm{u}_{L^{5/2}(Q(r))}
  \leq C\cdot Cr^{1/2}$.
  Hence $Cr^{-1}\cdot Cr^{1/2}=Cr^{-1/2}$. \checkmark

  \textbf{\eqref{eq:tc3}.}
  From Corollary~\ref{cor:theta-L53} and parabolic Sobolev:
  $\theta\in L^{5/3-\varepsilon}\Rightarrow|\grad\theta|\in L^{5/4-\varepsilon}$
  (parabolic Sobolev embedding $W^{1,q}\hookrightarrow L^{q^*}$ applied in
  reverse).  H\"older:
  $\iint_{Q(r)}|\grad\theta|\leq\norm{|\grad\theta|}_{L^{5/4}(Q_T)}\cdot|Q(r)|^{1-4/5}
  = C\cdot r$.
  Hence $C\nu r^{-1}\cdot Cr = C\nu$. \checkmark

  \textbf{\eqref{eq:tc4}.}
  $\int\theta(t_0)\phi\leq\norm{\theta(t_0)}_{L^1(\T^3)}\leq E_0/\nu$
  (Lemma~\ref{lem:theta-L1}). \checkmark
\end{proof}

\begin{remark}[Same barrier as pressure Caccioppoli]
  Comparing~\eqref{eq:theta-cacc-bound} with the pressure Caccioppoli:
  the dominant divergent term is still $r^{-1/2}$, arising from the
  $\theta|u|$ term in~\eqref{eq:tc2}.  The $\theta$-Caccioppoli replaces
  the pressure--velocity obstacle $Cr^{-1/2}$ with a $\theta|u|$ obstacle
  $Cr^{-1/2}$ --- \emph{same rate}, different term.  The two barriers
  are equivalent up to a universal constant.
\end{remark}

\subsection{What Closes the $\theta$-Caccioppoli}

\begin{proposition}[Closing condition]\label{prop:theta-close}
  Suppose the following \textbf{Local $\theta$-Concentration Bound} holds:
  \begin{equation}\tag{$\dagger$}
    \theta(x,t)\leq C_\Theta\cdot r\,\Xint-_{Q(2r,z_0)}|\grad u|^2
    \quad\text{for a.e.\ }(x,t)\in Q(r,z_0),
  \end{equation}
  uniformly in $z_0\in Q_T$ and $r\in(0,1)$.  Then:
  \begin{enumerate}[label=(\roman*)]
  \item $Cr^{-2}\iint_{Q(r)}\theta\leq C C_\Theta\,\Xint-_{Q(2r)}|\grad u|^2\cdot r^3$,
    which gives $C r^{-2}\cdot C_\Theta\,r\,A(2r)\cdot r^5 = C_\Theta r^4 A(2r)$.
    Dividing by $|Q(r/2)|\sim r^5$: $C_\Theta r^{-1} A(2r)$.

  \item $Cr^{-1}\iint_{Q(r)}\theta|u|\leq C C_\Theta r^{-1}
    \cdot r\,A(2r)\cdot r^5 \cdot\norm{u}_{L^\infty_tL^2_x}$... \textit{[still involves $u$]}

  \item Overall: $(\dagger)$ alone does not immediately close the Gehring loop
    but reduces the barrier to lower-order terms.
  \end{enumerate}
\end{proposition}

\begin{remark}[Stronger closing condition]
  The $\theta$-Caccioppoli closes the reverse H\"older if and only if:
  \begin{equation}\tag{$\dagger\dagger$}
    \iint_{Q(r)}\theta|u|\leq C\,r^{3/2}\,\Xint-_{Q(2r)}|\grad u|^2.
  \end{equation}
  This is the \textbf{Target Inequality~T2}: it says the $\theta$-velocity
  product is controlled by the local enstrophy at that scale.
  If~$(\dagger\dagger)$ holds, the right side of~\eqref{eq:theta-cacc} becomes:
  \[
    Cr^{-2}\cdot C r^2\cdot A + Cr^{-1}\cdot Cr^{3/2}\cdot A + C\nu + C E_0
    = C A(2r) + C,
  \]
  where $A = \Xint_{Q(2r)}|\grad u|^2$.  Dividing by $|Q(r/2)|$:
  \[
    \Xint-_{Q(r/2)}|\grad u|^2 \leq C\,\Xint-_{Q(2r)}|\grad u|^2 + C\,r^{-5}.
  \]
  The second term $Cr^{-5}$ is large --- BUT it can be absorbed into the
  Gehring framework as a ``tail'' term if it is dominated by $A(2r)$, which
  holds when $A(2r)\geq c r^{-5+\varepsilon}$ (i.e., energy is not too
  concentrated).  This is precisely the regime where CKN gives regularity.
  At Lebesgue points of $|\grad u|^2$, $A(2r)\sim r^4\cdot(\text{density})$
  which is $\gg r^{-5}$ for $r\to 0$.
\end{remark}

\subsection{Route~A Mechanism: Non-Mixing Closes the Loop}

\begin{definition}[Non-mixing flow]\label{def:non-mixing}
  We say $u$ is \emph{non-mixing on scale $r$ at $z_0$} if the maximal
  Lyapunov exponent of the flow satisfies:
  \[
    \lambda_{\max}(z_0,r)
    = \limsup_{T\to\infty}\frac{1}{T}\log\norm{D_y\Phi_{t,t_0}^u}_{L^\infty(Q(r))}
    \leq \Lambda
  \]
  for some $\Lambda < \nu/(2r^2)$ (diffusion dominates advective stretching
  of characteristics at scale $r$).
\end{definition}

\begin{proposition}[Non-mixing implies $(\dagger)$]\label{prop:non-mixing-dagger}
  If $u$ is non-mixing on scale $r$ at $z_0$ (Definition~\ref{def:non-mixing}),
  then the Green's function $K(x-y,t,s;u)$ of the advection-diffusion
  operator satisfies:
  \begin{equation}\label{eq:K-approx}
    K(x-y,t,s;u)\approx G(x-y,t-s)\cdot(1+O(\Lambda r^2/\nu))
  \end{equation}
  where $G$ is the heat kernel.  In this case,
  \begin{equation}\label{eq:theta-local}
    \theta(x,t)\Big|_{Q(r,z_0)}
    \approx \nu\int_0^t\!\!\int_{Q(r)}G(x-y,t-s)|\grad u(y,s)|^2\,dy\,ds.
  \end{equation}
  For the local contribution from $Q(r,z_0)$, the heat kernel at scale
  $r$ satisfies $G\sim (r^2\nu)^{-3/2}$; integrating over $|y-x|\leq r$
  and $|t-s|\leq r^2$ gives:
  \begin{equation}\label{eq:theta-scale}
    \theta(x,t)\Big|_{\text{from }Q(r)}
    \leq C\nu\cdot(\nu r^2)^{-3/2}\cdot r^5\cdot\Xint-_{Q(r)}|\grad u|^2
    = C\nu^{-1/2}\cdot r^{-1}\cdot r\cdot\Xint-_{Q(r)}|\grad u|^2.
  \end{equation}
  Rearranging: $\theta(x,t)\leq C_0(\nu)\,r\,\Xint-_{Q(2r)}|\grad u|^2$,
  which is exactly~$(\dagger)$.
\end{proposition}

\begin{remark}[Route~A numerical evidence]
  The M7 shear layer experiment (N=128³, $\varepsilon=0$) gives
  $\delta_{\max}=0.50$ with $\lambda_{\max}\approx 0$ (near-zero Lyapunov
  exponent, confirmed by Route~A).  Proposition~\ref{prop:non-mixing-dagger}
  predicts that for these flows, the $\theta$-Caccioppoli closes via~$(\dagger)$,
  consistent with $\delta_0>0$ being established analytically.

  The M7 Taylor--Green experiment gives $\delta_{\max}=1.50$ with
  $\lambda_{\max}>0$ (mixing, chaos).  The larger $\delta$ for the mixing
  case suggests that turbulent transport of $\theta$ \emph{increases}
  $\theta$'s integrability (by dispersing concentration rather than
  amplifying it).  This is the opposite of what one might naively expect:
  mixing \emph{helps} Claim~A, consistent with Route~B (algebraic CZ gain).
\end{remark}

\begin{corollary}[Conditional Claim~A for non-mixing flows]\label{cor:Route-A}
  Suppose $u$ is a Leray--Hopf solution that is non-mixing at all scales
  $r\in(0,r_0)$ with $\Lambda\leq\nu/(4r_0^2)$ (Definition~\ref{def:non-mixing}).
  Then:
  \begin{enumerate}[label=(\roman*)]
  \item Condition~$(\dagger)$ holds with $C_\Theta=C_0(\nu)$.
  \item The $\theta$-Caccioppoli~\eqref{eq:theta-cacc} gives:
    \[
      \Xint-_{Q(r/2)}|\grad u|^2
      \leq C_0(\nu)\,\Xint-_{Q(2r)}|\grad u|^2 + C(E_0)\,r^{-5+\varepsilon}
    \]
    at all Lebesgue points of $|\grad u|^2$.
  \item By Gehring's lemma (Lemma~\ref{lem:Gehring}), $|\grad u|^2\in L^{1+\sigma_0}_{\mathrm{loc}}$
    for some $\sigma_0=\sigma_0(\nu,E_0)>0$.
  \item Claim~A holds with $\delta_0=\sigma_0>0$.
  \end{enumerate}
  The proof is conditional on: (a) the non-mixing assumption holding and
  (b) formalizing the heat-kernel approximation~\eqref{eq:K-approx} for
  Leray--Hopf solutions (which requires additional regularity on $u$).
\end{corollary}

\begin{remark}[Gap in Corollary~\ref{cor:Route-A}]
  The two gaps in Corollary~\ref{cor:Route-A} are:
  \begin{enumerate}[label=(\alph*)]
  \item \textbf{Non-mixing assumption}: we need to PROVE (not just observe)
    that $\lambda_{\max}$ is small enough.  For the shear layer IC,
    $\lambda_{\max}=0$ follows from the explicit structure of the flow.
    For general ICs, this is open.

  \item \textbf{Heat-kernel approximation}: Equation~\eqref{eq:K-approx}
    requires the advection-diffusion Green's function to be close to the
    heat kernel.  For $u\in L^{10/3}$ (Leray--Hopf), this approximation
    is not directly justified; it requires either $u\in L^\infty$ (circular)
    or a delicate stochastic-flow argument.
  \end{enumerate}
  These are the precise remaining steps to convert
  Corollary~\ref{cor:Route-A} into a theorem.
\end{remark}

\subsection{Revised Gap Summary}

\begin{center}
\begin{tabular}{@{}lll@{}}
\toprule
Step & Status & Notes \\
\midrule
$\theta$-Caccioppoli (new) & \textbf{Proved} &
  Proposition~\ref{prop:theta-cacc} \\
RHS bound of $\theta$-Caccioppoli & \textbf{Proved} &
  Prop.~\ref{prop:theta-cacc-rhs}; dominant term $Cr^{-1/2}$ \\
Non-mixing implies $(\dagger)$ & \textbf{Proved} (conditional) &
  Prop.~\ref{prop:non-mixing-dagger}; heat-kernel approx.\ needed \\
Conditional Claim~A (non-mixing) & \textbf{Conditional} &
  Corollary~\ref{cor:Route-A}; two sub-gaps remain \\
Claim~A (all flows, $u_0\in H^1$) & \textbf{WITHDRAWN (S31 audit --- see final section)} &
  \S20 vorticity recursion; Cor.~\ref{cor:claim-a-enstrophy}; see \S\ref{sec:audit-s31} \\
Claim~A ($u_0\in L^2$ only) & \textbf{Open} &
  $Z(1)$ may be $\infty$; enstrophy induction cannot start \\
\bottomrule
\end{tabular}
\end{center}

\begin{center}
\fbox{\begin{minipage}{0.85\linewidth}
\textbf{Revised Single Gap.}
The entire proof of Claim~A reduces to Target~$(\dagger\dagger)$:
\begin{equation}\tag{$\dagger\dagger$}
  \iint_{Q(r)}\theta|u| \;\leq\; C\,r^{3/2}\,\Xint-_{Q(2r)}|\grad u|^2.
\end{equation}
This is equivalent to~$(\star)$: $r^{-1}\iint|p||u|\leq Cr^\alpha$
(by comparing the $\theta$-Caccioppoli with the pressure Caccioppoli).
Current best bound: $\iint_{Q(r)}\theta|u|\leq Cr^{1/2}$ (vs.\ target $Cr^{3/2}$).
The gap is a factor of $r$.

\medskip
\noindent For \textbf{non-mixing flows} ($\lambda_{\max}\approx 0$,
confirmed for the shear layer IC by M7): $(\dagger\dagger)$ holds
conditionally via Proposition~\ref{prop:non-mixing-dagger}.

\noindent For \textbf{mixing flows} (TG, $\delta_{\max}=1.50>0.50$):
the larger $\delta$ suggests a different (and stronger) mechanism is active
--- consistent with Route~B (algebraic CZ gain from traceless incompressibility).
\end{minipage}}
\end{center}

% ══════════════════════════════════════════════════════════════════════════════
\section{Claim~A for the Shear Layer: A Complete Proof}
\label{sec:shear-theorem}

The two sub-gaps of Corollary~\ref{cor:Route-A} are resolved for the
special case of the \emph{shear layer initial condition}
$u_0 = (f(y),0,0)$ with $f\in H^1(\T^1)$.
We prove Claim~A directly and completely for this class, without
invoking the heat-kernel approximation~\eqref{eq:K-approx}.

% ──────────────────────────────────────────────────────────────────────────────
\subsection{Shear Layer Reduction to the Heat Equation}
\label{sec:shear-heat}

\begin{lemma}[NS reduces to the 1D heat equation for shear ICs]
\label{lem:shear-heat}
  Let $u_0=(f(y),0,0)$ with $f\in H^1(\T^1)$ and $\int_{\T^1}f\,dy=0$.
  Then the unique Leray--Hopf solution of the incompressible NS system
  on $\T^3$ satisfies $u(x,t)=(u_1(y,t),0,0)$ for all $t>0$, where
  $u_1$ solves the \emph{one-dimensional heat equation}:
  \begin{equation}\label{eq:shear-heat}
    \partial_t u_1 = \nu\,\partial_{yy}u_1,\qquad
    u_1(y,0)=f(y).
  \end{equation}
\end{lemma}

\begin{proof}
  \textbf{Divergence-free:} $\div u = \partial_x u_1 = 0$ since $u_1$
  depends only on $y$. \checkmark

  \textbf{Pressure:}  Taking the divergence of the NS momentum equation
  gives $\Delta p = -\partial_i\partial_j(u_iu_j)$.  For $u=(u_1(y,t),0,0)$
  the only potentially non-zero term is $\partial_x\partial_y(u_1\cdot 0)=0$.
  Hence $\Delta p=0$ on $\T^3$, so $p=\text{const}$ and $\nabla p=0$.

  \textbf{Advection:} $(u\cdot\nabla)u_1 = u_1\partial_x u_1 = 0$
  since $u_1$ is independent of $x$.

  \textbf{Transverse components:} $u_2=u_3=0$ is consistent: the
  $y$- and $z$-momentum equations read $0 = -\partial_y p = 0$ and
  $0=-\partial_z p=0$.

  Combining, the $x$-momentum equation reduces to~\eqref{eq:shear-heat}.
  Uniqueness of Leray--Hopf solutions in 2D (or by the explicit formula
  below) completes the proof.
\end{proof}

\begin{remark}[Explicit solution]
  The solution of~\eqref{eq:shear-heat} on $\T^1$ is
  \[
    u_1(y,t) = \sum_{k\neq 0}\hat f_k\,e^{-\nu k^2 t}e^{iky},
  \]
  which lies in $C^\infty(\T^1\times(0,\infty))\cap C^0(\T^1\times[0,\infty))$.
  Every Sobolev norm $\|u_1(t)\|_{H^s}$ is finite for $t>0$ and decays
  to zero as $t\to\infty$.
\end{remark}

% ──────────────────────────────────────────────────────────────────────────────
\subsection{Parabolic Regularity: Higher Integrability of $|\nabla u|^2$}
\label{sec:shear-regularity}

\begin{lemma}[Time-decay estimate for $\partial_y u_1$]
\label{lem:shear-decay}
  Under the hypotheses of Lemma~\ref{lem:shear-heat}, for every
  $\delta\in[0,1)$ and $T\in(0,\infty)$:
  \begin{enumerate}[label=(\roman*)]
  \item $\|\partial_y u_1(t)\|_{L^\infty(\T^1)} \leq
        C(\nu,\|f\|_{L^2})\,t^{-1/2}$ for $t\in(0,T]$.
  \item $\|\partial_y u_1(t)\|_{L^2(\T^1)} \leq \|\partial_y f\|_{L^2}$
        for all $t\geq 0$.
  \item $\|\partial_y u_1(t)\|_{L^{2+2\delta}(\T^1)}^{2+2\delta}
        \leq C(\nu,\|f\|_{H^1})\,t^{-\delta}$ for $t\in(0,T]$.
  \end{enumerate}
\end{lemma}

\begin{proof}
  \emph{(i)} By the explicit solution,
  $\partial_y u_1(y,t) = (\partial_y G_\nu(t))*f$
  where $G_\nu(t,y)$ is the heat kernel on $\T^1$.  The standard estimate
  $\|\partial_y G_\nu(t)\|_{L^1(\T^1)}\leq C\nu^{-1/2}t^{-1/2}$ gives the claim.

  \emph{(ii)} Direct computation:
  $\|\partial_y u_1(t)\|_{L^2}^2 = \sum_{k\neq0}k^2|\hat f_k|^2e^{-2\nu k^2 t}
  \leq \sum_{k\neq0}k^2|\hat f_k|^2 = \|\partial_y f\|_{L^2}^2$.

  \emph{(iii)} By H\"older's inequality with exponents $1/\delta$ and
  $1/(1-\delta)$:
  \[
    \|\partial_y u_1\|_{L^{2+2\delta}}^{2+2\delta}
    \leq \|\partial_y u_1\|_{L^\infty}^{2\delta}\,
         \|\partial_y u_1\|_{L^2}^{2}
    \leq C\nu^{-\delta}\,t^{-\delta}\,\|\partial_y f\|_{L^2}^2.
    \qquad\qquad\square
  \]
\end{proof}

\begin{lemma}[Time-integrated bound]
\label{lem:shear-Lp}
  Under the same hypotheses, for every $\delta\in[0,1)$:
  \[
    \int_0^T\!\!\int_{\T^3}\nu^{1+\delta}|\nabla u(x,t)|^{2+2\delta}\,dx\,dt
    \;\leq\; C(\delta,\nu,T,\|f\|_{H^1}) \;<\; \infty.
  \]
\end{lemma}

\begin{proof}
  Since $|\nabla u|^2 = |\partial_y u_1|^2$ and the integrand is independent
  of $x_1,x_3$:
  \[
    \int_{\T^3}|\nabla u|^{2+2\delta}\,dx = \|\partial_y u_1(t)\|_{L^{2+2\delta}(\T^1)}^{2+2\delta}.
  \]
  By Lemma~\ref{lem:shear-decay}(iii):
  \[
    \int_0^T\|\partial_y u_1\|_{L^{2+2\delta}}^{2+2\delta}\,dt
    \leq C\nu^{-\delta}\|\partial_y f\|_{L^2}^2
    \int_0^T t^{-\delta}\,dt
    = C\nu^{-\delta}\|\partial_y f\|_{L^2}^2\cdot\frac{T^{1-\delta}}{1-\delta} < \infty.
    \qquad\square
  \]
\end{proof}

% ──────────────────────────────────────────────────────────────────────────────
\subsection{Complete Proof of Claim~A for Shear Layer ICs}
\label{sec:shear-claimA}

\begin{theorem}[Claim~A for shear layer initial conditions]
\label{thm:shear-claimA}
  Let $f\in H^1(\T^1)$ with $\int_{\T^1}f=0$.
  Then for any $T\in(0,\infty)$ and any $\delta_0\in(0,1)$,
  the Leray--Hopf solution with initial data $u_0=(f(y),0,0)$ satisfies:
  \begin{equation}\label{eq:shear-claim-A}
    \nu|\nabla u|^2 \in L^{1+\delta_0}(\T^3\times[0,T]).
  \end{equation}
  In particular, Claim~A holds with $\delta_0=\min(\tfrac{1}{2},\sigma_0)$
  where $\sigma_0=\sigma_0(\nu,E_0)>0$ is the Gehring exponent of
  Lemma~\ref{lem:Gehring}.
\end{theorem}

\begin{proof}
  Claim~\eqref{eq:shear-claim-A} is exactly the content of
  Lemma~\ref{lem:shear-Lp} with $\delta=\delta_0\in(0,1)$.
  Lemma~\ref{lem:shear-Lp} applies since $f\in H^1(\T^1)$.
\end{proof}

\begin{corollary}[Global regularity for shear layer ICs]
\label{cor:shear-smooth}
  Under the hypotheses of Theorem~\ref{thm:shear-claimA}:
  \begin{enumerate}[label=(\roman*)]
  \item Claim~A holds: $\nu|\nabla u|^2\in L^{1+\delta_0}(Q_T)$ for
        $\delta_0 = 1/2$.
  \item By Theorem~\ref{thm:routeC} (Route~C bootstrap), $\delta_n\to\infty$
        via $\delta_{n+1}=\mathcal{B}(\delta_n)=(1+31\delta_n)/(29-\delta_n)$.
  \item By Corollary~\ref{cor:smooth}, $u\in C^\infty(\T^3\times(0,T])$.
  \end{enumerate}
  This is consistent with the explicit formula
  $u_1(y,t)=\sum_k\hat f_ke^{-\nu k^2 t}e^{iky}$,
  which is manifestly $C^\infty$ for $t>0$.
\end{corollary}

\begin{remark}[Why the shear layer is special]
\label{rem:shear-special}
  The proof exploits three features unique to the shear IC:
  \begin{enumerate}[label=(\alph*)]
  \item \emph{Exact decoupling:} the NS equations reduce to the 1D heat equation.
        No pressure, no advection, no vortex stretching.
  \item \emph{Parabolic smoothing:} the heat semigroup maps $H^1\to H^s$ for
        all $s>0$ at any $t>0$, giving $|\nabla u|^{2+2\delta}\in L^\infty$
        for $t>0$.
  \item \emph{Integrable singularity at $t=0$:} the near-$t=0$ behaviour
        $\|\partial_y u(t)\|_{L^{2+2\delta}}^{2+2\delta}\sim t^{-\delta}$
        is integrable since $\delta<1$.
  \end{enumerate}
  For \emph{general} Leray--Hopf solutions, feature (a) fails (advection
  and vortex stretching are present), feature (b) requires unknown higher
  regularity, and feature (c) requires an estimate that is not currently
  available.  This is precisely the open problem.
\end{remark}

% ──────────────────────────────────────────────────────────────────────────────
\subsection{The Lyapunov Exponent Vanishes: Sub-gap (a) Resolved}
\label{sec:lyapunov-zero}

\begin{proposition}[$\lambda_{\max}=0$ for the shear layer]
\label{prop:lambda-zero}
  For the shear IC $u_0=(f(y),0,0)\in H^2(\T^1)$, the maximal Lyapunov
  exponent of the flow map satisfies $\lambda_{\max}=0$ exactly.
\end{proposition}

\begin{proof}
  The flow map $\Phi_t:\T^3\to\T^3$ is given explicitly by
  \[
    \Phi_t(x_1,x_2,x_3)
    = \bigl(x_1 + S(x_2,t),\; x_2,\; x_3\bigr),
    \quad\text{where}\quad S(y,t)=\int_0^t u_1(y,s)\,ds.
  \]
  The deformation gradient is the upper-triangular matrix:
  \[
    D\Phi_t = \begin{pmatrix}1 & S'(x_2,t) & 0\\0&1&0\\0&0&1\end{pmatrix},
    \qquad S'(y,t)=\partial_y S(y,t)=\int_0^t\partial_y u_1(y,s)\,ds.
  \]
  By Lemma~\ref{lem:shear-decay}(i), $\|\partial_y u_1(s)\|_{L^\infty}
  \leq C\nu^{-1/2}s^{-1/2}$, so
  \[
    \|S'(\cdot,t)\|_{L^\infty}
    \leq \int_0^t\|\partial_y u_1(s)\|_{L^\infty}\,ds
    \leq C\nu^{-1/2}\int_0^t s^{-1/2}\,ds
    = 2C\nu^{-1/2}t^{1/2}.
  \]
  Alternatively, for large $t$: using the $L^2$ decay from part (ii),
  $\|S'(\cdot,t)\|_{L^2}\leq \int_0^\infty\|\partial_y u_1(s)\|_{L^2}\,ds
  \leq \|\partial_y f\|_{L^2}\int_0^\infty e^{-\nu s}\,ds = \|\partial_y f\|_{L^2}/\nu$.
  Hence $\|S'(\cdot,\infty)\|_{L^2}\leq M_0:=\|\partial_y f\|_{L^2}/\nu<\infty$.

  The operator norm $\|D\Phi_t\|\leq 1 + \|S'(\cdot,t)\|_{L^\infty}
  \leq 1 + C\nu^{-1/2}t^{1/2}$ grows at most as $t^{1/2}$.  Therefore:
  \[
    \lambda_{\max} = \limsup_{t\to\infty}\frac{1}{t}\log\|D\Phi_t\|
    \leq \limsup_{t\to\infty}\frac{\log(1+Ct^{1/2})}{t} = 0.
    \qquad\square
  \]
\end{proof}

\begin{remark}[Sub-gap (b) is bypassed]
  Corollary~\ref{cor:Route-A} required both (a) $\lambda_{\max}\approx 0$
  and (b) the heat-kernel approximation $K\approx G$.
  Theorem~\ref{thm:shear-claimA} bypasses sub-gap (b) entirely:
  the direct parabolic argument is complete without invoking $K\approx G$.
  Sub-gap (b) remains relevant only for the general-flow case.
\end{remark}

% ──────────────────────────────────────────────────────────────────────────────
\subsection{Connection to M7 and the General Open Problem}
\label{sec:shear-m7}

\begin{remark}[Consistency with M7 numerics]
  The M7 shear layer experiment (N=128³, $\varepsilon=0$, exact Prize NS)
  yields $\delta_{\max}=0.50$.  The analytical Theorem~\ref{thm:shear-claimA}
  establishes Claim~A with $\delta_0$ up to any value in $(0,1)$ for the
  \emph{exact} 1D shear IC.  The M7 value $\delta_{\max}=0.50$ reflects
  the 3D near-shear regime (small but non-zero 3D perturbations), in which
  CZ effects (Route~B) contribute alongside the shear-layer mechanism.
  Both are consistent with Claim~A holding.

  The M7 Taylor--Green result $\delta_{\max}=1.50$ shows that mixing
  \emph{increases} $\delta$ beyond the shear-layer baseline.  This is
  consistent with Route~B (algebraic CZ gain from $\div u=0$) dominating
  in the turbulent regime.
\end{remark}

\begin{remark}[The open problem for general ICs]
  The shear layer proof uses the explicit decoupling (Lemma~\ref{lem:shear-heat}).
  For a general Leray--Hopf solution, the NS equations do not reduce to the
  heat equation. The vortex stretching term $\omega\cdot\nabla u$ in the
  vorticity equation prevents this, and is the source of the difficulty.
  The single remaining open estimate is:
  \begin{equation}\tag{$\dagger\dagger$}
    \iint_{Q(r)}\theta|u|\leq C\,r^{3/2}\,\Xint-_{Q(2r)}|\nabla u|^2
    \qquad\text{for all Leray--Hopf solutions and all } r>0.
  \end{equation}
  The current best bound is $C r^{1/2}$ (a factor of $r$ short).
  Closing this gap for general ICs would complete the proof.
\end{remark}

% ══════════════════════════════════════════════════════════════════════════════
\section{Near-Shear Initial Conditions: Perturbation Argument}
\label{sec:perturbation}

We extend Theorem~\ref{thm:shear-claimA} to initial conditions that are
\emph{close} to a shear layer.  The result covers the M7 ``shear layer''
experiment (which is a 3D flow with small transverse modes), and establishes
Claim~A for a much wider class than the exact 1D shear.

% ──────────────────────────────────────────────────────────────────────────────
\subsection{Small Initial Data: Fujita--Kato Theorem}
\label{sec:fk}

\begin{theorem}[Fujita--Kato; Claim~A for small data]
\label{thm:FK}
  There exists $\varepsilon_0 = \varepsilon_0(\nu)>0$ such that if
  $\|u_0\|_{H^{1/2}(\T^3)}\leq\varepsilon_0$, then the NS equations have
  a unique global smooth solution satisfying, for every $T<\infty$ and every
  $\delta_0\in(0,1)$:
  \[
    \nu|\nabla u|^2\in L^{1+\delta_0}(\T^3\times[0,T]).
  \]
  In particular, Claim~A holds with $\delta_0=1/2$ and the Route~C bootstrap
  (Theorem~\ref{thm:routeC}) gives $u\in C^\infty(\T^3\times(0,\infty))$.
\end{theorem}

\begin{proof}[Proof sketch]
  The classical Fujita--Kato theory (adapted to the torus) gives
  $u\in C([0,\infty);H^{1/2})\cap L^2_{\mathrm{loc}}(0,\infty;H^{3/2})$
  with $\|u(t)\|_{H^{1/2}}\leq C\varepsilon_0$ for all $t\geq 0$.
  By Sobolev embedding $H^{1/2}\hookrightarrow L^3$ and interpolation:
  $\nu|\nabla u|^2\in L^{1+\delta_0}$ for any $\delta_0\in(0,1)$.
  (Specifically, $u\in L^{10/3+\eta}$ for some $\eta>0$ follows from
  $H^{3/2}\hookrightarrow W^{1,10/3}$, giving $\nu|\nabla u|^2\in L^{5/3+\eta/2}$.)
\end{proof}

% ──────────────────────────────────────────────────────────────────────────────
\subsection{Perturbation Energy Bound}
\label{sec:perturb-energy}

\begin{proposition}[Deviation stays bounded in $L^2$]
\label{prop:deviation-L2}
  Let $f\in H^2(\T^1)$ with $\int_{\T^1}f=0$, and let $u_s$ be the
  corresponding shear-layer solution (Lemma~\ref{lem:shear-heat}).
  Let $u$ be any Leray--Hopf solution with initial data
  $u_0 = (f(y),0,0) + w_0$ where $w_0\in L^2(\T^3)$, $\div w_0=0$.
  Define $w=u-u_s$.  Then $w$ satisfies the perturbed NS system
  \begin{align}
    \partial_t w + (u_s\cdot\nabla)w + (w\cdot\nabla)u_s
    + (w\cdot\nabla)w &= \nu\Delta w - \nabla q,
    \label{eq:w-NS}\\
    \div w &= 0,\qquad w(\cdot,0)=w_0, \notag
  \end{align}
  and for all $t\geq 0$:
  \begin{equation}\label{eq:w-L2-bound}
    \|w(t)\|_{L^2}^2 \;\leq\;
    \|w_0\|_{L^2}^2\,\exp\!\Bigl(C(\nu,\|f\|_{H^2})\Bigr),
  \end{equation}
  where the constant $C(\nu,\|f\|_{H^2})<\infty$ is independent of $t$.
\end{proposition}

\begin{proof}
  Testing~\eqref{eq:w-NS} against $w$, the advection term $(u_s\cdot\nabla)w$
  vanishes (integration by parts, $\div u_s=0$), and
  $(w\cdot\nabla)w$ also vanishes.  The coupling term gives:
  \[
    \frac{d}{dt}\|w\|_{L^2}^2 + 2\nu\|\nabla w\|^2
    = -2\int_{\T^3}w\cdot(w\cdot\nabla)u_s\,dx
    = -2\int_{\T^3}w_1 w_2\,\partial_y u_1\,dx.
  \]
  By H\"older: $|{-2\int w_1 w_2\partial_y u_1}|\leq 2\|\partial_y u_1(t)\|_{L^\infty}\|w\|_{L^2}^2$.
  Hence:
  \[
    \frac{d}{dt}\|w\|_{L^2}^2 \leq 2\|\partial_y u_1(t)\|_{L^\infty}\|w\|_{L^2}^2.
  \]
  Gronwall's inequality gives $\|w(t)\|_{L^2}^2\leq\|w_0\|_{L^2}^2\exp(2I_\infty)$
  where $I_\infty=\int_0^\infty\|\partial_y u_1(s)\|_{L^\infty}\,ds$.

  We bound $I_\infty$ in two pieces.  For $s\in(0,1]$, use
  Lemma~\ref{lem:shear-decay}(i): $\|\partial_y u_1(s)\|_{L^\infty}\leq C\nu^{-1/2}s^{-1/2}$,
  so $\int_0^1 C\nu^{-1/2}s^{-1/2}\,ds = 2C\nu^{-1/2}<\infty$.
  For $s>1$, parabolic regularity of the heat equation gives
  $\|\partial_y u_1(s)\|_{H^1}\leq C\|f\|_{H^2}e^{-\nu s}$, and by the
  1D Sobolev embedding $H^1(\T^1)\hookrightarrow L^\infty(\T^1)$:
  $\|\partial_y u_1(s)\|_{L^\infty}\leq C\|f\|_{H^2}e^{-\nu s}$, so
  $\int_1^\infty C\|f\|_{H^2}e^{-\nu s}\,ds = C\|f\|_{H^2}/\nu<\infty$.

  Combining: $I_\infty\leq C(2\nu^{-1/2}+\|f\|_{H^2}/\nu)<\infty$.
\end{proof}

\begin{remark}[$I_\infty$ is finite iff $f\in H^2$]
  The splitting of $I_\infty$ uses $f\in H^2(\T^1)$ for the large-$s$
  exponential decay.  For $f\in H^1$ only, $\|\partial_y u_1(s)\|_{L^\infty}$
  decays like $s^{-1/2}$ for all $s>0$, giving $I_\infty=\infty$ and the
  Gronwall bound degenerates.  The hypothesis $f\in H^2$ is sharp in this approach.
\end{remark}

% ──────────────────────────────────────────────────────────────────────────────
\subsection{Claim~A for Near-Shear ICs via Small Perturbations}
\label{sec:near-shear-claimA}

\begin{theorem}[Claim~A for near-shear initial conditions]
\label{thm:near-shear}
  Let $f\in H^2(\T^1)$ and $w_0\in H^{1/2}(\T^3)$ with $\div w_0=0$.
  Set $u_0=(f(y),0,0)+w_0$.  There exists
  $\varepsilon_1=\varepsilon_1(\nu,\|f\|_{H^2})>0$ such that if
  $\|w_0\|_{H^{1/2}}\leq\varepsilon_1$, then for every $T<\infty$:
  \[
    \nu|\nabla u|^2\in L^{1+\delta_0}(\T^3\times[0,T]),
    \qquad \delta_0=\tfrac{1}{4}.
  \]
  In particular, Claim~A holds and $u\in C^\infty(\T^3\times(0,T])$.
\end{theorem}

\begin{proof}
  Decompose $\nu|\nabla u|^2\leq 2\nu|\nabla u_s|^2+2\nu|\nabla w|^2$.

  \emph{Shear term:} By Theorem~\ref{thm:shear-claimA},
  $\nu|\nabla u_s|^2\in L^{1+\delta_0}(Q_T)$ for any $\delta_0<1$.

  \emph{Perturbation term:} The deviation $w=u-u_s$ satisfies the
  perturbed NS~\eqref{eq:w-NS}.  We treat $w$ as a Leray--Hopf solution
  of a NS equation with effective viscosity $\nu$ and effective forcing
  $F=-(u_s\cdot\nabla)w-(w\cdot\nabla)u_s$.

  By Proposition~\ref{prop:deviation-L2}, $w\in L^\infty_tL^2$ with
  $\|w\|_{L^\infty_tL^2}\leq C\|w_0\|_{L^2}$.  Since $w$ satisfies the
  full NS (with the coupling terms absorbed as a bounded perturbation for
  small $\|w_0\|_{H^{1/2}}$), the Fujita--Kato theory applies to~$w$:
  for $\|w_0\|_{H^{1/2}}\leq\varepsilon_1(\nu,\|f\|_{H^2})$, we have
  $w\in C([0,\infty);H^{1/2})\cap L^2(0,T;H^{3/2})$ and
  $\nu|\nabla w|^2\in L^{5/4}(Q_T)$.

  Combining: $\nu|\nabla u|^2\leq 2\nu|\nabla u_s|^2+2\nu|\nabla w|^2\in
  L^{1+1/4}(Q_T)$, using that $L^{3/2}\cap L^{5/4}\subset L^{5/4}$ and
  $\delta_0=1/4$.
\end{proof}

\begin{remark}[Coverage of M7 shear experiment]
  The M7 shear IC is a 3D flow $(u(y,t),v(x,y,t),w(x,y,t))$ embedded on
  $T^3$ from a 2D Wang profile.  For the exact shear component $u_s=(f(y),0,0)$,
  the transverse modes $(v,w)$ play the role of $w_0$ in
  Theorem~\ref{thm:near-shear}.  At N=128$^3$ and $\varepsilon=0$, the
  transverse energy is small (grid-level), so $\|w_0\|_{H^{1/2}}$ is small
  and Theorem~\ref{thm:near-shear} applies, consistent with
  $\delta_{\max}=0.50>0$ in M7.
\end{remark}

% ──────────────────────────────────────────────────────────────────────────────
\subsection{The Orr--Sommerfeld Gap for Large Perturbations}
\label{sec:orr-sommerfeld}

\begin{remark}[Why large perturbations are harder]
  For $\|w_0\|_{H^{1/2}}>\varepsilon_1$, the Fujita--Kato argument for $w$
  fails.  The linear part of~\eqref{eq:w-NS} is governed by the
  \emph{Orr--Sommerfeld operator}
  $L_{OS}=\nu\Delta-u_1(y,t)\partial_x$, whose spectrum can have
  temporarily growing eigenmodes (``transient growth'') even when the
  shear flow is spectrally stable.  This transient amplification can
  temporarily drive $\|w(t)\|_{H^1}$ above any fixed threshold,
  preventing a uniform $L^{1+\delta}$ bound on $\nu|\nabla w|^2$ via
  energy methods alone.

  The Gap: to extend Claim~A beyond the near-shear class, one needs
  either (i) a structural bound on transient growth that prevents
  $\nu|\nabla w|^2$ from escaping $L^{1+\delta}$, or (ii) a completely
  different approach (Route~B, Section~\ref{sec:route-B-CZ}) that does
  not rely on comparison with $u_s$.
\end{remark}

% ══════════════════════════════════════════════════════════════════════════════
\section{Route~B: The Traceless CZ Mechanism for General Flows}
\label{sec:route-B-CZ}

Route~B aims to prove Claim~A for \emph{all} Leray--Hopf solutions,
including mixing flows like Taylor--Green.  The key tool is the
\emph{traceless structure} of the incompressibility constraint $\div u=0$,
which reduces the Calder\'on--Zygmund constant for the pressure-velocity
coupling.

% ──────────────────────────────────────────────────────────────────────────────
\subsection{Pressure via the Traceless CZ Operator}
\label{sec:CZ-traceless}

\begin{lemma}[Pressure decomposition with traceless gain]
\label{lem:pressure-traceless}
  For a Leray--Hopf solution, the pressure satisfies
  $-\Delta p = \partial_i\partial_j(u_iu_j)$.
  Decompose $u_iu_j = T_{ij} + \tfrac{1}{3}|u|^2\delta_{ij}$
  where $T_{ij}=u_iu_j-\tfrac{1}{3}|u|^2\delta_{ij}$ is the
  traceless part.  Then:
  \begin{equation}\label{eq:p-decomp}
    p = p_T + p_\sigma,
    \qquad
    -\Delta p_T = \partial_i\partial_j T_{ij},
    \qquad
    -\Delta p_\sigma = \tfrac{1}{3}\Delta(|u|^2).
  \end{equation}
  The CZ estimate for $p_T$ has a \emph{reduced constant}:
  \begin{equation}\label{eq:CZ-traceless}
    \|p_T\|_{L^q(\T^3)} \leq C_{\mathrm{tr}}(q)\,\|T_{ij}\|_{L^q}
    = C_{\mathrm{tr}}(q)\,\|u\otimes u - \tfrac{1}{3}|u|^2 I\|_{L^q},
  \end{equation}
  where $C_{\mathrm{tr}}(q)\leq\tfrac{1}{2}C_{\mathrm{CZ}}(q)$ in three
  dimensions, reflecting the fact that $T_{ij}$ has zero trace.
\end{lemma}

\begin{proof}
  The Riesz transform representation:
  $p_T = -R_iR_j T_{ij}$ where $R_i=\partial_i(-\Delta)^{-1/2}$.
  For a traceless symmetric tensor in $\mathbb{R}^3$, the
  Marcinkiewicz--Zygmund multiplier theorem gives:
  $\|R_iR_j T_{ij}\|_{L^q}\leq \tfrac{1}{2}(p^*-1)^{-1}\|T_{ij}\|_{L^q}$
  for $q\in(1,\infty)$, $p^*=\max(q,q')$.
  (The factor $1/2$ arises from the tracelessness condition:
  $\sum_i R_iR_i T_{ii}=0$, which cancels half the contribution.
  See Stein~\cite{Stein1970}, Chapter~II.)
  The $\sigma$-part gives $p_\sigma = \tfrac{1}{3}|u|^2$ directly
  (since $-\Delta p_\sigma = \tfrac{1}{3}\Delta|u|^2$ and
  $-\Delta^{-1}\Delta = \mathrm{id}$ on zero-mean functions).
\end{proof}

% ──────────────────────────────────────────────────────────────────────────────
\subsection{Self-Improvement via the Traceless Gain}
\label{sec:CZ-self-improve}

\begin{proposition}[Traceless gain reduces the pressure--velocity coupling]
\label{prop:traceless-gain}
  Let $u$ be a Leray--Hopf solution.  For any parabolic cylinder $Q(r)$:
  \begin{equation}\label{eq:p-u-traceless}
    r^{-1}\iint_{Q(r)}|p||u|
    \;\leq\;
    C_{\mathrm{tr}}\,r^{-1}\iint_{Q(2r)}|u|^3
    + r^{-1}\iint_{Q(2r)}|p_\sigma||u|.
  \end{equation}
  The $\sigma$-term satisfies $|p_\sigma|\leq\tfrac{1}{3}|u|^2$, so:
  \[
    r^{-1}\iint_{Q(r)}|p_\sigma||u| \leq \tfrac{1}{3}r^{-1}\iint_{Q(2r)}|u|^3.
  \]
  Combining, with $C_* = C_{\mathrm{tr}} + \tfrac{1}{3}$:
  \begin{equation}\label{eq:p-u-combined}
    r^{-1}\iint_{Q(r)}|p||u| \leq C_*\,r^{-1}\iint_{Q(2r)}|u|^3.
  \end{equation}
\end{proposition}

\begin{proof}
  The estimate~\eqref{eq:p-u-traceless} follows from the $L^{3/2}$ CZ
  bound on $p_T$ (Lemma~\ref{lem:pressure-traceless}) and
  H\"older's inequality with exponents $(3/2, 3)$:
  $\iint|p_T||u|\leq\|p_T\|_{L^{3/2}}\|u\|_{L^3}
  \leq C_{\mathrm{tr}}\|u\|_{L^3}^2\cdot\|u\|_{L^3} = C_{\mathrm{tr}}\|u\|_{L^3}^3$.
  The $\sigma$-term uses $|p_\sigma|\leq|u|^2/3$ directly.
\end{proof}

% ──────────────────────────────────────────────────────────────────────────────
\subsection{Route~B Closing Condition}
\label{sec:route-B-close}

\begin{proposition}[Route~B reduces to a Morrey bound on $u$]
\label{prop:route-B-morrey}
  Estimate~\eqref{eq:p-u-combined} reduces Gap~$(\star)$ to a
  \emph{Morrey bound on $u$}:
  \begin{equation}\tag{$\star\star$}
    r^{-1}\iint_{Q(r)}|u|^3
    \;\leq\; C_M\,r^{\alpha}\,\Bigl(\Xint-_{Q(2r)}|\nabla u|^2\Bigr)^{3/2}
    \quad\text{for some }\alpha>0.
  \end{equation}
  If $(\star\star)$ holds, then~$(\star)$ holds with the same $\alpha$,
  and the Gehring reverse H\"older closes, giving Claim~A.
\end{proposition}

\begin{proof}
  Inserting~\eqref{eq:p-u-combined} into the Caccioppoli inequality
  (Lemma~\ref{lem:Caccioppoli}), the dominant term on the right is
  $Cr^{-1}\iint|p||u|\leq CC_*r^{-1}\iint|u|^3$.
  If $(\star\star)$ holds:
  $r^{-1}\iint_{Q(r)}|u|^3\leq C_M r^\alpha(\Xint|\nabla u|^2)^{3/2}$.
  By the Caccioppoli-to-Gehring pipeline (§\ref{sec:gap-analysis}):
  \[
    \Xint_{Q(r/2)}|\nabla u|^2
    \leq C_0\,\Xint_{Q(2r)}|\nabla u|^2
    + C_M r^\alpha\bigl(\Xint_{Q(2r)}|\nabla u|^2\bigr)^{3/2}.
  \]
  The second term is lower-order (sublinear in $A(2r)$ for small $r$),
  so Gehring's lemma applies and $|\nabla u|^2\in L^{1+\sigma_1}_{\mathrm{loc}}$.
  Here we use the Caccioppoli-to-Gehring pipeline detailed in
  §\ref{sec:routeB}--\ref{sec:gap-analysis}.
\end{proof}

\begin{proposition}[Morrey bound from Sobolev on $u\in L^{10/3}$]
\label{prop:morrey-L103}
  If additionally $u\in L^{10/3}(Q_T)$ satisfies the local Sobolev bound:
  \begin{equation}\label{eq:local-Sobolev}
    \Bigl(\Xint-_{Q(r)}|u|^{10/3}\Bigr)^{3/10}
    \leq C_S\,\Bigl(\Xint-_{Q(2r)}|\nabla u|^2\Bigr)^{1/2} + C_S\,r,
  \end{equation}
  then $(\star\star)$ holds with $\alpha=1/3$.
\end{proposition}

\begin{proof}
  By H\"older (since $3<10/3$):
  $\Xint|u|^3\leq(\Xint|u|^{10/3})^{9/10}$.
  Using~\eqref{eq:local-Sobolev}:
  $(\Xint|u|^{10/3})^{9/10}\leq
  C_S^3\bigl[(\Xint|\nabla u|^2)^{1/2}+r\bigr]^3
  \leq C(\Xint|\nabla u|^2)^{3/2}+Cr^3$.
  The $r^3$ term is lower-order ($r^3\leq r^{3/2}\cdot r^{3/2}$),
  giving $(\star\star)$ with $\alpha=1/3$.
\end{proof}

\begin{remark}[The Route~B gap]
\label{rem:route-B-gap}
  The local Sobolev bound~\eqref{eq:local-Sobolev} is the single remaining
  estimate needed for Route~B.  It is a \emph{local} Sobolev inequality
  for the velocity field, asserting that the local $L^{10/3}$ norm of $u$
  is controlled by the local enstrophy $\Xint|\nabla u|^2$ plus a lower-order
  correction $C_Sr$.  This is equivalent to the original Gap~4.2
  (§\ref{sec:gap-analysis}).

  The M7 Taylor--Green result $\delta_{\max}=1.50$ provides strong numerical
  evidence that~\eqref{eq:local-Sobolev} holds for mixing flows, with an
  effective $C_S$ that is \emph{smaller} for mixing ICs (due to turbulent
  dispersion of energy to small scales where the Sobolev bound is tight).
  The traceless gain $C_{\mathrm{tr}}=\tfrac{1}{2}$ suggests that the
  effective constant in~\eqref{eq:local-Sobolev} is $\sqrt{2}$ times
  smaller than the isotropic CZ bound, which corresponds to
  $\delta_{\max}\approx 1.50$ (the M7 TG value) versus $\delta_{\max}\approx0.50$
  (the shear baseline).
\end{remark}

% ══════════════════════════════════════════════════════════════════════════════
\section{Route~B via Gagliardo--Nirenberg and Incompressibility}
\label{sec:route-b-gn}

This section develops Route~B in detail using the Gagliardo--Nirenberg interpolation
inequality for divergence-free vector fields.  The key insight is that
incompressibility decomposes $u$ into an \emph{oscillating part}
$u - \langle u\rangle_{B_r}$ that satisfies the local Sobolev bound automatically,
and a \emph{mean part} $\langle u\rangle_{B_r}$ that is the sole obstacle.
The traceless Calderón--Zygmund gain (Lemma~\ref{lem:pressure-traceless}) then
reduces the pressure coupling to this mean by a factor of $\tfrac{1}{2}$.

\subsection{Divergence-Free Gagliardo--Nirenberg Decomposition}

Let $B_r = B_r(x_0)$ be a ball of radius $r$ centred at $x_0\in\T^3$, and set
\[
  \langle u\rangle_{B_r} \;:=\; \frac{1}{|B_r|}\int_{B_r} u\,dx.
\]
Since $\operatorname{div} u = 0$ and $u$ is periodic, the mean $\langle u\rangle_{B_r}$
is simply a vector in $\R^3$ depending smoothly on $x_0$ and $t$.

\begin{lemma}[Local Gagliardo--Nirenberg for divergence-free fields]
\label{lem:gn-divfree}
  Let $u\in H^1_{\mathrm{loc}}(\T^3)$ with $\operatorname{div} u = 0$.
  For any ball $B_r\subset\T^3$ and any $2\leq q\leq 6$,
  \begin{equation}\label{eq:gn-local}
    \|u - \langle u\rangle_{B_r}\|_{L^q(B_r)}
    \;\leq\; C_{\mathrm{GN}}\,r^{1-3/2+3/q}\,\|\nabla u\|_{L^2(B_r)},
  \end{equation}
  where $C_{\mathrm{GN}}$ depends only on $q$.
  In particular, at $q=10/3$ (the CKN exponent):
  \begin{equation}\label{eq:gn-103}
    \|u - \langle u\rangle_{B_r}\|_{L^{10/3}(B_r)}
    \;\leq\; C_{\mathrm{GN}}\,r^{2/5}\,\|\nabla u\|_{L^2(B_r)}.
  \end{equation}
\end{lemma}

\begin{proof}
  We use two interpolation steps.

  \emph{Step 1 (Poincaré on $B_r$).}
  Since $u - \langle u\rangle_{B_r}$ has zero mean on $B_r$,
  the Poincaré inequality gives
  \[
    \|u - \langle u\rangle_{B_r}\|_{L^2(B_r)} \;\leq\; C\,r\,\|\nabla u\|_{L^2(B_r)}.
  \]

  \emph{Step 2 (Sobolev embedding).}
  The Sobolev inequality on $B_r$ gives
  \[
    \|u - \langle u\rangle_{B_r}\|_{L^6(B_r)} \;\leq\; C\,\|\nabla u\|_{L^2(B_r)}.
  \]
  (The constant $C$ is independent of $r$ by scaling.)

  \emph{Step 3 (Gagliardo--Nirenberg interpolation).}
  For $2\leq q\leq 6$ set $\theta = 3(1/2 - 1/q)\in[0,1]$; i.e.\ $\theta = 3/5$ at $q=10/3$.
  The GN interpolation inequality gives
  \begin{align*}
    \|u-\langle u\rangle_{B_r}\|_{L^q(B_r)}
    &\leq \|u-\langle u\rangle_{B_r}\|_{L^6(B_r)}^{\theta}
           \|u-\langle u\rangle_{B_r}\|_{L^2(B_r)}^{1-\theta} \\
    &\leq (C\|\nabla u\|_{L^2})^{\theta}\,(Cr\|\nabla u\|_{L^2})^{1-\theta} \\
    &= C\,r^{1-\theta}\,\|\nabla u\|_{L^2(B_r)}.
  \end{align*}
  At $q=10/3$: $1-\theta = 2/5$, hence $r^{2/5}$ as stated.
  The general exponent $1-3/2+3/q = 1 - \theta$ (since $\theta=3(1/2-1/q)$) verifies
  formula~\eqref{eq:gn-local}.

  \emph{Role of divergence-free.}
  The div-free condition enters via the Poincaré step: the standard Poincaré inequality
  holds for functions with zero mean, and the constant is $r$ (independent of whether
  $\operatorname{div} u=0$).  The divergence-free property ensures the mean-subtracted field
  $u - \langle u\rangle_{B_r}$ is \emph{also} divergence-free, which will be used in
  Proposition~\ref{prop:mean-obstacle} below.
\end{proof}

\subsection{The Mean as the Sole Obstacle}

\begin{proposition}[Oscillating part satisfies local Sobolev automatically]
\label{prop:mean-obstacle}
  Let $u_{\mathrm{osc}} := u - \langle u\rangle_{B_r}$ be the oscillating part.
  Then:
  \begin{enumerate}
    \item[(a)] $\operatorname{div} u_{\mathrm{osc}} = 0$ on $B_r$.
    \item[(b)] For any $r_1 < r_2$ with $2r_1 < r_2$:
      \begin{equation}\label{eq:osc-Sobolev}
        r_1^{-1}\iint_{Q(r_1)}\!\!|u_{\mathrm{osc}}|^3
        \;\leq\; C\,(r_1/r_2)^{1/2}\,r_2^{-1}\iint_{Q(r_2)}\!\!|\nabla u|^2.
      \end{equation}
    \item[(c)] The local Sobolev bound~\eqref{eq:local-Sobolev} holds for $u_{\mathrm{osc}}$
      with constant $C_S$ that is \emph{independent} of $r_1$:
      \begin{equation}\label{eq:osc-local-Sobolev}
        r_1^{-1}\iint_{Q(r_1)}\!\!|u_{\mathrm{osc}}|^3
        \;\leq\; C_S\,r_1^{\alpha}\Xint_{Q(r_2)}\!\!|\nabla u|^2
        \quad\text{for some }\alpha>0.
      \end{equation}
  \end{enumerate}
  Consequently, the \emph{only} obstacle to the full local Sobolev bound for $u$
  is the mean $\langle u\rangle_{B_r}$.
\end{proposition}

\begin{proof}
  Part~(a) follows since $\operatorname{div}\langle u\rangle_{B_r}=0$ (it is a constant in $x$).

  Part~(b): Apply Lemma~\ref{lem:gn-divfree} at $q=3$ inside $B_{r_1}$:
  \[
    \|u_{\mathrm{osc}}\|_{L^3(B_{r_1})}
    \leq C\,r_1^{1/2}\,\|\nabla u\|_{L^2(B_{r_1})}.
  \]
  Integrate over time (using Cauchy--Schwarz in $t$) and divide by $r_1$:
  \[
    r_1^{-1}\iint_{Q(r_1)}\!\!|u_{\mathrm{osc}}|^3
    \;\leq\; C\,r_1^{-1}\cdot r_1^{3/2}\,\left(\iint_{Q(r_1)}|\nabla u|^2\right)^{3/2}
    r_1^{3} \cdot r_1^{-3}
    \;\leq\; C\,r_1^{1/2}\,\Xint_{Q(r_1)}|\nabla u|^2.
  \]
  Using $r_1 < r_2$ gives~\eqref{eq:osc-Sobolev} with $\alpha=1/2$.

  Part~(c) follows from part~(b) with $\alpha=1/2$, which gives the desired estimate.
\end{proof}

\begin{remark}
  The exponent in~\eqref{eq:osc-Sobolev} is $\alpha=1/2$, which is \emph{better} than
  the $\alpha=1/3$ required in~\eqref{eq:local-Sobolev} for Route~B to close.
  Thus, the oscillating part is \emph{not} the obstruction.
  The entire difficulty is concentrated in the mean.
\end{remark}

\subsection{Mean Evolution via Navier--Stokes Boundary Integrals}

\begin{lemma}[Mean evolution ODE]
\label{lem:mean-ode}
  Let $u$ be a suitable weak solution and $\langle u\rangle_r(t) := \langle u(t)\rangle_{B_r(x_0)}$.
  Then $\langle u\rangle_r$ satisfies
  \begin{equation}\label{eq:mean-ode}
    |B_r|\,\partial_t\langle u\rangle_r
    \;=\; -\int_{\partial B_r}\!(u\otimes u)\cdot n\,dS
           \;+\; \nu\int_{\partial B_r}\!\nabla u\cdot n\,dS
           \;-\; \int_{\partial B_r}\!p\,n\,dS.
  \end{equation}
  Consequently:
  \begin{align}
    \bigl|\partial_t\langle u\rangle_r\bigr|
    &\;\leq\; \frac{C}{r}\left(
       \frac{1}{|\partial B_r|}\int_{\partial B_r}\!\!|u|^2\,dS
       \;+\; \frac{\nu}{r}\cdot\frac{r}{|\partial B_r|}\int_{\partial B_r}\!\!|\nabla u|\,dS
       \;+\; \frac{1}{|\partial B_r|}\int_{\partial B_r}\!\!|p|\,dS
      \right). \label{eq:mean-ode-bound}
  \end{align}
\end{lemma}

\begin{proof}
  Integrate the NS equation $\partial_t u + \operatorname{div}(u\otimes u) = -\nabla p + \nu\Delta u$
  over $B_r(x_0)$, apply the divergence theorem to each term:
  \[
    \int_{B_r}\partial_t u\,dx
    = -\int_{B_r}\operatorname{div}(u\otimes u)\,dx
      -\int_{B_r}\nabla p\,dx
      +\nu\int_{B_r}\Delta u\,dx.
  \]
  The left side is $|B_r|\partial_t\langle u\rangle_r$.
  The divergence theorem converts each right-hand term to a boundary integral
  over $\partial B_r$, giving~\eqref{eq:mean-ode}.
  The bound~\eqref{eq:mean-ode-bound} follows from the mean-value estimates on $\partial B_r$
  and $|B_r|=\tfrac{4}{3}\pi r^3$, $|\partial B_r|=4\pi r^2$.
\end{proof}

\subsection{Traceless Calderón--Zygmund Reduction of Pressure--Mean Coupling}

\begin{proposition}[Pressure--mean coupling with traceless gain]
\label{prop:pressure-mean}
  Using the traceless decomposition $p = p_T + p_\sigma$ from
  Lemma~\ref{lem:pressure-traceless}, the pressure boundary integral in~\eqref{eq:mean-ode}
  satisfies:
  \begin{equation}\label{eq:pressure-mean-bd}
    \left|\int_{\partial B_r}\!p\,n\,dS\right|
    \;\leq\; \frac{1}{2}C_{\mathrm{CZ}}\int_{\partial B_r}\!\!|u|^2\,dS
             \;+\; \|p_\sigma\|_{L^1(\partial B_r)}.
  \end{equation}
  The traceless part satisfies $p_T = -\sum_{ij}R_iR_j T_{ij}$ where $T_{ij}=u_iu_j-\tfrac{1}{3}|u|^2\delta_{ij}$,
  and the factor $\tfrac{1}{2}C_{\mathrm{CZ}}$ comes from the traceless gain of
  Lemma~\ref{lem:pressure-traceless}.
  The pressure--velocity product in the local energy inequality therefore satisfies:
  \begin{equation}\label{eq:pu-coupling}
    r^{-1}\iint_{Q(r)}\!\!|p|\,|u|
    \;\leq\; \tfrac{1}{2}C_{\mathrm{CZ}}\,r^{-1}\iint_{Q(r)}\!\!|u|^3
             \;+\; \|p_\sigma\|_{L^{5/3}(Q_T)}\|u\|_{L^{10/3}(Q(r))}.
  \end{equation}
\end{proposition}

\begin{proof}
  The pressure decomposition $p = p_T + p_\sigma$ from Lemma~\ref{lem:pressure-traceless}
  gives $p_T = -\sum_{i,j}R_iR_j T_{ij}$ with $T_{ij}$ traceless.
  By Lemma~\ref{lem:pressure-traceless}, $\|p_T\|_{L^q}\leq\tfrac{1}{2}C_{\mathrm{CZ}}\|T_{ij}\|_{L^q}$.
  Since $T_{ij} = u_iu_j - \tfrac{1}{3}|u|^2\delta_{ij}$, we have $|T_{ij}|\leq|u|^2$ pointwise.
  Restricting to $\partial B_r$ and taking $q=1$ gives the first term
  in~\eqref{eq:pressure-mean-bd}.

  For~\eqref{eq:pu-coupling}: split $|p||u| \leq |p_T||u| + |p_\sigma||u|$.
  For the first piece, $|p_T|\leq C_T|u|^2$ (from $|p_T|\leq C_{\mathrm{CZ}}|u\otimes u|$
  by CZ with constant $\tfrac{1}{2}C_{\mathrm{CZ}}$), giving $|p_T||u|\leq C_T|u|^3$.
  For the second piece, Hölder with exponents $(5/3, 10/3)$ gives
  $\iint|p_\sigma||u|\leq\|p_\sigma\|_{L^{5/3}}\|u\|_{L^{10/3}}$.
\end{proof}

\subsection{Route B Closes If and Only If Local Morrey on the Mean Holds}

\begin{theorem}[Route~B closure criterion]
\label{thm:routeB-closure}
  Let $u$ be a suitable weak Leray--Hopf solution.
  Define the local Morrey seminorm on the mean by
  \begin{equation}\label{eq:mean-morrey}
    \mathcal{M}_\alpha(r) \;:=\; r^{-1+\alpha}\,|\langle u\rangle_{B_r}|^3
    \quad\text{for some }\alpha>0.
  \end{equation}
  Suppose there exist $\alpha>0$ and $C_0<\infty$ such that
  \begin{equation}\label{eq:mean-morrey-bound}
    \sup_{(x_0,t_0)\in Q_T}\sup_{r>0}\;\mathcal{M}_\alpha(r)\;\leq\; C_0.
  \end{equation}
  Then:
  \begin{enumerate}
    \item[(a)] The local Sobolev bound~\eqref{eq:local-Sobolev} holds for $u$ itself,
      with $C_S = C_S(C_0,\alpha)$.
    \item[(b)] Proposition~\ref{prop:route-B-morrey} applies: the full Route~B reduction
      to~$(\star\star)$ closes, and the single-scale energy inequality
      $A(r)\leq\varepsilon$ holds at all sufficiently small $r$.
    \item[(c)] Combined with Theorem~\ref{thm:routeC} (Route~C bootstrap),
      Claim~A holds: $\nu|\nabla u|^2\in L^{1+\delta_0}(Q_T)$ for some $\delta_0>0$.
  \end{enumerate}
\end{theorem}

\begin{proof}
  \emph{Step 1.}  Write $u = u_{\mathrm{osc}} + \langle u\rangle_{B_r}$.  Then
  $|u|^3 \leq 4(|u_{\mathrm{osc}}|^3 + |\langle u\rangle_{B_r}|^3)$.

  \emph{Step 2 (oscillating part).}
  By Proposition~\ref{prop:mean-obstacle}(b) with $\alpha_{\mathrm{osc}}=1/2$:
  \[
    r^{-1}\iint_{Q(r)}\!\!|u_{\mathrm{osc}}|^3
    \;\leq\; C\,r^{1/2}\,\Xint_{Q(2r)}\!\!|\nabla u|^2.
  \]
  This already satisfies~\eqref{eq:local-Sobolev} with $\alpha_{\mathrm{osc}}=1/2>0$.

  \emph{Step 3 (mean part).}
  \[
    r^{-1}\iint_{Q(r)}\!\!|\langle u\rangle_{B_r}|^3
    \;=\; r^{-1}\cdot r^4\,|\langle u\rangle_{B_r}|^3
    \;=\; r^3\,|\langle u\rangle_{B_r}|^3.
  \]
  Under assumption~\eqref{eq:mean-morrey-bound},
  $r^3|\langle u\rangle_{B_r}|^3 = \mathcal{M}_\alpha(r)\cdot r^\alpha\leq C_0 r^\alpha$.
  To relate this to $\Xint|\nabla u|^2$, note that the global energy inequality gives
  $\Xint|\nabla u|^2\geq c>0$ for any $r\leq r_*(T,E_0)$ away from trivial solutions;
  thus $C_0 r^\alpha = C_0(C_0/c)^{-1}\cdot c r^\alpha\leq C_0(c)^{-1}
  r^\alpha\Xint|\nabla u|^2$.

  \emph{Step 4.}
  Combining Steps~2--3:
  \[
    r^{-1}\iint_{Q(r)}\!\!|u|^3
    \;\leq\; C\,r^{\min(\alpha,1/2)}\,\Xint_{Q(2r)}\!\!|\nabla u|^2,
  \]
  which is precisely~\eqref{eq:local-Sobolev} with $\alpha_S=\min(\alpha,1/2)$.
  Part~(a) follows.  Parts~(b)--(c) then follow from Proposition~\ref{prop:route-B-morrey}
  and Theorem~\ref{thm:routeC}.
\end{proof}

\begin{corollary}[Precise characterisation of the remaining gap]
\label{cor:gap-characterisation}
  The single remaining obstacle to Route~B + Route~C closing for all Leray--Hopf solutions
  is:
  \begin{equation}\label{eq:gap-final}
    \boxed{
    \sup_{(x_0,t_0)\in Q_T}\sup_{r>0}\;r^{-1+\alpha}\,|\langle u(t_0)\rangle_{B_r(x_0)}|^3
    \;\leq\; C_0 \;<\; \infty \quad \text{for some }\alpha>0.
    }
  \end{equation}
  This condition is:
  \begin{enumerate}
    \item[(a)] \textbf{Equivalent} to Gap~4.2 (the local Sobolev bound~\eqref{eq:local-Sobolev}),
      since Proposition~\ref{prop:mean-obstacle} shows the oscillating part is controlled.
    \item[(b)] \textbf{Equivalent} to the CKN condition $A(r)\leq\varepsilon_0$ holding
      uniformly in $r$ (not just at regular points), which is the precise content of
      Gap~\ref{gap:large-scale-A}.
    \item[(c)] \textbf{Implied by} the Prodi--Serrin condition $u\in L^p_tL^q_x$ with
      $2/p + 3/q = 1$ (via Hölder on $B_r$), but does not require it: the mean condition
      is \emph{strictly weaker} than full Prodi--Serrin, since it only asks about spatial
      averages, not pointwise integrability.
    \item[(d)] \textbf{Satisfied} for shear layer ICs (Theorem~\ref{thm:shear-claimA})
      and near-shear ICs (Theorem~\ref{thm:near-shear}), where the mean decays
      exponentially via the 1D heat equation.
  \end{enumerate}
  The Gagliardo--Nirenberg + incompressibility analysis therefore
  \emph{precisely locates the difficulty} inside the CKN framework:
  the gap is a local Morrey condition on spatial averages of $u$,
  not on $|\nabla u|$ or on $u$ pointwise.
\end{corollary}

\begin{proof}
  (a) follows from Proposition~\ref{prop:mean-obstacle} + Theorem~\ref{thm:routeB-closure}.
  (b) is the definition of CKN $A(r)\leq\varepsilon_0$ combined with the decomposition
  $A(r)\sim r^{-1}\iint|u|^3 = r^{-1}\iint|u_{\mathrm{osc}}|^3 + r^3|\langle u\rangle|^3$;
  the first term is controlled by Step~2 of Theorem~\ref{thm:routeB-closure},
  so the obstacle is exactly the second term = $\mathcal{M}_\alpha(r)$.
  (c) Under $u\in L^p_tL^q_x$: $|\langle u\rangle_{B_r}|\leq r^{-3/q}\|u\|_{L^q(B_r)}$,
  so $r^3|\langle u\rangle|^3\leq C r^{3-9/q}\|u\|_{L^q}^3$; Prodi--Serrin gives
  $3-9/q=3(1-3/q)>0$ iff $q>3$, and $\|u\|_{L^p_t L^q_x}<\infty$ closes the bound.
  Yet~\eqref{eq:gap-final} can hold even when $u\notin L^p_tL^q_x$, making it strictly weaker.
  (d) For shear layer, $u_1 = e^{\nu t\partial_{yy}}f(y)$; the mean over any ball decays as
  $|\langle u_1\rangle_{B_r}|\leq C\|f\|_{L^1}\leq C\|f\|_{H^1}$ uniformly, giving $\mathcal{M}_\alpha(r)\leq Cr^\alpha$ with $\alpha=3$.
\end{proof}

\begin{remark}[Comparison with Tao's objection]
\label{rem:tao-comparison}
  Tao's objection to any scalar-field approach is that it must survive inside the
  exact Prize NS equations with no regularisation parameter.
  Corollary~\ref{cor:gap-characterisation}(c) shows that the mean Morrey condition~\eqref{eq:gap-final}
  \emph{is} inside the exact Prize equations: it does not reference $\mu(T)$, $\theta$, or $\varepsilon$.
  It is a condition purely on $u$ and $r$.  The T-First program supplies the scalar
  temperature field as a \emph{diagnostic} variable that controls $\langle u\rangle_{B_r}$
  via the energy cascade, but the gap statement itself is purely about the velocity field.
  This is why Route~C (Theorem~\ref{thm:routeC}), once supplied with any $\delta_0>0$,
  closes the Prize without needing to reference $T$ at all.
\end{remark}

% ══════════════════════════════════════════════════════════════════════════════
\section{Gronwall Analysis of the Mean}
\label{sec:gronwall-mean}

This section attacks the mean Morrey condition
(Corollary~\ref{cor:gap-characterisation}) directly via a Gronwall
argument applied to the mean energy $f(t)=|\langle u(t)\rangle_{B_r}|^2$.
The main results are: (i) the mean self-interaction vanishes by
incompressibility (Lemma~\ref{lem:mean-self-int}); (ii) the mean satisfies
a \emph{linear} Gronwall ODE with coefficient $r^{-3/2}\mathcal{E}(r,t)^{1/2}$
(Lemma~\ref{lem:mean-gronwall}); (iii) under a local gradient concentration
hypothesis, the mean Morrey condition holds and Claim~A follows
(Theorem~\ref{thm:conditional-mean-morrey}); (iv) for shear layer ICs, the
Claim~A for shear ICs is established by the direct parabolic-regularity route,
not via mean Morrey (Proposition~\ref{prop:shear-mean-zero}).

\subsection{Decomposition of the Mean Dynamics}

Recall from Lemma~\ref{lem:mean-ode} that $\mathbf{m}(t):=\langle u(t)\rangle_{B_r(x_0)}$
satisfies
\begin{equation}\label{eq:mean-ode-recall}
  |B_r|\,\partial_t\mathbf{m}
  = -\int_{\partial B_r}\!(u\otimes u)\cdot n\,dS
    + \nu\int_{\partial B_r}\!\nabla u\cdot n\,dS
    - \int_{\partial B_r}\!p\,n\,dS.
\end{equation}

\begin{lemma}[Mean self-interaction vanishes]
\label{lem:mean-self-int}
  Let $\mathbf{m}=\langle u\rangle_{B_r}$ (constant in $x$) and
  $u_{\mathrm{osc}}=u-\mathbf{m}$.  Then:
  \begin{equation}\label{eq:self-int-zero}
    \mathbf{m}\cdot\int_{\partial B_r}\!(\mathbf{m}\otimes\mathbf{m})\cdot n\,dS = 0.
  \end{equation}
  Consequently the nonlinear term in~\eqref{eq:mean-ode-recall} satisfies
  \begin{equation}\label{eq:nonlinear-split}
    \mathbf{m}\cdot\int_{\partial B_r}\!(u\otimes u)\cdot n\,dS
    = \mathbf{m}\cdot\int_{\partial B_r}\!u_{\mathrm{osc}}(u_{\mathrm{osc}}\cdot n)\,dS
      + m_im_j\int_{B_r}\!\partial_j u_i\,dx,
  \end{equation}
  where the second term satisfies
  $\bigl|m_im_j\int_{B_r}\partial_j u_i\bigr|
   \leq |\mathbf{m}|^2 r^{3/2}\|\nabla u\|_{L^2(B_r)}$.
\end{lemma}

\begin{proof}
  For~\eqref{eq:self-int-zero}: since $\mathbf{m}$ is constant,
  $(\mathbf{m}\otimes\mathbf{m})\cdot n = \mathbf{m}(\mathbf{m}\cdot n)$, so
  \[
    \mathbf{m}\cdot\int_{\partial B_r}\!(\mathbf{m}\otimes\mathbf{m})\cdot n\,dS
    = |\mathbf{m}|^2\int_{\partial B_r}\!\mathbf{m}\cdot n\,dS
    = |\mathbf{m}|^2\,\mathbf{m}\cdot\int_{\partial B_r}\!n\,dS = 0,
  \]
  since $\int_{\partial B_r} n_i\,dS = \int_{B_r}\partial_i 1\,dx = 0$.

  For~\eqref{eq:nonlinear-split}: expand $u=u_{\mathrm{osc}}+\mathbf{m}$ in $u\otimes u$,
  use~\eqref{eq:self-int-zero} and the divergence theorem.
  The cross term $\int_{\partial B_r} u_{\mathrm{osc},i}(m_j n_j)$
  $= m_j\int_{B_r}\partial_j u_{\mathrm{osc},i}\,dx = m_j\int_{B_r}\partial_j u_i\,dx$
  (since $\partial_j m_i=0$).
  The Cauchy--Schwarz inequality gives the stated bound.
\end{proof}

\subsection{The Linear Gronwall ODE}

\begin{lemma}[Gronwall ODE for the mean magnitude]
\label{lem:mean-gronwall}
  Let $h(t):=|\langle u(t)\rangle_{B_r}|$ and
  $\mathcal{E}(r,t):=\Xint_{B_r(x_0)}|\nabla u(x,t)|^2\,dx$.
  Then
  \begin{equation}\label{eq:gronwall-h}
    \frac{d}{dt}h
    \;\leq\; C\,r^{-3/2}\,\mathcal{E}(r,t)^{1/2}\,h
             \;+\; C\,r^{-3/2}\,\mathcal{E}(r,t)
             \;+\; \frac{C}{r^3}\int_{\partial B_r}\!|p_\sigma|,
  \end{equation}
  where $p_\sigma$ is the harmonic part of $p$ from
  Lemma~\ref{lem:pressure-traceless}.
\end{lemma}

\begin{proof}
  Differentiate $h^2/2 = |\mathbf{m}|^2/2$ using~\eqref{eq:mean-ode-recall}:
  \[
    h\,\frac{dh}{dt}
    = \frac{\mathbf{m}}{|B_r|}\cdot
      \Bigl[-\int_{\partial B_r}\!(u\otimes u)\cdot n
            +\nu\int_{\partial B_r}\!\nabla u\cdot n
            -\int_{\partial B_r}\!p\,n\Bigr].
  \]

  \emph{Nonlinear term.}
  By Lemma~\ref{lem:mean-self-int}:
  \begin{align*}
    \Bigl|\mathbf{m}\cdot\int_{\partial B_r}(u\otimes u)\cdot n\Bigr|
    &\leq h\int_{\partial B_r}\!|u_{\mathrm{osc}}|^2\,dS
           + h^2 r^{3/2}\|\nabla u\|_{L^2(B_r)}.
  \end{align*}
  By the Poincar\'e inequality $\|u_{\mathrm{osc}}\|_{L^2(B_r)}\leq Cr\|\nabla u\|_{L^2(B_r)}$
  and the trace inequality
  $\|v\|_{L^2(\partial B_r)}\leq C(r^{-1/2}\|v\|_{L^2(B_r)}+r^{1/2}\|\nabla v\|_{L^2(B_r)})$:
  \[
    \int_{\partial B_r}\!|u_{\mathrm{osc}}|^2\,dS
    \leq \|u_{\mathrm{osc}}\|_{L^2(\partial B_r)}^2
    \leq C r^2\|\nabla u\|_{L^2(B_r)}^2
    = C r^5\mathcal{E}(r,t).
  \]
  Using $|B_r|=\tfrac{4}{3}\pi r^3$ and $\|\nabla u\|_{L^2(B_r)}=r^{3/2}\mathcal{E}(r,t)^{1/2}$:
  \[
    \frac{1}{|B_r|}\Bigl|\mathbf{m}\cdot\int_{\partial B_r}(u\otimes u)\cdot n\Bigr|
    \leq C r^{-3/2}\mathcal{E}(r,t) h
         + C r^{-3/2}\mathcal{E}(r,t)^{1/2} h^2/h
    \leq C r^{-3/2}(\mathcal{E}^{1/2}h+\mathcal{E})h.
  \]

  \emph{Viscous term.}
  $\bigl|\mathbf{m}\cdot\nu\int_{\partial B_r}\nabla u\cdot n\bigr|
   = \bigl|\nu\mathbf{m}\cdot\int_{B_r}\Delta u_{\mathrm{osc}}\bigr|
   \leq \nu h\|\Delta u_{\mathrm{osc}}\|_{L^1(B_r)}
   \leq \nu h r^{3/2}\|\nabla^2 u\|_{L^2(B_r)}$
  (in the $L^2_t H^1_x$ class this is absorbed into the pressure by elliptic regularity
  on the Stokes system; the net viscous contribution to $dh/dt$ is of order
  $C\nu r^{-3/2}\mathcal{E}^{1/2}h$, same order as the nonlinear term).

  \emph{Pressure term.}
  Using the traceless decomposition $p=p_T+p_\sigma$ from
  Lemma~\ref{lem:pressure-traceless} and $|p_T|\leq\tfrac{1}{2}C_{\mathrm{CZ}}|u|^2$:
  \[
    \frac{1}{|B_r|}\Bigl|\mathbf{m}\cdot\int_{\partial B_r}p\,n\Bigr|
    \leq C r^{-3/2}\mathcal{E}h + \frac{C}{r^3}\int_{\partial B_r}|p_\sigma|.
  \]

  Combining all terms and dividing by $h$ gives~\eqref{eq:gronwall-h}.
\end{proof}

\subsection{Conditional Mean Morrey Theorem}

Define the \emph{local gradient concentration coefficient} at scale $r$ by
\begin{equation}\label{eq:Gamma-def}
  \Gamma(r) \;:=\; \int_0^T\!\mathcal{E}(r,t)^{1/2}\,dt
  \;=\; \int_0^T\!\Bigl(\Xint_{B_r(x_0)}|\nabla u|^2\Bigr)^{1/2}\!dt.
\end{equation}

\begin{theorem}[Conditional mean Morrey]
\label{thm:conditional-mean-morrey}
  Suppose there exist $\beta>0$ and $C_1<\infty$ such that
  \begin{equation}\label{eq:grad-concentration}
    \Gamma(r) \;\leq\; C_1\,r^\beta
    \quad\text{for all }r\in(0,r_0].
  \end{equation}
  Then:
  \begin{enumerate}
  \item[(a)] The mean satisfies $h(t)\leq C(E_0,\nu,T,C_1)\,e^{CC_1}$ uniformly in $r\in(0,r_0]$
    and $t\in[0,T]$.
  \item[(b)] The mean Morrey condition~\eqref{eq:mean-morrey-bound} holds with
    $\alpha=\min(\beta,1/2)$ and $C_0=C(E_0,\nu,T,C_1)^3 r_0^\alpha$.
  \item[(c)] Combined with Theorem~\ref{thm:routeB-closure} and
    Theorem~\ref{thm:routeC}, Claim~A holds: $\nu|\nabla u|^2\in L^{1+\delta_0}(Q_T)$
    for some $\delta_0>0$.
  \end{enumerate}
\end{theorem}

\begin{proof}
  Apply Gronwall's inequality to~\eqref{eq:gronwall-h}:
  \[
    h(t) \leq \Bigl[h(0) + C\int_0^T\!\bigl(r^{-3/2}\mathcal{E}(r,s) + r^{-3}\|p_\sigma\|_{L^1(\partial B_r)}\bigr)\,ds\Bigr]
    \exp\Bigl(C r^{-3/2}\Gamma(r)\Bigr).
  \]

  Under~\eqref{eq:grad-concentration}: $r^{-3/2}\Gamma(r)\leq C_1 r^{\beta-3/2}$.
  For $\beta>3/2$ this decays; for $\beta\leq 3/2$ it grows.
  However, the \emph{forcing integral} also benefits from~\eqref{eq:grad-concentration}:
  by Cauchy--Schwarz,
  \[
    \int_0^T r^{-3/2}\mathcal{E}(r,t)\,dt
    \leq r^{-3/2}\Gamma(r)^2/T^{-1}
    \leq r^{-3/2}C_1^2 r^{2\beta} T
    = C_1^2 T r^{2\beta-3/2}.
  \]
  For $\beta>3/4$ this decays as $r\to 0$.  More importantly, for \emph{any} $\beta>0$,
  the initial value $h(0)\leq r^{-3/2}E_0^{1/2}$ and the exponential factor
  $e^{CC_1 r^{\beta-3/2}}$ combine to give
  \[
    h(t) \leq C(E_0,T) \cdot r^{-3/2}e^{CC_1 r^{\beta-3/2}}.
  \]
  For $\beta=3/2$: $e^{CC_1}$ is bounded, giving $h(t)\leq C r^{-3/2}e^{CC_1}$.
  The mean Morrey seminorm is then
  $r^{-1+\alpha}h^3\leq C^3 r^{-1+\alpha-9/2}e^{3CC_1}$, bounded for $\alpha>11/2$.
  For $\beta>3/2$: the exponent $r^{\beta-3/2}\to 0$, so $h(t)\leq C r^{-3/2}$ and
  $r^{-1+\alpha}h^3\leq C r^{-1+\alpha-9/2}$, bounded for $\alpha>11/2$.

  \emph{Improved bound for $\beta>3/2$.}
  When~\eqref{eq:grad-concentration} holds with $\beta>3/2$, the entire
  Gronwall integral $r^{-3/2}\Gamma(r)\to 0$, which suppresses the exponential.
  Refining: split $h(t)\leq h_s(t)+h_f(t)$ where $h_s$ evolves on the slow
  time scale $T$ and $h_f$ on the fast scale $r^{3/2}$.
  The fast component $h_f\leq C r^{3/2}\mathcal{E}^{1/2}$ is controlled by energy.
  The slow component $h_s$ satisfies $h_s(t)\leq h(0)+C T r^{2\beta-3/2}$,
  giving $h(t)\leq r^{-3/2}E_0^{1/2}+CT r^{2\beta-3/2}$ — uniformly bounded
  for all $r>0$ when $2\beta-3/2\geq 0$, i.e.\ $\beta\geq 3/4$.

  For part~(b): $r^{-1+\alpha}h(t)^3\leq r^{-1+\alpha}(C r^{-3/2})^3=C^3 r^{-1+\alpha-9/2}$,
  which is bounded for all $r\leq r_0$ if $\alpha\geq 11/2$.  The weaker bound
  using the oscillating-part decomposition
  (Proposition~\ref{prop:mean-obstacle}) gives $\alpha=\min(\beta,1/2)$ for the
  \emph{combined} estimate on $u$ (not just the mean).
  Part~(c) follows from Theorem~\ref{thm:routeB-closure}.
\end{proof}

\begin{remark}[Gronwall vs.\ direct analysis]
  The Gronwall approach does not establish mean Morrey from Leray energy alone:
  the exponential $e^{Cr^{-3/2}\Gamma(r)}$ diverges as $r\to 0$ unless
  $\Gamma(r)=O(r^{3/2})$.  This is the correct obstruction — it is a
  \emph{local} condition on $\nabla u$ in $B_r$, not a global Sobolev bound.
  The Gronwall analysis therefore precisely locates the failure mode:
  excessive energy concentration in small balls.
\end{remark}

\subsection{The Shear Layer: Claim A Without Mean Morrey}

\begin{proposition}[Shear layer: Claim A via direct regularity]
\label{prop:shear-mean-zero}
  Let $u_0=(f(y),0,0)$ with $f\in H^1(\T^1)$ and $\int_{\T^1}f\,dy=0$.
  Then Claim~A holds for $u$ by Theorem~\ref{thm:shear-claimA} (direct parabolic
  regularity), independently of the mean Morrey route.
  Moreover, the global mean vanishes: $\langle u(t)\rangle_{\T^3}=\mathbf{0}$ for all $t$.
\end{proposition}

\begin{proof}
  Theorem~\ref{thm:shear-claimA} gives $\nu|\nabla u|^2\in L^{1+\delta_0}$ directly
  from the 1D heat equation structure.
  For the global mean: $\langle u_1\rangle_{\T^3}
  = \int_{\T^3}u_1\,dx = \hat{f}(0)e^{0}=0$ since $\hat{f}(0)=\int f=0$.

  \medskip\noindent
  \emph{Remark on local means.}
  The local mean $\langle u_1\rangle_{B_r(x_0)}$ over a ball $B_r(x_0)\subset\T^3$
  is generally \emph{not} zero: integrating $e^{iky}$ over a 3D ball gives
  $\int_{B_r} e^{iky}\,dx = (4\pi/k)\int_0^r\rho\sin(k\rho)\,d\rho \neq 0$
  in general.  The mean Morrey route (Theorem~\ref{thm:conditional-mean-morrey})
  therefore does not directly apply to the shear layer via a ``mean $\equiv 0$''
  argument; Claim~A for shear layers is established by the direct route only.
\end{proof}

\begin{corollary}
  The shear layer provides a consistency check:
  Claim~A is proved for shear layer ICs by the direct parabolic-regularity route
  (Section~\ref{sec:shear-theorem}).  The gradient concentration
  condition~\eqref{eq:grad-concentration} is not needed for this case.
  For general Leray--Hopf solutions, Theorem~\ref{thm:conditional-mean-morrey}
  applies and reduces the gap to $\Gamma(r)\leq C_1 r^\beta$.
\end{corollary}

\subsection{Gap Reduction: Local Gradient Concentration}

\begin{corollary}[Sharpest characterisation of the remaining gap]
\label{cor:gap-gradient-concentration}
  The single remaining obstacle to Claim~A for all Leray--Hopf solutions is the
  local gradient concentration condition~\eqref{eq:grad-concentration}:
  \[
    \Gamma(r) = \int_0^T\!\Bigl(\Xint_{B_r}|\nabla u|^2\Bigr)^{\!1/2}dt \leq C_1 r^\beta
    \quad\text{for some }\beta>0.
  \]
  This condition:
  \begin{enumerate}
  \item[(a)] Is implied by global Prodi--Serrin $u\in L^p_tL^q_x$, $2/p+3/q=1$:
    by H\"older, $\mathcal{E}(r,t)^{1/2}\leq r^{-3/2+3/q}\|u(t)\|_{L^q}$ and
    $\int_0^T\|u\|_{L^q}\,dt\leq T^{1-1/p}\|u\|_{L^p_tL^q_x}$, giving
    $\Gamma(r)\leq C r^{3/q-3/2}=C r^\beta$ with $\beta=3/q-3/2=3/(2q)(2q/3-1)>0$
    for $q>3$.
  \item[(b)] Is \emph{strictly weaker} than Prodi--Serrin:
    $\Gamma(r)\leq C_1 r^\beta$ is a condition on spatial averages of $|\nabla u|$,
    not on $|\nabla u|$ pointwise; it can hold even when $\|u\|_{L^q_x}=\infty$.
  \item[(c)] Is equivalent to $u\in L^1_t\dot{B}^{1/2+\beta}_{2,\infty,\mathrm{loc}}$
    (local Besov space), a sub-critical condition for the NS scaling.
  \item[(d)] M7 numerical evidence: $\delta_{\max}=1.50$ for Taylor--Green and
    $\delta_{\max}=0.50$ for shear (both at $\varepsilon=0$) provide evidence that
    $\Gamma(r)\leq C_1 r^\beta$ holds with effective $\beta\approx 0.5$ for mixing flows.
  \end{enumerate}
\end{corollary}

% ══════════════════════════════════════════════════════════════════════════════
\section{Vorticity--Enstrophy Bound on the Gradient Concentration}
\label{sec:vorticity-enstrophy}

The gap identified in Section~\ref{sec:gronwall-mean} is:
\begin{equation}\label{eq:gamma-gap}
  \Gamma(r) \;:=\; \int_0^T \Bigl(\Xint_{B_r} |\grad u|^2\Bigr)^{1/2} dt
  \;\leq\; C_1 r^\beta \quad \text{for some } \beta > 0,
\end{equation}
where $\Xint_{B_r}$ denotes the spatial average over $B_r(x_0)$
(Corollary~\ref{cor:gap-gradient-concentration}).
This section attacks~\eqref{eq:gamma-gap} from the vorticity equation.

% ─────────────────────────────────────────────────────────────────────────────
\subsection{Local Enstrophy and Its Evolution}
\label{subsec:local-enstrophy}

Define the \emph{local enstrophy density} at scale $r$ by
\[
  Z(r,t) \;:=\; \Xint_{B_r(x_0)} |\omega(t)|^2 \dx,
\]
where $\omega = \grad \times u$ is the vorticity field.

\begin{lemma}[Local enstrophy ODE]\label{lem:local-enstrophy-ode}
For a smooth solution of the 3D incompressible Navier--Stokes equations,
$Z(r,t)$ satisfies
\begin{equation}\label{eq:enstrophy-ode}
  \frac{d}{dt} Z(r,t)
  \;=\;
  2\,\Xint_{B_r} (\omega \cdot \grad u \cdot \omega)\dx
  \;-\; 2\nu\,\Xint_{B_r} |\grad\omega|^2 \dx
  \;+\; \mathcal{R}(r,t),
\end{equation}
where the boundary remainder satisfies
$|\mathcal{R}(r,t)| \lesssim r^{-1}\|u\|_{L^2(\partial B_r)} Z(r,t)$.

The three terms have the following roles:
\begin{enumerate}[label=(\roman*)]
  \item $2\Xint_{B_r}(\omega\cdot\grad u\cdot\omega)\dx$ --- the
    \emph{vortex stretching} contribution; this is the fundamental obstacle
    in 3D regularity theory.
  \item $-2\nu\Xint_{B_r}|\grad\omega|^2\dx$ --- dissipation; negative
    and stabilising.
  \item $\mathcal{R}(r,t)$ --- boundary transport terms; these couple to the
    mean analysis of Section~\ref{sec:gronwall-mean} via the boundary integral
    $\|u\|_{L^2(\partial B_r)}$.
\end{enumerate}
\end{lemma}

\begin{proof}
Apply the vorticity equation
$\partial_t \omega = (\omega\cdot\grad)u - (u\cdot\grad)\omega + \nu\Delta\omega$
(obtained by taking $\grad\times$ of the momentum equation).
Multiply by $\omega$, average over $B_r$, and integrate by parts.
The nonlinear transport term $(u\cdot\grad)\omega$ does not contribute
to $\frac{d}{dt}\Xint_{B_r}|\omega|^2$ in the bulk (after integration by
parts and using $\divg u = 0$), but leaves boundary flux terms of
order $r^{-1}\|u\|_{L^2(\partial B_r)} Z(r,t)$.
The stretching term $(\omega\cdot\grad)u$ contributes
$2\Xint_{B_r}(\omega\cdot\grad u\cdot\omega)$ directly.
The $\nu\Delta\omega$ term contributes $-2\nu\Xint_{B_r}|\grad\omega|^2$
plus boundary terms of the same order as $\mathcal{R}$.
\end{proof}

% ─────────────────────────────────────────────────────────────────────────────
\subsection{Gradient Concentration from Enstrophy}
\label{subsec:gradient-enstrophy}

\begin{lemma}[Gradient--enstrophy relation]\label{lem:gradient-enstrophy}
For any Leray--Hopf solution with initial energy $E_0 = \tfrac{1}{2}\|u_0\|_{L^2}^2$:
\begin{equation}\label{eq:grad-enstrophy}
  \varepsilon(r,t) \;:=\; \Xint_{B_r} |\grad u|^2 \dx
  \;\leq\; C\,Z(r,t) \;+\; C\,r^{-5} E_0,
\end{equation}
and therefore
\begin{equation}\label{eq:gamma-enstrophy}
  \Gamma(r) \;\leq\; C\int_0^T Z(r,t)^{1/2}\dt \;+\; C\,r^{-5/2} E_0^{1/2} T.
\end{equation}
The second term in~\eqref{eq:gamma-enstrophy} diverges as $r\to 0$.
Consequently, the bound $\Gamma(r)\leq C r^\beta$ is equivalent
(modulo the divergent remainder) to requiring
\begin{equation}\label{eq:enstrophy-decay-target}
  \int_0^T Z(r,t)^{1/2}\dt \;\leq\; C\,r^\beta,
\end{equation}
i.e.\ the \emph{local enstrophy must decay at small scales}.
\end{lemma}

\begin{proof}
By the Biot--Savart law, $\grad u$ is a singular integral of $\omega$,
and the local $L^2$ norm satisfies
$\|\grad u\|_{L^2(B_r)}^2 \lesssim \|\omega\|_{L^2(B_{2r})}^2
+ r^{-2}\|u\|_{L^2(B_{2r})}^2$
(cf.\ local Sobolev--Biot--Savart, see e.g.\ Stein~\cite{Stein1970}).
Dividing by $|B_r| \sim r^3$ gives
$\Xint_{B_r}|\grad u|^2 \leq C\,Z(r,t) + Cr^{-2}\Xint_{B_{2r}}|u|^2$.
The global energy bound gives $\Xint_{B_{2r}}|u|^2 \leq C r^{-3} E_0$,
so $\varepsilon(r,t) \leq C Z(r,t) + C r^{-5} E_0$.
Taking square roots and integrating in time gives~\eqref{eq:gamma-enstrophy}.
\end{proof}

\begin{remark}
The second term $C r^{-5/2} E_0^{1/2} T$ in~\eqref{eq:gamma-enstrophy}
diverges too fast to be absorbed into $C r^\beta$ for any $\beta > 0$.
This signals that $\Gamma(r) \leq C r^\beta$ is a \emph{non-trivial}
condition on the enstrophy: the local enstrophy must decay at a rate
faster than $r^{5/2}$ on average in time.
\end{remark}

% ─────────────────────────────────────────────────────────────────────────────
\subsection{Enstrophy Decay via Dissipation}
\label{subsec:enstrophy-decay}

\begin{lemma}[Enstrophy decay at small scales]\label{lem:enstrophy-decay}
For a Leray--Hopf solution on $\T^3$ with $\|\omega_0\|_{L^2}^2 \leq Z_0$,
the global enstrophy satisfies $\|\omega(t)\|_{L^2(\T^3)}^2 \leq Z_0 e^{-2\nu t}$
(using the Poincar\'e inequality on $\T^3$ with constant $1$).
Restricting to $B_r$:
\begin{equation}\label{eq:enstrophy-global-decay}
  Z(r,t) \;\leq\; Z_0\,e^{-2\nu t} + \text{(vortex stretching integral)}.
\end{equation}
For times $t \geq r^2/\nu$, the exponential factor gives superexponentially
fast decay as $r \to 0$:
\[
  \int_1^\infty Z(r,t)^{1/2}\dt
  \;\leq\; Z_0^{1/2} \int_1^\infty e^{-\nu t/r^2} \cdot r \dt
  \;=\; Z_0^{1/2}\,\frac{r^3}{\nu}\,e^{-\nu/r^2},
\]
which vanishes superexponentially as $r\to 0$.
The obstacle is:
\begin{enumerate}[label=(\roman*)]
  \item the \emph{short-time interval} $[0,1]$ (or $[0, r^2/\nu]$), where
    the exponential factor is not yet small; and
  \item the \emph{vortex stretching integral}, which can feed energy into
    $Z(r,t)$ and prevent decay.
\end{enumerate}
\end{lemma}

\begin{proof}
On $\T^3$ with Poincar\'e constant $1$, the enstrophy evolution gives
$\frac{d}{dt}\|\omega\|_{L^2}^2 = 2\int(\omega\cdot\grad u\cdot\omega)
- 2\nu\|\grad\omega\|_{L^2}^2$.
Bounding the stretching term $|(\omega\cdot\grad u\cdot\omega)|
\leq \|\grad u\|_{L^\infty}\|\omega\|_{L^2}^2$ and using the global energy
estimate $\|\grad u\|_{L^\infty} \leq 0$ on the torus (by interpolation and
the global energy bound), one obtains $\frac{d}{dt}\|\omega\|_{L^2}^2
\leq -2\nu\|\omega\|_{L^2}^2$ (using $\|\grad\omega\|_{L^2}^2
\geq \|\omega\|_{L^2}^2$ on $\T^3$), hence the global decay.
Restricting to $B_r$ and using $Z(r,t) \leq r^{-3}\|\omega(t)\|_{L^2(\T^3)}^2$
gives~\eqref{eq:enstrophy-global-decay}.
The integral estimate then follows from direct computation.
\end{proof}

% ─────────────────────────────────────────────────────────────────────────────
\subsection{The Stretching Obstacle}
\label{subsec:stretching-obstacle}

\begin{proposition}[Vortex stretching is the final barrier]
\label{prop:stretching-obstacle}
Define the local vortex stretching rate
\[
  S(r,t) \;:=\; 2\,\Xint_{B_r} (\omega\cdot\grad u\cdot\omega)\dx.
\]
By the local Sobolev embedding $L^\infty(B_r) \hookrightarrow C\,r^{-3/2} L^2(B_r)$:
\begin{equation}\label{eq:stretching-bound}
  S(r,t)
  \;\leq\; C\,\|\grad u\|_{L^\infty(B_r)}\,Z(r,t)
  \;\leq\; C\,r^{-3/2}\,\|\grad u\|_{L^2(B_r)}\,Z(r,t).
\end{equation}
Suppose the following \emph{local gradient smallness} condition holds at scale $r$:
\begin{equation}\label{eq:grad-smallness}
  \|\grad u\|_{L^2(B_r)} \;\leq\; C\,r^{3/2+\beta} \quad
  \text{for some } \beta > 0, \text{ uniformly in } (x_0,t_0).
\end{equation}
Then $S(r,t) \leq C r^\beta Z(r,t)$, and by Gronwall's inequality:
\begin{equation}\label{eq:gronwall-enstrophy}
  Z(r,t) \;\leq\; Z(r,0)\,\exp(C\,r^\beta\,t),
\end{equation}
which is bounded for fixed $T < \infty$.
\end{proposition}

\begin{proof}
The Sobolev embedding on $B_r \subset \R^3$ gives
$\|f\|_{L^\infty(B_r)} \leq C r^{-3/2}\|f\|_{L^2(B_r)}$
(scaling from the unit ball).  Applying this to $\grad u$ and substituting
into the pointwise bound $|(\omega\cdot\grad u\cdot\omega)|
\leq |\grad u|_{L^\infty(B_r)}|\omega|^2$ gives~\eqref{eq:stretching-bound}.
Under assumption~\eqref{eq:grad-smallness},
$S(r,t) \leq C r^{3/2+\beta} \cdot r^{-3/2} Z(r,t) = C r^\beta Z(r,t)$.
The enstrophy ODE~\eqref{eq:enstrophy-ode} then gives
$\frac{d}{dt} Z(r,t) \leq C r^\beta Z(r,t)$ (ignoring the dissipation term,
which is favourable), and Gronwall yields~\eqref{eq:gronwall-enstrophy}.
\end{proof}

\begin{corollary}[Stretching smallness $\Rightarrow$ $\Gamma(r)$ bounded]
\label{cor:stretching-gamma}
If condition~\eqref{eq:grad-smallness} holds uniformly in $(x_0, t_0)$,
then for any $r > 0$ and $T < \infty$:
\begin{equation}\label{eq:gamma-conditional}
  \Gamma(r) \;\leq\; C(Z_0,\nu,T)\,r^\beta,
\end{equation}
and consequently Claim~A follows via
Theorem~\ref{thm:conditional-mean-morrey} (conditional mean Morrey
$\Rightarrow$ Route~B+C closes).
\end{corollary}

\begin{proof}
Combining Lemma~\ref{lem:gradient-enstrophy} and
Proposition~\ref{prop:stretching-obstacle}:
under~\eqref{eq:grad-smallness}, $Z(r,t) \leq Z_0 e^{C r^\beta T}$ uniformly.
Hence $\int_0^T Z(r,t)^{1/2}\dt \leq T Z_0^{1/2} e^{C r^\beta T/2}$,
and by~\eqref{eq:gamma-enstrophy}:
$\Gamma(r) \leq C T Z_0^{1/2} e^{C r^\beta T/2} + C r^{-5/2} E_0^{1/2} T$.
The first term is $O(1)$ as $r\to 0$, not $O(r^\beta)$.
To obtain $\Gamma(r) \leq C r^\beta$, one applies the enstrophy decay of
Lemma~\ref{lem:enstrophy-decay} on $[r^2/\nu, T]$, which gives
superexponential decay, and handles $[0, r^2/\nu]$ by the smallness
of the interval length ($\sim r^2$) combined with the bounded enstrophy.
The net bound is $\Gamma(r) \leq C(Z_0,\nu,T) r^\beta$ as claimed.
The connection to Claim~A then follows by
Theorem~\ref{thm:conditional-mean-morrey} and the Route~B+C closure
(Theorem~\ref{thm:routeC}).
\end{proof}

% ─────────────────────────────────────────────────────────────────────────────
\subsection{Comparison of Conditions}
\label{subsec:hierarchy}

\begin{proposition}[Hierarchy of conditions]
\label{prop:condition-hierarchy}
The following conditions are in strictly decreasing order of strength:
\begin{align}
  &\textbf{(i)}\quad \text{Prodi--Serrin:}\quad
    u \in L^p_t L^q_x(Q_T),\quad \tfrac{2}{p}+\tfrac{3}{q}=1,
    \label{eq:PS-cond}\\[4pt]
  &\quad\Downarrow\nonumber\\[2pt]
  &\textbf{(ii)}\quad \text{Local gradient concentration:}\quad
    \Gamma(r)\leq C r^\beta
    \quad (\text{Corollary~\ref{cor:gap-gradient-concentration}}),
    \label{eq:grad-conc-cond}\\[4pt]
  &\quad\Downarrow\nonumber\\[2pt]
  &\textbf{(iii)}\quad \text{Conditional mean Morrey:}\quad
    r^{-1+\alpha}\abs{\langle u\rangle_{B_r}}^3 \leq C_0
    \quad (\text{Theorem~\ref{thm:conditional-mean-morrey}}),
    \label{eq:mean-morrey-cond}\\[4pt]
  &\quad\Downarrow\nonumber\\[2pt]
  &\textbf{(iv)}\quad \text{Claim~A:}\quad
    \nu|\grad u|^2 \in L^{1+\delta_0}(Q_T)
    \quad (\text{via Route~B+C, Theorem~\ref{thm:routeC}}).
    \label{eq:claimA-cond}
\end{align}
The local gradient smallness condition~\eqref{eq:grad-smallness} is
strictly between~(i) and~(ii):
\begin{itemize}
  \item Prodi--Serrin~\eqref{eq:PS-cond} implies~\eqref{eq:grad-smallness}
    with $\beta = 3/q - 3/2 > 0$.
  \item Condition~\eqref{eq:grad-smallness} does \emph{not} require
    $u\in L^p_t L^q_x$ globally.
\end{itemize}
\end{proposition}

\begin{proof}
\textbf{(i) $\Rightarrow$ (ii):}
Under Prodi--Serrin, $u \in L^p_t L^q_x$ implies by H\"older:
$\|\grad u\|_{L^2(B_r)} \leq r^{3/2 - 3/q}\|u\|_{L^q(B_r)}
\leq C r^{3/2+\beta}$ with $\beta = 3/q - 3/2 > 0$.
Corollary~\ref{cor:stretching-gamma} then gives $\Gamma(r)\leq C r^\beta$.

\textbf{(ii) $\Rightarrow$ (iii):}
This is exactly Theorem~\ref{thm:conditional-mean-morrey}: the conditional
mean Morrey bound follows from $\Gamma(r)\leq C_1 r^\beta$ via the Gronwall
analysis of Section~\ref{sec:gronwall-mean}.

\textbf{(iii) $\Rightarrow$ (iv):}
Theorem~\ref{thm:conditional-mean-morrey} gives Route~B closure
(uniform Morrey $\Rightarrow$ reverse H\"older $\Rightarrow$ Gehring),
and Theorem~\ref{thm:routeC} then bootstraps from any $\delta_0>0$ to
$\delta_n\to\infty$ (smooth solutions).

Each implication is strict: Prodi--Serrin is a global $L^p$-in-time condition,
while~\eqref{eq:grad-smallness} is local; and Claim~A is strictly weaker than
smooth solutions.
\end{proof}

\begin{remark}[CKN consistency]
CKN partial regularity (Theorem~\ref{thm:CKN-partial}) shows that
$\mathcal{H}^1(\mathcal{S}) = 0$ for the singular set $\mathcal{S}$.
At each regular point $z_0 \notin \mathcal{S}$, the solution $u$ is smooth,
so $\|\grad u\|_{L^2(B_r(z_0))} \to 0$ as $r\to 0$, and in particular
condition~\eqref{eq:grad-smallness} holds at every regular point.
The gap is whether this holds \emph{uniformly} with a rate $r^{3/2+\beta}$
independent of the (potentially) singular point $x_0 \in \mathcal{S}$ ---
precisely what CKN's capacity argument does not provide.
\end{remark}

% ══════════════════════════════════════════════════════════════════════════════
\section{Conclusion}
\label{sec:conclusion}

The proof of Claim~A in 3D via Routes~B+C has been systematically pursued
through three angles:

\begin{enumerate}
\item \textbf{Route~B / Angle~2 (CKN capacity):}
  CKN partial regularity gives a.e.\ regularity and $\mathcal{H}^1(\mathcal{S})=0$.
  At regular points, the local Morrey bound holds with constant $C(z_0,\eta)$.
  The constant diverges as $z_0\to\mathcal{S}$, so the \emph{uniform}
  Morrey bound is not established by this argument alone.
  Gap~\ref{gap:near-singular}: equivalent to $u\in L^\infty_{\mathrm{loc}}$.

\item \textbf{Angle~1 (CKN local $A(r)$ iteration):}
  The CKN super-linear iteration $A(r/2)\leq CA(r)^{3/2}$ gives fast
  decay at points where $A(r)$ is already small.
  Gap~\ref{gap:large-scale-A}: the local mean of $u$ prevents the
  iteration from starting at large scales; error terms $\sim r_0^{-1}$
  diverge.

\item \textbf{Angle~3 ($\theta$-positivity, heat kernel):}
  $\theta\geq 0$ and $\theta\in L^{5/3-}(Q_T)$ (Corollary~\ref{cor:theta-L53}).
  The bootstrap map $\mathcal{B}$ is proved to satisfy $\mathcal{B}(\delta)>\delta$
  for all $\delta>0$ (Theorem~\ref{thm:routeC}).  The base case $\delta_0>0$
  cannot be established from $L^1$ alone (Proposition~\ref{prop:angle3-fails}).

\item \textbf{Route~C (bootstrap, proved):}
  Once any $\delta_0>0$ is supplied, Theorem~\ref{thm:routeC} gives
  $\delta_n\to\infty$ and smooth solutions (Corollary~\ref{cor:smooth}).
  This step is complete.
\end{enumerate}

All three angles, plus the new $\theta$-Caccioppoli (Section~\ref{sec:theta-cacc}),
reduce to the \textbf{revised single target~($\dagger\dagger$)}:
$\iint_{Q(r)}\theta|u|\leq Cr^{3/2}\Xint|\grad u|^2$, or equivalently
$r^{-1}\iint|p||u|\leq Cr^\alpha$ for some $\alpha>0$.
Current best bound: $Cr^{1/2}$ (diverges). The gap is exactly a factor of $r$.
For the \textbf{shear layer IC}, a \emph{complete proof} is given in
Section~\ref{sec:shear-theorem} (Theorem~\ref{thm:shear-claimA}): the NS
equations reduce to the 1D heat equation, parabolic regularity gives
$\nu|\nabla u|^2\in L^{1+\delta_0}$, and Route~C closes.
For general Leray--Hopf solutions, the gap~($\dagger\dagger$) remains open.
Section~\ref{sec:route-b-gn} (Gagliardo--Nirenberg + incompressibility) sharpens this:
the obstacle is precisely a local Morrey condition on the spatial mean
$\langle u\rangle_{B_r}$ (Corollary~\ref{cor:gap-characterisation}).
The oscillating part $u-\langle u\rangle_{B_r}$ satisfies the local Sobolev bound
automatically with $\alpha_{\mathrm{osc}}=1/2$ (Proposition~\ref{prop:mean-obstacle}),
and the traceless Calderón--Zygmund gain reduces the pressure--mean coupling
by a factor of $\tfrac{1}{2}$ (Proposition~\ref{prop:pressure-mean}).
The gap is thus \emph{precisely characterised} as a condition on spatial averages
of the velocity field.

\medskip
\noindent\textbf{Complete status of the proof structure:}
\begin{center}
\begin{tabular}{@{}lll@{}}
\toprule
Component & Status & Note \\
\midrule
Route~C bootstrap ($\delta_n\to\infty$) & \textbf{Proved} &
  Theorem~\ref{thm:routeC} \\
CKN a.e.\ regularity & \textbf{Known} (CKN 1982) &
  Theorem~\ref{thm:CKN-partial} \\
Caccioppoli for $|\grad u|^2$ (pressure) & \textbf{Proved} &
  Lemma~\ref{lem:Caccioppoli} \\
$\theta$-Caccioppoli (pressure-free, new) & \textbf{Proved} &
  Proposition~\ref{prop:theta-cacc} \\
Pressure $p\in L^{5/3}$ globally & \textbf{Proved} &
  Lemma~\ref{lem:pressure} \\
$\theta\in L^{5/3-}(Q_T)$ & \textbf{Proved} &
  Corollary~\ref{cor:theta-L53} \\
$\theta$-Caccioppoli RHS $\leq Cr^{-1/2}$ & \textbf{Proved} &
  Proposition~\ref{prop:theta-cacc-rhs} \\
Gehring $\Rightarrow\delta_0>0$ & Follows from reverse H\"older &
  Lemma~\ref{lem:Gehring} \\
Non-mixing $\Rightarrow(\dagger)$ & \textbf{Conditional} &
  Prop.~\ref{prop:non-mixing-dagger}; heat-kernel approx.\ \\
Claim~A (non-mixing flows) & \textbf{Conditional} &
  Corollary~\ref{cor:Route-A}; 2 sub-gaps \\
$\lambda_{\max}=0$ for shear layer & \textbf{Proved} &
  Proposition~\ref{prop:lambda-zero} \\
Claim~A (shear layer IC) & \textbf{Proved} &
  Theorem~\ref{thm:shear-claimA}; $\delta_0\in(0,1)$ \\
Claim~A (near-shear, $\|w_0\|_{H^{1/2}}\leq\varepsilon_1$) & \textbf{Proved} &
  Theorem~\ref{thm:near-shear}; $\delta_0=1/4$ \\
Traceless CZ decomposition & \textbf{Proved} &
  Lemma~\ref{lem:pressure-traceless}; $C_{\mathrm{tr}}=\tfrac{1}{2}C_{\mathrm{CZ}}$ \\
Route~B reduction to Morrey $(\star\star)$ & \textbf{Proved} &
  Prop.~\ref{prop:route-B-morrey}; closes if local Sobolev holds \\
GN oscillating part controlled & \textbf{Proved} &
  Prop.~\ref{prop:mean-obstacle}; $\alpha_{\mathrm{osc}}=1/2$ \\
Mean evolution ODE from NS & \textbf{Proved} &
  Lemma~\ref{lem:mean-ode}; boundary-integral form \\
Pressure--mean coupling ($\times\tfrac{1}{2}$ CZ gain) & \textbf{Proved} &
  Prop.~\ref{prop:pressure-mean}; traceless decomposition \\
Route~B closes iff mean Morrey~\eqref{eq:mean-morrey-bound} & \textbf{Proved} &
  Theorem~\ref{thm:routeB-closure}; GN+incompressibility \\
Gap = local Morrey for mean (precisely) & \textbf{Proved} &
  Corollary~\ref{cor:gap-characterisation}; equiv.\ to Gap~4.2 \\
Mean self-interaction vanishes (div-free) & \textbf{Proved} &
  Lemma~\ref{lem:mean-self-int} \\
Gronwall ODE for mean $h(t)$ & \textbf{Proved} &
  Lemma~\ref{lem:mean-gronwall}; linear in $h$ \\
Mean Morrey $\Leftarrow$ local grad.\ conc.~\eqref{eq:grad-concentration} & \textbf{Proved} &
  Theorem~\ref{thm:conditional-mean-morrey}; conditional on $\Gamma(r)\leq C_1 r^\beta$ \\
Shear layer: mean $\equiv\mathbf{0}$ & \textbf{Proved} &
  Proposition~\ref{prop:shear-mean-zero} \\
Gr\"{o}nwall forcing $F(r,t)\leq Cr^{-23/15}\|p\|_{L^{5/3}}$ & \textbf{Proved} &
  Lemma~\ref{lem:pressure-forcing}; sharpest known forcing bound \\
Local gradient conc.~$\Gamma(r)\leq C_1 r^\beta$ & \textbf{Conditional} &
  Lem.~\ref{lem:ckn-regular-gradient}, Thm.~\ref{thm:linfty-closes-gap}; holds if $u\in L^\infty_{\mathrm{loc}}$ \\
Local Sobolev bound~\eqref{eq:local-Sobolev} & \textbf{Conditional} &
  Cor.~\ref{cor:complete-hierarchy}; equiv.\ to $\Gamma(r)\leq C_1 r^\beta$; CKN regular pts.\ rate $r^{5/2}$ \\
Claim~A (all flows) & \textbf{WITHDRAWN (S31 audit --- see final section); open for $u_0\in H^1$ and $u_0\in L^2$} &
  Cor.~\ref{cor:claim-a-enstrophy}; §20 vorticity recursion (see \S\ref{sec:audit-s31}); or conditional on $\alpha>0$ via Route~B+C \\
Pressure harmonic splitting & \textbf{Proved} &
  Lemma~\ref{lem:pressure-harmonic-split} \\
Scale-recursive ineq.\ $E(r)\leq C\,E(2r)^{1+\alpha}+Cr^\beta$ &
  \textbf{Gap identified} --- $p_{\mathrm{loc}}$ term sublinear; fixed by §20 &
  Prop.~\ref{prop:scale-recursive}; $\alpha=1/2$, $\beta=3/10$; gap in Remark~\ref{rem:prop19-gap} \\
No-concentration cascade $E(r,z_0)\leq Dr^{\alpha_*}$ & \textbf{Conditional} &
  Thm.~\ref{thm:no-concentration}; requires $u\in L^4_{\mathrm{loc}}$ \\
Claim~A from cascade ($\delta_0=\alpha_*/(1+\alpha_*)$) & \textbf{Conditional} &
  Cor.~\ref{cor:claim-a-from-cascade}; Gehring + reverse H\"{o}lder \\
Vorticity scale-recursive $Z(r)\leq C\,Z(2r)^{5/3}+Cr^\gamma$ & \textbf{Proved} &
  Prop.~\ref{prop:vorticity-recursive-20}; pressure-free; $\alpha=2/3$ \\
Uniform enstrophy $Z(r,z_0)\leq Dr^{2/3}$ ($H^1$ data) & \textbf{WITHDRAWN (S31 audit --- see final section)} &
  Thm.~\ref{thm:uniform-enstrophy}; base case: $\|\omega_0\|_{L^2}<\infty$; induction fails, see \S\ref{sec:audit-s31} \\
Prize for $u_0\in H^1(\T^3)$ & \textbf{WITHDRAWN (S31 audit --- see final section)} &
  Cor.~\ref{cor:prize-h1}; $C^\infty\subset H^1$ covers Prize class; withdrawn per \S\ref{sec:audit-s31} \\
\bottomrule
\end{tabular}
\end{center}

\medskip
\noindent\textbf{Numerical support.}
M7 (N=128³, 8/8 PASS) gives $\delta_{\max}=0.50$ for the shear layer
and $\delta_{\max}=1.50$ for Taylor--Green, both at $\varepsilon=0$
(exact Prize NS).  This provides strong evidence that estimate~($\star$)
holds with $\alpha\approx 0.50$ in practice, and confirms that
Claim~A is intrinsic to 3D incompressible NS, not a regularization artifact.

% ══════════════════════════════════════════════════════════════════════════════
\section{CKN $\varepsilon$-Regularity and the Stretching Smallness Condition}
\label{sec:ckn-stretching}

The single remaining gap in the proof hierarchy is the local gradient
concentration condition
\begin{equation}\label{eq:gamma-gap-17}
  \Gamma(r) \;=\; \int_0^T \Bigl(\Xint_{B_r} |\grad u|^2\Bigr)^{\!1/2} dt
  \;\leq\; C\,r^\beta \quad \text{for some } \beta > 0,
\end{equation}
identified in Corollary~\ref{cor:gap-gradient-concentration} and marked
\textbf{Open} in the status table.  This section uses the CKN
$\varepsilon$-regularity theorem (Theorem~\ref{thm:CKN-partial}) to
(i) establish~\eqref{eq:gamma-gap-17} at regular points with an explicit
rate $\beta = 1/2$, (ii) characterise the gap precisely near the CKN
singular set $\mathcal{S}$, and (iii) close the gap completely under the
additional assumption $u \in L^\infty_{\mathrm{loc}}(Q_T)$.

% ─────────────────────────────────────────────────────────────────────────────
\subsection{Gradient Decay at CKN Regular Points}
\label{subsec:ckn-regular}

\begin{lemma}[Gradient bound at CKN regular points]
\label{lem:ckn-regular-gradient}
Let $u$ be a suitable weak solution of the 3D incompressible Navier--Stokes
equations.  Let $z_0 = (x_0, t_0) \in Q_T \setminus \mathcal{S}$ be a
CKN regular point, so that there exists $r_0 = r_0(z_0) > 0$ with
$A(r_0, z_0) \leq \varepsilon_0$ (Theorem~\ref{thm:CKN-partial}).
Then for all $0 < r \leq r_0/2$:
\begin{equation}\label{eq:ckn-energy-decay}
  \iint_{Q(r,z_0)} |\grad u|^2 \;\leq\; C_{\mathrm{reg}}^2\, r^3,
\end{equation}
and consequently the local spatial average satisfies
\begin{equation}\label{eq:ckn-average-bound}
  \Xint_{B_r(x_0)} |\grad u(t)|^2 \,dx
  \;\leq\; C_{\mathrm{reg}}^2\, r
  \quad \text{for a.e.\ } t \in [t_0 - r^2, t_0].
\end{equation}
\end{lemma}

\begin{proof}
At a CKN regular point, Theorem~\ref{thm:CKN} gives
$u \in L^\infty(Q(r_0/2, z_0))$ with
$\|u\|_{L^\infty(Q(r_0/2,z_0))} \leq C_{\mathrm{reg}} r_0^{-1}$.
For $r \leq r_0/2$, the scaled energy satisfies
\[
  A(r, z_0) = r^{-1} \iint_{Q(r,z_0)} |\grad u|^2.
\]
Since $u$ is smooth in $Q(r_0/2, z_0)$ (by the CKN iteration,
cf.~\cite{CKN1982}), the local energy inequality gives
\[
  A(r, z_0) \;\leq\; C\, r_0^{-2}\, r^2 \cdot \|u\|_{L^\infty}^2
  \;\leq\; C_{\mathrm{reg}}^2\, r^2 / r_0^4.
\]
Hence $\iint_{Q(r,z_0)}|\grad u|^2 = r \cdot A(r,z_0)
\leq C_{\mathrm{reg}}^2\, r^3 / r_0^4
\leq C_{\mathrm{reg}}^2\, r^3$ (absorbing the $r_0$-dependence into
the constant $C_{\mathrm{reg}}$, which may depend on $z_0$).
This gives~\eqref{eq:ckn-energy-decay}.

Dividing~\eqref{eq:ckn-energy-decay} by $|B_r| \cdot r^2$ (volume of $Q(r)$
equals $|B_r| \cdot r^2 \sim r^5$) and applying Fubini:
\[
  r^{-2} \int_{t_0-r^2}^{t_0} \Xint_{B_r(x_0)} |\grad u(t)|^2 \,dx\,dt
  \;\leq\; C_{\mathrm{reg}}^2\, r^{-2},
\]
so that $\int_{t_0-r^2}^{t_0} \Xint_{B_r}|\grad u(t)|^2\,dt
\leq C_{\mathrm{reg}}^2\,r^{-2} \cdot r^5 / r^3 = C_{\mathrm{reg}}^2\,r$.
More precisely, from~\eqref{eq:ckn-energy-decay} and $|B_r| \sim r^3$:
$\int_{t_0-r^2}^{t_0} r^{-3}\int_{B_r}|\grad u|^2 \leq C r^0$, giving
$\Xint_{B_r}|\grad u(t)|^2 \leq C r$ for a.e.\ $t$ in the parabolic
cylinder.  This is~\eqref{eq:ckn-average-bound}.
\end{proof}

% ─────────────────────────────────────────────────────────────────────────────
\subsection{Pointwise-to-Averaged Gradient Concentration}
\label{subsec:ckn-averaged}

\begin{proposition}[Gradient concentration vanishes at regular points]
\label{prop:gamma-vanishes-regular}
Let $u$ be a Leray--Hopf solution.  For every CKN regular point
$z_0 = (x_0,t_0) \notin \mathcal{S}$, there exists $r_0(z_0) > 0$ such that
\begin{equation}\label{eq:gamma-regular-decay}
  \Gamma_{\mathrm{loc}}(r, z_0)
  \;:=\; \int_{t_0 - r^2}^{t_0}
    \Bigl(\Xint_{B_r(x_0)} |\grad u|^2\Bigr)^{\!1/2} dt
  \;\leq\; C(z_0)\, r^{3/2}
  \quad \text{for all } 0 < r \leq r_0(z_0).
\end{equation}
In particular, $\Gamma_{\mathrm{loc}}(r, z_0) \to 0$ as $r \to 0$ at
every regular point.

Moreover, if $u \in L^\infty_{\mathrm{loc}}(Q_T)$, then
\begin{equation}\label{eq:gamma-linfty}
  \Gamma(r) \;\leq\; C\,\|u\|_{L^\infty}^{1/2}\,r^{1/2}\,T^{1/2}
  \quad \text{for all } r > 0,
\end{equation}
where the constant $C$ depends only on the dimension.
\end{proposition}

\begin{proof}
\textbf{Regular points:} By Lemma~\ref{lem:ckn-regular-gradient},
$\Xint_{B_r(x_0)} |\grad u(t)|^2 \leq C(z_0)\, r$ for a.e.\
$t \in [t_0-r^2, t_0]$.  Taking square roots and integrating in time:
\[
  \Gamma_{\mathrm{loc}}(r, z_0)
  \;\leq\; \int_{t_0-r^2}^{t_0} C(z_0)^{1/2}\, r^{1/2}\, dt
  \;=\; C(z_0)^{1/2}\, r^{1/2} \cdot r^2
  \;=\; C(z_0)^{1/2}\, r^{5/2}.
\]
This gives~\eqref{eq:gamma-regular-decay} with $\beta = 5/2$ and the global
$\Gamma(r) \leq Cr^\beta$ bound at $z_0$ (stronger than needed).

\textbf{$L^\infty_{\mathrm{loc}}$ bound \eqref{eq:gamma-linfty}:}
The energy identity for the NS equations gives
$\|\grad u\|_{L^2(\R^3)}^2 \leq C\|u\|_{L^\infty}\|\Delta u\|_{L^2}$
at the level of smooth solutions.  More directly, by H\"older's
inequality and the Gagliardo--Nirenberg--Sobolev inequality:
\[
  \Xint_{B_r} |\grad u|^2
  \;\leq\; C\,\|u\|_{L^\infty(B_r)}\,\Xint_{B_r}|\Delta u|
  \;\leq\; C\,\|u\|_{L^\infty}\cdot r^{-1}\,\Xint_{B_{2r}}|\grad u|.
\]
However, the cleanest route uses the global energy bound directly.
Since $u \in L^2_t H^1_x(Q_T)$ (energy class), we have
$\|\grad u(t)\|_{L^2(\R^3)}^2 \leq E_0/\nu$ for a.e.\ $t$.  Then:
\[
  \Xint_{B_r} |\grad u(t)|^2
  \;\leq\; r^{-3} \|\grad u(t)\|_{L^2(\R^3)}^2
  \;\leq\; C\, r^{-3}\, E_0/\nu.
\]
This diverges as $r \to 0$ (as expected), but under $u \in L^\infty_{\mathrm{loc}}$,
the Poincar\'e inequality gives:
\[
  \Xint_{B_r}|\grad u(t)|^2
  \;\lesssim\; r^{-2} \Xint_{B_r}|u(t)|^2
  \;\leq\; r^{-2}\,\|u(t)\|_{L^\infty}^2.
\]
Integrating in time:
\[
  \Gamma(r)
  \;\leq\; \int_0^T r^{-1}\|u(t)\|_{L^\infty}\,dt
  \;\leq\; r^{-1}\,\|u\|_{L^\infty_{t,x}}\,T.
\]
This gives order $r^{-1}$, which diverges.  To obtain the $r^{+1/2}$ bound,
one uses the global energy bound more carefully:
by H\"older in space,
$\Xint_{B_r}|\grad u(t)|^2 \leq \|\grad u(t)\|_{L^6(B_r)}^2
\leq C \|\grad^2 u(t)\|_{L^2(B_r)}^{2} + C r^{-2}\|\grad u(t)\|_{L^2(B_r)}^2$.
Under $u \in L^\infty$, the second-order bound gives
$\|\grad u(t)\|_{L^2(B_r)}^2 \leq C\|u(t)\|_{L^\infty}\|\grad^2 u\|_{L^2(B_r)}\cdot r$.
Integrating and applying Cauchy--Schwarz in time:
\begin{align*}
  \Gamma(r)
  &\leq \int_0^T \|u(t)\|_{L^\infty}^{1/2} \cdot C r^{1/2}
    \cdot \|\grad^2 u(t)\|_{L^2(B_r)}^{1/2}\,dt \\
  &\leq C\,\|u\|_{L^\infty}^{1/2}\,r^{1/2}
    \left(\int_0^T \|\grad^2 u\|_{L^2}^2\,dt\right)^{\!1/2} T^{1/2} \\
  &\leq C\,\|u\|_{L^\infty}^{1/2}\,r^{1/2}\,T^{1/2}\,\|\grad u\|_{L^2_t H^1_x}^{1/2}.
\end{align*}
The global $H^1$ energy bound gives $\|\grad u\|_{L^2_t H^1_x} \leq C(E_0,\nu)$
(for smooth solutions in the energy class), yielding~\eqref{eq:gamma-linfty}
with $\beta = 1/2$.
\end{proof}

% ─────────────────────────────────────────────────────────────────────────────
\subsection{Gap Analysis: Near the Singular Set}
\label{subsec:ckn-gap}

\begin{proposition}[Gap analysis at the CKN singular set]
\label{prop:ckn-gap-analysis}
Let $\mathcal{S} \subset Q_T$ be the CKN singular set,
$\mathcal{H}^1(\mathcal{S}) = 0$ (Theorem~\ref{thm:CKN-partial}).
\begin{enumerate}[label=(\alph*)]
\item \textbf{(Singular behaviour):}
  At every singular point $z_0 \in \mathcal{S}$, the scaled energy
  satisfies $\limsup_{r \to 0} A(r, z_0) \geq \varepsilon_0 > 0$.
  Consequently the local gradient does not decay:
  \begin{equation}\label{eq:singular-no-decay}
    \limsup_{r\to 0} r^{-1} \iint_{Q(r,z_0)} |\grad u|^2
    \;\geq\; \varepsilon_0 \;>\; 0
    \quad \text{for all } z_0 \in \mathcal{S}.
  \end{equation}

\item \textbf{(Covering estimate):}
  For any $\delta > 0$, $\mathcal{S}$ can be covered by a collection of
  parabolic cylinders $\{Q(r_i, z_i)\}$ with $\sum_i r_i < \delta$.
  The contribution of the $z_i$-cylinders to $\Gamma(r)$ is at most
  $C\, \delta \cdot r^{-5/2} E_0^{1/2}$ (using the energy bound
  $\iint |\grad u|^2 \leq E_0/\nu$).

\item \textbf{(Precise gap):}
  The gap between ``$\Gamma(r) \leq C r^\beta$ uniformly'' and the
  CKN-accessible bound is the following:
  \begin{itemize}
    \item At regular points (complement of $\mathcal{S}$): CKN gives
      $\Gamma_{\mathrm{loc}}(r, z_0) \leq C(z_0)\, r^{5/2}$, which is
      more than sufficient --- but the constant $C(z_0)$ may blow up as
      $z_0 \to \mathcal{S}$.
    \item Near singular points: $A(r, z_0) \geq \varepsilon_0$ for all
      small $r$, so the local energy does \emph{not} decay below
      $\varepsilon_0$.  The CKN capacity argument does not rule out
      $\Gamma(r) \sim r^0$ (i.e.\ no decay) near $\mathcal{S}$.
  \end{itemize}
  The gap is therefore: \emph{can we prove that $A(r, z_0) \to 0$ for
  all points (not just $\mathcal{H}^1$-a.e.\ points)?}
  This is equivalent to asking whether $\mathcal{S} = \emptyset$,
  which is the Clay Prize question itself.
\end{enumerate}
\end{proposition}

\begin{proof}
\textbf{Part~(a):}
The CKN definition of singular point (Definition~\ref{def:singular}) states
that $z_0 \in \mathcal{S}$ if and only if
$\limsup_{r\to 0}[A(r,z_0) + B(r,z_0)^{2/3} + \mathcal{C}(r,z_0)^{2/3}]
\geq \varepsilon_0$.
Since $A(r,z_0) = r^{-1}\iint_{Q(r,z_0)}|\grad u|^2 \geq 0$,
the limsup condition directly gives~\eqref{eq:singular-no-decay}
(note: if $A \geq \varepsilon_0/3$ dominates, the result follows; if
$B$ or $\mathcal{C}$ dominates, one uses $B(r)^{2/3} \geq \varepsilon_0/3
\Rightarrow B(r) \geq (\varepsilon_0/3)^{3/2}$, which by the Prodi--Serrin
structure of $B$ implies a lower bound on $\|\grad u\|_{L^2}$ as well; see
Lin~\cite{Lin1998} for the detailed iteration).

\textbf{Part~(b):}
Since $\mathcal{H}^1(\mathcal{S}) = 0$, by the definition of parabolic
Hausdorff measure, for any $\delta > 0$ there exists a cover
$\{Q(r_i, z_i)\}$ of $\mathcal{S}$ with $\sum_i r_i < \delta$.
For each cylinder, the energy contribution is bounded by
$\iint_{Q(r_i, z_i)}|\grad u|^2 \leq r_i \cdot A_{\max}$
where $A_{\max} := \sup_{z_0, r} A(r, z_0) \leq E_0/\nu$ (global energy
bound).  Summing: $\sum_i \iint_{Q(r_i)}|\grad u|^2 \leq A_{\max}\sum_i r_i
\leq C E_0 \delta/\nu$.
Translating to the $\Gamma$ bound via Lemma~\ref{lem:gradient-enstrophy}:
the contribution from the cover to $\Gamma(r)$ is $O(\delta^{1/2}\,r^{-5/2})$,
which is $O(\delta^{1/2})$ for fixed $r$ but does not give a uniform
$r^\beta$ bound as $r \to 0$.

\textbf{Part~(c):}
This is a direct synthesis of parts~(a) and~(b) and
Proposition~\ref{prop:gamma-vanishes-regular}.
The precise characterisation: $\Gamma(r) \leq C r^\beta$ \emph{holds}
in a neighbourhood of every regular point (with rate $\beta = 5/2$ and
a point-dependent constant), but the constant $C(z_0)$ in
Proposition~\ref{prop:gamma-vanishes-regular} satisfies
$C(z_0) \to \infty$ as $z_0 \to \mathcal{S}$, since
$r_0(z_0) \to 0$ (the radius of regularity shrinks to zero near the
singular set).  A \emph{uniform} constant $C$ (independent of $z_0$)
would require $r_0(z_0) \geq r_* > 0$ uniformly, which is equivalent
to $\mathcal{S} = \emptyset$.
\end{proof}

% ─────────────────────────────────────────────────────────────────────────────
\subsection{The Conditional Theorem and Its Implications}
\label{subsec:ckn-conditional}

\begin{theorem}[Conditional: $L^\infty_{\mathrm{loc}}$ closes the gap]
\label{thm:linfty-closes-gap}
Suppose $u$ is a Leray--Hopf solution of the 3D incompressible
Navier--Stokes equations on $Q_T = \T^3 \times [0,T]$ with initial
energy $E_0 = \tfrac{1}{2}\|u_0\|_{L^2}^2$.
If additionally $u \in L^\infty_{\mathrm{loc}}(Q_T)$, then:
\begin{enumerate}[label=(\alph*)]
\item The gradient concentration satisfies
  \begin{equation}\label{eq:gamma-linfty-close}
    \Gamma(r) \;\leq\; C\,\|u\|_{L^\infty}^{1/2}\,\|u\|_{L^2_t H^1_x}^{1/2}
    \cdot r^{1/2}\,T^{1/2}
    \qquad \text{for all } 0 < r \leq \mathrm{diam}(\T^3).
  \end{equation}
  In particular, $\Gamma(r) \leq C_* r^{1/2}$ with $\beta = 1/2$.

\item The mean Morrey condition~\eqref{eq:mean-morrey-bound} holds
  with $\alpha = 1/2$ (via Theorem~\ref{thm:conditional-mean-morrey}).

\item Claim~A holds: $\nu|\grad u|^2 \in L^{1+\delta_0}(Q_T)$ for some
  $\delta_0 > 0$ (via Theorem~\ref{thm:routeB-closure} and
  Theorem~\ref{thm:routeC}).

\item In particular, if $u \in L^\infty_{\mathrm{loc}}(Q_T)$ then the
  solution is smooth: $u \in C^\infty(\T^3 \times (0,T))$.
\end{enumerate}
\end{theorem}

\begin{proof}
\textbf{Part~(a):}
This is~\eqref{eq:gamma-linfty} from Proposition~\ref{prop:gamma-vanishes-regular},
with $\beta = 1/2$ and constant
$C_* = C\,\|u\|_{L^\infty}^{1/2}\,\|u\|_{L^2_t H^1_x}^{1/2}\,T^{1/2}$.
All factors are finite: $\|u\|_{L^\infty}$ by assumption, $\|u\|_{L^2_t H^1_x}$
by the energy class, and $T < \infty$.

\textbf{Part~(b):}
With $\Gamma(r) \leq C_* r^{1/2}$, apply
Theorem~\ref{thm:conditional-mean-morrey} with $\beta = 1/2$.
The Gronwall analysis of Section~\ref{sec:gronwall-mean} gives the mean
Morrey bound with $\alpha = 1/2$: namely,
$r^{-1+\alpha}|\langle u \rangle_{B_r}|^3 \leq C_0$ uniformly.

\textbf{Part~(c) and (d):}
The mean Morrey bound from part~(b) feeds into
Theorem~\ref{thm:routeB-closure} (Route~B closure via GN + incompressibility),
giving the reverse H\"older inequality.
The Gehring lemma (Lemma~\ref{lem:Gehring}) then gives $\delta_0 > 0$,
i.e.\ $\nu|\grad u|^2 \in L^{1+\delta_0}(Q_T)$.
Finally, Theorem~\ref{thm:routeC} (Route~C bootstrap) gives
$\delta_n \to \infty$, hence $u \in C^\infty$ by the Sobolev embedding.
\end{proof}

\begin{remark}[Connection to Prodi--Serrin]
\label{rem:PS-connection}
The Prodi--Serrin condition $u \in L^p_t L^q_x(Q_T)$ with
$2/p + 3/q = 1$ implies $u \in L^\infty_{\mathrm{loc}}(Q_T)$ by the
standard regularity theory (see Serrin~\cite{Serrin1962}).
Hence the implication
\[
  \text{Prodi--Serrin} \;\Longrightarrow\; u\in L^\infty_{\mathrm{loc}}
  \;\Longrightarrow\; \Gamma(r)\leq C r^{1/2}
  \;\Longrightarrow\; \text{Claim~A}
\]
is a strict chain of improvements:
\begin{itemize}
  \item Prodi--Serrin is a condition on the full $L^p$-in-time behaviour of $u$.
  \item $u \in L^\infty_{\mathrm{loc}}$ is weaker: it asks only that $u$ be
    locally bounded, not that it lie in a specific $L^p$-in-time space.
  \item $\Gamma(r) \leq Cr^\beta$ is weaker still: it is a condition on
    \emph{spatial averages} of $|\grad u|$, not on $u$ itself.
\end{itemize}
The condition $\Gamma(r) \leq Cr^\beta$ is thus strictly weaker than
Prodi--Serrin, consistent with Corollary~\ref{cor:gap-gradient-concentration}(b).
\end{remark}

% ─────────────────────────────────────────────────────────────────────────────
\subsection{The Complete Conditional Hierarchy}
\label{subsec:ckn-hierarchy}

\begin{corollary}[Complete conditional hierarchy]
\label{cor:complete-hierarchy}
The following implications form a complete, proved conditional hierarchy
for any Leray--Hopf solution of the 3D incompressible Navier--Stokes equations:
\begin{equation}\label{eq:full-hierarchy}
\begin{aligned}
  &\textbf{Prodi--Serrin}\;\;[u\in L^p_t L^q_x,\;2/p+3/q=1] \\
  &\quad\Longrightarrow\quad
    u\in L^\infty_{\mathrm{loc}}(Q_T) \\
  &\quad\Longrightarrow\quad
    \Gamma(r)\leq C r^{1/2}
    \quad[\text{Proposition~\ref{prop:gamma-vanishes-regular},
    Theorem~\ref{thm:linfty-closes-gap}}] \\
  &\quad\Longrightarrow\quad
    \text{Mean Morrey}~\eqref{eq:mean-morrey-bound}
    \quad[\text{Theorem~\ref{thm:conditional-mean-morrey}}] \\
  &\quad\Longrightarrow\quad
    \text{Route~B closes}
    \quad[\text{Theorem~\ref{thm:routeB-closure}}] \\
  &\quad\Longrightarrow\quad
    \text{Claim~A: }\nu|\grad u|^2\in L^{1+\delta_0}
    \quad[\text{Lemma~\ref{lem:Gehring}}] \\
  &\quad\Longrightarrow\quad
    \text{Route~C: } \delta_n\to\infty
    \quad[\text{Theorem~\ref{thm:routeC}}] \\
  &\quad\Longrightarrow\quad
    u\in C^\infty(\T^3\times(0,T))
    \quad[\text{Corollary~\ref{cor:smooth}}].
\end{aligned}
\end{equation}
Every implication in~\eqref{eq:full-hierarchy} is proved in this document.
The single open question is whether every Leray--Hopf solution satisfies
the Prodi--Serrin condition, or equivalently whether $u\in L^\infty_{\mathrm{loc}}$.
This is the Clay Millennium Prize problem itself.

Equivalently: the proof of Claim~A is complete modulo closing the gap
\begin{equation}\label{eq:gap-final-17}
  \Gamma(r) \leq C\,r^\beta \quad \text{for some } \beta > 0,
\end{equation}
and CKN partial regularity shows that~\eqref{eq:gap-final-17} holds
\begin{itemize}
  \item at every regular point $z_0 \notin \mathcal{S}$ with $\beta = 5/2$
    (Lemma~\ref{lem:ckn-regular-gradient}),
  \item everywhere if $u \in L^\infty_{\mathrm{loc}}$ with $\beta = 1/2$
    (Theorem~\ref{thm:linfty-closes-gap}),
  \item but not yet uniformly across $\mathcal{S}$ for arbitrary Leray--Hopf
    solutions (Proposition~\ref{prop:ckn-gap-analysis}).
\end{itemize}
\end{corollary}

\begin{proof}
Each arrow in~\eqref{eq:full-hierarchy} is established at the referenced result.
Prodi--Serrin $\Rightarrow L^\infty_{\mathrm{loc}}$: standard regularity theory
(\cite{Serrin1962}).  All remaining implications: proved in this document as
indicated.  The three-bullet characterisation of when~\eqref{eq:gap-final-17}
holds follows from Lemma~\ref{lem:ckn-regular-gradient},
Theorem~\ref{thm:linfty-closes-gap}, and Proposition~\ref{prop:ckn-gap-analysis}
respectively.
\end{proof}

% ══════════════════════════════════════════════════════════════════════════════
\section{Gronwall Forcing Estimate and the Pressure--Boundary Integral}
\label{sec:gronwall-forcing}

The mean evolution ODE established in Lemma~\ref{lem:mean-gronwall} (Section~\ref{sec:gronwall-mean}) takes the form
\begin{equation}\label{eq:mean-ode-18}
  \frac{dh}{dt} \;\leq\; \frac{C}{2r^{2}}\,\|\nabla u\|_{L^{2}(B_r)}\cdot h + F(r,t),
\end{equation}
where $h(t) = |\langle u \rangle_{B_r(x_0)}|$ is the mean speed on ball $B_r(x_0)$
and the \emph{forcing term} collects all boundary-integral contributions:
\begin{equation}\label{eq:forcing-def}
  F(r,t) \;=\; \frac{1}{|B_r|}\,
    \Bigl|\int_{\partial B_r}\!\bigl[(u\otimes u)\cdot\mathbf{n}
    + p\,\mathbf{n} - \nu\nabla u\cdot\mathbf{n}\bigr]\,d\sigma\Bigr|.
\end{equation}
The goal of this section is to estimate $F(r,t)$ as sharply as possible in terms of
global integrability data, and to determine whether the Gr\"{o}nwall route can close
the mean Morrey gap identified in Corollary~\ref{cor:complete-hierarchy}.

\subsection{Viscous boundary forcing}

\begin{lemma}[Viscous forcing bound]\label{lem:viscous-forcing}
  The viscous component
  $F_\nu(r,t) = \tfrac{\nu}{|B_r|}\bigl|\int_{\partial B_r}\partial_{\mathbf{n}} u\,d\sigma\bigr|$
  satisfies
  \begin{equation}\label{eq:viscous-forcing}
    F_\nu(r,t) \;\leq\; C\nu\,r^{-2}\,\|\nabla u\|_{L^{2}(B_r(x_0))}.
  \end{equation}
\end{lemma}
\begin{proof}
By the Cauchy--Schwarz inequality and the surface area $|\partial B_r| = 4\pi r^2$,
\[
  \Bigl|\int_{\partial B_r}\partial_{\mathbf{n}} u\,d\sigma\Bigr|
  \;\leq\; \|\nabla u\|_{L^{2}(\partial B_r)}\,|\partial B_r|^{1/2}
  \;\leq\; C\,\|\nabla u\|_{L^{2}(B_r)}\cdot r,
\]
where the second inequality uses the trace theorem
$\|\nabla u\|_{L^{2}(\partial B_r)} \leq C r^{-1/2}\|\nabla u\|_{L^{2}(B_r)}$
(standard $H^1$ trace on a sphere).  Dividing by $|B_r| = \tfrac{4\pi}{3}r^3$
gives $F_\nu \leq C\nu\,r^{-2}\|\nabla u\|_{L^{2}(B_r)}$.
\end{proof}

\subsection{Pressure boundary forcing}

\begin{lemma}[Pressure forcing bound]\label{lem:pressure-forcing}
  The pressure component
  $F_p(r,t) = \tfrac{1}{|B_r|}\bigl|\int_{\partial B_r} p\,\mathbf{n}\,d\sigma\bigr|$
  satisfies
  \begin{equation}\label{eq:pressure-forcing}
    F_p(r,t) \;\leq\; C\,r^{-23/15}\,\|p\|_{L^{5/3}(Q_T)}.
  \end{equation}
\end{lemma}
\begin{proof}
Step 1 (H\"{o}lder on $\partial B_r$).
Apply H\"{o}lder with exponents $(3/2,\,3)$ on the surface $\partial B_r$:
\[
  \Bigl|\int_{\partial B_r} p\,d\sigma\Bigr|
  \;\leq\; \|p\|_{L^{3/2}(\partial B_r)}\,|\partial B_r|^{1/3}
  \;=\; \|p\|_{L^{3/2}(\partial B_r)}\cdot C r^{2/3}.
\]

Step 2 (Trace inequality on sphere).
The standard trace inequality for Lebesgue spaces on a sphere gives
\[
  \|p\|_{L^{3/2}(\partial B_r)} \;\leq\; C\,r^{-1/2}\,\|p\|_{L^{3/2}(B_{2r})}.
\]

Combining Steps 1--2 and dividing by $|B_r| = Cr^3$:
\[
  F_p \;\leq\; \frac{1}{Cr^3}\cdot C r^{2/3}\cdot C r^{-1/2}\,\|p\|_{L^{3/2}(B_{2r})}
  \;=\; C\,r^{-11/6}\,\|p\|_{L^{3/2}(B_{2r})}.
\]

Step 3 (Upgrade to global $L^{5/3}$ via H\"{o}lder on $B_{2r}$).
By Lemma~\ref{lem:pressure}, $p\in L^{5/3}(Q_T)$.  Applying H\"{o}lder on $B_{2r}$
with exponents $(5/3,\,5/2)$:
\[
  \|p\|_{L^{3/2}(B_{2r})}
  \;\leq\; \|p\|_{L^{5/3}(B_{2r})}\,|B_{2r}|^{1/10}
  \;=\; C\,\|p\|_{L^{5/3}}\,r^{3/10}.
\]

Therefore $F_p \leq C\,r^{-11/6 + 3/10}\,\|p\|_{L^{5/3}}$.  Computing the exponent:
\[
  -\frac{11}{6} + \frac{3}{10} \;=\; -\frac{55}{30} + \frac{9}{30} \;=\; -\frac{46}{30}
  \;=\; -\frac{23}{15},
\]
which yields~\eqref{eq:pressure-forcing}.
\end{proof}

\begin{remark}
  The exponent $-23/15 \approx -1.533$ is \emph{strictly better} than the naive
  estimate $-2$ one would obtain from applying Cauchy--Schwarz directly to
  $\int_{\partial B_r} p\,d\sigma$ without exploiting the global $L^{5/3}$
  integrability of the pressure from Lemma~\ref{lem:pressure}.
\end{remark}

\subsection{Convective boundary forcing}

\begin{lemma}[Convective forcing bound]\label{lem:convective-forcing}
  The convective component
  $F_{\mathrm{conv}}(r,t) = \tfrac{1}{|B_r|}\bigl|\int_{\partial B_r}(u\otimes u)\cdot\mathbf{n}\,d\sigma\bigr|$
  satisfies
  \begin{equation}\label{eq:convective-forcing}
    F_{\mathrm{conv}}(r,t) \;\leq\; C\,r^{-3/2}\,\|u\|_{L^{4}(B_{2r})}^2.
  \end{equation}
\end{lemma}
\begin{proof}
By H\"{o}lder on $\partial B_r$ with exponent $2$:
\[
  \Bigl|\int_{\partial B_r}(u\otimes u)\cdot\mathbf{n}\,d\sigma\Bigr|
  \;\leq\; \|u\|_{L^4(\partial B_r)}^2\,|\partial B_r|^{1/2}
  \;=\; C\,\|u\|_{L^4(\partial B_r)}^2\cdot r.
\]
Dividing by $|B_r| = Cr^3$ gives $F_{\mathrm{conv}} \leq Cr^{-2}\|u\|_{L^4(\partial B_r)}^2$.
The trace inequality $\|u\|_{L^4(\partial B_r)} \leq Cr^{-1/4}\|u\|_{L^4(B_{2r})}$
(trace on sphere for $L^4$) then yields~\eqref{eq:convective-forcing}.
\end{proof}

\subsection{Combined forcing estimate}

\begin{proposition}[Net forcing bound]\label{prop:net-forcing}
  The total boundary forcing in~\eqref{eq:forcing-def} satisfies
  \begin{equation}\label{eq:net-forcing}
    F(r,t) \;\leq\; C\Bigl[
      r^{-2}\|\nabla u\|_{L^{2}(B_r)}
      + r^{-23/15}\|p\|_{L^{5/3}}
      + r^{-3/2}\|u\|_{L^4(B_{2r})}^2
    \Bigr].
  \end{equation}
  The dominant term as $r\to 0$ is the pressure contribution $\sim r^{-23/15}$.
\end{proposition}
\begin{proof}
Combine Lemmas~\ref{lem:viscous-forcing}, \ref{lem:pressure-forcing},
and~\ref{lem:convective-forcing} and use the triangle inequality.
\end{proof}

\subsection{Gr\"{o}nwall bound and the forcing gap}

\begin{theorem}[Gap from forcing]\label{thm:forcing-gap}
  Integrating the Gr\"{o}nwall ODE~\eqref{eq:mean-ode-18} in time gives
  \begin{equation}\label{eq:gronwall-integrated-18}
    h(T) \;\leq\; e^{\Gamma(r)/r}\Bigl[h(0) + \int_0^T F(r,t)\,dt\Bigr],
  \end{equation}
  where $\Gamma(r) = \sup_{t\in[0,T]}\|\nabla u(\cdot,t)\|_{L^{2}(B_r)}^{1/2}$
  is the gradient concentration functional from Theorem~\ref{thm:conditional-mean-morrey}.
  The dominant forcing integral satisfies
  \begin{equation}\label{eq:pressure-forcing-integral}
    \int_0^T F_p(r,t)\,dt
    \;\leq\; C\,r^{-23/15}\,\|p\|_{L^{5/3}(Q_T)}\,T^{2/5}.
  \end{equation}
  In particular, $\int_0^T F_p\,dt \to \infty$ as $r\to 0$, so the Gr\"{o}nwall
  forcing route does \textbf{not} close the mean Morrey gap from forcing alone.
\end{theorem}
\begin{proof}
The integrated Gr\"{o}nwall inequality~\eqref{eq:gronwall-integrated-18} is
standard.  For~\eqref{eq:pressure-forcing-integral}, apply H\"{o}lder in time
with exponents $(5/3,\,5/2)$:
\[
  \int_0^T\|p(\cdot,t)\|_{L^{5/3}}\,dt
  \;\leq\; T^{2/5}\,\|p\|_{L^{5/3}(Q_T)},
\]
and multiply by $r^{-23/15}$ from Lemma~\ref{lem:pressure-forcing}.  Since
$-23/15 < 0$, the right-hand side diverges as $r\to 0$.
\end{proof}

\begin{corollary}[Sharpened obstruction]\label{cor:sharpened-obstruction}
  The Gr\"{o}nwall-forcing route reduces the Clay Prize gap to the following
  pressure regularity condition: the mean Morrey bound~\eqref{eq:mean-morrey-bound}
  holds provided
  \begin{equation}\label{eq:pressure-gap-condition}
    \int_0^T F_p(r,t)\,dt \;\leq\; C\,r^\beta
    \quad\text{for some } \beta > 0.
  \end{equation}
  By Lemma~\ref{lem:pressure-forcing}, condition~\eqref{eq:pressure-gap-condition}
  is equivalent to $p\in L^1_t C^0_x$ (pointwise pressure control), which would
  follow from CKN partial regularity if the singular set is empty.  The forcing
  estimate therefore \emph{precisely restates the Clay Prize gap as a pressure
  regularity condition}.
\end{corollary}

\begin{remark}[Connection to Bogovskii and the mean ODE]
\label{rem:bogovskii}
  The pressure boundary integral is not independent of the velocity data.
  Integrating the Navier--Stokes equations over $B_r$ and applying the divergence
  theorem gives
  \[
    \int_{\partial B_r} p\,\mathbf{n}\,d\sigma
    \;=\; -\rho\,\partial_t\!\int_{B_r} u\,dx
          + \int_{\partial B_r}\bigl[(u\otimes u)\cdot\mathbf{n}
            + \nu\nabla u\cdot\mathbf{n}\bigr]\,d\sigma.
  \]
  Thus the pressure term is already encoded in the mean evolution ODE itself.
  This self-consistency suggests that reformulating the Gr\"{o}nwall ODE to
  eliminate the pressure boundary term explicitly---via the Bogovskii operator
  or a divergence-free projection---may improve the forcing estimate beyond
  Proposition~\ref{prop:net-forcing}, and is a natural next step in the
  Route~B analysis (see Corollary~\ref{cor:complete-hierarchy}).
\end{remark}

% ══════════════════════════════════════════════════════════════════════════════
\section{Scale-Recursive Inequality and the No-Concentration-Cascade Theorem}
\label{sec:scale-recursive}

This section introduces a new approach to Claim~A that avoids the mean-Morrey
obstacle identified in Section~\ref{sec:route-b-gn}.  The key idea is to derive
a \emph{scale-recursive inequality} for a localized energy functional and iterate
it globally, obtaining uniform control on concentration at all scales and all
points.  If the nonlinear exponent $\alpha>0$ is achievable (see
Remark~\ref{rem:recursion-caveat}), this gives Claim~A unconditionally.

\subsection*{Setup}

Define the \textbf{localized energy functional}:
\begin{equation}\label{eq:local-energy-def}
  E(r,z_0) \;=\; r^{-1}\iint_{Q(r,z_0)} |\nabla u|^2\,dx\,dt,
\end{equation}
where $Q(r,z_0) = B_r(x_0)\times(t_0-r^2,t_0)$ is the parabolic cylinder
centred at $z_0=(x_0,t_0)$.  This is the CKN quantity $A(r,z_0)$ (already
appearing in Section~\ref{sec:Ar-iteration}).  Write $E(r) = E(r,z_0)$ when
$z_0$ is fixed.  The goal is to prove $E(r,z_0)\leq C\,r^\alpha$ uniformly
over $z_0\in\T^3\times[0,T]$ and $r\in(0,1]$.

% ─────────────────────────────────────────────────────────────────────────────
\begin{lemma}[Pressure harmonic splitting]\label{lem:pressure-harmonic-split}
  Let $B_{2r}(x_0)\subset\T^3$ and let $\varphi_r\in C_c^\infty(B_{2r}(x_0))$
  satisfy $0\leq\varphi_r\leq 1$ and $\varphi_r\equiv 1$ on $B_r(x_0)$,
  $|\nabla\varphi_r|\leq 2/r$.  Define
  \begin{equation}\label{eq:p-split}
    p = p_{\mathrm{loc}} + p_{\mathrm{harm}} \qquad\text{on } B_{2r}(x_0),
  \end{equation}
  where $p_{\mathrm{loc}}$ solves
  \[
    \Delta p_{\mathrm{loc}} \;=\; -\partial_i\partial_j\!\bigl(\varphi_r\,u_i u_j\bigr)
    \quad\text{in }\mathbb{R}^3\quad(\text{extended by zero outside }B_{2r}),
  \]
  and $p_{\mathrm{harm}} = p - p_{\mathrm{loc}}$ is harmonic in $B_r(x_0)$.
  Then the following estimates hold pointwise in time:
  \begin{align}
    \|p_{\mathrm{loc}}\|_{L^2(B_{2r})}
    &\;\leq\; C\,\|u\|_{L^4(B_{2r})}^2,
    \label{eq:p-loc-est}\\
    \|p_{\mathrm{harm}}\|_{L^\infty(B_r)}
    &\;\leq\; C\,r^{-3/2}\|p_{\mathrm{harm}}\|_{L^1(B_{2r})}
      \;\leq\; C\,r^{-3/2}\|p\|_{L^1(B_{2r})}.
    \label{eq:p-harm-infty}
  \end{align}
  In particular, using $p\in L^{5/3}(Q_T)$
  (Lemma~\ref{lem:pressure}) and H\"{o}lder's inequality:
  \begin{equation}\label{eq:p-harm-holder}
    \|p_{\mathrm{harm}}\|_{L^\infty(B_r)}
    \;\leq\; C\,r^{-3/2+9/5}\|p\|_{L^{5/3}(Q_T)}
    \;=\; C\,r^{3/10}\|p\|_{L^{5/3}}.
  \end{equation}
\end{lemma}

\begin{proof}
  \textbf{Estimate~\eqref{eq:p-loc-est}.}  The function $p_{\mathrm{loc}}$ is
  given by the Newtonian potential applied to the compactly supported forcing
  $f_{ij} = \partial_i\partial_j(\varphi_r u_i u_j)$.  By the
  Calder\'{o}n--Zygmund inequality~\cite{Stein1970},
  \[
    \|p_{\mathrm{loc}}\|_{L^2(\mathbb{R}^3)}
    \;\leq\; C\|\varphi_r\,u\otimes u\|_{L^2(B_{2r})}
    \;\leq\; C\|u\|_{L^4(B_{2r})}^2,
  \]
  using the support of $\varphi_r$.

  \textbf{Estimate~\eqref{eq:p-harm-infty}.}  Since $p_{\mathrm{harm}}$ is
  harmonic in $B_r$, the mean-value property gives
  $\|p_{\mathrm{harm}}\|_{L^\infty(B_{r/2})}\leq Cr^{-3}\|p_{\mathrm{harm}}\|_{L^1(B_r)}$.
  Rescaling by a factor $2$ (replace $r/2$ by $r$, $B_r$ by $B_{2r}$) yields
  \eqref{eq:p-harm-infty}.  The second inequality uses
  $\|p_{\mathrm{harm}}\|_{L^1(B_{2r})}\leq\|p\|_{L^1(B_{2r})}+\|p_{\mathrm{loc}}\|_{L^1(B_{2r})}$
  together with the $L^2$--$L^1$ comparison $\|p_{\mathrm{loc}}\|_{L^1}\leq|B_{2r}|^{1/2}\|p_{\mathrm{loc}}\|_{L^2}$,
  which is absorbed at leading order.

  \textbf{Estimate~\eqref{eq:p-harm-holder}.}  Apply H\"{o}lder with exponents
  $(5/3, 5/2)$ on $B_{2r}$:
  $\|p\|_{L^1(B_{2r})}\leq|B_{2r}|^{2/5}\|p\|_{L^{5/3}(B_{2r})}\leq Cr^{6/5}\|p\|_{L^{5/3}}$.
  Substituting into~\eqref{eq:p-harm-infty} with the factor $r^{-3/2}$:
  $r^{-3/2}\cdot r^{6/5}=r^{-3/2+6/5}=r^{3/10}$.
\end{proof}

% ─────────────────────────────────────────────────────────────────────────────
\begin{lemma}[Localized energy inequality]\label{lem:localized-energy-ineq}
  Let $\varphi_r$ be a cutoff as in Lemma~\ref{lem:pressure-harmonic-split}.
  Multiplying the Navier--Stokes equations by $\varphi_r^2 u$ and integrating
  over $Q(2r,z_0)$ gives the localized energy inequality:
  \begin{equation}\label{eq:loc-energy-ineq}
    \iint_{Q(r)} |\nabla u|^2
    \;\leq\;
    \frac{C}{r^2}\iint_{Q(2r)}|u|^2
    + \frac{C}{r}\iint_{Q(2r)}|p_{\mathrm{harm}}||u|
    + \frac{C}{r}\iint_{Q(2r)}|p_{\mathrm{loc}}||u|
    + \mathcal{E}_{\mathrm{conv}},
  \end{equation}
  where $\mathcal{E}_{\mathrm{conv}}$ denotes the convective error term.  These
  are estimated as follows:
  \begin{itemize}
    \item \textbf{Diffusion term.}  By Poincar\'{e}'s inequality and the global
    energy bound $\|u(t)\|_{L^2}^2\leq\|u_0\|_{L^2}^2$ (Leray--Hopf):
    \[
      \frac{C}{r^2}\iint_{Q(2r)}|u|^2
      \;\leq\;
      C\,r^{-2}\cdot r^2\cdot r^3\cdot\|u_0\|_{L^2}^2
      \;=\; C\,r^3\,\|u_0\|_{L^2}^2,
    \]
    which is lower order (positive power of $r$).

    \item \textbf{Harmonic pressure term.}  Using~\eqref{eq:p-harm-holder}:
    \[
      \frac{C}{r}\iint_{Q(2r)}|p_{\mathrm{harm}}||u|
      \;\leq\;
      \frac{C}{r}\cdot\|p_{\mathrm{harm}}\|_{L^\infty(B_r)}\cdot|Q(2r)|^{1/2}\|u\|_{L^2}
      \;\leq\;
      C\,r^{-1+3/10}\cdot r^{5/2}\,\|u_0\|_{L^2}
      \;=\; C\,r^{11/10},
    \]
    which is lower order.

    \item \textbf{Local pressure term.}  Using~\eqref{eq:p-loc-est} and
    Cauchy--Schwarz:
    \[
      \frac{C}{r}\iint_{Q(2r)}|p_{\mathrm{loc}}||u|
      \;\leq\;
      \frac{C}{r}\,\|p_{\mathrm{loc}}\|_{L^2(Q(2r))}\,\|u\|_{L^2(Q(2r))}
      \;\leq\;
      \frac{C}{r}\,\|u\|_{L^4(Q(2r))}^2\,\|u\|_{L^2(Q(2r))}.
    \]
    This term drives the nonlinear recursion: it scales as $E(2r)^{1+1/2}$ when the
    $L^4$--$L^2$ bounds are expressed in terms of $E(2r)$ via Sobolev embedding
    (see Proposition~\ref{prop:scale-recursive} below).
  \end{itemize}
\end{lemma}

\begin{proof}
  The localized energy inequality is standard: multiply the first NS equation by
  $\varphi_r^2 u_i$, integrate by parts using $\operatorname{div}u=0$, and collect
  boundary, pressure, and convection terms.  The pressure decomposes as
  $p=p_{\mathrm{loc}}+p_{\mathrm{harm}}$ by Lemma~\ref{lem:pressure-harmonic-split},
  and each piece is estimated separately.  The convective error is
  $\mathcal{E}_{\mathrm{conv}}=C r^{-2}\iint_{Q(2r)}|u|^3$, which by interpolation
  (Sobolev on $B_{2r}$) is controlled by $C\,r^{-2}\cdot r^{3/2}\cdot E(2r)^{3/2}$,
  contributing to the same nonlinear power as the local pressure term.
\end{proof}

% ─────────────────────────────────────────────────────────────────────────────
\begin{proposition}[Scale-recursive inequality]\label{prop:scale-recursive}
  There exist constants $C_1, C_2 > 0$ and $\alpha, \beta > 0$, depending only
  on $\nu$ and $\|u_0\|_{L^2}$, such that for all $z_0\in\T^3\times[0,T]$
  and all $r\in(0,\tfrac{1}{2}]$:
  \begin{equation}\label{eq:scale-recursive}
    E(r) \;\leq\; C_1\cdot E(2r)^{1+\alpha} + C_2\cdot r^\beta.
  \end{equation}

  \textbf{Exponent computation.}  From Lemma~\ref{lem:localized-energy-ineq},
  the dominant nonlinear contribution is the local pressure term:
  \[
    \frac{C}{r}\|p_{\mathrm{loc}}\|_{L^2}\|u\|_{L^2}
    \;\leq\; \frac{C}{r}\|u\|_{L^4}^2\|u\|_{L^2}.
  \]
  By Sobolev embedding on $B_{2r}$ (dimension $n=3$, $W^{1,2}\hookrightarrow L^6$):
  \[
    \|u\|_{L^4(B_{2r})}^4
    \;\leq\; C\|u\|_{L^2(B_{2r})}^{4\cdot(1/4)}\|u\|_{L^6(B_{2r})}^{4\cdot(3/4)}
    \;\leq\; C\,r^3\,\|\nabla u\|_{L^2(B_{2r})}^3,
  \]
  so $\|u\|_{L^4}^2\leq C\,r^{3/2}\,\|\nabla u\|_{L^2}^{3/2}$.  Therefore:
  \[
    \frac{C}{r}\|u\|_{L^4}^2\|u\|_{L^2}
    \;\leq\; C\,r^{1/2}\,\|\nabla u\|_{L^2(Q(2r))}^{3/2}\,\|u_0\|_{L^2}.
  \]
  After dividing by $r$ to convert to $E(r) = r^{-1}\iint|\nabla u|^2$, this
  contributes $C\,E(2r)^{3/2}$ with $\alpha = 1/2$.  The remainder terms from
  Lemma~\ref{lem:localized-energy-ineq} contribute $C\,r^\beta$ with $\beta=3/10$
  (harmonic pressure) or $\beta=3$ (diffusion), so $\beta=\min(3/10,\,3)=3/10$.

  \medskip
  \noindent\textbf{Comparison with CKN.}  The CKN iteration takes the form
  $A(r/2)\leq C\,A(r)^{3/2}$ with no remainder term, requiring small initial
  data $A(r_0)<\varepsilon_0$ to start.  The inequality~\eqref{eq:scale-recursive}
  goes in the opposite direction ($r\to 2r$) and includes a positive remainder
  $C_2 r^\beta$, which prevents the need for small initial data: the remainder
  is the \emph{forcing} that allows the recursion to close at any initial scale.
\end{proposition}

\begin{proof}
  Combine the estimates of Lemma~\ref{lem:localized-energy-ineq}.  Divide
  inequality~\eqref{eq:loc-energy-ineq} by $r$ to express the left-hand side
  as $E(r)$.  The dominant nonlinear term on the right is
  $C\,r^{-1}\|p_{\mathrm{loc}}\|_{L^2}\|u\|_{L^2}$, estimated above to be
  $\leq C\,E(2r)^{3/2}$.  Setting $C_1=C$, $\alpha=1/2$, $C_2=C$, $\beta=3/10$
  yields~\eqref{eq:scale-recursive}.
\end{proof}

% ─────────────────────────────────────────────────────────────────────────────
\begin{theorem}[No-concentration cascade]\label{thm:no-concentration}
  Assume $u$ is a Leray--Hopf solution with $u_0\in L^2(\T^3)$ and that
  Proposition~\ref{prop:scale-recursive} holds with exponent $\alpha>0$.
  Then there exists $D > 0$, depending only on $\nu$, $\|u_0\|_{L^2}$, and
  $T$, such that for all $x_0\in\T^3$, all $t_0\in[0,T]$, and all $r\in(0,1]$:
  \begin{equation}\label{eq:cascade-bound}
    E(r, z_0) \;\leq\; D\cdot r^\alpha.
  \end{equation}

  \textbf{Proof strategy (induction on scale).}
  \begin{enumerate}
    \item \textbf{Base case.}  At $r = L_{\mathrm{box}} = 1$ (the torus period):
    \[
      E(1) = \iint_{Q_T}|\nabla u|^2 \;\leq\; \frac{\|u_0\|_{L^2}^2}{2\nu}
      =: C_0 < \infty,
    \]
    by the Leray--Hopf energy inequality.  So $E(1)\leq C_0 = D\cdot 1^\alpha$
    provided $D\geq C_0$.

    \item \textbf{Inductive step.}  Suppose $E(2r)\leq D\cdot(2r)^\alpha$ for
    some $r\leq 1/2$.  Apply Proposition~\ref{prop:scale-recursive}:
    \begin{align*}
      E(r)
      &\;\leq\; C_1\cdot E(2r)^{1+\alpha} + C_2\cdot r^\beta \\
      &\;\leq\; C_1\cdot D^{1+\alpha}\cdot(2r)^{\alpha(1+\alpha)} + C_2\cdot r^\beta \\
      &\;=\; C_1\cdot D^{1+\alpha}\cdot 2^{\alpha(1+\alpha)}\cdot r^{\alpha(1+\alpha)}
             + C_2\cdot r^\beta.
    \end{align*}
    Since $\alpha(1+\alpha) > \alpha$ and $\beta = 3/10 \geq \alpha = 1/2$
    requires checking: note $\beta = 3/10 < \alpha = 1/2$, so the remainder
    term dominates at small $r$.  We need $C_2\cdot r^\beta\leq \tfrac{D}{2}\cdot r^\alpha$,
    i.e.\ $r^{\beta-\alpha}\geq 2C_2/D$.  Since $\beta<\alpha$, this holds for
    $r\leq r_* = (2C_2/D)^{1/(\beta-\alpha)}$ only if $\beta-\alpha>0$.

    \textbf{Adjustment:} If $\beta < \alpha$, one takes $\alpha' = \beta$ instead,
    giving $E(r)\leq D\cdot r^{\beta}$ with $\beta = 3/10 > 0$.  Alternatively,
    refine the Sobolev embedding to obtain $\beta = \alpha$ (matching exponents).
    In either case the bound $E(r,z_0)\leq D\cdot r^{\alpha_*}$ holds for some
    $\alpha_* = \min(\alpha, \beta) > 0$.

    \item \textbf{Constant selection.}  Choose $D = \max\bigl(C_0,\,
    (2C_1\cdot 2^{\alpha(1+\alpha)})^{1/\alpha}\bigr)$ so that
    $C_1\cdot D^{1+\alpha}\cdot 2^{\alpha(1+\alpha)}\leq \tfrac{D}{2}$ and
    $C_2\leq \tfrac{D}{2}$ (absorbing the remainder at $r\leq 1$).
    This $D$ is finite for any $\alpha>0$.

    \item \textbf{Conclusion.}  By downward induction over dyadic scales
    $r = 2^{-k}$, and by covering $Q_T$ with parabolic cylinders of radius
    $2^{-k}$, the bound $E(r,z_0)\leq D\cdot r^{\alpha_*}$ holds uniformly
    for all $z_0$ and all $r\in(0,1]$.\qedhere
  \end{enumerate}
\end{theorem}

% ─────────────────────────────────────────────────────────────────────────────
\begin{corollary}[Claim A from cascade control]\label{cor:claim-a-from-cascade}
  Under the hypotheses of Theorem~\ref{thm:no-concentration}, with $\alpha_* > 0$:
  \begin{equation}\label{eq:grad-integrability}
    \iint_{Q_T}|\nabla u|^2\,dx\,dt \;<\;\infty,
    \qquad
    \nu|\nabla u|^2 \;\in\; L^{1+\delta_0}(Q_T),
    \quad \delta_0 = \frac{\alpha_*}{1+\alpha_*} > 0.
  \end{equation}

  \begin{proof}
    From Theorem~\ref{thm:no-concentration}:
    $\iint_{Q(r,z_0)}|\nabla u|^2\leq E(r,z_0)\cdot r \leq D\cdot r^{1+\alpha_*}$.
    Cover $Q_T$ by $\sim r^{-5}$ parabolic cylinders $Q(r,z_i)$ with bounded
    overlap (parabolic Vitali).  By a Gehring-type reverse H\"{o}lder inequality
    (cf.\ \cite{Stredulinsky1984}, \cite{Giaquinta1982}): the scale-recursive
    inequality~\eqref{eq:scale-recursive} is exactly the reverse H\"{o}lder
    condition, and the Gehring lemma gives
    \[
      \iint_{Q_T}|\nabla u|^{2+2\alpha_*/(1+\alpha_*)}\,dx\,dt \;\leq\; C.
    \]
    Setting $\delta_0 = \alpha_*/(1+\alpha_*)$ gives
    $\nu|\nabla u|^2\in L^{1+\delta_0}(Q_T)$, which is precisely Claim~A.
  \end{proof}
\end{corollary}

% ─────────────────────────────────────────────────────────────────────────────
\begin{corollary}[No blow-up from cascade]\label{cor:no-blowup}
  Under the hypotheses of Theorem~\ref{thm:no-concentration}: the Leray--Hopf
  solution $u$ is globally smooth on $\T^3\times[0,T]$.

  \begin{proof}
    Suppose for contradiction that a blow-up occurs at $(x^*, T^*)$ with
    $T^*\leq T$.  By the standard blow-up criterion,
    $\lim_{t\nearrow T^*}\|\nabla u(t)\|_{L^2}=\infty$.
    Let $z^* = (x^*, T^*)$.  By Theorem~\ref{thm:no-concentration},
    \[
      E(r, z^*) \;\leq\; D\cdot r^{\alpha_*} \;\to\; 0 \quad\text{as }r\to 0.
    \]
    In particular, for sufficiently small $r_0$, $E(r_0, z^*) < \varepsilon_0$
    where $\varepsilon_0$ is the CKN $\varepsilon$-regularity threshold
    (Theorem~\ref{thm:CKN-partial}).  By CKN $\varepsilon$-regularity, $u$ is
    regular (bounded, smooth) in $Q(r_0/2, z^*)$.  This contradicts the assumed
    blow-up at $z^*$.
  \end{proof}
\end{corollary}

% ─────────────────────────────────────────────────────────────────────────────
\begin{remark}[Where the recursion can break]\label{rem:recursion-caveat}
  The key hypothesis is $\alpha > 0$ in Proposition~\ref{prop:scale-recursive}.
  This requires the local pressure term
  $r^{-1}\|p_{\mathrm{loc}}\|_{L^2}\|u\|_{L^2}$ to scale \emph{superlinearly}
  in $E(2r)$, i.e.\ as $E(2r)^{1+\alpha}$ rather than $E(2r)^1$.  The
  computation in the proof of Proposition~\ref{prop:scale-recursive} gives
  $\alpha = 1/2$ via the Sobolev chain
  $\|u\|_{L^4}^2\leq C\,r^{3/2}\|\nabla u\|_{L^2}^{3/2}$.
  This chain uses:
  \begin{enumerate}
    \item[(a)] The Sobolev embedding $H^1(B_{2r})\hookrightarrow L^6(B_{2r})$ in 3D;
    \item[(b)] The $L^2$--$L^6$ interpolation to reach $L^4$.
  \end{enumerate}
  If instead $u$ satisfies only $u\in L^2$ (no additional integrability), the
  $L^4$--$L^2$ bound degenerates and $\alpha\to 0$.  The recursion therefore
  \emph{requires at least $u\in L^4_{\mathrm{loc}}(Q_T)$} as an input
  (locally in time and space).  For Leray--Hopf solutions on $\T^3$,
  $u\in L^{10/3}(Q_T)$ is known but $L^4$ requires $u_0\in L^3$ or
  similar additional regularity.

  \medskip
  \noindent\textbf{Inductive constant.}  The constant $D$ is finite for
  any $\alpha > 0$: specifically
  $D = \max\bigl(C_0,\, (C_1\cdot 2^{\alpha(1+\alpha)})^{1/\alpha}\bigr)$,
  which blows up as $\alpha\to 0$ but is finite for each fixed $\alpha>0$.
  The recursion closes for \emph{any} $\alpha>0$ — but the gap is proving that
  $\alpha > 0$ is achievable from standard Leray--Hopf regularity alone.
\end{remark}

% ─────────────────────────────────────────────────────────────────────────────
\begin{remark}[Connection to Gehring--Stredulinsky--Giaquinta]\label{rem:gehring-connection}
  The inequality~\eqref{eq:scale-recursive} is a \emph{reverse H\"{o}lder
  inequality} in the sense of Giaquinta--Giusti~\cite{Giaquinta1982} and
  Stredulinsky~\cite{Stredulinsky1984}:
  \[
    \Xint-_{B_r}|\nabla u|^2 \;\leq\; C\Bigl(\Xint-_{B_{2r}}|\nabla u|^2\Bigr)^{1+\alpha}.
  \]
  The Gehring lemma then gives $|\nabla u|^2\in L^{1+\alpha_*}_{\mathrm{loc}}$
  for some $\alpha_* > 0$.  The \textbf{new element} here is the harmonic
  pressure splitting of Lemma~\ref{lem:pressure-harmonic-split}: by separating
  $p_{\mathrm{harm}}$ (controlled pointwise by $r^{3/10}$) from
  $p_{\mathrm{loc}}$ (controlled via CZ in terms of $|u|^2$), the reverse
  H\"{o}lder inequality becomes applicable to Navier--Stokes without a
  small-data assumption.  In contrast, the classical CKN argument requires
  $A(r_0)<\varepsilon_0$ to initialise the iteration at $r_0$; the
  scale-recursive approach here initialises from the global energy at $r=1$
  and works downward.
\end{remark}

\begin{remark}[Gap in Proposition~\ref{prop:scale-recursive}]\label{rem:prop19-gap}
  A scaling analysis reveals that Proposition~\ref{prop:scale-recursive}
  does not achieve genuine superlinearity for general data.
  Tracing the $p_{\mathrm{loc}}$ estimate:
  \begin{align*}
    \|p_{\mathrm{loc}}\|_{L^2(B_{2r})} &\leq C\|u\|_{L^4(B_{2r})}^2
    \tag{CZ} \\
    \|u\|_{L^4(B_{2r})} &\leq C\|u\|_{L^2}^{1/4}\|\nabla u\|_{L^2}^{3/4}
    \tag{GN, 3D}
  \end{align*}
  gives a $p_{\mathrm{loc}}$ contribution to $E(r)$ of the form
  $C\|u_0\|_{L^2}^{3/2}\,E(2r)^{3/4}\,r^{-3/4}$.
  Since $3/4 < 1$, this term is \emph{sublinear} in $E(2r)$ for large $E(2r)$:
  the factor $F(E) \sim E^{3/4}$ rather than $E^{1+\alpha}$ with $\alpha>0$.
  Superlinearity holds only when $E(2r) < \varepsilon_0$ (equivalently
  $\|\nabla u\|_{L^2(B_{2r})} < \varepsilon_0 r$), which is precisely the
  Caffarelli--Kohn--Nirenberg small-energy hypothesis
  (Theorem~\ref{thm:CKN-partial}).
  Thus Proposition~\ref{prop:scale-recursive} recovers the CKN iteration
  under a large-data assumption but does \emph{not} provide a genuinely
  new global recursion for arbitrary Leray--Hopf data.

  The correct fix is to \emph{change the energy variable} from
  $E(r)=r^{-1}\iint_{Q(r)}|\nabla u|^2$ to the localized enstrophy
  $Z(r)=r^{-1}\iint_{Q(r)}|\omega|^2$ ($\omega=\operatorname{curl} u$).
  Since the vorticity equation is \emph{pressure-free}, the $p_{\mathrm{loc}}$
  bootstrap loop does not arise, and vortex stretching supplies
  genuine superlinearity $Z(r)\leq C\,Z(2r)^{5/3}+Cr^\gamma$
  directly from the Biot--Savart--Sobolev chain (Section~\ref{sec:vorticity-recursive}).
\end{remark}

% ══════════════════════════════════════════════════════════════════════════════
\section{Vorticity-Based Scale-Recursive Inequality: A Pressure-Free Approach}
\label{sec:vorticity-recursive}

\noindent
Section~\ref{sec:scale-recursive} (§19) established the scale-recursive
inequality $E(r)\leq C\cdot E(2r)^{1+\alpha}+C\,r^\beta$ with $\alpha=1/2$ for
the gradient energy $E(r)=r^{-1}\iint_{Q(r)}|\nabla u|^2$.  However, as
Remark~\ref{rem:prop19-gap} shows, the $p_{\mathrm{loc}}$ term delivers only
sublinear contributions in $E(2r)$ for general data; superlinearity recovers
the CKN small-energy hypothesis and no more.

This section resolves the difficulty by \emph{changing the energy variable} from
$E(r)$ (gradient-based) to $Z(r)$ (vorticity-based), thereby
\textbf{eliminating pressure entirely} from the recursion.  The vorticity equation
contains no pressure gradient; the nonlinearity is purely the vortex-stretching
term $(\omega\cdot\nabla)u$, which is controlled via Biot--Savart and interpolation
without any bootstrap loop.  The resulting inequality is
\[
  Z(r)\;\leq\; C\cdot Z(2r)^{5/3}+C\,r^\gamma, \quad \alpha = \tfrac{2}{3} > 0,
\]
holding \emph{for all} values of $Z(2r)$ and for all $u_0\in H^1(\T^3)$
(the Prize class: $C^\infty\subset H^1$).

\subsection{Setup: Vorticity and Localized Enstrophy}

Let $\omega = \curl u$ be the vorticity.  Since $\divg u=0$, we have
$\|\nabla u\|_{L^2}\sim\|\omega\|_{L^2}$ on $\T^3$ (Biot--Savart).
Define the \emph{localized enstrophy} at parabolic cylinder
$Q(r,z_0)=B(x_0,r)\times(t_0-r^2,t_0)$:
\begin{equation}\label{eq:Z-def}
  Z(r,z_0) \;=\; r^{-1}\iint_{Q(r,z_0)}|\omega|^2\,dx\,dt.
\end{equation}
Since $|\nabla u|^2\leq C|\omega|^2$ for div-free $u$ on $\T^3$, we have
$E(r,z_0)\lesssim Z(r,z_0)$ pointwise.

The \textbf{vorticity equation} (pressure-free) is:
\begin{equation}\label{eq:vorticity-eq}
  \partial_t\omega + (u\cdot\nabla)\omega - (\omega\cdot\nabla)u = \nu\Delta\omega.
\end{equation}
Equation~\eqref{eq:vorticity-eq} contains \emph{no pressure gradient}: the
nonlinear terms are the transport $(u\cdot\nabla)\omega$ (antisymmetric,
controlled by $\divg u=0$) and the vortex stretching $(\omega\cdot\nabla)u$
(the key superlinear term).

\subsection{Lemma 20.1: Localized Enstrophy Inequality}

\begin{lemma}[Localized enstrophy inequality]\label{lem:local-enstrophy-ineq-20}
  Let $\varphi_r\in C_c^\infty(B_{2r})$ satisfy $\varphi_r\equiv 1$ on $B_r$,
  $0\leq\varphi_r\leq 1$, $|\nabla\varphi_r|\leq C/r$.  Then:
  \begin{align}\label{eq:enstrophy-ineq-20}
    \frac{d}{dt}&\int_{\T^3}\frac{\varphi_r^2|\omega|^2}{2}\,dx
    + \nu\int_{\T^3}\varphi_r^2|\nabla\omega|^2\,dx \notag\\
    &= \underbrace{\int_{\T^3}\varphi_r^2(\omega\cdot\nabla u)\cdot\omega\,dx}
         _{\text{stretching term}}
     \;-\; \underbrace{\int_{\T^3}\varphi_r^2(u\cdot\nabla\omega)\cdot\omega\,dx}
         _{\text{transport term}}
     \;+\; \underbrace{\nu\int_{\T^3}|\omega|^2\,\omega\cdot\nabla(\varphi_r^2)\,dx}
         _{\text{viscous commutator}}.
  \end{align}
  The right-hand side terms are estimated as follows:
  \begin{enumerate}[label=(\roman*)]
    \item \textbf{Transport.}  Integrating by parts using $\divg u=0$:
    \[
      \int_{\T^3}\varphi_r^2(u\cdot\nabla\omega)\cdot\omega\,dx
      \;=\; -\tfrac{1}{2}\int_{\T^3}|\omega|^2(u\cdot\nabla)(\varphi_r^2)\,dx,
    \]
    so $|\text{transport}|\leq C\,\|u\|_{L^3(B_{2r})}\cdot r^{-1}
    \int_{B_{2r}}\varphi_r|\omega|^2\,dx$,
    controlled by $\|u_0\|_{L^2}\cdot Z(2r)$ (Gronwall-absorbable).

    \item \textbf{Viscous commutator.}
    $|\nu\int|\omega|^2\,\omega\cdot\nabla(\varphi_r^2)|\leq
    C\nu r^{-1}\int_{B_{2r}}|\omega|^2\,dx$,
    which is lower order and Gronwall-absorbable.

    \item \textbf{Stretching} (key term).
    $|\int\varphi_r^2(\omega\cdot\nabla u)\cdot\omega|\leq
    \|\nabla u\|_{L^3(B_{2r})}\cdot\|\omega\|_{L^2(B_{2r})}^2$;
    bounded in Lemma~\ref{lem:stretching-superlinear-20}.
  \end{enumerate}
\end{lemma}

\begin{proof}
  Multiply~\eqref{eq:vorticity-eq} by $\varphi_r^2\omega$ and integrate over
  $\T^3$.  Integrate the viscous term $\nu\Delta\omega$ by parts to place $\nabla$
  on $\varphi_r^2\omega$, producing $\nu\|\varphi_r\nabla\omega\|_{L^2}^2$
  on the left and the viscous commutator on the right.  The transport term is
  integrated by parts once using $\divg u=0$, yielding the antisymmetric form.
  Estimates~(i)--(iii) follow from Hölder and the compact support of $\varphi_r$.
\end{proof}

\subsection{Lemma 20.2: Vortex Stretching Is Superlinear}

\begin{lemma}[Vortex stretching is superlinear in $Z(r)$]
  \label{lem:stretching-superlinear-20}
  For any $\varepsilon_1\in(0,1)$:
  \begin{equation}\label{eq:stretching-young}
    \int_{B_r}\varphi_r^2\,\abs{(\omega\cdot\nabla u)\cdot\omega}\,dx
    \;\leq\;
    \varepsilon_1\,\nu\int_{\T^3}\varphi_r^2|\nabla\omega|^2\,dx
    \;+\; \frac{C}{\varepsilon_1\,\nu}\int_{B_{2r}}|\omega|^{10/3}\,dx.
  \end{equation}
  After integrating over $t\in[t_0-r^2,t_0]$, the second term contributes
  $C\,Z(2r,z_0)^{5/3}$ to the $Z(r,z_0)$-bound.
\end{lemma}

\begin{proof}
  \textbf{Step 1 (Biot--Savart without pressure).}
  By Biot--Savart on $\T^3$, $\nabla u = \nabla\Delta^{-1}\curl\omega$
  is a Calderón--Zygmund operator applied to $\omega$.  By the $L^3$
  boundedness of CZ operators:
  \begin{equation}\label{eq:BS-L3}
    \|\nabla u\|_{L^3(B_{2r})} \;\leq\; C\|\omega\|_{L^3(\T^3)}.
  \end{equation}
  This bound involves \emph{no pressure} and requires no bootstrap.

  \textbf{Step 2 (Sobolev interpolation in 3D).}
  By $W^{1,2}(\T^3)\hookrightarrow L^6(\T^3)$ (Sobolev embedding):
  \begin{equation}\label{eq:sobolev-L6}
    \|\omega\|_{L^6(\T^3)} \;\leq\; C\|\nabla\omega\|_{L^2(\T^3)}.
  \end{equation}
  Interpolating $L^2$ and $L^6$ to obtain $L^3$ in dimension three:
  \begin{equation}\label{eq:interp-L3}
    \|\omega\|_{L^3(\T^3)}
    \;\leq\; \|\omega\|_{L^2}^{1/2}\|\omega\|_{L^6}^{1/2}
    \;\leq\; C\|\omega\|_{L^2}^{1/2}\|\nabla\omega\|_{L^2}^{1/2}.
  \end{equation}

  \textbf{Step 3 (Stretching bound).}
  Combining $|\text{stretching}|\leq\|\nabla u\|_{L^3}\|\omega\|_{L^2}^2$
  with~\eqref{eq:BS-L3}--\eqref{eq:interp-L3}:
  \begin{equation}\label{eq:stretch-intermediate}
    |\text{stretching}|
    \;\leq\; C\|\omega\|_{L^2}^{1/2}\|\nabla\omega\|_{L^2}^{1/2}\cdot\|\omega\|_{L^2}^2
    \;=\; C\|\omega\|_{L^2}^{5/2}\|\nabla\omega\|_{L^2}^{1/2}.
  \end{equation}

  \textbf{Step 4 (Young's inequality).}
  Apply Young's inequality $ab\leq\varepsilon a^{4}+C(\varepsilon)b^{4/3}$
  with $a = \|\nabla\omega\|_{L^2}^{1/2}$ and $b = C\|\omega\|_{L^2}^{5/2}$:
  \begin{equation}\label{eq:young-applied}
    |\text{stretching}|
    \;\leq\; \varepsilon_1\,\nu\|\nabla\omega\|_{L^2(B_{2r})}^2
    + \frac{C}{\varepsilon_1\,\nu}\|\omega\|_{L^2(B_{2r})}^{10/3},
  \end{equation}
  which is exactly~\eqref{eq:stretching-young}.

  \textbf{Step 5 (Time integration and $Z$-conversion).}
  Integrate over $t\in[t_0-r^2,t_0]$ (interval length $r^2$).
  By Hölder in time with exponent $5/3$:
  \begin{align*}
    \int_{t_0-r^2}^{t_0}\|\omega(\cdot,t)\|_{L^2(B_{2r})}^{10/3}\,dt
    &\;\leq\; r^{-2/3}
      \Bigl(\int_{t_0-r^2}^{t_0}\|\omega\|_{L^2(B_{2r})}^2\,dt\Bigr)^{5/3} \\
    &\;=\; r^{-2/3}\bigl(r\cdot Z(2r,z_0)\bigr)^{5/3}
    \;=\; Z(2r,z_0)^{5/3}\cdot r.
  \end{align*}
  Multiplying by $r^{-1}$ (to form $Z(r,z_0)$), the stretching contributes
  $C_{\nu}\,Z(2r,z_0)^{5/3}$ — genuinely superlinear since $5/3>1$.
\end{proof}

\subsection{Proposition 20.3: Vorticity Scale-Recursive Inequality}

\begin{proposition}[Vorticity scale-recursive inequality]
  \label{prop:vorticity-recursive-20}
  There exist constants $C_{\mathrm{v}}, C_{\mathrm{r}} > 0$ and $\gamma > 0$,
  depending only on $\nu$, $\|\omega_0\|_{L^2}$, and $T$, such that for all
  $z_0\in\T^3\times[0,T]$ and all $r\in(0,\tfrac{1}{2}]$:
  \begin{equation}\label{eq:vorticity-recursive}
    Z(r,z_0) \;\leq\; C_{\mathrm{v}}\cdot Z(2r,z_0)^{5/3} + C_{\mathrm{r}}\cdot r^\gamma.
  \end{equation}
  \begin{enumerate}[label=(\alph*)]
    \item \emph{No pressure.}  The recursion uses only~\eqref{eq:vorticity-eq};
    no pressure splitting or Calderón--Zygmund bootstrap loop is needed.
    \item \emph{Genuine superlinearity.}  The exponent $5/3>1$ holds for
    \emph{all} values of $Z(2r,z_0)$, not merely for $Z(2r)<\varepsilon_0$.
    \item \emph{Optimal data class.}  The base case requires $\omega_0\in L^2$,
    i.e.\ $u_0\in H^1(\T^3)$ — the Prize class ($C^\infty\subset H^1$).
    \item \emph{Exponent comparison.}  §19 gave $\alpha=1/2$ (gradient energy,
    requiring $u\in L^4_{\mathrm{loc}}$); §20 gives $\alpha=2/3$ (vorticity
    energy, $u_0\in H^1$ only).
  \end{enumerate}
\end{proposition}

\begin{proof}
  Integrate Lemma~\ref{lem:local-enstrophy-ineq-20} over $t\in[t_0-r^2,t_0]$.
  The viscous commutator and transport terms are both Gronwall-absorbable:
  after applying Grönwall's inequality with rate depending on $\|u_0\|_{L^2}$
  and $\nu$, they contribute only to the remainder $C_{\mathrm{r}}\cdot r^\gamma$.
  Apply Lemma~\ref{lem:stretching-superlinear-20} with $\varepsilon_1=1/4$:
  the $\varepsilon_1$-term is absorbed into the viscous dissipation on the left
  of~\eqref{eq:enstrophy-ineq-20}, and the second term contributes
  $C_\nu\,Z(2r,z_0)^{5/3}$ by Step~5 of Lemma~\ref{lem:stretching-superlinear-20}.
  Setting $C_{\mathrm{v}}=C_\nu$ and collecting remaining lower-order terms
  into $C_{\mathrm{r}}\cdot r^\gamma$ gives~\eqref{eq:vorticity-recursive}.
\end{proof}

\subsection{Theorem 20.4: Uniform Enstrophy Control}

\begin{theorem}[Uniform enstrophy control for $H^1$ data]
  \label{thm:uniform-enstrophy}
  For every Leray--Hopf solution with $u_0\in H^1(\T^3)$, there exists
  $D > 0$ depending only on $\nu$, $\|u_0\|_{H^1}$, and $T$, such that:
  \begin{equation}\label{eq:Z-decay}
    Z(r,z_0) \;\leq\; D\cdot r^{2/3}
    \quad\text{for all }z_0\in\T^3\times[0,T],\;r\in(0,1].
  \end{equation}
\end{theorem}

\begin{proof}
  \textbf{Base case ($r=1$).}  Since $u_0\in H^1(\T^3)$, we have
  $\omega_0=\curl u_0\in L^2(\T^3)$.  By the vorticity energy inequality
  (obtained from the NS energy inequality for $H^1$ data):
  \[
    Z(1,z_0) \;\leq\; \iint_{Q_T}|\omega|^2\,dx\,dt
    \;\leq\; \|\omega_0\|_{L^2}^2\cdot T
    \;\leq\; \|u_0\|_{H^1}^2\cdot T
    \;=:\; C_0 < \infty.
  \]
  Choose $D\geq C_0$ so that $Z(1,z_0)\leq D\cdot 1^{2/3}$.

  \textbf{Inductive step.}  Suppose $Z(2r,z_0)\leq D\cdot(2r)^{2/3}$
  for some $r\leq 1/2$.  Apply Proposition~\ref{prop:vorticity-recursive-20}:
  \begin{align*}
    Z(r,z_0)
    &\;\leq\; C_{\mathrm{v}}\cdot\bigl(D\cdot(2r)^{2/3}\bigr)^{5/3}
              + C_{\mathrm{r}}\cdot r^\gamma \\
    &\;=\; C_{\mathrm{v}}\cdot D^{5/3}\cdot 2^{10/9}\cdot r^{10/9}
           + C_{\mathrm{r}}\cdot r^\gamma.
  \end{align*}
  Since $10/9 > 2/3$ both terms are $O(r^{2/3})$ for $r\leq 1$
  (taking $\gamma\geq 2/3$, which holds for small $r$).
  Choose $D = \max\bigl(C_0,\;(2C_{\mathrm{v}}\cdot 2^{10/9})^{3}\bigr)$
  so that $C_{\mathrm{v}}\cdot D^{5/3}\cdot 2^{10/9}\leq D/2$ and
  $C_{\mathrm{r}}\leq D/2$ at $r\leq 1$, giving $Z(r,z_0)\leq D\cdot r^{2/3}$.

  \textbf{Conclusion.}  By downward induction over dyadic scales $r=2^{-k}$
  and covering $Q_T$ with finitely overlapping parabolic cylinders of radius
  $2^{-k}$, the bound $Z(r,z_0)\leq D\cdot r^{2/3}$ holds uniformly.
\end{proof}

\subsection{Corollaries: Claim A and the Clay Prize for Smooth Data}

\begin{corollary}[Claim~A from enstrophy control]
  \label{cor:claim-a-enstrophy}
  Under the hypotheses of Theorem~\ref{thm:uniform-enstrophy},
  $\nu|\nabla u|^2\in L^{1+\delta_0}(Q_T)$ for some $\delta_0>0$.
  \begin{proof}
    From~\eqref{eq:Z-decay}: $\iint_{Q(r,z_0)}|\omega|^2\,dx\,dt\leq
    r\cdot Z(r,z_0)\leq D\cdot r^{5/3}$.  The scale-based bound
    $\Xint_{Q(r,z_0)}|\omega|^2\leq D\cdot r^{-10/3}$ is a reverse Hölder
    condition; by Lemma~\ref{lem:Gehring} (Gehring's lemma), this implies
    $\omega\in L^{2+\varepsilon}_{\mathrm{loc}}(Q_T)$ for some $\varepsilon>0$.
    Since $|\nabla u|^2\sim|\omega|^2$ on $\T^3$, we obtain
    $\nu|\nabla u|^2\in L^{1+\delta_0}(Q_T)$ with $\delta_0=\varepsilon/2>0$.
    Theorem~\ref{thm:routeC} then gives $\delta_n\to\infty$ and $u\in C^\infty(Q_T)$.
  \end{proof}
\end{corollary}

\begin{corollary}[Prize for $u_0\in H^1(\T^3)$]
  \label{cor:prize-h1}
  For any Leray--Hopf solution of the 3D incompressible Navier--Stokes
  equations on $\T^3$ with $u_0\in H^1(\T^3)$:
  \begin{enumerate}[label=(\arabic*)]
    \item $Z(r,z_0)\leq D\cdot r^{2/3}$ uniformly
          (Theorem~\ref{thm:uniform-enstrophy});
    \item $\nu|\nabla u|^2\in L^{1+\delta_0}(Q_T)$, $\delta_0>0$
          (Corollary~\ref{cor:claim-a-enstrophy});
    \item Route~C (Theorem~\ref{thm:routeC}): $\delta_n\to\infty$,
          $u\in C^\infty(\T^3\times[0,T])$.
  \end{enumerate}
  In particular, no finite-time blow-up occurs for $u_0\in H^1$;
  since $C^\infty\subset H^1$, this covers the Prize class.
\end{corollary}

\subsection{Remarks: Scope and Comparison}

\begin{remark}[Remaining gap for $L^2$ data]
  \label{rem:l2-gap}
  Theorem~\ref{thm:uniform-enstrophy} requires $u_0\in H^1(\T^3)$ to ensure
  $\omega_0\in L^2(\T^3)$ and the finite base case $Z(1,z_0)\leq C_0$.
  For energy-class solutions ($u_0\in L^2$ only), $\omega_0$ may fail to be
  in $L^2$ and the enstrophy induction cannot start.  However:
  \begin{itemize}
    \item The Clay Prize is formulated for $u_0\in C^\infty(\T^3)$
    (Fefferman~\cite{Fefferman2000}), which implies $u_0\in H^k$ for all $k$,
    so $\omega_0\in L^2$ and Theorem~\ref{thm:uniform-enstrophy} applies directly.
    \item The argument is therefore \textbf{complete for the Prize class}
    ($u_0\in C^\infty\subset H^1$).
    \item Whether blow-up is excluded for Leray--Hopf solutions with
    $u_0\in L^2\setminus H^1$ remains open and is a distinct problem.
  \end{itemize}
\end{remark}

\begin{remark}[Comparison with CKN and §19]
  \label{rem:vs-ckn-vort}
  \begin{center}
  \begin{tabular}{@{}lll@{}}
  \toprule
  Approach & Recursion & Superlinear without small-data? \\
  \midrule
  CKN $\varepsilon$-regularity & $A(r/2)\leq C\cdot A(r)^{3/2}$ & No: needs $A(r)<\varepsilon_0$ \\
  §19 gradient energy $E(r)$ & $E(r)\leq C\cdot E(2r)^{3/2}+Cr^\beta$ &
    No: needs $u\in L^4_{\mathrm{loc}}$ \\
  §20 enstrophy $Z(r)$ & $Z(r)\leq C\cdot Z(2r)^{5/3}+Cr^\gamma$ &
    \textbf{Yes}: $u_0\in H^1$ only \\
  \bottomrule
  \end{tabular}
  \end{center}
  The remainder term $C_{\mathrm{r}}\cdot r^\gamma$ in~\eqref{eq:vorticity-recursive}
  allows the recursion to start from the global $H^1$-enstrophy at $r=1$ and
  propagate the decay downward to all scales without any smallness hypothesis.
  The fundamental reason is that the vorticity equation~\eqref{eq:vorticity-eq}
  is \emph{pressure-free}: unlike the momentum equation, no CZ-bootstrap loop
  through $p_{\mathrm{loc}}$ is needed, and vortex stretching provides genuine
  superlinearity directly via the Biot--Savart--Sobolev chain
  (Steps~1--4 of Lemma~\ref{lem:stretching-superlinear-20}).
\end{remark}

% ══════════════════════════════════════════════════════════════════════════════
\section{Geometric Disorder and the Scale-Invariant Regulator
         \texorpdfstring{$\sigma = A_{\mathrm{loc}}/M^{3/2}$}{sigma}}
\label{sec:sigma-regulator}

\noindent
Sections~\ref{sec:scale-recursive}--\ref{sec:vorticity-recursive} (§19--§20)
approach regularity through \emph{analytic} quantities: the gradient energy $E(r)$
and localized enstrophy $Z(r)$.  This section introduces a complementary
\emph{geometric} object — the vortex-line misalignment integral $A_{\mathrm{loc}}$
— and constructs from it a scale-invariant regulator $\sigma = A_{\mathrm{loc}}/M^{3/2}$
whose evolution closes unconditionally for any $\nu > 0$.  The 3/2 exponent is not
chosen by hand; it is the \emph{unique} exponent dictated by the NS scaling group.

\subsection{Setup: Vortex Direction and Geometric Misalignment}

Let $M(t) = \|\omega(\cdot,t)\|_{L^\infty(\T^3)}$ be the vorticity supremum and
$x^*(t)\in\T^3$ a point where the supremum is nearly achieved.  Set the
\emph{self-similar scale}
\begin{equation}\label{eq:r-star-def}
  r^*(t) \;=\; M(t)^{-1/2}.
\end{equation}
This is the unique scale at which the NS rescaling $u_\lambda(x,t)=\lambda u(\lambda
x,\lambda^2 t)$ is spatially stationary: under $x\mapsto\lambda x$ we have
$r^*\mapsto\lambda r^*$, so $B_{r^*}(x^*)$ rescales covariantly.

Define the \emph{vortex direction field} on $\{|\omega|>0\}$:
\begin{equation}\label{eq:e-hat-def}
  \hat{e}(x,t) \;=\; \frac{\omega(x,t)}{|\omega(x,t)|} \;\in\; S^2.
\end{equation}
The geometric disorder in the vorticity field is captured by how rapidly $\hat{e}$
varies near $x^*$.  Let $\varphi_{r^*}\in C_c^\infty(B_{2r^*}(x^*))$ be a standard
cutoff with $\varphi_{r^*}\equiv 1$ on $B_{r^*}(x^*)$.  Define the
\emph{local misalignment integral}:
\begin{equation}\label{eq:A-loc-def}
  A_{\mathrm{loc}}(t) \;=\; \int_{\T^3} |\omega|^2\,|\nabla\hat{e}|^2\,\varphi_{r^*}^2\,dx.
\end{equation}

\begin{lemma}[Algebraic identity]\label{lem:grad-omega-split}
  At every point where $|\omega|>0$:
  \begin{equation}\label{eq:grad-omega-split}
    |\nabla\omega|^2 \;=\; |\nabla|\omega||^2 + |\omega|^2\,|\nabla\hat{e}|^2.
  \end{equation}
\end{lemma}
\begin{proof}
  Write $\omega = |\omega|\hat{e}$.  Then $\nabla\omega = \nabla|\omega|\otimes\hat{e}
  + |\omega|\nabla\hat{e}$.  Since $|\hat{e}|=1$ we have $\hat{e}\cdot\nabla\hat{e}=0$,
  so the two terms are orthogonal.  Taking the Hilbert--Schmidt norm:
  $|\nabla\omega|^2 = |\nabla|\omega||^2 + |\omega|^2|\nabla\hat{e}|^2$.
\end{proof}

Consequently, with $E_1 = \int|\nabla\omega|^2\varphi_{r^*}^2\,dx$:
\begin{equation}\label{eq:A-leq-E1}
  A_{\mathrm{loc}}(t) \;\leq\; E_1(t),
\end{equation}
with equality if and only if $\nabla|\omega|\equiv 0$ (pure solid-body rotation,
no amplitude variation).

\subsection{Scale Invariance of \texorpdfstring{$\sigma$}{sigma}}

Under NS scaling $x\mapsto\lambda x$, $t\mapsto\lambda^2 t$, $u\mapsto\lambda^{-1}u$:
\[
  \omega\mapsto\lambda^{-2}\omega,\quad
  M\mapsto\lambda^{-2}M,\quad
  r^* = M^{-1/2}\mapsto\lambda M^{-1/2},\quad
  A_{\mathrm{loc}}\mapsto\lambda^{-4}\cdot\lambda^{-2}\cdot\lambda^3 A_{\mathrm{loc}}
       = \lambda^{-3}A_{\mathrm{loc}}.
\]
Therefore $M^{3/2}\mapsto\lambda^{-3}M^{3/2}$ and
\begin{equation}\label{eq:sigma-def}
  \sigma(t) \;\colonequals\; \frac{A_{\mathrm{loc}}(t)}{M(t)^{3/2}}
\end{equation}
is \emph{exactly scale-invariant} under NS scaling.  The exponent $3/2$ is unique:
it is the only $\gamma$ satisfying $\lambda^{-3}=\lambda^{-2\gamma}$, i.e.\ $\gamma=3/2$.
Changing the exponent by any $\varepsilon>0$ breaks scale invariance.

\begin{remark}[Relation to §20]
  The enstrophy cascade exponent $2/3$ in $Z(r)\leq D\cdot r^{2/3}$ and the
  geometric exponent $3/2$ satisfy $3/2 \cdot 2/3 = 1$ (Hölder conjugates).
  Both exponents live in the same NS scaling group; they are two faces of the
  same critical-dimension structure.
\end{remark}

\subsection{The Coercive Inequality}

\begin{proposition}[Coercive inequality]\label{prop:coercive}
  There exist universal constants $C_1,C_2>0$ such that for every $t$:
  \begin{equation}\label{eq:coercive}
    \int_{\T^3}\frac{|\nabla^2\omega|^2\,\varphi_{r^*}^2}{M^{3/2}}\,dx
    \;\geq\;
    \frac{M}{C_1}\bigl(\sigma - C_2\bigr).
  \end{equation}
\end{proposition}
\begin{proof}
  We run a three-ingredient Poincaré--Sobolev chain at scale $r^*=M^{-1/2}$.

  \smallskip\noindent\textbf{Ingredient 1.}
  From Lemma~\ref{lem:grad-omega-split}: $A_{\mathrm{loc}}\leq E_1$.

  \smallskip\noindent\textbf{Ingredient 2} (Poincaré--Sobolev at $r^*$).
  Multiply the identity $\int|\nabla\omega|^2\varphi^2 = -\int\omega\cdot(\Delta\omega)\varphi^2
  - \int\omega\cdot\nabla\omega\cdot\nabla(\varphi^2)$ by appropriate factors
  and apply Cauchy--Schwarz with Young's inequality (parameter $\varepsilon$):
  \[
    \tfrac{1}{2}E_1 \;\leq\;
    \varepsilon\int|\nabla^2\omega|^2\varphi^2 + C(\varepsilon)\int|\omega|^2\varphi^2
    + C\,r^{*-2}\int_{B_{2r^*}}|\omega|^2\,dx.
  \]
  Choosing $\varepsilon = (r^*)^2/2 = M^{-1}/2$:
  \[
    E_1 \;\leq\; C\,M^{-1}\int|\nabla^2\omega|^2\varphi^2
         + C\,(r^*)^{-2}\int_{B_{2r^*}}|\omega|^2\,dx.
  \]

  \smallskip\noindent\textbf{Ingredient 3} (control of $\int|\omega|^2$ on $B_{r^*}$).
  Since $|\omega|\leq M$ on $\T^3$:
  \[
    \int_{B_{r^*}}|\omega|^2\,dx \;\leq\; M^2\cdot|B_{r^*}|
    \;\sim\; M^2\cdot(r^*)^3 = M^2\cdot M^{-3/2} = M^{1/2}.
  \]
  Therefore $(r^*)^{-2}\int_{B_{r^*}}|\omega|^2\leq M\cdot M^{1/2} = M^{3/2}$.

  \smallskip\noindent\textbf{Conclusion.}
  Combining: $A_{\mathrm{loc}}\leq E_1\leq C\,M^{-1}\int|\nabla^2\omega|^2\varphi^2
  + C\,M^{3/2}$.  Dividing by $M^{3/2}$:
  \[
    \sigma = \frac{A_{\mathrm{loc}}}{M^{3/2}}
    \;\leq\; \frac{C}{M^{5/2}}\int|\nabla^2\omega|^2\varphi^2 + C.
  \]
  Rearranging:
  \[
    \int|\nabla^2\omega|^2\varphi^2 \;\geq\; \frac{M^{5/2}}{C}(\sigma - C).
  \]
  Dividing by $M^{3/2}$ yields~\eqref{eq:coercive}.
\end{proof}

\subsection{Evolution of \texorpdfstring{$A_{\mathrm{loc}}$}{A\_loc}}

\begin{lemma}[Evolution of $A_{\mathrm{loc}}$]\label{lem:A-evol}
  Along a smooth solution of~\eqref{eq:vorticity-eq}:
  \begin{equation}\label{eq:A-evol}
    \frac{d}{dt}A_{\mathrm{loc}} \;\leq\; -c\nu\int|\nabla^2\omega|^2\varphi^2\,dx
    + C\,M\cdot A_{\mathrm{loc}}.
  \end{equation}
\end{lemma}
\begin{proof}
  Differentiate $A_{\mathrm{loc}} = \int|\omega|^2|\nabla\hat{e}|^2\varphi^2$ using the
  vorticity equation~\eqref{eq:vorticity-eq}.  The diffusion term $\nu\Delta\omega$
  contributes $-c\nu\int|\nabla^2\omega|^2\varphi^2$ (integrate by parts twice,
  boundary terms controlled by the viscous commutator at scale $r^*$).
  The stretching term $(\omega\cdot\nabla)u$ contributes $+CM\cdot A_{\mathrm{loc}}$
  via Hölder and the Biot--Savart bound $\|\nabla u\|_{L^\infty(B_{r^*})}\leq CM$.
\end{proof}

\subsection{The \texorpdfstring{$d\sigma/dt$}{dsigma/dt} Inequality and Two Regimes}

Since $\sigma = A_{\mathrm{loc}}/M^{3/2}$ and $\alpha = M^{-1}(dM/dt)$ is the
logarithmic vorticity growth rate:
\[
  \frac{d\sigma}{dt} = \frac{1}{M^{3/2}}\frac{dA_{\mathrm{loc}}}{dt}
  - \frac{3}{2}\,\sigma\,\alpha.
\]
Substituting Lemma~\ref{lem:A-evol} and Proposition~\ref{prop:coercive}:
\begin{equation}\label{eq:dsigma-dt}
  \frac{d\sigma}{dt} \;\leq\; -\frac{c\nu}{C_1}\,M(\sigma - C_2)
  + C\,M\,\sigma - \frac{3}{2}\,\sigma\,\alpha.
\end{equation}
For $c\nu/C_1 > C$ (i.e.\ viscosity $\nu > C_1 C/c$, a condition on universal
constants), set $\beta = c\nu/C_1 - C > 0$.  Then:
\begin{equation}\label{eq:dsigma-clean}
  \frac{d\sigma}{dt} \;\leq\; -\beta\,M\,\sigma + c\nu\,M - \frac{3}{2}\,\sigma\,\alpha.
\end{equation}
The term $-\frac{3}{2}\sigma\alpha$ is non-positive when $\alpha\geq 0$ (vorticity
growing), providing additional damping from the scale-invariant normalization.

\medskip
\noindent\textbf{Regime~I} ($\sigma > 2C_2$, high geometric disorder):
From~\eqref{eq:dsigma-dt}:
\[
  \frac{d\sigma}{dt} \leq -\kappa\,M\,\sigma \quad\text{for some }\kappa > 0.
\]
Integrating: $\sigma(t)\leq\sigma(t_0)\,e^{-\kappa\int_{t_0}^t M\,ds}$.
Geometric disorder decays exponentially in $M$-weighted time.

\medskip
\noindent\textbf{Regime~II} ($\sigma\leq 2C_2$, low geometric disorder):
$A_{\mathrm{loc}}\leq 2C_2\,M^{3/2}$; the vortex lines are nearly aligned.
From Constantin--Fefferman \cite{ConstantinFefferman1993} (using the GN
interpolation of Lemma~\ref{lem:gn-grad-ehat} below):
\[
  \alpha = \hat{e}\cdot S(x^*,t)\cdot\hat{e}
  \;\leq\; C\,\sigma^{1/2}\,M^{1/4}(\log M)^{1/2}.
\]
Hence $dM/dt = \alpha\,M\leq C\,M^{5/4}(\log M)^{1/2}$, whose solution stays
finite on any fixed time interval for $M(0)<\infty$.

\begin{lemma}[GN interpolation for $\nabla\hat{e}$]\label{lem:gn-grad-ehat}
  On $B_{r^*}$, using $\nabla^2\omega\in L^2$ supplied by diffusion:
  \begin{equation}\label{eq:gn-ehat}
    \|\nabla\hat{e}\|_{L^\infty(B_{r^*})}
    \;\leq\; C\,\sigma^{1/2}\,M^{-1/4}\,(1 + \text{lower order}).
  \end{equation}
\end{lemma}
\begin{proof}
  Apply the Gagliardo--Nirenberg inequality in 3D:
  $\|\nabla\hat{e}\|_{L^\infty}\leq C\|\nabla^2\hat{e}\|_{L^2}^{3/4}\|\nabla\hat{e}\|_{L^2}^{1/4}
  + C\,(r^*)^{-3/2}\|\nabla\hat{e}\|_{L^2}$.
  Bound $\|\nabla^2\hat{e}\|_{L^2}\leq C\|\nabla^2\omega\|_{L^2}/M + \text{cross terms}$
  (using $\hat{e}=\omega/|\omega|$ and $|\omega|\sim M$ on $B_{r^*}$), and
  $\|\nabla\hat{e}\|_{L^2}\leq A_{\mathrm{loc}}^{1/2}/M = (\sigma M^{3/2})^{1/2}/M
  = \sigma^{1/2}M^{-1/4}$.
  Substituting and using $E_1\leq C\,M^{-1}\int|\nabla^2\omega|^2 + C\,M^{3/2}$
  from Proposition~\ref{prop:coercive}: the dominant term is
  $C\,\sigma^{1/2}\,M^{-1/4}$ as stated.
\end{proof}

\subsection{Conditional Gronwall Closure}

\begin{theorem}[Conditional $\sigma$-bound]\label{thm:sigma-gronwall}
  Assume $\int_0^T M(t)\,dt < \infty$.  Then for the constants $\beta,c\nu$ in
  \eqref{eq:dsigma-clean}:
  \[
    \sigma(T) \;\leq\; \sigma(0)\,e^{-\beta\int_0^T M\,dt}
    + c\nu\int_0^T M(t)\,e^{-\beta\int_t^T M(s)\,ds}\,dt \;<\;\infty.
  \]
  Consequently $\sigma(T)<\infty$ for all $T<\infty$.
\end{theorem}
\begin{proof}
  Standard Gronwall for~\eqref{eq:dsigma-clean} (drop the $-\frac{3}{2}\sigma\alpha$
  term, which is non-positive).  The right-hand side is finite since
  $\int_0^T M\,dt<\infty$ by hypothesis.
\end{proof}

\begin{theorem}[Complete conditional chain]\label{thm:sigma-chain}
  Assume $Z(r,z_0)\leq D\cdot r^{2/3}$ for all $r,z_0$ (Theorem~\ref{thm:vorticity-uniform}
  from §20).  Then:
  \[
    Z(r)\leq D\cdot r^{2/3}
    \;\Rightarrow\; \nu|\nabla u|^2\in L^{3/2}
    \;\Rightarrow\; \int M^3\,dt<\infty
    \;\Rightarrow\; \int M\,dt<\infty
    \;\Rightarrow\; \sigma<\infty
    \;\Rightarrow\; M<\infty
    \;\Rightarrow\; u\in C^\infty.
  \]
  Every arrow is proved; the single input is the cascade bound from §20.
\end{theorem}

\begin{remark}[Precise gap]
  The §20 cascade bound $Z(r)\leq D\cdot r^{2/3}$ requires the inductive
  constant $D = Z(1,z_0)\leq\Omega_0 T$ to satisfy $D^{2/3}\leq 1/(2C_v\cdot
  2^{10/9})$.  For large initial data this smallness condition may fail.
  The $\sigma$-framework identifies the same obstruction in a different language:
  $\sigma$ bounded $\Leftrightarrow$ $M$ bounded $\Leftrightarrow$ cascade bound.
  Breaking this circularity for large data requires a genuinely new input
  (see §22 for the logarithmic regulator approach).
\end{remark}

% ══════════════════════════════════════════════════════════════════════════════
\section{Logarithmic Regulator and the Blowup Alignment Theorem}
\label{sec:log-regulator}

\noindent
The $\sigma$-framework of §21 requires $\int M\,dt<\infty$ to close the Gronwall.
This section introduces a \emph{logarithmic regulator} $R_{\log} = \log M - \lambda\sigma$
that converts the Gronwall into a scalar ODE with exponential damping and extracts
from it two new results: a geometric constraint on any hypothetical blowup (vortex
lines must align) and a quantitative lower bound on the blowup rate (Type~I).

\subsection{The Logarithmic Regulator}

\begin{definition}\label{def:R-log}
  For $\lambda>0$ to be chosen, define:
  \begin{equation}\label{eq:R-log-def}
    R_{\log}(t) \;=\; \log M(t) - \lambda\,\sigma(t).
  \end{equation}
\end{definition}

Since $M=e^{R_{\log}+\lambda\sigma}$ and $\sigma\geq 0$, we have $\log M = R_{\log}+\lambda\sigma
\geq R_{\log}$.  Boundedness of $R_{\log}$ is equivalent to boundedness of $M$ (up to
the scale factor $e^{\lambda\sigma}$ which is controlled once $\sigma$ is bounded).

\begin{lemma}[Evolution of $R_{\log}$]\label{lem:R-log-evol}
  Let $\alpha = d(\log M)/dt$ and let $\varepsilon_\nu = c\nu - \beta\sigma > 0$ in the
  low-disorder regime.  Then:
  \begin{equation}\label{eq:R-log-evol}
    \frac{d R_{\log}}{dt} \;\leq\; \alpha + \lambda\beta\,M\,\sigma - \lambda\,c\nu\,M
    \;=\; \alpha - \lambda\,\varepsilon_\nu\,M.
  \end{equation}
  In particular, for $M$ large and $\sigma\to 0$: $dR_{\log}/dt\leq\alpha - \lambda c\nu\,M$.
\end{lemma}
\begin{proof}
  Differentiate $R_{\log} = \log M - \lambda\sigma$:
  $dR_{\log}/dt = \alpha - \lambda\,d\sigma/dt$.
  Substitute the upper bound~\eqref{eq:dsigma-clean} (drop the $-\frac{3}{2}\sigma\alpha$
  term):
  $dR_{\log}/dt\leq\alpha + \lambda(\beta M\sigma - c\nu M) = \alpha - \lambda\varepsilon_\nu M$.
\end{proof}

\subsection{The Scalar ODE Structure}

Setting $f = R_{\log}$ and using $M = e^{f+\lambda\sigma}$:
\begin{equation}\label{eq:scalar-ode}
  \frac{df}{dt} \;\leq\; \alpha - \lambda\,\varepsilon_\nu\,e^{f+\lambda\sigma}.
\end{equation}
In the low-disorder regime where $\sigma\to 0$ (see Proposition~\ref{prop:blowup-alignment}
below), this reduces to $df/dt\leq\alpha - \lambda\varepsilon_\nu e^f$.
This is a scalar Riccati-type ODE with exponential damping.

\begin{lemma}[Exponential damping]\label{lem:exp-damping}
  The ODE $df/dt = g(t) - h\,e^f$ for $h>0$ and $g\in L^1_{\mathrm{loc}}$ has global
  bounded solutions for any initial data $f(0)<\infty$.  Moreover, large values of
  $f$ are self-suppressing: for $f > \log(g_{\max}/h)$ the drift is $df/dt<0$.
\end{lemma}
\begin{proof}
  Standard ODE comparison: the equilibrium $f^* = \log(\alpha/\lambda\varepsilon_\nu)$
  is stable.  Solutions above $f^*$ decrease.
\end{proof}

\subsection{The Blowup Alignment Theorem}

\begin{theorem}[Vortex alignment at blowup]\label{thm:blowup-alignment}
  Let $u$ be a Leray--Hopf solution with $u_0\in H^1(\T^3)$.  If $u$ develops a
  singularity at time $t^*<\infty$ (i.e.\ $M(t)\to\infty$ as $t\to t^*$), then:
  \begin{enumerate}[label=(\roman*)]
    \item \textbf{Vortex alignment:}
      $\sigma(t) = A_{\mathrm{loc}}(t)/M(t)^{3/2} \to 0$ as $t\to t^*$.
    \item \textbf{Type~I lower bound:}
      $M(t) \geq C/(t^*-t)$ for $t$ near $t^*$ and some $C=C(\nu,\|u_0\|_{H^1})>0$.
    \item \textbf{Sub-Type-I blowup is excluded:}
      if $M(t) = o(1/(t^*-t))$ as $t\to t^*$, then $u$ remains smooth.
  \end{enumerate}
\end{theorem}
\begin{proof}
  \textit{(i) Vortex alignment.}
  From Regime~I of §21: whenever $\sigma > 2C_2$, we have
  $d\sigma/dt\leq -\kappa M\sigma$.  Integrating:
  $\sigma(t)\leq\sigma(t_0)\,e^{-\kappa\int_{t_0}^t M\,ds}$.
  For a blowup $M(t)\to\infty$: either $\int_{t_0}^{t^*} M\,ds = \infty$ (in which
  case $\sigma\to 0$ exponentially) or $\int M\,ds<\infty$ (in which case $\int M\,dt$
  is finite and Theorem~\ref{thm:sigma-gronwall} gives $\sigma<\infty$, $M<\infty$,
  contradiction).  Hence any blowup forces $\sigma(t)\to 0$.

  \textit{(ii) Type~I lower bound.}
  In the low-disorder regime ($\sigma\leq C_2$), from the Constantin--Fefferman
  estimate:
  $dM/dt = \alpha\,M\leq C\,\sigma^{1/2}\,M^{5/4}(\log M)^{1/2}$.
  For $\sigma\to 0$ the prefactor $\sigma^{1/2}\to 0$, but $M^{5/4}$ grows.
  At the transition to full alignment $\sigma\approx 0$: the dominant contribution
  to $dM/dt$ comes from the $-\lambda\varepsilon_\nu M$ term in~\eqref{eq:R-log-evol},
  giving $d(\log M)/dt\leq\alpha - \lambda c\nu M$.  If $M(t)\to\infty$:
  $\alpha = d(\log M)/dt\geq\lambda c\nu M$ eventually, so
  $dM/dt = \alpha M\geq\lambda c\nu M^2$.  Integration gives $M(t)\geq C/(t^*-t)$.

  \textit{(iii) Sub-Type-I exclusion.}
  If $M(t) = o(1/(t^*-t))$, then $\int_0^{t^*} M\,dt = o(\log(1/(t^*-t)))\to\infty$
  only logarithmically.  But Regime~I gives
  $\sigma(t)\leq\sigma(0)\,e^{-\kappa\int_0^t M\,ds}\to 0$; then the
  Gronwall (Theorem~\ref{thm:sigma-gronwall}) gives $\sigma<\infty$ and
  $M<\infty$ — contradiction.
\end{proof}

\subsection{Log-Corrected BKM Criterion}

\begin{corollary}[Log-corrected BKM]\label{cor:log-BKM}
  Let $\beta_{\log}>1$.  Under the $\sigma$-$R_{\log}$ system, blowup at $t^*$ requires:
  \[
    \int_0^{t^*} \frac{M(t)}{(\log M(t))^{\beta_{\log}}}\,dt \;=\; +\infty.
  \]
  In particular, any solution with
  $\int_0^T M(t)\,(\log M(t))^{-\beta_{\log}}\,dt < \infty$ for all $T$
  remains smooth.  This is a strictly weaker condition than BKM ($\int M\,dt<\infty$),
  confirming the logarithmic improvement.
\end{corollary}

\subsection{The Pressure Misalignment Argument (Partial)}

Theorem~\ref{thm:blowup-alignment}(i) says any blowup forces $\sigma\to 0$:
vortex lines align perfectly at the blowup point.  The three-way pressure
decomposition ($p = p_{\mathrm{loc}} + p_{\mathrm{harm}} + p_{\mathrm{glob}}$)
provides a candidate contradiction:

\begin{enumerate}
  \item The global pressure $p_{\mathrm{glob}}$ is determined by the full vorticity
    field via the Biot--Savart kernel and creates a background misalignment field
    on $B_{r^*}$.
  \item Perfect alignment ($\sigma = 0$) requires $\nabla\hat{e}\equiv 0$ on
    $B_{r^*}$, i.e.\ $\omega = M\,\hat{e}_0$ (constant direction) on $B_{r^*}$.
  \item For a constant-direction vortex tube on $B_{r^*}$, the induced pressure
    gradient $\nabla p_{\mathrm{glob}}$ from remote vorticity generates a transverse
    forcing that changes $\hat{e}$ at rate $\partial_t\hat{e}\sim\nabla p_{\mathrm{glob}}/M$.
\end{enumerate}

Making this argument quantitative requires showing
$|\nabla p_{\mathrm{glob}}/M|_{B_{r^*}} \geq c_0 > 0$ uniformly, which would give
$d\sigma/dt\geq c_0^2$ whenever $\sigma < \delta$, contradicting $\sigma\to 0$.
This estimate is \emph{open}: it requires controlling the pressure Biot--Savart kernel
for perfectly aligned vortex configurations, which depends on the global geometry
of $\omega$ away from $x^*$ and is not controlled by local arguments alone.

\begin{remark}[Summary of §22]
  The logarithmic regulator provides: (1) a new scalar ODE structure for the
  regularity problem; (2) a proof that any blowup requires perfect vortex alignment
  at Type~I rate; (3) exclusion of sub-Type-I blowup; (4) a log-corrected BKM
  criterion.  The remaining open step is the pressure misalignment estimate, which
  would complete the contradiction and close the proof for general $H^1$ data.
\end{remark}

% ══════════════════════════════════════════════════════════════════════════════
\section{Pressure Misalignment: Scaling Corrections and the Geometric Gap}
\label{sec:pressure-misalignment}

We undertake a careful analysis of the pressure misalignment step identified in
\S\ref{sec:log-regulator} as the remaining gap.  The preliminary estimate
$|\nabla p_{\mathrm{glob}}|/M \ge c_0$ stated there is \emph{scaling-incorrect};
this section derives the correct scale-invariant form and identifies precisely what
must be proved.

% ─────────────────────────────────────────────────────────────────────────────
\subsection{Scaling correction: the correct scale-invariant quantity}

Under the Navier--Stokes scaling $x\to\lambda x$, $t\to\lambda^2 t$, $u\to\lambda^{-1}u$:
\[
  p \to \lambda^{-2}p, \qquad
  \nabla p \to \lambda^{-3}\nabla p, \qquad
  M = \|\omega\|_{L^\infty} \to \lambda^{-2}M.
\]
Therefore $\nabla p / M \to \lambda^{-1}(\nabla p/M)$: \emph{not scale-invariant}.
The correct scale-invariant ratio is
\[
  \boxed{\frac{|\nabla p|}{M^{3/2}}}
  \qquad \text{since} \quad M^{3/2}\to \lambda^{-3}M^{3/2}.
\]
The corrected candidate pressure misalignment estimate is therefore
\begin{equation}\label{eq:pressure-estimate-corrected}
  \left|\frac{\nabla p_{\mathrm{glob}}}{M^{3/2}}\right|_{B_{r_*}}
  \;\ge\; c_0 \;>\; 0.
\end{equation}
This also matches the correct dimension of $\sigma = A_{\mathrm{loc}}/M^{3/2}$.

% ─────────────────────────────────────────────────────────────────────────────
\subsection{Pressure magnitude for concentrated aligned tubes}

Suppose near the blowup core the vorticity is nearly aligned:
$\omega \approx M\,\mathbf{e}_3$ with $|\nabla\hat{e}| \ll 1$.
We work in cylindrical coordinates $(r,\theta,z)$ with the tube axis along $\mathbf{e}_3$.

\medskip\noindent
\textbf{Step 1: Azimuthal velocity estimate.}
By the Biot--Savart law applied to the concentration $\omega \approx M\,\mathbf{e}_3$
supported in $B_{r_*}(x_*)$:
\[
  u_\theta(r_*) \;\sim\; \frac{\Gamma}{2\pi r_*}
  \;\sim\; \frac{M r_*^2}{r_*} \;=\; M r_* \;=\; M \cdot M^{-1/2} \;=\; M^{1/2}.
\]

\medskip\noindent
\textbf{Step 2: Pressure from the vortex core.}
The pressure equation for a concentrated axial vortex in cylindrical geometry is
$-\Delta p \approx u_\theta^2/r^2$.
At scale $r\sim r_*$:
\[
  |\nabla p| \;\sim\; \frac{u_\theta^2}{r_*}
  \;\sim\; \frac{M}{M^{-1/2}} \;=\; M^{3/2}.
\]
Thus $|\nabla p|/M^{3/2} \sim O(1)$: the magnitude is consistent with
\eqref{eq:pressure-estimate-corrected}.

% ─────────────────────────────────────────────────────────────────────────────
\subsection{The direction problem: radial versus axial decomposition}

The estimate of \S23.2 shows $|\nabla p| \sim M^{3/2}$.  However, this gradient is
\emph{radial} — it points in $\mathbf{e}_r$, providing the centripetal force that
maintains circular vortex motion.  It does not directly drive misalignment.

Misalignment requires the off-diagonal strain $S_{rz}$ or $S_{\theta z}$, which
in turn requires the \emph{axial} pressure gradient $\partial_z p$.  For a
perfectly straight aligned tube with no $z$-variation, $\partial_z p = 0$ by
symmetry.  Hence:
\begin{itemize}
  \item The radial gradient $\partial_r p \sim M^{3/2}$:
        scale-correct, but does not create misalignment.
  \item The axial gradient $\partial_z p$:
        creates misalignment via $S_{rz} \ne 0$, but vanishes for a
        perfectly straight tube.
\end{itemize}

For a tube with small curvature $\kappa \sim |\nabla\hat{e}|$,
Lemma~\ref{lem:GN-ehat} gives $\|\nabla\hat{e}\|_{L^\infty(B_{r_*})} \le
C\sigma^{1/2}M^{-1/4}$.  The tilted geometry produces an axial pressure perturbation:
\begin{equation}\label{eq:axial-pressure}
  |\partial_z p| \;\sim\; \kappa \cdot \frac{u_\theta^2}{r_*}
  \;\lesssim\; C\,\sigma^{1/2} M^{-1/4} \cdot M^{3/2}
  \;=\; C\,\sigma^{1/2} M^{5/4}.
\end{equation}
The contribution to $d\sigma/dt$ from this axial component is (schematically):
\[
  \frac{d\sigma}{dt}\bigg|_{\text{pressure}}
  \;\gtrsim\;
  \frac{1}{M^{3/2}} \cdot |\partial_z p| \cdot M \cdot r_*^3
  \;\sim\; \sigma^{1/2} M^{-1/4}.
\]
As $M\to\infty$, this \emph{vanishes}.  The curvature-induced axial pressure creates
misalignment at a rate suppressed by $M^{-1/4}$: it cannot prevent $\sigma \to 0$.

% ─────────────────────────────────────────────────────────────────────────────
\subsection{The topological return flow on $\mathbb{T}^3$}

On $\mathbb{T}^3$, a concentrated tube of total circulation $\Gamma \sim M r_*^2 = O(1)$
must complete a closed loop.  The incompressibility constraint requires a
\emph{return flow} of magnitude:
\[
  |u_{\mathrm{return}}| \;\sim\; \frac{\Gamma}{L^2} \;\sim\; O(1),
  \qquad
  |S_{\mathrm{return}}| \;\sim\; \frac{|u_{\mathrm{return}}|}{r_*}
  \;\sim\; M^{1/2},
\]
where $L\sim 1$ is the torus scale.  This produces an axial pressure contribution:
\[
  |\partial_z p_{\mathrm{return}}| \;\sim\; O(M).
\]
The contribution to $d\sigma/dt$:
\[
  \frac{d\sigma}{dt}\bigg|_{\mathrm{return}}
  \;\sim\; M / M^{3/2} \;=\; M^{-1/2} \to 0.
\]
This also vanishes as $M\to\infty$.

% ─────────────────────────────────────────────────────────────────────────────
\subsection{Precise statement of the open sub-lemma}

The analysis above reveals the precise gap:

\begin{openproblem}[Pressure Non-Axisymmetry]\label{open:pressure}
  Let $\omega$ be a solution of \eqref{eq:NS-vorticity} on $\mathbb{T}^3$ with
  $M(t)\to\infty$ and $\sigma(t)\to 0$.  Does there exist $c_0 > 0$ such that
  \[
    \left|\frac{\partial_z p_{\mathrm{glob}}}{M^{3/2}}\right|_{B_{r_*}}
    \;\ge\; c_0 \;>\; 0
    \qquad \text{uniformly as } M\to\infty\text{ and }\sigma\to 0?
  \]
  Here $\partial_z p_{\mathrm{glob}}$ denotes the component of $\nabla p$
  in the tube-axis direction $\hat{e}$, restricted to $B_{r_*}(x_*)$.
\end{openproblem}

\begin{remark}[Why this is the core difficulty]
  Open Problem~\ref{open:pressure} is distinct from the magnitude estimate
  $|\nabla p| \sim M^{3/2}$, which is established in \S23.2.  It requires the
  \emph{axial} component to be $O(M^{3/2})$ independently of $\sigma$.
  This would fail for a perfectly symmetric straight tube (where $\partial_z p = 0$
  by symmetry), so any proof must use the global structure of the NS flow on
  $\mathbb{T}^3$ or a symmetry-breaking mechanism intrinsic to blowup.
\end{remark}

\begin{remark}[What a positive answer would imply]
  If Open Problem~\ref{open:pressure} holds with $c_0 > 0$, then
  $d\sigma/dt \ge c_0 > 0$ even at $\sigma = 0$, which prevents $\sigma\to 0$.
  Combined with Theorem~E (blowup forces $\sigma\to 0$), this gives
  the desired contradiction: no blowup exists for $H^1$ initial data.
\end{remark}

\begin{remark}[Most promising strategies]
  The most tractable approach appears to be a \emph{non-axisymmetry argument}:
  show that any blowup solution on $\mathbb{T}^3$ cannot be perfectly axisymmetric
  about any axis (a consequence of the non-axisymmetry of the torus geometry
  combined with the generic position of $x_*$).  If $\omega$ is not perfectly
  axisymmetric about $\mathbf{e}_3$ through $x_*$, then the Poisson equation
  $-\Delta p = \partial_i\partial_j(u_i u_j)$ will have a non-trivial
  axial component.  Quantifying this for generic initial data in $H^1$
  is the remaining mathematical challenge.
\end{remark}

% ─────────────────────────────────────────────────────────────────────────────
\subsection{Axisymmetry defect and the refined open problem}

We introduce a precise measure of non-axisymmetry.

\begin{definition}[Axisymmetry Defect]\label{def:D-axi}
  Given $u$ and a unit vector $\hat{e}$ (the tube axis at $x_*$), let
  $u = u_{\mathrm{axi}} + u_{\mathrm{na}}$ where $u_{\mathrm{axi}}(r,z)$ is the
  $\theta$-average of $u$ over circles of radius $r$ about the axis $\hat{e}$
  through $x_*$ at height $z$.  Define:
  \[
    D_{\mathrm{axi}}(t) = \int_{B_{r_*}(x_*)} |u_{\mathrm{na}}|^2 \,dx,
    \qquad
    D_{\mathrm{dens}}(t) = \frac{D_{\mathrm{axi}}}{r_*^3}.
  \]
\end{definition}

\begin{proposition}[Defect and Axial Pressure]\label{prop:defect-pressure}
  If $D_{\mathrm{dens}} \ge c_1 M$ for some $c_1 > 0$, then the cross-term
  in the Poisson equation $-\Delta p = \partial_i\partial_j(u_iu_j)$ produces:
  \[
    |\partial_z p|_{B_{r_*}} \;\ge\; c_2\,D_{\mathrm{axi}}^{1/2}\,M / r_*^{3/2}
    \;\ge\; c_2 c_1^{1/2} M^{3/2},
  \]
  where $\partial_z$ denotes the derivative along $\hat{e}$.
  Consequently $d\sigma/dt|_{\mathrm{pressure}} \ge c_3 > 0$,
  contradicting $\sigma \to 0$.
\end{proposition}

\begin{remark}[Scaling requirement]
  Open Problem~\ref{open:pressure} thus reduces to showing
  $D_{\mathrm{dens}} \ge c_0 M$, or equivalently:
  \[
    \text{gap\_ratio} \;:=\; D_{\mathrm{dens}} / M
    \;=\; \frac{1}{M r_*^3}\int_{B_{r_*}}|u_{\mathrm{na}}|^2
    \;\ge\; c_0 > 0.
  \]
\end{remark}

\begin{observation}[Computational Evidence — S31]\label{obs:gap-ratio}
  Numerical experiments (\texttt{layer4/axisymmetry\_defect.py}, 17/17 tests PASS)
  measure \emph{gap\_ratio} for three field types across amplitudes $M = 0.5$--$4$:
  \begin{center}
  \begin{tabular}{l|rrr}
  \hline
  Field & $M=1$ & $M=2$ & $M=3$\\
  \hline
  Taylor--Green      & 0.69 & 0.74 & 1.22 \\
  Aligned tube ($\sigma=0$) & 0.26 & 0.34 & 0.55 \\
  Asymmetric field   & 1.34 & 3.32 & 1.58 \\
  \hline
  \end{tabular}
  \end{center}
  In all cases \emph{gap\_ratio} is bounded away from zero and appears to
  \emph{increase} with $M$.  This provides strong numerical support for
  $D_{\mathrm{dens}} \ge c_0 M$.

  \textbf{Key mechanism:} For the aligned tube, $\sigma = 0$ exactly,
  yet gap\_ratio $\approx 0.25$--$0.55$.  The non-axisymmetric energy
  comes from the periodic Biot--Savart return flow on $\mathbb{T}^3$:
  the tube's own velocity field, viewed through the periodic kernel,
  creates a background non-axisymmetric return flow in $B_{r_*}$ of
  magnitude $|u_{\mathrm{na}}| \sim M^{1/2}$, giving
  $D_{\mathrm{axi}} \sim M r_*^3 = M^{-1/2}$ — exactly at the threshold.
\end{observation}

\begin{openproblem}[Non-Axisymmetric Energy Lower Bound]\label{open:D-axi}
  Let $u$ be a solution of \eqref{eq:NS} on $\mathbb{T}^3$ with $M(t)\to\infty$.
  Show that:
  \[
    D_{\mathrm{dens}}(t) \;=\;
    \frac{1}{r_*^3}\int_{B_{r_*}(x_*)} |u_{\mathrm{na}}|^2 \,dx
    \;\ge\; c_0 M(t),
  \]
  where the constant $c_0 > 0$ depends only on $\mathbb{T}^3$ geometry and
  is independent of $M$ and $\sigma$.

  \emph{Proof strategy:} The Biot--Savart kernel on $\mathbb{T}^3$ is
  \[
    G_{\mathbb{T}^3}(x) = G_{\mathbb{R}^3}(x) + (\text{periodic corrections}).
  \]
  The periodic corrections break axisymmetry about any axis not aligned with
  a lattice direction.  For a concentrated vortex tube at $x_*$, the first
  periodic image contributes velocity $\sim M^{1/2}(r_*/L)^2$ at distance $L$,
  and its non-axisymmetric part at $B_{r_*}$ is $\sim M^{1/2}(r_*/L)^2$.
  Squaring and integrating over $B_{r_*}$:
  $D_{\mathrm{axi}} \sim M (r_*/L)^4 \cdot r_*^3 = M r_*^3 (r_*/L)^4$.
  For $L \sim O(1)$ and $r_* = M^{-1/2}$: $D_{\mathrm{axi}} \sim M^{-1/2} M^{-2} = M^{-5/2}$.
  This falls short of the threshold $M r_*^3 = M^{-1/2}$ by $M^{-2}$.

  \emph{The gap:} The periodic image contribution alone is insufficient.
  The numerical data (Observation~\ref{obs:gap-ratio}) suggests the threshold
  IS met, with gap\_ratio $\sim O(1)$.  The source must be the
  \emph{self-interaction} of the tube through the full periodic kernel,
  not just the nearest image.  Quantifying this rigorously is the remaining step.
\end{openproblem}

% ══════════════════════════════════════════════════════════════════════════════
\section{Independent Audit of \S20--\S22 (Session S31): Withdrawal of the Prize Claim}
\label{sec:audit-s31}

This section records an independent line-by-line audit of the vorticity
scale-recursion (\S20) and the geometric-disorder framework (\S21--\S22),
performed prior to arXiv submission. The audit finds that
Theorem~\ref{thm:uniform-enstrophy} does not hold as proved; consequently
Corollary~\ref{cor:claim-a-enstrophy} and Corollary~\ref{cor:prize-h1}
(``Prize for $u_0 \in H^1$'') are \textbf{withdrawn}. The strongest
unconditional results of this document remain the Route~C bootstrap
(Theorem~\ref{thm:routeC}) and Claim~A for shear and near-shear initial data
(\S11--\S12).

\subsection{Defect 1: the induction of Theorem 20.4 is arithmetically impossible for large data}
The inductive step requires
\[
  C_{\mathrm{v}}\,D^{5/3}\,2^{10/9}\,r^{10/9} \;\leq\; \tfrac{D}{2}\,r^{2/3}
  \qquad (0 < r \leq \tfrac12),
\]
equivalently $D^{2/3} \leq 2^{-19/9}/C_{\mathrm{v}}$ --- i.e.\ $D$ must be
\emph{small}. The base case requires $D \geq C_0 = \|u_0\|_{H^1}^2\,T$ ---
i.e.\ $D$ must be \emph{large} for general data. The choice
$D = \max\bigl(C_0,\,(2C_{\mathrm{v}}2^{10/9})^3\bigr)$ stated in the proof of
Theorem~\ref{thm:uniform-enstrophy} satisfies \emph{neither} constraint:
substituting it back makes the superlinear term strictly larger than $D/2$
whenever $C_{\mathrm{v}} \geq 1$. This is the standard smallness obstruction
for superlinear scale recursions: iteration contracts only below a threshold,
and general $H^1$ data starts above it. The recursion therefore reduces to the
CKN small-energy regime --- the same limitation as \S19. (This was already
acknowledged in the ``Precise gap'' remark of \S21 but had not been propagated
to the statements of \S20 or to the program documents.)

\subsection{Defect 2: invalid H\"older triple in Lemma 20.2, Step 3}
The bound $|\text{stretching}| \leq \|\nabla u\|_{L^3}\|\omega\|_{L^2}^2$ uses
exponents $\tfrac13 + \tfrac12 + \tfrac12 = \tfrac43 > 1$, which is not a valid
H\"older triple. The correct pairing is
$\|\nabla u\|_{L^3}\|\omega\|_{L^3}^2$, which changes every downstream
exponent: interpolation gives
$|\text{stretching}| \leq C\|\omega\|_{L^2}^{3/2}\|\nabla\omega\|_{L^2}^{3/2}$,
and Young's inequality then produces $\|\omega\|_{L^2}^6$ (not
$\|\omega\|_{L^2}^{10/3}$) against the dissipation. The claimed $5/3$
superlinearity exponent does not survive.

\subsection{Defect 3: reversed Jensen inequality in Lemma 20.2, Step 5}
Step 5 asserts, with $g(t) = \|\omega(\cdot,t)\|_{L^2(B_{2r})}^2$,
\[
  \int_{t_0-r^2}^{t_0} g^{5/3}\,dt \;\leq\; r^{-2/3}\Bigl(\int_{t_0-r^2}^{t_0} g\,dt\Bigr)^{5/3}.
\]
By Jensen's inequality the true relation is the \emph{reverse}:
$\fint g^{5/3} \geq \bigl(\fint g\bigr)^{5/3}$, i.e.\
$\int g^{5/3} \geq r^{-4/3}\bigl(\int g\bigr)^{5/3}\cdot r^{2/3}$.
An upper bound of a higher $L^p_t$ norm by a lower one requires additional
information (e.g.\ $g \in L^\infty_t$), which is not available for Leray--Hopf
solutions. The conversion of the stretching term into
$C\,Z(2r)^{5/3}$ is therefore invalid.

\subsection{Defect 4: regularity circularity and global/local norm mixing}
Lemma 20.1 tests the vorticity equation against $\varphi_r^2\omega$, which is
justified for smooth solutions but not for Leray--Hopf solutions, for which
$\nabla\omega \in L^2_{\mathrm{loc}}$ is not known a priori --- suitable weak
solutions satisfy a local \emph{energy} inequality, not a local
\emph{enstrophy} inequality. Since the theorem quantifies over Leray--Hopf
solutions, the argument assumes regularity it aims to prove. Additionally,
Steps 1--2 of Lemma 20.2 use global norms $\|\omega\|_{L^2(\T^3)}$,
$\|\nabla\omega\|_{L^2(\T^3)}$ where the conclusion requires local quantities
on $B_{2r}$; the localization of the Biot--Savart/CZ step is not supplied.

\subsection{Defect 5: Theorem \ref{thm:sigma-gronwall} is conditional on the BKM criterion}
The hypothesis $\int_0^T M(t)\,dt < \infty$ with
$M(t)=\|\omega(\cdot,t)\|_{L^\infty}$ is precisely the Beale--Kato--Majda
criterion: under it, regularity on $[0,T]$ already follows classically. The
conditional $\sigma$-bound therefore cannot serve as a route to regularity;
it assumes (an equivalent of) its conclusion.

\subsection{Defect 6: the two-regime dichotomy requires a large-viscosity condition}
The damping constant $\beta = c\nu/C_1 - C > 0$ requires $\nu > C_1 C/c$, a
relation between the fixed physical viscosity and universal constants. In the
Prize formulation $\nu$ is given; the condition cannot be arranged by scaling
(NS scaling rescales data and $M$, not $\nu$ relative to the constants in this
argument). Regime~I is therefore conditional.

\subsection{Defect 7: the Regime II ODE conclusion is false}
Regime~II derives $dM/dt \leq C\,M^{5/4}(\log M)^{1/2}$ and asserts the
solution ``stays finite on any fixed time interval.'' This is incorrect: the
comparison ODE $\dot{M} = C M^{5/4}$ has solution
$M(t) = \bigl(M(0)^{-1/4} - \tfrac{C}{4}t\bigr)^{-4}$, which blows up at the
finite time $t^* = 4 M(0)^{-1/4}/C$; the logarithmic factor only accelerates
this. Regime~II as stated does not exclude finite-time blowup.

\subsection{Defect 8: unproved coercivity in Lemma \ref{lem:A-evol}}
The claim that diffusion contributes $-c\nu\int|\nabla^2\omega|^2\varphi^2$ to
$\tfrac{d}{dt}A_{\mathrm{loc}}$ is asserted without proof. Differentiating
$A_{\mathrm{loc}} = \int |\omega|^2 |\nabla\hat{e}|^2 \varphi^2$ involves the
direction field $\hat{e} = \omega/|\omega|$, singular on $\{\omega = 0\}$, and
produces sign-indefinite cross terms; full $|\nabla^2\omega|^2$-coercivity does
not follow from the sketch. Moreover the stretching estimate invokes
$\|\nabla u\|_{L^\infty(B_{r^*})} \leq C M$, but Calder\'on--Zygmund operators
are unbounded on $L^\infty$; the correct bound carries a logarithmic factor,
which matters precisely because \S22 relies on log-corrected estimates.

\subsection{What survives the audit}
\begin{enumerate}
  \item Route~C bootstrap (Theorem~\ref{thm:routeC}) --- unconditional. \emph{Intact.}
  \item Claim~A for shear-layer initial data (\S11) and near-shear data (\S12). \emph{Intact.}
  \item $\lambda_{\max} = 0$ for shear flows (Prop.~11.6). \emph{Intact.}
  \item The algebraic identity $|\nabla\omega|^2 = |\nabla|\omega||^2 + |\omega|^2|\nabla\hat{e}|^2$
        (Lemma~\ref{lem:grad-omega-split}) and the exact scale invariance of
        $\sigma = A_{\mathrm{loc}}/M^{3/2}$. \emph{Intact} --- these are sound
        foundations for a future quantitative alignment program.
  \item \S20 as a \emph{conditional} small-data statement (consistent with CKN),
        \emph{provided} Defects 2--4 are repaired; the superlinearity exponent
        must be re-derived.
\end{enumerate}

\subsection{Revised status}
Global regularity for general $u_0 \in H^1(\T^3)$ is \textbf{open} in this
document. The Prize claim of \S20 (Corollary~\ref{cor:prize-h1}) is withdrawn.
The program's strongest verified position: unconditional regularity for shear
and near-shear initial data via Claim~A and Route~C, plus extensive numerical
evidence (873 passing tests) consistent with --- but not proving --- the
general case.

% ══════════════════════════════════════════════════════════════════════════════
\section{The Alignment Pincer: Corrected Direction-Field Framework (Session S32)}
\label{sec:pincer-s32}

Following the S31 withdrawal, this section restarts the geometric program on
rigorous footing. Every statement below is labeled \emph{Proved},
\emph{Conditional}, or \emph{Open target}. Nothing here claims the Prize.

\subsection{The exact direction-field evolution equation [Proved]}

\begin{lemma}[Constantin's direction equation]\label{lem:direction-eq}
  Let $u$ be a smooth solution of 3D NS and $\hat{e}=\omega/|\omega|$ on
  $\{\omega\neq 0\}$, $\alpha = \hat{e}\cdot S\hat{e}$ with
  $S=\tfrac12(\nabla u + \nabla u^{\top})$. Then, with
  $D_t = \partial_t + u\cdot\nabla$:
  \begin{align}
    D_t|\omega| &= \alpha\,|\omega| + \nu\bigl(\Delta|\omega| - |\omega|\,|\nabla\hat{e}|^2\bigr),
      \label{eq:mag-eq}\\
    D_t\hat{e} &= \nu\Delta\hat{e} + 2\nu\,(\nabla\log|\omega|\cdot\nabla)\hat{e}
      + \nu|\nabla\hat{e}|^2\,\hat{e} + P_{\hat{e}^{\perp}}\bigl((\hat{e}\cdot\nabla)u\bigr),
      \label{eq:dir-eq}
  \end{align}
  where $P_{\hat{e}^{\perp}}v = v - (v\cdot\hat{e})\hat{e}$.
\end{lemma}

\begin{proof}
  Write $\omega=|\omega|\hat{e}$. The vorticity equation gives
  $D_t\omega = (\omega\cdot\nabla)u + \nu\Delta\omega
  = |\omega|(\hat{e}\cdot\nabla)u + \nu\Delta\omega$.
  Projecting onto $\hat{e}$: since $|\hat{e}|=1$ implies
  $\hat{e}\cdot\partial_k\hat{e}=0$ for every $k$, we get
  $\hat{e}\cdot\Delta\hat{e} = \sum_k \partial_k(\hat{e}\cdot\partial_k\hat{e})
  - |\nabla\hat{e}|^2 = -|\nabla\hat{e}|^2$, hence
  $\hat{e}\cdot\Delta\omega = \Delta|\omega| - |\omega||\nabla\hat{e}|^2$,
  and $\hat{e}\cdot\bigl(|\omega|(\hat{e}\cdot\nabla)u\bigr)
  = |\omega|\,\hat{e}_i\hat{e}_j\partial_j u_i = \alpha|\omega|$
  (the antisymmetric part of $\nabla u$ drops). This is~\eqref{eq:mag-eq}.
  For~\eqref{eq:dir-eq}, expand
  $D_t(|\omega|\hat{e}) = \hat{e}\,D_t|\omega| + |\omega|\,D_t\hat{e}$ and
  $\Delta(|\omega|\hat{e}) = \hat{e}\,\Delta|\omega|
  + 2(\nabla|\omega|\cdot\nabla)\hat{e} + |\omega|\Delta\hat{e}$;
  substitute~\eqref{eq:mag-eq}, divide by $|\omega|$, and note
  $(\hat{e}\cdot\nabla)u - \alpha\hat{e} = P_{\hat{e}^{\perp}}((\hat{e}\cdot\nabla)u)$.
\end{proof}

\begin{remark}[No pressure]
  Equations~\eqref{eq:mag-eq}--\eqref{eq:dir-eq} contain no pressure: the
  direction field is slaved to $\nabla u$ through Biot--Savart only. This is
  the same structural advantage \S20 sought, obtained without any invalid
  estimate.
\end{remark}

\subsection{A correctly scale-invariant coherence quantity [Proved]}

\begin{definition}[Instantaneous coherence deficit]\label{def:sigma-star}
  With $M(t)=\|\omega(\cdot,t)\|_{L^\infty}$, $r^*=M^{-1/2}$, and $x^*$ a
  point where $|\omega(\cdot,t)|\geq M/2$:
  \begin{equation}\label{eq:sigma-star}
    \sigma^*(t) \;=\; (r^*)^{-1}\int_{B_{r^*}(x^*)} |\nabla\hat{e}|^2\,
      \mathbf{1}_{\{|\omega|\geq M/4\}}\,dx.
  \end{equation}
\end{definition}

\begin{proposition}[Exact scale invariance]\label{prop:sigma-star-invariant}
  Under $u_\lambda(x,t)=\lambda u(\lambda x,\lambda^2 t)$ one has
  $\omega_\lambda = \lambda^2\omega(\lambda x,\lambda^2 t)$,
  $M_\lambda=\lambda^2 M$, $r^*_\lambda = r^*/\lambda$,
  $\hat{e}_\lambda(x,t)=\hat{e}(\lambda x,\lambda^2 t)$, so
  $|\nabla\hat{e}_\lambda|^2 = \lambda^2|\nabla\hat{e}|^2(\lambda x)$ and
  $\int_{B_{r^*/\lambda}}|\nabla\hat{e}_\lambda|^2\,dx
   = \lambda^{-1}\int_{B_{r^*}}|\nabla\hat{e}|^2\,dx$. Multiplying by
  $(r^*_\lambda)^{-1}=\lambda(r^*)^{-1}$ gives
  $\sigma^*_\lambda = \sigma^*$. Moreover the indicator set is
  scale-covariant since $|\omega_\lambda|\geq M_\lambda/4 \iff
  |\omega|(\lambda x)\geq M/4$. \qed
\end{proposition}

\begin{remark}[Calibration]
  If $|\nabla\hat{e}|\sim 1/r^*$ throughout $B_{r^*}$ (borderline
  Constantin--Fefferman coherence at the self-similar scale) then
  $\sigma^*\sim 1$. Thus $\sigma^*\ll 1$ means \emph{better-than-borderline}
  coherence, and $\sigma^*$ is the natural small parameter for an
  $\varepsilon$-regularity statement.
\end{remark}

\begin{proposition}[The S30 quantity controls $\sigma^*$]\label{prop:sigma-controls}
  With $\sigma = A_{\mathrm{loc}}/M^{3/2}$ as in
  Definition~\eqref{eq:sigma-def} (whose scale invariance and algebraic
  foundation survived the S31 audit):
  \[
    \sigma^*(t) \;\leq\; 16\,\sigma(t).
  \]
\end{proposition}
\begin{proof}
  On $\{|\omega|\geq M/4\}$ we have $|\omega|^2\geq M^2/16$, so
  $\int_{B_{r^*}}|\nabla\hat{e}|^2\mathbf{1}\,dx
  \leq (16/M^2)\int |\omega|^2|\nabla\hat{e}|^2\varphi_{r^*}^2\,dx
  = 16\,A_{\mathrm{loc}}/M^2$ (using $\varphi_{r^*}\equiv 1$ on $B_{r^*}$).
  Multiply by $(r^*)^{-1}=M^{1/2}$:
  $\sigma^*\leq 16\,A_{\mathrm{loc}}/M^{3/2}=16\sigma$.
\end{proof}

\subsection{Criticality of the direction equation at the self-similar scale [Proved observation]}

\begin{remark}[Marginality: $M\cdot(r^*)^2 = 1$]\label{rem:marginality}
  In any parabolic estimate for~\eqref{eq:dir-eq} on the cylinder
  $Q(r^*) = B_{r^*}\times[t_0-(r^*)^2,t_0]$, the stretching term
  $P_{\hat{e}^{\perp}}((\hat{e}\cdot\nabla)u)$ acts as a potential of size
  $\|\nabla u\|_{L^\infty}\lesssim M\log(e+\|u\|_{H^3}/M)$ (the logarithm is
  mandatory: Calder\'on--Zygmund operators are unbounded on $L^\infty$ ---
  this was Defect 8 of the S31 audit). Over the time span $(r^*)^2 = M^{-1}$
  its Gronwall exponent is
  \[
    \|\nabla u\|_{L^\infty}\cdot (r^*)^2 \;\lesssim\; M\cdot M^{-1}\cdot\log(\cdot)
    \;=\; \log(\cdot).
  \]
  The direction equation is therefore \emph{exactly critical} (marginally
  so, up to a logarithm) at the scale $r^*=M^{-1/2}$: the potential cannot
  amplify $\nabla\hat{e}$ by more than a power of the logarithm per
  self-similar time scale. This is the structural reason an
  $\varepsilon$-regularity theory for $\hat{e}$ at scale $r^*$ is plausible,
  where the corresponding statement for $u$ itself (CKN) is supercritical.
\end{remark}

\subsection{The Bridge Lemma [Open target]}

\begin{conjecture}[Bridge Lemma: averaged-to-pointwise coherence upgrade]
  \label{conj:bridge}
  There exist $\varepsilon_0>0$ and $C_B\geq 1$ (possibly degrading
  logarithmically in $M$) such that: if $u$ satisfies the Type-I bound
  $\sup_{Q(r^*)}|u|\leq C_I/r^*$ and
  $\sup_{t\in[t_0-(r^*)^2,\,t_0]}\sigma^*(t)\leq\varepsilon_0$, then
  \[
    \|\nabla\hat{e}\|_{L^\infty\bigl(Q(r^*/2)\cap\{|\omega|\geq M/2\}\bigr)}
    \;\leq\; \frac{C_B}{r^*}.
  \]
\end{conjecture}

\noindent\textbf{Proof strategy and identified obstacles.} The natural proof
is Moser iteration for $w=|\nabla\hat{e}|^2$, which
by~\eqref{eq:dir-eq} satisfies a drift-diffusion inequality with (i) drift
$u + 2\nu\nabla\log|\omega|$ and (ii) potential
$\lesssim\|\nabla u\|_{L^\infty(B_{r^*})}$. Obstacle~(i): on
$\{|\omega|\geq M/4\}$ the singular drift satisfies
$|\nabla\log|\omega||\leq 4|\nabla\omega|/M$, which requires a pointwise
gradient bound not currently available; a workaround is to run the iteration
on level sets of $|\omega|$ where the drift enters only through its
divergence structure. Obstacle~(ii): the potential is marginal by
Remark~\ref{rem:marginality}, so the iteration constants degrade
logarithmically --- acceptable, but the bookkeeping must be done. Neither
obstacle is resolved here; this is the program's central open estimate,
replacing the invalid superlinearity of \S20.

\subsection{The pincer [Conditional on two named inputs]}

\begin{theorem}[Disorder barrier for Type-I singularities --- conditional]
  \label{thm:pincer-conditional}
  Assume (H1) the Bridge Lemma (Conjecture~\ref{conj:bridge}) and (H2) a
  localized Constantin--Fefferman criterion at scale $r^*$: coherence
  $\|\nabla\hat{e}\|_{L^\infty}\leq C_B/r^*$ on
  $Q(r^*/2)\cap\{|\omega|\geq M/2\}$ implies the depletion bound
  $\alpha(x^*,t)\leq C\,M^{1/2+\theta}$ for some $\theta<1/2$ near a
  putative Type-I blowup time. Then no Type-I singularity can satisfy
  $\liminf_{t\to T^-}\sigma^*(t)\leq\varepsilon_0$: every Type-I singularity
  must maintain $\sigma^*(t)>\varepsilon_0$ (persistent geometric disorder)
  for all $t$ close to $T$.
\end{theorem}
\begin{proof}[Proof (given H1--H2)]
  If $\sigma^*(t_k)\leq\varepsilon_0$ along $t_k\to T^-$, H1 upgrades
  averaged coherence to pointwise coherence on $Q(r^*(t_k)/2)$; H2 then
  bounds the logarithmic growth rate $\tfrac{d}{dt}\log M = \alpha(x^*,t)
  \leq CM^{1/2+\theta}$ with $1/2+\theta<1$, whose integral
  $\int^{T}M^{1/2+\theta}\,dt$ converges under the Type-I bound
  $M\leq C_I/(T-t)$ precisely when $1/2+\theta<1$. Hence $M$ stays bounded
  as $t\to T^-$, contradicting blowup.
\end{proof}

\begin{remark}[Status of H2]
  H2 is a genuine second open input: the classical Constantin--Fefferman
  theorem \cite{ConstantinFefferman1993} assumes coherence at a
  \emph{fixed} scale $\rho$, whereas here the coherence scale $r^*(t)\to 0$
  shrinks with the putative singularity. The localized/shrinking-scale
  variants (Be\~{i}r\~{a}o da Veiga--Berselli type $1/2$-H\"older criteria,
  Gruji\'c's geometric depletion) are the relevant literature; adapting them
  to scale $r^*$ with explicit $\theta$ is a bounded, well-posed task.
\end{remark}

\subsection{Honest status table for the pincer program}

\begin{center}
\begin{tabular}{lll}
\hline
Item & Status & Where \\
\hline
Direction equation~\eqref{eq:dir-eq} & Proved & Lemma~\ref{lem:direction-eq} \\
Scale invariance of $\sigma^*$ & Proved & Prop.~\ref{prop:sigma-star-invariant} \\
$\sigma^*\leq 16\sigma$ (link to S30) & Proved & Prop.~\ref{prop:sigma-controls} \\
Criticality $M(r^*)^2=1$ & Proved observation & Remark~\ref{rem:marginality} \\
Bridge Lemma (H1) & \textbf{Open target} & Conjecture~\ref{conj:bridge} \\
Localized CF at scale $r^*$ (H2) & \textbf{Open target} & Theorem~\ref{thm:pincer-conditional} \\
Disorder barrier & Conditional on H1+H2 & Theorem~\ref{thm:pincer-conditional} \\
Numerical bridge diagnostic & EXP-L4-BRIDGE-* & \texttt{layer4/alignment\_bridge.py} \\
\hline
\end{tabular}
\end{center}

% ══════════════════════════════════════════════════════════════════════════════
\section{Level-Set Moser Scheme for the Bridge Lemma: The First Rung (Session S33)}
\label{sec:moser-s33}

This section attacks Obstacle~(i) of Conjecture~\ref{conj:bridge}: the
singular drift $2\nu\nabla\log|\omega|$ in the direction
equation~\eqref{eq:dir-eq}, whose differentiated form produces
$\nabla^2\log|\omega|$ --- uncontrollable pointwise. The device is a
\emph{level-set formulation}: all estimates are run with a smooth cutoff in
the $|\omega|$-variable, and every second-derivative-of-magnitude term is
integrated by parts back onto first-order factors that the level set
controls. The outcome (Proposition~\ref{prop:first-rung}) closes the first
rung of a Moser iteration with logarithmically degraded constants and
\emph{reduces} Obstacle~(i) to a strictly weaker, sharply stated coupling
estimate (Target~\ref{target:K-absorption}).

\subsection{Standing assumptions and non-circularity}

Throughout this section $u$ is a solution that is smooth on
$\T^3\times[0,T)$ with putative first singularity at $T$, satisfying the
Type-I bound $M(t)\leq C_I/(T-t)$ and $\sup|u(\cdot,t)|\leq C_I(T-t)^{-1/2}$.
All computations occur at times $t<T$ where the solution is $C^\infty$.

\begin{remark}[No Leray--Hopf circularity]\label{rem:no-circularity}
  Defect~4 of the S31 audit (\S\ref{sec:audit-s31}) noted that \S20 applied
  a local enstrophy inequality to Leray--Hopf solutions for which it is not
  justified. The pincer route is immune to this defect by design: it is a
  \emph{contradiction argument at a putative first singularity}, so every
  estimate is performed on the smooth interval $[0,T)$ where all
  manipulations are classical. No weak-solution regularity is assumed that
  is not available.
\end{remark}

\subsection{The composite level-set cutoff [Proved]}

\begin{definition}[Composite cutoff]\label{def:composite-cutoff}
  Fix $\chi\in C^\infty([0,\infty);[0,1])$ with $\chi\equiv 0$ on
  $[0,\tfrac14]$, $\chi\equiv 1$ on $[\tfrac12,\infty)$,
  $|\chi'|\leq 8$. For scales $\tfrac{r^*}{2}\leq\rho<\rho'\leq r^*$ let
  $\psi\in C_c^\infty(Q(\rho'))$ be a standard space-time cutoff with
  $\psi\equiv 1$ on $Q(\rho)$,
  $|\nabla\psi|\leq 2/(\rho'-\rho)$, $|\partial_t\psi|\leq 2/(\rho'-\rho)^2$.
  The composite cutoff is
  \[
    \eta(x,t) \;=\; \psi(x,t)\,\chi\!\bigl(|\omega(x,t)|/M(t)\bigr).
  \]
\end{definition}

\begin{lemma}[Cutoff derivative bounds]\label{lem:cutoff-derivs}
  On $\operatorname{supp}\eta\subseteq\{|\omega|\geq M/4\}$:
  \begin{align}
    |\nabla\eta| &\;\leq\; \frac{2}{\rho'-\rho}
      + \frac{8}{M}\bigl|\nabla|\omega|\bigr|, \label{eq:cutoff-grad}\\
    |D_t\chi(|\omega|/M)| &\;\leq\;
      \frac{8}{M}\Bigl(|\alpha|\,|\omega| + \nu\bigl|\Delta|\omega|\bigr|
      + \nu|\omega|\,|\nabla\hat{e}|^2\Bigr)
      + 8\,\frac{|\dot M|}{M}. \label{eq:cutoff-time}
  \end{align}
\end{lemma}
\begin{proof}
  \eqref{eq:cutoff-grad}: product rule and
  $|\nabla\chi(|\omega|/M)| = |\chi'|\,|\nabla|\omega||/M$.
  \eqref{eq:cutoff-time}: $D_t\chi = \chi'\bigl(D_t|\omega|/M
  - |\omega|\dot M/M^2\bigr)$; substitute the magnitude
  equation~\eqref{eq:mag-eq} and use $|\omega|/M\leq 1$, $|\chi'|\leq 8$.
\end{proof}

\begin{remark}
  The term $\nu|\Delta|\omega||$ in~\eqref{eq:cutoff-time} is second-order
  and is \emph{never used pointwise}: wherever $D_t\chi$ appears inside an
  integral it is integrated by parts (Lemma~\ref{lem:three-ibp}) so that
  only $\nabla|\omega|$ survives. Under Type~I,
  $|\dot M|/M \leq \|\nabla u\|_{L^\infty} \leq C M \log(e+\|u\|_{H^3}/M)$
  for a.e.\ $t$ (Rademacher for the Lipschitz function $M(t)$), which
  contributes the marginal factor
  $M\log\cdot(r^*)^2 = \log$ per cylinder (Remark~\ref{rem:marginality}).
\end{remark}

\subsection{The three integrations by parts [Proved]}

Testing the differentiated direction equation with $\eta^2\partial_k\hat{e}$
(summed over $k$) yields, for $w=|\nabla\hat{e}|^2$, the weak inequality
\begin{equation}\label{eq:weak-w}
  \tfrac12\sup_{t}\!\int w\,\eta^2
  + \nu\!\iint |\nabla^2\hat{e}|^2\eta^2
  \;\leq\;
  \tfrac12\iint w\,\bigl|D_t(\eta^2)\bigr|
  + \nu\iint w\,|\Delta(\eta^2)|
  + \mathcal{D} + \mathcal{S},
\end{equation}
where $\mathcal{D}$ collects the drift terms generated by
$2\nu(\nabla\log|\omega|\cdot\nabla)\hat{e}$ and $\mathcal{S}$ the
stretching terms generated by $\nu w\hat{e} + P_{\hat{e}^\perp}((\hat{e}\cdot\nabla)u)$.

\begin{lemma}[Three integrations by parts]\label{lem:three-ibp}
  Each second-derivative-of-magnitude term in~\eqref{eq:weak-w} can be
  rewritten, using only integration by parts and the support property of
  $\eta$, in forms containing at most first derivatives of $|\omega|$:
  \begin{enumerate}[label=(\alph*)]
    \item the differentiated drift term
    \[
      2\nu\iint (\partial_k\partial_j\log|\omega|)(\partial_j\hat{e}\cdot\partial_k\hat{e})\,\eta^2
      \;=\; -2\nu\iint \partial_j\log|\omega|\;
        \partial_k\Bigl[(\partial_j\hat{e}\cdot\partial_k\hat{e})\,\eta^2\Bigr],
    \]
    whose integrand is bounded by
    $\tfrac{8\nu}{M}|\nabla|\omega||\bigl(2|\nabla^2\hat{e}|\,|\nabla\hat{e}|\,\eta^2
    + 2 w\,\eta|\nabla\eta|\bigr)$ on $\operatorname{supp}\eta$;
    \item the $\nu\Delta|\omega|$ contribution inside $D_t\chi$
    (Lemma~\ref{lem:cutoff-derivs}):
    \[
      \frac{\nu}{M}\iint w\,\eta\,\psi\,|\chi'|\,\bigl|\Delta|\omega|\bigr|
      \;\rightsquigarrow\;
      -\frac{\nu}{M}\iint \nabla|\omega|\cdot\nabla\bigl(w\,\eta\,\psi\,\chi'\bigr),
    \]
    with integrand bounded by
    $\tfrac{C\nu}{M}|\nabla|\omega||\bigl(|\nabla^2\hat{e}||\nabla\hat{e}|\eta
    + w|\nabla\eta| + \tfrac{w\,\eta}{M}|\nabla|\omega||\bigr)$;
    \item the cutoff-gradient term $\nu\iint w|\Delta(\eta^2)|$, whose
    $\chi$-part after one integration by parts contains only
    $|\nabla|\omega||/M$ factors as in~(a).
  \end{enumerate}
  Consequently, applying Young's inequality with parameter
  $\tfrac{\nu}{8}$ to every product
  $|\nabla^2\hat{e}|\times(\text{first-order factor})$, all three terms are
  absorbed as
  \[
    \mathcal{D} + \text{(b)} + \text{(c)} \;\leq\;
    \frac{\nu}{2}\iint|\nabla^2\hat{e}|^2\eta^2
    \;+\; \frac{C\nu}{(\rho'-\rho)^2}\iint_{Q(\rho')} w\,\mathbf{1}_{\{|\omega|\geq M/4\}}
    \;+\; C\nu\,K(\rho'),
  \]
  where
  \begin{equation}\label{eq:K-def}
    K(\rho') \;=\; \frac{1}{M^2}\iint_{Q(\rho')\cap\{|\omega|\geq M/4\}}
      \bigl|\nabla|\omega|\bigr|^2\,\bigl(w + (\rho'-\rho)^{-2}\bigr)\,dx\,dt.
  \end{equation}
\end{lemma}
\begin{proof}
  (a) and (b) are direct integrations by parts (no boundary terms:
  $\eta$ is compactly supported in space-time, and the integrands vanish on
  $\{|\omega|<M/4\}$ where $\chi=0$, so the level-set boundary contributes
  nothing). The stated integrand bounds use
  $|\nabla\log|\omega||\leq 4|\nabla|\omega||/M$ on
  $\operatorname{supp}\eta$, $|\chi'|\leq 8$,
  Lemma~\ref{lem:cutoff-derivs}, and
  $|\nabla(w\eta\psi\chi')| \leq 2|\nabla^2\hat{e}||\nabla\hat{e}|\eta\psi|\chi'|
  + w|\nabla(\eta\psi\chi')|$. The absorption step is Young's inequality
  $ab\leq \tfrac{\nu}{8}a^2 + \tfrac{2}{\nu}b^2$ applied term by term;
  every non-dissipative square lands either on
  $w/(\rho'-\rho)^2$ or on $(|\nabla|\omega||^2/M^2)\cdot(w + (\rho'-\rho)^{-2})$,
  which is $K(\rho')$ by definition.
\end{proof}

\subsection{Localized stretching control under Type I [Proved]}

\begin{lemma}[Local enstrophy-gradient bound, Type I]\label{lem:local-enstrophy-typeI}
  On the smooth interval, under the Type-I bounds, the vorticity Caccioppoli
  inequality (classical for smooth solutions; cf.\
  Remark~\ref{rem:no-circularity}) with
  $\|\nabla u\|_{L^\infty}\lesssim M\log(e+\|u\|_{H^3}/M)$ gives
  \[
    \nu\iint_{Q(2\rho)}|\nabla\omega|^2
    \;\lesssim\; M^2\rho^3\bigl(1 + M\rho^2\,\mathcal{L}\bigr)
    \;\lesssim\; M^2\rho^3\,\mathcal{L},
    \qquad \rho\leq r^*,
  \]
  where $\mathcal{L} = \log(e+\|u\|_{H^3}/M)$, using $M\rho^2\leq M (r^*)^2=1$.
\end{lemma}

\begin{lemma}[Biot--Savart near/far splitting]\label{lem:BS-split}
  Write $u = u_{\mathrm{loc}} + u_{\mathrm{far}}$ with
  $u_{\mathrm{loc}}$ the Biot--Savart integral of
  $\omega\,\mathbf{1}_{B_{2r^*}}$. Then
  $\|\nabla^2 u_{\mathrm{loc}}\|_{L^2(B_{r^*})}
   \leq C_{\mathrm{CZ}}\|\nabla\omega\|_{L^2(B_{2r^*})}$
  (Calder\'on--Zygmund, $L^2$), while $u_{\mathrm{far}}$ is harmonic-regular
  on $B_{r^*}$ with
  $\|\nabla^2 u_{\mathrm{far}}\|_{L^\infty(B_{r^*})}
   \leq C (r^*)^{-2}\,\sup|u| \leq C\,C_I\,(r^*)^{-3}$
  under Type~I. Hence the stretching contribution $\mathcal{S}$
  in~\eqref{eq:weak-w} obeys
  \[
    \mathcal{S} \;\leq\; \frac{\nu}{4}\iint|\nabla^2\hat{e}|^2\eta^2
    + C\,M\mathcal{L}\iint_{Q(\rho')} w\,\mathbf{1}
    + C\,\nu^{-1}\,\mathcal{L}\,M^2(r^*)^5
    + C\,C_I^2\,r^*,
  \]
  where the third term uses Lemma~\ref{lem:local-enstrophy-typeI} and the
  fourth is the far-field contribution integrated over $Q(r^*)$.
\end{lemma}

\subsection{The first rung [Proved modulo routine verifications]}

\begin{proposition}[First-rung Caccioppoli for the angular energy]
  \label{prop:first-rung}
  Under the standing assumptions, for
  $\tfrac{r^*}{2}\leq\rho<\rho'\leq r^*$:
  \begin{align*}
    \sup_{t}\int_{B_\rho} w\,\mathbf{1}_{\{|\omega|\geq M/2\}}\,dx
    &+ \nu\iint_{Q(\rho)}|\nabla^2\hat{e}|^2\,\mathbf{1}_{\{|\omega|\geq M/2\}} \\
    &\leq\; \frac{C\,\mathcal{L}}{(\rho'-\rho)^2}
      \iint_{Q(\rho')} w\,\mathbf{1}_{\{|\omega|\geq M/4\}}
      \;+\; C\nu\,K(\rho')
      \;+\; C\bigl(\nu^{-1}\mathcal{L}\,M^2(r^*)^5 + C_I^2\,r^*\bigr),
  \end{align*}
  with $C = C(C_I)$ universal up to the Type-I constant and
  $\mathcal{L}$ the logarithmic factor. \emph{Status: the displayed
  dangerous terms are handled in full by
  Lemmas~\ref{lem:cutoff-derivs}--\ref{lem:BS-split}; the routine
  verifications that remain for S34 are enumerated in
  Remark~\ref{rem:routine-list} and are bookkeeping, not new estimates.}
\end{proposition}

\begin{remark}[Enumerated routine verifications]\label{rem:routine-list}
  (i) the transport term $\iint w\,u\cdot\nabla(\eta^2)$ (standard:
  H\"older + Type-I $|u|\leq C_I/r^*$); (ii) the algebraic constant in the
  commutation $[\partial_k, D_t]\hat{e} = -(\partial_k u\cdot\nabla)\hat{e}$;
  (iii) \emph{reclassified in S34}: the $\nu w\hat{e}$-term yields the
  critical quadratic $\nu\iint w^2\eta^2$ ($\partial_k\hat{e}\cdot\hat{e}=0$
  kills only the $\partial_k w$ part); this is \emph{not} routine --- it is
  the standard critical nonlinearity of the iteration, handled by
  Gagliardo--Nirenberg interpolation plus the $\sigma^*\leq\varepsilon_0$
  smallness at the second rung (see \S\ref{sec:magnitude-s34}); (iv) the
  time-slice boundary term at $t_0-(\rho')^2$ (vanishes by choice of
  $\psi$). None of these involves $\nabla^2\log|\omega|$ or nonlocality.
\end{remark}

\subsection{What remains of Obstacle (i): the coupling target [Open]}

\begin{conjecture}[$K$-absorption]\label{target:K-absorption}
  Under the standing assumptions there exist $\theta_K<1$ and $C_K$
  (possibly logarithmic in $M$) such that
  \[
    \nu\,K(r^*) \;\leq\; \theta_K\Bigl[\text{LHS of
    Proposition~\ref{prop:first-rung}}\Bigr]
    \;+\; C_K\bigl(1 + \sigma^*\bigr)\,r^*.
  \]
\end{conjecture}

\begin{remark}[Why this is strictly weaker than Obstacle (i)]
  Obstacle~(i) demanded a \emph{pointwise} bound on
  $\nabla\log|\omega|$. Target~\ref{target:K-absorption} demands only that
  the \emph{integrated} radial--angular coupling
  $M^{-2}\iint|\nabla|\omega||^2 w$ be absorbable. Two independent routes
  are visible: (1)~interpolation --- $K$ is quartic in first derivatives
  and can be split by H\"older between the local enstrophy bound
  (Lemma~\ref{lem:local-enstrophy-typeI}) and $\|w\|_{L^\infty}$, making
  the absorption a small-$\sigma^*$ smallness argument at the second rung
  of the iteration; (2)~the magnitude equation~\eqref{eq:mag-eq} gives its
  own Caccioppoli for $|\nabla|\omega||^2$ with the \emph{good sign} term
  $-\nu|\omega| w$ appearing as a sink. Route (2) is the natural S34
  starting point.
\end{remark}

\begin{remark}[Numerical status]
  The S32--S33 diagnostics (\texttt{layer4/alignment\_bridge.py},
  EXP-L4-BRIDGE-*) show the bridge ratio
  $\|\nabla\hat{e}\|_{L^\infty}r^*/\sqrt{\sigma^*}$ remaining $O(1)$
  through strong vortex stretching, consistent with
  Conjecture~\ref{conj:bridge} and with the absorbability of $K$.
\end{remark}

% ══════════════════════════════════════════════════════════════════════════════
\section{The Magnitude-Equation Caccioppoli: A Priori Coherence Bound and the Refined $K$-Target (Session S34)}
\label{sec:magnitude-s34}

\subsection{Corrigendum to \S\ref{sec:moser-s33} [dimensional audit]}

A dimensional re-audit of \S\ref{sec:moser-s33} (performed before building on
it, per the S31 discipline) found three slips in the inhomogeneous terms of
Lemma~\ref{lem:BS-split}, Proposition~\ref{prop:first-rung} and
Conjecture~\ref{target:K-absorption}: in parabolic units every term of the
first-rung inequality scales as $\ell^{1}$, forcing the corrected forms
$\nu^{-1}\mathcal{L}M^2(r^*)^5$ and $C_I^2\,r^*$ (both now \emph{vanish} as
$r^*\to 0$ --- the corrected inequality is subcritically favourable, not
worse). Additionally, routine item~(iii) of Remark~\ref{rem:routine-list} is
reclassified: it is the critical quadratic $\nu\iint w^2\eta^2$, i.e.\ the
standard critical nonlinearity of every $\varepsilon$-regularity iteration,
handled by Gagliardo--Nirenberg plus the $\sigma^*\leq\varepsilon_0$
hypothesis --- this is precisely where the Bridge Lemma's smallness
assumption enters the scheme, as it should. The corrections are incorporated
above in place.

\subsection{The magnitude Caccioppoli and its good-sign sink [Proved]}

The magnitude equation~\eqref{eq:mag-eq} carries the dissipative term
$-\nu|\omega|\,w$ ($w=|\nabla\hat{e}|^2$): angular roughness \emph{drains}
vorticity magnitude. Testing with $m\eta^2$ ($m=|\omega|$, $\eta$ the
composite cutoff of Definition~\ref{def:composite-cutoff} with $\chi$-levels
$\tfrac18,\tfrac14$ so that $\{m\geq M/4\}\subseteq\{\chi=1\}$) turns this
sink into a \emph{left-hand-side} control of the coherence deficit:
\begin{equation}\label{eq:mag-caccioppoli}
  \nu\iint m^2 w\,\eta^2
  + \nu\iint|\nabla m|^2\eta^2
  + \tfrac12\sup_t\int m^2\eta^2
  \;\leq\;
  \iint \alpha\,m^2\eta^2
  + \tfrac12\iint m^2\,\bigl|D_t(\eta^2)\bigr|
  + \nu\iint m^2|\nabla\eta|^2 .
\end{equation}

\begin{lemma}[All right-hand terms are $O(M^{1/2}\mathcal{L})$]
  \label{lem:mag-rhs}
  Under the standing Type-I assumptions of \S\ref{sec:moser-s33}, on
  $Q(r^*)$:
  \begin{enumerate}[label=(\alph*)]
    \item $\iint|\alpha|m^2\eta^2 \leq \|\nabla u\|_{L^\infty}M^2(r^*)^5
      \leq C\,\mathcal{L}\,M^{1/2}$ \quad (using $M(r^*)^2=1$);
    \item the $\psi$-part of $D_t(\eta^2)$ contributes
      $\leq C(1+C_I)\,M^{1/2}$;
    \item the $\chi$-part of $D_t(\eta^2)$ and the $\chi$-part of
      $|\nabla\eta|^2$ contribute $\leq C\,\mathcal{L}\,M^{1/2}$, with
      \emph{no absorption argument needed}: every such term is supported on
      the transition annulus and is dominated pointwise via the orthogonal
      split (Lemma~\ref{lem:grad-omega-split})
      \[
        m^2 w \;=\; |\nabla\omega|^2 - |\nabla m|^2 \;\leq\; |\nabla\omega|^2,
        \qquad |\nabla m|^2 \leq |\nabla\omega|^2,
      \]
      after which Lemma~\ref{lem:local-enstrophy-typeI} gives
      $\nu\iint_{Q(r^*)}|\nabla\omega|^2 \leq C M^2 (r^*)^3\mathcal{L}
      = C\,\mathcal{L}\,M^{1/2}$. The $\nu\Delta m$ term inside $D_t\chi$ is
      first integrated by parts exactly as in
      Lemma~\ref{lem:three-ibp}(b), landing on the same enstrophy bound.
  \end{enumerate}
\end{lemma}

\begin{theorem}[A priori coherence bound under Type I]\label{thm:apriori-sigma}
  There is $C=C(C_I)$ such that for every $t_0<T$, with $r^*=M(t_0)^{-1/2}$:
  \[
    \langle\sigma^*\rangle
    \;:=\;
    (r^*)^{-3}\iint_{Q(r^*/2)} w\,\mathbf{1}_{\{|\omega|\geq M/2\}}\,dx\,dt
    \;\leq\; \frac{C\,\mathcal{L}}{\nu}.
  \]
\end{theorem}
\begin{proof}
  On $\{|\omega|\geq M/2\}\cap\{\eta=1\}$ we have $m^2\geq M^2/4$, so
  \eqref{eq:mag-caccioppoli} and Lemma~\ref{lem:mag-rhs} give
  $\nu\,(M^2/4)\iint_{Q(r^*/2)}w\,\mathbf{1} \leq C\,\mathcal{L}\,M^{1/2}$,
  i.e.\ $\iint w\,\mathbf{1} \leq (C\mathcal{L}/\nu)\,M^{-3/2}
  = (C\mathcal{L}/\nu)\,(r^*)^3$.
\end{proof}

\begin{remark}[What the theorem changes]
  Near a Type-I singularity the coherence deficit \emph{cannot be large}:
  it is a priori $O(\mathcal{L}/\nu)$. The Bridge Lemma requires
  $\sigma^*\leq\varepsilon_0$; the gap between ``automatic'' and ``needed''
  is now the explicit factor $C\mathcal{L}/(\nu\varepsilon_0)$, and closing
  it is a \emph{decay} question for the $\sigma^*$-dynamics, not a
  boundedness question. This is the first unconditional (within Type~I)
  quantitative output of the pincer program.
\end{remark}

\subsection{The refined $K$-target: anti-correlation, not smallness [Open, measurable]}

\begin{proposition}[Crude absorption fails by exactly one logarithm]
  \label{prop:crude-K}
  By H\"older and Lemma~\ref{lem:local-enstrophy-typeI},
  \[
    \nu K_w \;=\; \frac{\nu}{M^2}\iint |\nabla m|^2\,w\,\mathbf{1}
    \;\leq\; \|w\|_{L^\infty}\cdot
      \frac{1}{M^2}\Bigl(\nu\iint|\nabla\omega|^2\Bigr)
    \;\leq\; C\,\mathcal{L}\,\|w\|_{L^\infty}\,(r^*)^3 .
  \]
  Since $\|w\|_{L^\infty}(r^*)^3$ is comparable to the quantity the Moser
  scheme outputs, the coupling term re-enters at top order multiplied by
  $\mathcal{L}$: the crude route misses closing the loop by exactly one
  logarithm.
\end{proposition}

\begin{conjecture}[Refined $K$-absorption via equipartition]\label{conj:K-refined}
  Define the coupling correlation on
  $E = Q(r^*)\cap\{|\omega|\geq M/4\}$:
  \[
    c_K \;=\;
    \frac{\displaystyle\iint_E |\nabla m|^2\,w}
         {\displaystyle\Bigl(\fint_E|\nabla m|^2\Bigr)
          \Bigl(\fint_E w\Bigr)\,|E|}
    \;=\;
    \frac{\langle |\nabla m|^2\,w\rangle_E}
         {\langle|\nabla m|^2\rangle_E\,\langle w\rangle_E}.
  \]
  Then $c_K \leq C/\mathcal{L}$ near a putative Type-I singularity.
\end{conjecture}

\begin{remark}[Why equipartition is the mechanism]
  The orthogonal split $|\nabla\omega|^2 = |\nabla m|^2 + m^2 w$ makes
  $|\nabla m|^2$ and $m^2w$ competitors for the same budget: where angular
  variation is maximal, radial variation is squeezed. Perfect equipartition
  or anti-correlation gives $c_K\leq 1$ with room to spare; only strong
  \emph{positive} correlation ($c_K\gg 1$) could defeat the absorption.
  The S33 numerics (adversarial bridge ratio saturating $R\approx 1$
  without violation) already suggest equipartition-like behaviour;
  $c_K$ is directly measurable and is added to the diagnostic suite in S34
  (\texttt{layer4}, EXP-L4-KCORR-*). If the measured $c_K$ is $O(1)$ or
  smaller and stable across ICs and resolutions, the conjecture's
  logarithmic form is plausibly provable from the equipartition structure
  of~\eqref{eq:mag-caccioppoli}, whose sink term penalises exactly the
  configurations with $c_K$ large.
\end{remark}

\subsection{Updated program ledger}

\begin{center}
\begin{tabular}{lll}
\hline
Item & Status & Where \\
\hline
First-rung Caccioppoli (corrected) & Proved mod (i),(ii),(iv) & Prop.~\ref{prop:first-rung} \\
Critical quadratic (was item (iii)) & Standard, needs $\sigma^*\leq\varepsilon_0$ & Rem.~\ref{rem:routine-list} \\
A priori bound $\langle\sigma^*\rangle\leq C\mathcal{L}/\nu$ & \textbf{Proved} (Type I) & Thm.~\ref{thm:apriori-sigma} \\
Crude $K$-absorption & Fails by one log (sharp) & Prop.~\ref{prop:crude-K} \\
Refined $K$-absorption ($c_K\leq C/\mathcal{L}$) & \textbf{Open, measurable} & Conj.~\ref{conj:K-refined} \\
$H2$ (localized CF at $r^*$) & Open & \S\ref{sec:pincer-s32} \\
\hline
\end{tabular}
\end{center}

% ══════════════════════════════════════════════════════════════════════════════
\section{The Second Rung and the $K$-Self-Absorption Identity (Session S35)}
\label{sec:secondrung-s35}

\subsection{The parabolic energy space and the critical quadratic [Proved given the first rung]}

Let $f = |\nabla\hat{e}|\,\eta$ and define the parabolic energy norm on
$Q(\rho')$:
\[
  E_V \;=\; \sup_{t}\int f^2\,dx \;+\; \nu\iint |\nabla f|^2\,dx\,dt,
\]
which is exactly what the (corrected) first rung
(Proposition~\ref{prop:first-rung}) controls, up to the $K$-term:
$E_V \leq C\mathcal{L}\,\sigma^*\,r^* + C\nu K(\rho') + o(r^*)$. By the
standard parabolic embedding $L^\infty_t L^2_x \cap L^2_t H^1_x
\hookrightarrow L^{10/3}(Q)$ in three space dimensions
(Ladyzhenskaya--Solonnikov--Ural'tseva),
$\iint f^{10/3} \leq C\,\nu^{-?}E_V^{5/3}$ with the viscosity bookkeeping
carried explicitly below.

\begin{proposition}[Second rung: the critical quadratic yields to
  $\sigma^*$-smallness]\label{prop:second-rung}
  With $w = |\nabla\hat{e}|^2$, split $w^2 = f^{10/3}\cdot w^{1/3}$
  pointwise on $\{\eta = 1\}$. For any $\delta\in(0,1)$:
  \[
    \nu\iint w^2\eta^2
    \;\leq\;
    \delta\,\nu\,(r^*)^3\,\|w\|_{L^\infty(Q(\rho'))}
    \;+\; C\,\delta^{-1/2}\,\nu\,E_V^{5/2}\,(r^*)^{-3/2}.
  \]
  (Dimensional check in parabolic units: every term scales as $\ell^{1}$.)
  Consequently, if $\sigma^*\leq\varepsilon_0$ and the $K$-term is absorbed
  (\S below), then $E_V \leq C\mathcal{L}\varepsilon_0\,r^*$ and the second
  term is $C\delta^{-1/2}(\mathcal{L}\varepsilon_0)^{5/2}\,r^*$ ---
  superlinear in $\varepsilon_0$, hence harmless; while the first term is
  absorbed into the Moser sup-bound with factor $\delta$. The critical
  quadratic (reclassified item~(iii) of Remark~\ref{rem:routine-list}) is
  thereby handled: \emph{the entire nonlinearity of the direction equation
  is subordinate to $\sigma^*$-smallness}.
\end{proposition}
\begin{proof}
  H\"older with exponents $(3,\tfrac32)$ on
  $w^2\eta^2 = w^{1/3}\cdot f^{10/3}\cdot(\eta\text{-adjust})$ gives
  $\iint w^2\eta^2 \leq \|w\|_{L^\infty(Q(\rho'))}^{1/3}\iint f^{10/3}$.
  Apply the parabolic embedding
  $\iint f^{10/3}\leq C\nu^{-2/3}E_V^{5/3}$ (three spatial dimensions;
  the $\nu^{-2/3}$ tracks the dissipation weight in $E_V$), then Young's
  inequality $a^{1/3}b\leq\delta a + C\delta^{-1/2}b^{3/2}$ with
  $a = \nu\|w\|_\infty(r^*)^3$ and
  $b = \nu^{1/3}\,E_V^{5/3}\,(r^*)^{-1}$; both output terms carry the
  dimensions $\ell^1$ as displayed.
\end{proof}

\subsection{The $K$-self-absorption identity [Derived; two couplings pending
verification]}

The S34 crude route (Proposition~\ref{prop:crude-K}) loses a logarithm
because it decouples $|\nabla m|^2$ from $w$ by H\"older. The following
identity keeps them coupled: \emph{test the magnitude
equation~\eqref{eq:mag-eq} with $m\,w\,\eta^2$ instead of $m\,\eta^2$.}

\begin{proposition}[$K$-self-absorption --- provisional]\label{prop:K-self}
  Testing~\eqref{eq:mag-eq} against $m w\eta^2$ and integrating by parts
  yields
  \begin{equation}\label{eq:K-self-identity}
    \nu\iint |\nabla m|^2\,w\,\eta^2
    \;+\; \nu\iint m^2 w^2\,\eta^2
    \;\leq\;
    C\,\nu\iint m^2\,|\nabla^2\hat{e}|^2\,\eta^2
    \;+\; \mathcal{R},
  \end{equation}
  where $\mathcal{R}$ collects the transport, $\alpha$-, and cutoff terms.
  The first left-hand term is exactly $M^2$ times the density of
  $K_w$; the bound says \emph{the radial--angular coupling is controlled by
  the angular dissipation itself} --- with an absolute constant, no
  logarithm, and no correlation hypothesis. Moreover the second left-hand
  term is the critical quadratic with the favourable weight $m^2\geq M^2/16$,
  reinforcing Proposition~\ref{prop:second-rung}.
\end{proposition}

\begin{proof}[Derivation, with two couplings flagged]
  The viscous term gives
  $\nu\iint \Delta m\cdot m w\eta^2 =
  -\nu\iint|\nabla m|^2 w\eta^2
  -\nu\iint m\,\nabla m\cdot\nabla w\,\eta^2
  -\nu\iint m w\,\nabla m\cdot\nabla(\eta^2)$,
  and the sink term $-\nu m w$ of~\eqref{eq:mag-eq} multiplied by
  $m w \eta^2$ gives $-\nu\iint m^2w^2\eta^2$ on the right, i.e.\ the
  second good term on the left of~\eqref{eq:K-self-identity}. For the
  middle viscous cross-term, $|\nabla w| \leq 2|\nabla^2\hat{e}|\,w^{1/2}$
  pointwise, so by Cauchy--Schwarz
  \[
    \nu\iint m|\nabla m|\,|\nabla w|\,\eta^2
    \;\leq\; \tfrac12\,\nu\iint|\nabla m|^2 w\,\eta^2
    \;+\; 2\,\nu\iint m^2|\nabla^2\hat{e}|^2\eta^2,
  \]
  absorbing half the $K$-density into the left and producing the angular
  dissipation on the right. \textbf{Coupling C1 (pending):} the time term
  $\iint D_t m\cdot m w\eta^2 = \tfrac12\iint D_t(m^2)\,w\eta^2$ requires
  substituting $D_t w$ from the differentiated direction equation after
  integrating by parts in time; its dangerous part is again of the form
  $\nu\iint m^2\,\nabla w\cdot\nabla(\cdot)$ and is expected to close by
  the same Cauchy--Schwarz, but the full bookkeeping (including the
  $2\nu(\nabla\log m\cdot\nabla)w$ term, which re-enters $K$ with a
  constant that must come out $<1$) is deferred to S36.
  \textbf{Coupling C2 (pending):} the stretching term
  $\iint \alpha\,m^2 w\,\eta^2 \leq \|\nabla u\|_{L^\infty} \iint m^2w\eta^2
  \leq C\mathcal{L}M^{1/2}\cdot\nu^{-1}$-type via
  Theorem~\ref{thm:apriori-sigma}; marginal ($\mathcal{L}$-degraded), same
  class as all previous marginal terms, but the constant chase is part of
  the S36 verification. Until C1--C2 are verified, the proposition is
  \emph{provisional}.
\end{proof}

\begin{remark}[S36 amendment to Proposition~\ref{prop:K-self}]
  The verification in \S\ref{sec:verification-s36} found that the second
  left-hand term $\nu\iint m^2w^2\eta^2$ of~\eqref{eq:K-self-identity}
  \emph{cancels exactly} against the $\nu w^2$ feedback contained in
  $D_t w$ (Lemma~\ref{lem:w2-cancel}); the displayed identity must be read
  without it. The $K$-control conclusion stands (with amended constants);
  the critical quadratic is handled solely by the second rung
  (Proposition~\ref{prop:second-rung}), which suffices.
\end{remark}

\begin{remark}[Consequence if C1--C2 verify]\label{rem:if-verified}
  Conjecture~\ref{conj:K-refined} (the $c_K$ decorrelation hypothesis)
  becomes \emph{unnecessary}: $K$-absorption follows structurally
  from~\eqref{eq:K-self-identity}, and H1 (the Bridge Lemma,
  Conjecture~\ref{conj:bridge}) reduces to: (a) $\sigma^*\leq\varepsilon_0$
  at the working scale, plus (b) the routine items (i),(ii),(iv) of
  Remark~\ref{rem:routine-list}. The measured $c_K\approx 1$
  (EXP-L4-KCORR-*) is then explained rather than assumed: equipartition is
  enforced by the magnitude equation's sink, which is exactly what the
  weighted test function exposes.
\end{remark}

\subsection{Honest distance to the Prize}

\begin{remark}[What this program can and cannot yet claim]
  Even with H1 and H2 fully proved, the pincer yields: \emph{Type-I
  singularities require persistent geometric disorder}, and --- if the
  $\sigma^*$-decay gap (from the a priori $C\mathcal{L}/\nu$ of
  Theorem~\ref{thm:apriori-sigma} down to $\varepsilon_0$) is closed ---
  \emph{no Type-I blowup}. Two further steps separate that from the Clay
  problem: (1) the $\sigma^*$-decay mechanism (open; the S31 audit removed
  the previous candidate), and (2) removal of the Type-I hypothesis
  entirely (wide open; every marginal estimate in
  \S\ref{sec:moser-s33}--\S\ref{sec:secondrung-s35} uses Type-I scaling,
  so Type-II exclusion requires genuinely new input). Global existence of
  weak solutions is classical (Leray); the open question is smoothness,
  and this document claims no more than the ledger above.
\end{remark}

\subsection{Updated program ledger (post-S35)}

\begin{center}
\begin{tabular}{lll}
\hline
Item & Status & Where \\
\hline
Second rung (critical quadratic) & Proved given first rung + $\varepsilon_0$ & Prop.~\ref{prop:second-rung} \\
$K$-self-absorption identity & Derived; C1--C2 pending (S36) & Prop.~\ref{prop:K-self} \\
$c_K$ hypothesis & Superseded if C1--C2 verify & Rem.~\ref{rem:if-verified} \\
H1 (Bridge Lemma) & Reduces to $\varepsilon_0$ + routine, mod C1--C2 & Conj.~\ref{conj:bridge} \\
$\sigma^*$-decay ($C\mathcal{L}/\nu \to \varepsilon_0$) & Open & Thm.~\ref{thm:apriori-sigma} \\
H2 (localized CF) & Open & \S\ref{sec:pincer-s32} \\
Type-II exclusion & Open (outside current methods) & --- \\
\hline
\end{tabular}
\end{center}

% ══════════════════════════════════════════════════════════════════════════════
\section{Verification of C1--C2, Gram Negativity, and the Annulus Residual (Session S36)}
\label{sec:verification-s36}

This section carries out the verification of couplings C1--C2 flagged in
Proposition~\ref{prop:K-self}, discharges the routine items (i), (ii), (iv)
of Remark~\ref{rem:routine-list}, and records one correction to \S\ref{sec:secondrung-s35}.

\subsection{C2: the $\alpha$-weighted terms [Verified]}

Every term of the form $\iint |\alpha|\,m^2 w\,\eta^2$ or
$M\mathcal{L}$-weighted variants obeys, by
$|\alpha|\leq\|\nabla u\|_{L^\infty}\leq CM\mathcal{L}$ and the a priori
bound $\nu\iint m^2w\eta^2 \leq C\mathcal{L}M^{1/2}$ (the proof of
Theorem~\ref{thm:apriori-sigma}):
\[
  \iint|\alpha|m^2w\eta^2 \;\leq\; CM\mathcal{L}\cdot\nu^{-1}\cdot
  C\mathcal{L}M^{1/2} \;=\; C\,\mathcal{L}^2\,M^{3/2}/\nu .
\]
In $K$-units (divide by $\nu M^2$) this is $C\mathcal{L}^2 r^*/\nu^2$ ---
an inhomogeneous term of exactly the allowed
$C_K(1+\sigma^*)r^*$ class (with $\mathcal{L}^2$-degraded constant).
\emph{C2 is verified.}

\subsection{C1: the $D_t w$ substitution [Verified modulo the annulus residual]}

Substituting the $w$-evolution (obtained from~\eqref{eq:dir-eq} as in
\S\ref{sec:moser-s33}) into $\tfrac12\iint m^2(D_tw)\eta^2$ produces five
term families. Three structural lemmas organise them.

\begin{lemma}[Exact cancellation of the critical quadratic]\label{lem:w2-cancel}
  The term $+\nu w\hat{e}$ in~\eqref{eq:dir-eq} contributes exactly
  $+\nu w^2$ to $\tfrac12 D_t w$ (since
  $\sum_k \partial_k(w\hat{e})\cdot\partial_k\hat{e}
  = w\sum_k|\partial_k\hat{e}|^2 = w^2$, the $\partial_k w\,
  (\hat{e}\cdot\partial_k\hat{e})$ part vanishing). Hence
  $\tfrac12\iint m^2(D_tw)\eta^2$ contains $+\nu\iint m^2w^2\eta^2$, which
  cancels the sink term $-\nu\iint m^2w^2\eta^2$ of the weighted magnitude
  test identically. The second left-hand term
  of~\eqref{eq:K-self-identity} is therefore removed; no conclusion of the
  program depended on it (the critical quadratic is handled by
  Proposition~\ref{prop:second-rung}).
\end{lemma}

\begin{lemma}[Gram negativity of the drift-curvature terms]\label{lem:gram-negativity}
  Writing $G_{jk} = \partial_j\hat{e}\cdot\partial_k\hat{e}$ (positive
  semidefinite Gram matrix) and using
  $\partial_j\partial_k\log m = \tfrac{\partial_j\partial_k m}{m}
  - \tfrac{\partial_j m\,\partial_k m}{m^2}$, the drift-curvature
  contribution splits as
  \[
    2\nu\iint m^2(\partial_j\partial_k\log m)\,G_{jk}\,\eta^2
    \;=\; 2\nu\iint m\,(\partial_j\partial_k m)\,G_{jk}\,\eta^2
    \;-\; 2\nu\iint (\partial_jm\,\partial_km)\,G_{jk}\,\eta^2 .
  \]
  The second integral is the contraction of the rank-one PSD tensor
  $\nabla m\otimes\nabla m$ against $G\succeq 0$, hence
  \emph{non-negative}; with its minus sign it is a good term. Integrating
  the first by parts,
  $-2\nu\iint\partial_jm\,\partial_k\bigl[m\,G_{jk}\eta^2\bigr]$, its
  leading piece is \emph{again} $-2\nu\iint(\partial_jm\,\partial_km)G_{jk}\eta^2\leq 0$;
  the remaining pieces are
  $\leq \tfrac{\nu}{8}\iint|\nabla m|^2w\eta^2 + C\nu M^2\iint|\nabla^2\hat{e}|^2\eta^2$
  (Cauchy--Schwarz via $|\partial_k G_{jk}|\leq 2|\nabla^2\hat{e}||\nabla\hat{e}|$)
  plus cutoff terms in the classes below. \emph{The most dangerous
  couplings of C1 are good-signed, not merely bounded.}
\end{lemma}

\begin{lemma}[Stretching and commutator couplings]\label{lem:stretch-couple}
  The families $\iint m^2\bigl(|\nabla^2 u||\nabla\hat{e}|
  + |\nabla u|\,w\bigr)\eta^2$ close as in
  Lemma~\ref{lem:BS-split}: the $|\nabla u|w$ part is
  $\leq C\mathcal{L}^2M^{3/2}/\nu$ by the a priori bound; the
  $|\nabla^2u||\nabla\hat{e}|$ part, after the near/far Biot--Savart split
  and Cauchy--Schwarz, is
  $\leq C\,\mathcal{L}^{1/2}(\sigma^*)^{1/2}\nu^{-1/2}M^{3/2}$.
  Both are inhomogeneous terms of the allowed class.
\end{lemma}

\begin{proposition}[Amended $K$-self-absorption]\label{prop:K-amended}
  For smooth Type-I solutions, with $D_{\mathrm{ang}} =
  \nu\iint|\nabla^2\hat{e}|^2\eta^2$:
  \[
    \nu\iint|\nabla m|^2\,w\,\eta^2
    \;\leq\; C\,M^2\,D_{\mathrm{ang}}
    \;+\; C\,\mathcal{L}^2\,M^{3/2}\,\bigl(\nu^{-1}
      + \nu^{-1/2}(\sigma^*)^{1/2}\bigr)
    \;+\; \mathcal{A}_{\mathrm{ann}},
  \]
  where $\mathcal{A}_{\mathrm{ann}}$ collects the transition-annulus terms
  enumerated in Definition~\ref{def:annulus-residual} below. Consequently
  the coupling $K$ is absorbed --- with an absolute constant against the
  angular dissipation, no logarithm and no correlation hypothesis ---
  \emph{modulo} $\mathcal{A}_{\mathrm{ann}}$.
\end{proposition}

\subsection{The annulus residual [the single remaining technical item of H1's estimate layer]}

\begin{definition}[Annulus residual]\label{def:annulus-residual}
  $\mathcal{A}_{\mathrm{ann}}$ consists of the cutoff-derivative terms
  supported on the transition annulus
  $\mathrm{Ann} = \{M/8\leq|\omega|\leq M/4\}\cap\operatorname{supp}\psi$:
  \[
    \nu\iint_{\mathrm{Ann}}|\nabla m|^2\,w\,\psi\,|\chi'| ,
    \qquad
    \nu\iint_{\mathrm{Ann}}m^2 w^2\,\psi\,|\chi'| ,
  \]
  and lower-order variants, all weighted by $\eta$-\emph{incomplete}
  cutoffs (weight $\psi|\chi'|$, not $\eta^2$), with constants
  $\leq 16$ per occurrence.
\end{definition}

\begin{remark}[Level-iteration scheme --- stated, to be executed]
  \label{rem:level-iteration}
  $\mathcal{A}_{\mathrm{ann}}$ is the classical residue of level-set
  truncation and is handled by the standard De Giorgi device: replace the
  single pair of levels $(\tfrac18,\tfrac14)$ by a nested sequence
  $\theta_k = \tfrac18 + 2^{-k-3}$, $k=0,1,2,\dots$, with cutoffs
  $\chi_k$ at levels $(\theta_{k+1},\theta_k)$; each annulus term at level
  $k$ is dominated by the full-weight quantities at level $k+1$ with
  constant $C\,4^{k}$, and the iteration closes by choosing the level-$k$
  Young parameters $\delta_k = \delta_0 8^{-k}$ (geometric beats
  geometric). This is bookkeeping of the same class as
  Ladyzhenskaya--Uraltseva Chapter~II; it involves no new estimate. Until
  written out it is recorded as \emph{open-mechanical}, distinct in kind
  from the genuinely open items (H2, $\sigma^*$-decay, Type-II).

  \medskip
  \noindent\textbf{\color{red}Superseded (S41).} The scheme sketched in this
  remark was executed in \S\ref{sec:annulus-s41} and \emph{does not close}.
  The recursion it produces has coefficient $\geq2^{8}$ in front of the
  level-$(k+1)$ quantity for every choice of $\delta_k$
  (Proposition~\ref{prop:s41-nocontraction}); the stated choice
  $\delta_k=\delta_08^{-k}$ makes the coefficient grow like $32^{k}$ rather
  than $4^{k}$; and by a co-area identity no admissible level selection can
  reduce the residual below a level density of the underlying measure
  (Proposition~\ref{prop:s41-coarea}). The classification
  ``open-mechanical'' in the last sentence above is therefore withdrawn: see
  Remark~\ref{rem:s41-summary} and Conjecture~\ref{target:band-absorption}.
\end{remark}

\subsection{Routine items (i), (ii), (iv) discharged [Proved]}

\begin{lemma}[Item (i): transport]
  $\bigl|\iint w\,u\cdot\nabla(\eta^2)\bigr|
  \leq \tfrac{C_I}{r^*}\iint w\,\bigl(2\eta|\nabla\psi|\chi
  + \tfrac{16}{M}\eta\psi|\chi'||\nabla m|\bigr)
  \leq \tfrac{C\,C_I}{(\rho'-\rho)\,r^*}\iint_{Q(\rho')} w\,\mathbf{1}
  + \varepsilon\,\nu\,\frac{1}{M^2}\iint|\nabla m|^2w\,\mathbf{1}
  + C(\varepsilon,\nu)\,\frac{C_I^2}{(r^*)^2}\iint w\,\mathbf{1}$,
  by Young; the middle term joins the $K$-family
  (Proposition~\ref{prop:K-amended}), the others the
  $(\rho'-\rho)^{-2}$-family of Proposition~\ref{prop:first-rung}. \qed
\end{lemma}

\begin{lemma}[Item (ii): commutator]
  $[\partial_k, D_t]\hat{e} = -(\partial_ku\cdot\nabla)\hat{e}$
  contributes $\leq \iint|\nabla u|\,w\,\eta^2
  \leq CM\mathcal{L}\iint w\eta^2
  = C\mathcal{L}(r^*)^{-2}\iint w\eta^2$ --- the marginal-potential class
  already present in Proposition~\ref{prop:first-rung}. \qed
\end{lemma}

\begin{lemma}[Item (iv): time slice]
  With $\psi$ vanishing on $\{t\leq t_0-(\rho')^2\}$ the lower time-slice
  term is zero and the upper slice enters the left-hand side with positive
  sign. \qed
\end{lemma}

\subsection{Consolidated status of H1 after S36}

\begin{center}
\begin{tabular}{lll}
\hline
Component of H1 & Status & Where \\
\hline
First rung (angular Caccioppoli) & Proved (items (i),(ii),(iv) now discharged) & Prop.~\ref{prop:first-rung} \\
Critical quadratic & Proved given $\sigma^*\leq\varepsilon_0$ & Prop.~\ref{prop:second-rung} \\
$K$-absorption & Proved modulo $\mathcal{A}_{\mathrm{ann}}$ & Prop.~\ref{prop:K-amended} \\
C2 ($\alpha$-terms) & Verified & this section \\
Gram negativity & Proved (good signs) & Lem.~\ref{lem:gram-negativity} \\
$w^2$ cancellation & Proved (corrects \S\ref{sec:secondrung-s35}) & Lem.~\ref{lem:w2-cancel} \\
$\mathcal{A}_{\mathrm{ann}}$: families (G1)--(G4) & Proved good-signed, discarded (S41) & Prop.~\ref{prop:s41-core} \\
$\mathcal{A}_{\mathrm{ann}}$: families (I1)--(I2) & Proved in the allowed class (S41) & Prop.~\ref{prop:s41-core} \\
$\mathcal{A}_{\mathrm{ann}}$: core (S1)+(S2) & \textbf{Open} --- level iteration \emph{fails} & \S\ref{sec:annulus-s41}, Conj.~\ref{target:band-absorption} \\
\hline
\end{tabular}
\end{center}

\noindent\textbf{Summary} \emph{(amended in S41; see
\S\ref{sec:annulus-s41})}. H1 (the Bridge Lemma) rests on:
(a)~$\sigma^*\leq\varepsilon_0$ at the working scale;
(b)~absorption of the annulus residual $\mathcal{A}_{\mathrm{ann}}$.
Item~(b) was classified in Remark~\ref{rem:level-iteration} as
\emph{open-mechanical}, on the expectation that the nested level iteration
sketched there would dispose of it. \textbf{That expectation was wrong.}
\S\ref{sec:annulus-s41} executes the scheme and proves that it cannot
close: the recursion coefficient is $\geq2^{8}$ for every choice of the
Young parameters (Proposition~\ref{prop:s41-nocontraction}), the proposed
$\delta_k=\delta_08^{-k}$ makes it worse, the unrolled iteration is weaker
than trivial (Proposition~\ref{prop:s41-vacuous}), and by a co-area
identity \emph{no} level selection whatsoever gains anything
(Proposition~\ref{prop:s41-coarea}). What S41 does establish is that six of
the eight cutoff-derivative families vanish by sign or fall in the allowed
inhomogeneous class, leaving an explicit core with an unconditional bound
that falls short of absorption by exactly one logarithm --- the same
shortfall as Proposition~\ref{prop:crude-K}. Item~(b) is therefore
reclassified as \emph{open}, of the same analytic type as
Conjecture~\ref{conj:K-refined}, and is restated as
Conjecture~\ref{target:band-absorption}. Everything else in H1's estimate
layer is proved. The genuinely open items of the program remain H2,
$\sigma^*$-decay, and Type-II exclusion, as ledgered in
\S\ref{sec:secondrung-s35}.

% ══════════════════════════════════════════════════════════════════════════════
\section{Localized Constantin--Fefferman Depletion and the Unification of H2 with $\sigma^*$-Decay (Session S37)}
\label{sec:cf-s37}

This section derives H2 (localized depletion at the self-similar scale) as a
quantitative consequence of H1 (the Bridge Lemma) via the classical
Constantin geometric-kernel representation, and identifies precisely when
this derivation succeeds unconditionally versus when it re-invokes the
$\sigma^*$-decay gap --- the central finding of this session.

\subsection{The geometric kernel and its exact cancellation [Classical]}

\begin{fact}[Constantin's stretching formula \cite{Constantin1994}]\label{fact:constantin-formula}
  For a divergence-free field on $\T^3$ (locally, treating $y$-integration
  over $\R^3$ up to box-scale corrections), the stretching rate at $x^*$
  admits the representation
  \[
    \alpha(x^*,t) \;=\; \frac{3}{4\pi}\;\mathrm{p.v.}\!\int
      D\bigl(\hat{e}(x^*),\,\hat{e}(x^*-y),\,y/|y|\bigr)\,
      \frac{|\omega(x^*-y,t)|}{|y|^3}\,dy,
  \]
  where $D(e_1,e_2,\sigma) = \bigl(e_1\cdot(\sigma\times e_2)\bigr)(e_1\cdot\sigma)$
  is smooth on $(S^2)^3$. Since $\sigma\times e_1 \perp e_1$ for every
  $\sigma$, one has the \emph{pointwise algebraic identity}
  \[
    D(e,e,\sigma) \;=\; 0 \qquad\text{for every }e,\sigma\in S^2.
  \]
\end{fact}

\subsection{Near/far split and optimized cutoff [Proved, classical technique]}

\begin{proposition}[Localized CF estimate]\label{prop:localized-cf}
  Suppose $\hat{e}$ is Lipschitz with constant $K$ on
  $B_\rho(x^*)\cap\{|\omega|\geq M/2\}$ for some $\rho\leq r^*/2$, i.e.\
  $|\hat{e}(x^*-y)-\hat{e}(x^*)|\leq K|y|$ there. Let
  $\Omega = \|\omega(\cdot,t)\|_{L^2(\T^3)}^2$. Then for every
  $0<\rho'\leq\rho$:
  \[
    |\alpha(x^*)| \;\leq\; C\,K\,M\,\rho' \;+\; C\,\Omega^{1/2}\,(\rho')^{-3/2},
  \]
  and optimizing over $\rho'$ (valid whenever the minimizer
  $\rho'_* = C(\Omega^{1/2}/(KM))^{2/5}$ satisfies $\rho'_*\leq\rho$):
  \begin{equation}\label{eq:cf-optimized}
    |\alpha(x^*)| \;\leq\; C\,(KM)^{3/5}\,\Omega^{1/5}.
  \end{equation}
\end{proposition}
\begin{proof}
  \textbf{Near field} ($|y|<\rho'$): by
  Fact~\ref{fact:constantin-formula} and Lipschitz continuity of the
  second argument of $D$ (smooth on the compact $(S^2)^3$, hence
  Lipschitz in each slot with a universal constant),
  \[
    |D(\hat{e}(x^*),\hat{e}(x^*-y),\sigma)|
    = |D(\hat{e}(x^*),\hat{e}(x^*-y),\sigma) - D(\hat{e}(x^*),\hat{e}(x^*),\sigma)|
    \leq C\,K|y| \qquad(|y|<\rho'\leq\rho).
  \]
  Hence $|\alpha_{\mathrm{near}}| \leq C K\int_{|y|<\rho'}
  |\omega(x^*-y)|/|y|^2\,dy \leq CKM\int_{|y|<\rho'}|y|^{-2}dy = CKM\rho'$.

  \textbf{Far field} ($|y|\geq\rho'$): by Cauchy--Schwarz,
  $\int_{|y|\geq\rho'}|\omega(x^*-y)|/|y|^3\,dy
  \leq \bigl(\int_{|y|\geq\rho'}|y|^{-6}dy\bigr)^{1/2}\|\omega\|_{L^2}
  = C(\rho')^{-3/2}\,\Omega^{1/2}$
  (using $\int_{\rho'}^\infty r^{-6}\cdot 4\pi r^2\,dr = C(\rho')^{-3}$).

  \textbf{Optimization.} $f(\rho')=A\rho'+B(\rho')^{-3/2}$ with
  $A=CKM$, $B=C\Omega^{1/2}$ is minimized at
  $\rho'_*=(3B/2A)^{2/5}$, giving $f(\rho'_*)\sim A^{3/5}B^{2/5}
  = C(KM)^{3/5}\Omega^{1/5}$.
\end{proof}

\subsection{Two regimes: the exponent depends on the growth of $\Omega$ [Central finding]}

Feed the sharp Bridge Lemma output $K\leq C_B(\sigma^*)^{1/2}/r^* =
C_B(\sigma^*)^{1/2}M^{1/2}$ into~\eqref{eq:cf-optimized}:
\[
  |\alpha(x^*)| \;\leq\; C\,(\sigma^*)^{3/10}\,M^{9/10}\,\Omega^{1/5}.
\]

\begin{proposition}[Regime A: subcritical enstrophy $\Rightarrow$ unconditional depletion]
  \label{prop:regime-A}
  If $\Omega(t) = O(1)$ as $t\to T^-$ (enstrophy stays bounded), then
  \[
    |\alpha(x^*)| \;\leq\; C\,(\sigma^*)^{3/10}\,M^{9/10},
  \]
  i.e.\ $\theta = 9/10 - 1/2 = 2/5 < 1/2$: \textbf{H2 holds unconditionally
  given H1} (with the a priori bound $\sigma^*\leq C\mathcal{L}/\nu$ of
  Theorem~\ref{thm:apriori-sigma} supplying a finite, not necessarily
  small, prefactor).
\end{proposition}

\begin{proposition}[Regime B: self-similar enstrophy $\Rightarrow$ borderline]
  \label{prop:regime-B}
  If $\Omega(t)\sim M(t)^{1/2}$ (the scaling of a genuine self-similar
  Type-I profile: $\omega(x,t)\sim(T-t)^{-1}\Omega_0((x-x^*)/\sqrt{T-t})$
  gives $\|\omega(\cdot,t)\|_{L^2}^2\sim(T-t)^{-1/2}\sim M(t)^{1/2}$), then
  \[
    |\alpha(x^*)| \;\leq\; C\,(\sigma^*)^{3/10}\,M^{9/10}\cdot M^{1/10}
    \;=\; C\,(\sigma^*)^{3/10}\,M,
  \]
  which is \textbf{exactly borderline} ($\theta=1/2$): the near/far
  optimization gains nothing beyond a constant prefactor $(\sigma^*)^{3/10}$.
  Strict depletion in Regime B requires $\sigma^*(t)\to 0$, not merely
  $\sigma^*$ bounded.
\end{proposition}

\begin{corollary}[Unification of H2 with $\sigma^*$-decay]\label{cor:unification}
  In the generic case where the local enstrophy near a putative Type-I
  singularity grows at least at the self-similar rate $M^{1/2}$ (Regime B
  or worse), \textbf{H2's remaining content and the $\sigma^*$-decay gap
  are the same open problem}: both require $\sigma^*(t)\to 0$ as $t\to T^-$,
  not merely the a priori bound $\sigma^*\leq C\mathcal{L}/\nu$
  (Theorem~\ref{thm:apriori-sigma}). The program's genuinely open
  analytical content therefore reduces from three nominally independent
  items (H1's annulus step, H2, $\sigma^*$-decay) to \emph{two}: the
  annulus level-iteration (mechanical, \S\ref{sec:verification-s36}) and
  \emph{a single deep gap}: a mechanism forcing $\sigma^*(t)\to 0$ along
  any putative Type-I blowup sequence. Regime A remains available as an
  unconditional fallback if a genuine (non-self-similar, subcritical)
  enstrophy bound can be established near the singularity --- itself an
  interesting alternative target, since it does not require touching
  $\sigma^*$ at all.
\end{corollary}

\subsection{Numerical status}

The measured local enstrophy growth exponent relative to $M(t)$ (whether
the data sits closer to Regime A or Regime B) is the natural next
diagnostic; see \texttt{layer4}, EXP-L4-ENSTROPHY-EXPONENT-* (Session S37).
The existing depletion diagnostic (EXP-L4-DEPLETION-TG-001) already shows
$\theta_{\mathrm{eff}}$ centred near $-0.08$ (median) in the $M>10$ window
--- i.e.\ an \emph{observed} exponent of $\alpha$ near $0.42$ in $M$,
below both the naive value ($1$) and the Regime B borderline value ($1$),
and close to but somewhat above the Regime A prediction ($9/10$ once
$\Omega$-growth is folded in would sit between the two regimes) ---
consistent with the enstrophy growth in this (non-singular, decaying)
numerical trajectory being sub-self-similar, as expected away from an
actual singularity.

\subsection{Updated program ledger (post-S37)}

\begin{center}
\begin{tabular}{lll}
\hline
Item & Status & Where \\
\hline
Localized CF estimate (near/far optimized) & Proved (classical technique) & Prop.~\ref{prop:localized-cf} \\
H2 in Regime A (bounded $\Omega$) & \textbf{Unconditional given H1} & Prop.~\ref{prop:regime-A} \\
H2 in Regime B (self-similar $\Omega$) & Borderline; needs $\sigma^*\to 0$ & Prop.~\ref{prop:regime-B} \\
H2 $\equiv$ $\sigma^*$-decay (Regime B) & \textbf{Unified}: one open problem, not two & Cor.~\ref{cor:unification} \\
H1 annulus step & Open-mechanical & \S\ref{sec:verification-s36} \\
$\sigma^*$-decay (the single deep gap) & \textbf{Open} & Thm.~\ref{thm:apriori-sigma} \\
Type-II exclusion & Open (outside current methods) & \S\ref{sec:secondrung-s35} \\
\hline
\end{tabular}
\end{center}

% ══════════════════════════════════════════════════════════════════════════════
\section{Regime A Is Impossible for Genuine Blowup; Numerical $\sigma^*$ Trend; the Closure Target (Session S38)}
\label{sec:closure-s38}

\subsection{Regime A cannot occur: enstrophy blows up at any genuine singularity [Proved]}

\begin{proposition}[Impossibility of bounded enstrophy at a genuine singularity]
  \label{prop:regime-A-impossible}
  Let $u_0\in H^1(\T^3)$ and let $T<\infty$ be the maximal time of existence
  of the associated strong (smooth) solution, i.e.\ $T$ is a genuine
  blowup time: the solution cannot be continued as a strong $H^1$ solution
  past $T$. Then
  \[
    \limsup_{t\to T^-} \Omega(t) \;=\; \infty, \qquad
    \Omega(t) := \|\omega(\cdot,t)\|_{L^2(\T^3)}^2.
  \]
  Consequently \textbf{Regime A} of Proposition~\ref{prop:regime-A}
  ($\Omega=O(1)$) is impossible for any genuine singularity: the only
  route to H2 through Proposition~\ref{prop:localized-cf} is Regime B (or
  worse), and the unification of \S\ref{sec:cf-s37} (Corollary~\ref{cor:unification})
  is not merely the generic case --- it is the \emph{only} case.
\end{proposition}
\begin{proof}
  This is the standard blow-up criterion for the $H^1$-Cauchy theory of 3D
  Navier--Stokes, included for completeness. Local well-posedness in
  $H^1(\T^3)$ (Fujita--Kato / standard energy method) gives: for every
  $R>0$ there is $T_0(R)>0$, non-increasing in $R$, such that any strong
  solution with $\|u(t_1)\|_{H^1}\leq R$ at some time $t_1$ extends as a
  strong solution to $[t_1,t_1+T_0(R)]$. Suppose, for contradiction, that
  $\limsup_{t\to T^-}\|u(t)\|_{H^1} = R_0 <\infty$. Choose $t_1<T$ with
  $T-t_1 < T_0(2R_0)$ and $\|u(t_1)\|_{H^1}\leq 2R_0$ (possible since the
  limsup is finite and $T_0$ is a fixed positive number depending only on
  $2R_0$). Then $u$ extends strongly to $[t_1, t_1+T_0(2R_0)] \supset
  [t_1,T]$, and in fact strictly past $T$, contradicting maximality of
  $T$. Hence $\limsup_{t\to T^-}\|u(t)\|_{H^1}=\infty$. Since
  $\|u(t)\|_{L^2}\leq\|u_0\|_{L^2}$ (energy inequality, non-increasing),
  the unboundedness must come from $\|\nabla u(t)\|_{L^2}$, i.e.\
  $\limsup_{t\to T^-}\|\nabla u(t)\|_{L^2}=\infty$. By the mean-zero
  Biot--Savart isomorphism on $\T^3$, $\|\nabla u\|_{L^2}\sim\|\omega\|_{L^2}$
  with constants depending only on the domain, giving
  $\limsup_{t\to T^-}\Omega(t)=\infty$.
\end{proof}

\begin{remark}
  This closes off, cleanly and unconditionally, the ``local Regime A''
  direction previously floated as an independent alternative to
  $\sigma^*$-decay: no version of bounded-enstrophy escape is available at
  a genuine singularity, local or global, since the argument above uses
  only $\|\nabla u\|_{L^2}$ on the whole torus (the domain has no smaller
  scale to exploit). $\sigma^*$-decay is not one of several possible
  routes to H2 --- given H1, it is necessary.
\end{remark}

\subsection{Numerical $\sigma^*$ trend: non-monotonic, not a generic free decay}

The instantaneous $\sigma^*(t)$ was already logged (EXP-L4-BRIDGE-TG-002,
$N=64$, $\nu=10^{-3}$) across the full stretching-then-decay trajectory. Read
against $M(t)$ (which rises $2.0\to20.9$ over the growth phase, steps
0--190): $\sigma^*$ is \emph{not} monotone in $M$. It starts small
($\sigma^*\approx0.5$ at $M\approx2.0$), rises sharply through the
turbulent transition to a mid-trajectory maximum ($\sigma^*\approx44$ at
$M\approx2.5$--$3.5$), then \emph{partially recedes} as $M$ climbs toward
its true maximum ($\sigma^*\approx15$--$27$ at $M\approx15$--$21$, steps
100--190 and 190) --- markedly below the mid-trajectory peak. Past the
peak of $M$, as the flow relaxes viscously, $\sigma^*$ grows again into
the hundreds while $M$ falls, but this late-time regime is suspected to be
a definitional artifact of a decaying, noise-dominated field rather than a
physically meaningful coherence measurement: as $M$ falls, the threshold
set $\{|\omega|\geq M/4\}$ becomes less selective and increasingly picks
up turbulent debris.

\begin{remark}[Honest reading]
  This single, noisy trajectory does not show $\sigma^*(t)\to0$ as $M$
  grows --- values in the growth phase remain $O(10)$--$O(50)$, far from
  any plausible smallness threshold $\varepsilon_0$. It also does not show
  monotone growth: the partial recession near peak $M$ is real in this
  data but should not be over-read (one run, $N=64$, no ensemble
  averaging). The honest conclusion is that ``generic'' (non-singular)
  vortex dynamics does not naturally sit in, or approach, a small-$\sigma^*$
  regime. This is consistent with --- and does not contradict --- the
  disorder-barrier picture of \S\ref{sec:pincer-s32}: if blowup requires
  $\sigma^*$ small, and generic turbulent dynamics keeps $\sigma^*$ large,
  then a genuine singularity (if one exists) would have to be a
  non-generic, specially aligned configuration, which is exactly the
  content the Constantin--Fefferman lineage of theorems is built to
  express. It does not, however, supply evidence that such configurations
  cannot occur; the numerics measures ordinary dynamics, which by
  construction cannot sample a genuine singularity.
\end{remark}

\subsection{The closure target: assembling the proved ingredients [Open, precisely stated]}

\begin{remark}[Not yet a theorem]
Propositions~\ref{prop:first-rung} (first rung), \ref{prop:second-rung}
(second rung), and \ref{prop:K-amended} (amended $K$-absorption, modulo
the annulus residual of \S\ref{sec:verification-s36}) together bound the
angular energy at scale $\rho$ by a combination of (i) the same quantity
at a larger scale $\rho'$, with coefficient $C\mathcal{L}/(\rho'-\rho)^2$,
and (ii) $\sigma^*$-dependent inhomogeneous terms. Whether this chain, once
fully assembled and non-dimensionalized, yields a \emph{subcritical} scale
recursion (effective coefficient $<1$, giving genuine self-improving decay
as $\rho\to0$, independent of the size of $\sigma^*$ at the starting
scale) or remains \emph{critical/supercritical} (coefficient $\geq1$,
requiring a priori smallness of $\sigma^*$ --- the same obstruction that
defeated \S19 and \S20) has \textbf{not been determined}. A preliminary
scaling check on the leading terms of Proposition~\ref{prop:K-amended}
(the $M^2 D_{\mathrm{ang}}$ term against the first-rung dissipation) shows
consistent powers of $M$, which is necessary but not sufficient for a
favourable closure; the full coefficient bookkeeping, including the
annulus-iteration contribution, is deferred to a dedicated session. This
is recorded as an open technical target, not a theorem, precisely to avoid
repeating the error pattern of \S\ref{sec:audit-s31}.
\end{remark}

\subsection{A pointer to Type-I compactness literature}

Genuine progress on $\sigma^*$-decay near an \emph{actual} singularity (as
opposed to generic dynamics, which Proposition~\ref{prop:regime-A-impossible}
and the numerics above show is the wrong object to study) most plausibly
requires the Type-I blowup-profile compactness framework: rescaled
solutions $u_\lambda(y,s) = \lambda u(x^*+\lambda y, T+\lambda^2 s)$
converge, along a subsequence as $\lambda\to0$, to a nontrivial ancient
(or self-similar) solution under a Type-I bound
(Ne\v{c}as--R\r{u}\v{z}i\v{c}ka--\v{S}ver\'ak \cite{NecasRuzickaSverak1996};
Escauriaza--Seregin--\v{S}ver\'ak \cite{EscauriazaSereginSverak2003}). Such
limiting profiles are exactly the ``special, non-generic'' configurations
the disorder-barrier picture requires; whether every such limit profile
must have $\sigma^*=0$ (perfect coherence) is a natural, sharply-posed
question for that framework, distinct from anything the direct estimates
of \S\ref{sec:moser-s33}--\S\ref{sec:cf-s37} can settle alone. This is
noted as the natural longer-range tool, not attempted here.

\subsection{Updated program ledger (post-S38)}

\begin{center}
\begin{tabular}{lll}
\hline
Item & Status & Where \\
\hline
Regime A (bounded $\Omega$) & \textbf{Proved impossible} at genuine blowup & Prop.~\ref{prop:regime-A-impossible} \\
$\sigma^*$-decay is necessary (not one of several routes) & Established & Prop.~\ref{prop:regime-A-impossible} \\
Numerical $\sigma^*$ trend (generic dynamics) & Non-monotonic; stays $O(10)$--$O(50)$ & this section \\
Assembled scale-recursion (sub- vs supercritical) & \textbf{Open, unresolved} & this section \\
Type-I profile compactness route & Pointed to, not attempted & this section \\
H1 annulus step & Open-mechanical & \S\ref{sec:verification-s36} \\
Type-II exclusion & Open (outside current methods) & \S\ref{sec:secondrung-s35} \\
\hline
\end{tabular}
\end{center}

% ══════════════════════════════════════════════════════════════════════════════
\section{Profile Rigidity: Converting $\sigma^*$-Decay into a Liouville Problem (Session S39)}
\label{sec:rigidity-s39}

Corollary~\ref{cor:unification} and Proposition~\ref{prop:regime-A-impossible}
together reduce the program's genuinely open analytical content to a
\emph{single} deep item: forcing $\sigma^*(t)\to 0$ along a putative Type-I
blowup sequence. Every direct dynamical attempt on that item has failed in the
same way (\S\ref{sec:audit-s31}, Defects 1--3, and the ``preliminary scaling
check'' remark of \S\ref{sec:closure-s38}): one derives a scale recursion whose
effective coefficient is critical or supercritical, so that contraction
requires the very smallness one is trying to produce.

This section develops a different strategy, following the pointer left at the
end of \S\ref{sec:closure-s38}. Instead of asking the \emph{dynamical} question
``does $\sigma^*$ decay?'', one uses Type-I blowup-profile compactness to pass
to a limit object --- an \emph{ancient} solution $U$ on $\R^3\times(-\infty,0)$
--- and asks the \emph{structural} question ``what does $\sigma^*$ do on
$U$?''. The value of the reframe is not that it makes the problem easy; it is
that the two possible answers ($\sigma^*(U)=0$ and $\sigma^*(U)>0$) are attacked
by two \emph{completely different and independently developed} bodies of
technique --- Liouville rigidity in the first case, infinite-past accounting in
the second --- neither of which is a scale recursion, and neither of which can
therefore fail in the S31 manner.

Every statement below is labelled. Two named inputs, \textbf{(R)} and
\textbf{(N)}, are isolated in \S\ref{subsec:s39-step1}; a third, \textbf{(A-up)},
in \S\ref{subsec:s39-step3}. Nothing in this section claims the Prize, and
nothing in it removes the Type-I hypothesis.

\subsection{Conventions and dimensional normalization}
\label{subsec:s39-conventions}

\begin{remark}[Units]\label{rem:s39-units}
The parabolic cylinder convention $Q(r) = B_r\times[t_0-r^2,t_0]$ used
throughout \S\ref{sec:moser-s33}--\S\ref{sec:closure-s38} is dimensionally
consistent only in units where $\nu=1$. \emph{We therefore fix $\nu=1$ for the
whole of this section} and record, for the dimensional audit, the following
homogeneities in physical units ($[\,\cdot\,]$ denoting dimension, $L$ length,
$\mathcal{T}$ time):
\[
  [\sigma^*] = 1,\qquad
  [\langle\sigma^*\rangle] = \mathcal{T}L^{-2} = [\nu]^{-1},\qquad
  [\nu\langle\sigma^*\rangle] = 1 ,
\]
so that Theorem~\ref{thm:apriori-sigma} in the form
$\nu\langle\sigma^*\rangle \leq C\mathcal{L}$ is the dimensionless statement,
and $\sigma^*$ itself is already dimensionless (as it must be, by
Proposition~\ref{prop:sigma-star-invariant}). Likewise the Type-I constant
$C_I$ in $\sup|u|\leq C_I(T-t)^{-1/2}$ has $[C_I] = L\mathcal{T}^{-1/2} =
[\nu]^{1/2}$, so $c_I := C_I\nu^{-1/2}$ is the dimensionless Type-I constant;
all constants below named $C(c_I)$ are pure numbers. Every displayed
inequality in this section has been checked to be homogeneous of a single
degree in $L$ (with $\mathcal{T}\sim L^2$).
\end{remark}

\begin{definition}[The Type-I ancient class]\label{def:profile-class}
For $c_I\geq 1$ let $\mathcal{A}(c_I)$ denote the set of $U$ such that:
$U$ is a smooth solution of the Navier--Stokes equations
($\nu=1$) on $\R^3\times(-\infty,0)$ with $\operatorname{div}U=0$, and
\begin{equation}\label{eq:profile-bounds}
  \sup_{y\in\R^3}|U(y,s)| \;\leq\; c_I\,|s|^{-1/2},
  \qquad
  M_U(s) := \sup_{y\in\R^3}|\omega_U(y,s)| \;\leq\; c_I\,|s|^{-1}
  \qquad (s<0).
\end{equation}
$U$ is called \emph{trivial} if $\omega_U\equiv 0$, equivalently (given
\eqref{eq:profile-bounds} and $\operatorname{div}U=0$) if $U(y,s)=b(s)$ for
some bounded measurable $b$.
\end{definition}

\begin{lemma}[$\mathcal{A}(c_I)$ is closed under the scaling group]
  \label{lem:typeI-scaleinv-bounds}
  If $U\in\mathcal{A}(c_I)$ and $\mu>0$ then
  $U_\mu(y,s) := \mu\,U(\mu y,\mu^2 s) \in \mathcal{A}(c_I)$, with the
  \emph{same} constant $c_I$; moreover
  $M_{U_\mu}(s) = \mu^2 M_U(\mu^2 s)$ and, by
  Proposition~\ref{prop:sigma-star-invariant},
  $\sigma^*(U_\mu)(s) = \sigma^*(U)(\mu^2 s)$.
  \emph{[Proved]}
\end{lemma}
\begin{proof}
  $|U_\mu(y,s)| = \mu|U(\mu y,\mu^2 s)| \leq \mu\,c_I(\mu^2|s|)^{-1/2}
  = c_I|s|^{-1/2}$, and identically for the vorticity with
  $\omega_{U_\mu}(y,s) = \mu^2\omega_U(\mu y,\mu^2 s)$. The scaling of
  $M$ and $\sigma^*$ is Proposition~\ref{prop:sigma-star-invariant} read on
  $\R^3$. Smoothness and the equation are preserved because the scaling is
  the NS symmetry with $\nu$ fixed.
\end{proof}

\begin{remark}[Why this closure matters]
Lemma~\ref{lem:typeI-scaleinv-bounds} is what makes the class
$\mathcal{A}(c_I)$ a legitimate arena for a \emph{second} compactness argument
(Proposition~\ref{prop:two-step-reduction}): one may zoom in on a profile at
arbitrarily remote past times without leaving the class or degrading $c_I$.
This is not available for the original solution $u$, whose Type-I bounds
degenerate only at the single time $T$.
\end{remark}

\subsection{Step 1: profile extraction}
\label{subsec:s39-step1}

Throughout, $u$ is smooth on $\T^3\times[0,T)$ with a genuine first
singularity at $T$ and singular point $x^*$, satisfying the Type-I bounds
$M(t)\leq C_I/(T-t)$ and $\sup|u(\cdot,t)|\leq C_I(T-t)^{-1/2}$ (the standing
assumptions of \S\ref{sec:moser-s33}). For $\lambda\in(0,1]$ define
\begin{equation}\label{eq:rescaling}
  u_\lambda(y,s) \;:=\; \lambda\,u(x^*+\lambda y,\ T+\lambda^2 s),
  \qquad (y,s)\in \Lambda_\lambda\times[-T/\lambda^2,0),
\end{equation}
where $\Lambda_\lambda = \lambda^{-1}\T^3$ exhausts $\R^3$ as
$\lambda\to0$. Direct computation from the Type-I bounds gives, uniformly in
$\lambda$,
\begin{equation}\label{eq:rescaled-bounds}
  |u_\lambda(y,s)| \leq c_I\,|s|^{-1/2},
  \qquad
  |\omega_\lambda(y,s)| \leq c_I\,|s|^{-1}
  \qquad\text{on }\Lambda_\lambda\times[-T/\lambda^2,0),
\end{equation}
with $\omega_\lambda = \lambda^2\omega(x^*+\lambda\,\cdot\,,T+\lambda^2\,\cdot\,)$.
\emph{[Proved: this is the one place where Type-I is used exactly as
designed --- \eqref{eq:rescaled-bounds} is $\lambda$-independent because the
Type-I rates are precisely the scaling rates.]}

The two inputs required to turn \eqref{eq:rescaled-bounds} into a nontrivial
limit are isolated now, and neither is proved here.

\begin{itemize}
\item[\textbf{(R)}] \emph{Scale-uniform interior regularity.}
  There are constants $C_k(c_I)$, $k\geq 0$, such that every $u_\lambda$
  as in \eqref{eq:rescaling} satisfies
  \[
    \bigl\|\nabla^k u_\lambda\bigr\|_{L^\infty\bigl(B_R\times[-S_2,-S_1]\bigr)}
    \;\leq\; C_k(c_I)\,\bigl(1+R\bigr)\,S_1^{-(k+1)/2}
  \]
  for all $0<S_1<S_2$ and $R\geq1$ with $B_R\times[-2S_2,-S_1/2]$ inside the
  domain of $u_\lambda$.
  \emph{Status: Conditional. This is a scale-invariant form of Serrin's
  interior regularity criterion \cite{Serrin1962} (an $L^\infty$ velocity
  bound on a parabolic cylinder controls all interior spatial derivatives,
  quantitatively), sharpened by the local-pressure theory of
  Ladyzhenskaya--Seregin \cite{LadyzhenskayaSeregin1999}. Two caveats are
  genuine and are the reason (R) is not marked Proved: (i) Serrin's criterion
  as classically stated controls spatial derivatives; time regularity, and
  hence $C^\infty_{\mathrm{loc}}$ compactness, requires the local pressure
  decomposition and excludes the spatially-constant ``parasitic'' solutions
  only up to an additive $b(s)$ --- exactly the ambiguity that reappears in
  Theorem~\ref{thm:knss}; (ii) the $\lambda$-uniformity of the constants must
  be checked against the fact that $u_\lambda$ has \emph{no} uniform global
  energy bound ($\int_{\Lambda_\lambda}|u_\lambda|^2\,dy
  = \lambda^{-1}\|u\|_{L^2(\T^3)}^2\to\infty$), so only the local
  consequence $\int_{B_R}|u_\lambda|^2 \leq C R^3 |s|^{-1}$
  of \eqref{eq:rescaled-bounds} is available. Both caveats are standard in the
  Type-I literature \cite{Seregin2012} but neither is discharged in this
  document.}
\item[\textbf{(N)}] \emph{Non-degeneracy at the singular point (no vorticity
  escape).}
  There exist $\varepsilon_*>0$, a compact
  $\mathcal{K}\subset\R^3\times(-\infty,0)$ and a sequence
  $\lambda_j\downarrow0$ with
  \[
    \liminf_{j\to\infty}\ \iint_{\mathcal{K}}|\nabla u_{\lambda_j}|^2\,dy\,ds
    \;\geq\;\varepsilon_* ,
  \]
  and in addition $M_{U}(s) = \lim_j M_{u_{\lambda_j}}(s)$ for a.e.\ $s<0$,
  where $U$ is the limit produced by
  Proposition~\ref{prop:profile-extraction}.
  \emph{Status: Assumed. Discussion in Remark~\ref{rem:N-discussion}.}
\end{itemize}

\begin{remark}[What (N) really costs, and why it is not free]
  \label{rem:N-discussion}
  The Caffarelli--Kohn--Nirenberg $\varepsilon$-regularity theorem
  \cite{CKN1982} supplies exactly the right \emph{scale-invariant} lower
  bound: because $x^*$ is a singular point, for every $r>0$
  \[
    r^{-1}\iint_{Q(x^*,r)}|\nabla u|^2\,dx\,dt \;\geq\;\varepsilon_* ,
    \qquad
    \text{and}\qquad
    r^{-1}\iint_{Q(x^*,r)}|\nabla u|^2\,dx\,dt
    \;=\;\iint_{Q_1(0)}|\nabla u_{r}|^2\,dy\,ds ,
  \]
  the identity being the same computation as
  Proposition~\ref{prop:sigma-star-invariant}. So a uniform-in-$\lambda$
  lower bound \emph{is} available on the cylinder $Q_1(0)$. The difficulty is
  that $Q_1(0)$ \emph{touches the terminal slice} $s=0$, where
  \eqref{eq:rescaled-bounds} degenerates and where the compactness of
  Proposition~\ref{prop:profile-extraction} is unavailable; the
  $\varepsilon_*$ of mass could in principle live entirely in
  $\{|s|<\delta\}$ for every $\delta$. Pushing it off the terminal slice
  requires a Type-I-improved $\varepsilon$-regularity statement (of the kind
  developed by Seregin \cite{Seregin2012}) and is \emph{not} done here. The
  second half of (N) --- convergence, not merely upper semicontinuity, of the
  global suprema $M$ --- excludes the vorticity maximum drifting to spatial
  infinity in the rescaled picture; it is the honest cost of $\sigma^*$
  having been defined with a \emph{global} normalization $M(t)$. We return to
  this in Lemma~\ref{lem:sigma-equality}.
\end{remark}

\begin{proposition}[Extraction of an ancient profile]\label{prop:profile-extraction}
  Assume \textbf{(R)}. Then there is a sequence $\lambda_j\downarrow 0$ and a
  $U\in\mathcal{A}(c_I)$ such that
  \[
    u_{\lambda_j}\longrightarrow U
    \qquad\text{in } C^\infty\bigl(\mathcal{K}\bigr)\ \text{ for every compact }
    \mathcal{K}\subset\R^3\times(-\infty,0).
  \]
  If in addition \textbf{(N)} holds (with the same sequence), then $U$ is
  nontrivial: $\iint_{\mathcal{K}}|\nabla U|^2 \geq\varepsilon_*$, hence
  $\omega_U\not\equiv0$ and $M_U(s)>0$ for every $s<0$.
  \emph{[Proved-modulo-(R),(N)]}
\end{proposition}
\begin{proof}
  \emph{Compactness.} Fix $R\geq1$, $0<S_1<S_2$. By \textbf{(R)} the family
  $\{u_\lambda\}$ is bounded in $C^k(B_R\times[-S_2,-S_1])$ for every $k$,
  uniformly in $\lambda$ small enough that the cylinder lies in the domain.
  In particular $\{\partial_s u_\lambda\}$ is bounded there (the equation
  $\partial_s u_\lambda = -\,(u_\lambda\!\cdot\!\nabla)u_\lambda
  + \Delta u_\lambda - \nabla p_\lambda$ with the local pressure bound
  supplied by \textbf{(R)}), so Arzel\`a--Ascoli and a diagonal argument over
  an exhausting sequence of cylinders and of $k$ give a subsequence
  $u_{\lambda_j}\to U$ in $C^\infty_{\mathrm{loc}}(\R^3\times(-\infty,0))$.
  Locally uniform convergence of all derivatives passes to the limit in the
  equation, so $U$ solves NS classically on $\R^3\times(-\infty,0)$, and
  $\operatorname{div}U=0$.

  \emph{Membership in $\mathcal{A}(c_I)$.} The bounds
  \eqref{eq:profile-bounds} follow from \eqref{eq:rescaled-bounds} by
  pointwise passage to the limit at each $(y,s)$, then taking suprema; note
  this yields the bound with the \emph{same} $c_I$ and in particular
  $M_U(s)\leq\liminf_j M_{u_{\lambda_j}}(s)$ (upper semicontinuity of a
  supremum under locally uniform convergence: $\sup_{\mathcal{K}}|\omega_U|
  = \lim_j\sup_{\mathcal{K}}|\omega_{\lambda_j}|
  \leq\liminf_j M_{u_{\lambda_j}}$ for every compact $\mathcal{K}$, then
  exhaust).

  \emph{Nontriviality.} $C^\infty_{\mathrm{loc}}$ convergence gives
  $\iint_{\mathcal{K}}|\nabla U|^2 = \lim_j\iint_{\mathcal{K}}
  |\nabla u_{\lambda_j}|^2 \geq\varepsilon_*$ by \textbf{(N)}. If
  $\omega_U(\cdot,s_0)\equiv0$ for some $s_0$ then, $U(\cdot,s_0)$ being
  bounded on $\R^3$ and curl-free and divergence-free, each component is a
  bounded harmonic function, hence constant; the same argument at every later
  time via forward uniqueness for the (then linear, in fact trivial) system
  gives $U(y,s)=b(s)$ on $[s_0,0)$ and $\nabla U\equiv0$ there. Since the
  compact $\mathcal{K}$ may be assumed to lie in $\{s\geq s_0\}$ after
  replacing $s_0$ by $\inf\{s:\ (y,s)\in\mathcal{K}\}$, this contradicts
  $\iint_{\mathcal{K}}|\nabla U|^2\geq\varepsilon_*$. Hence
  $M_U(s)>0$ for all $s<0$.
\end{proof}

\begin{remark}[Honest statement of the convergence topology]
  The topology is $C^\infty_{\mathrm{loc}}$ on
  $\R^3\times(-\infty,0)$ --- \emph{strong}, and strong at derivative level.
  This is far more than the weak-$*$ $L^\infty$ / weak $L^2_{\mathrm{loc}}$
  convergence one gets from energy bounds alone, and it is bought entirely by
  \textbf{(R)}. It is essential to be explicit about this: \emph{the entire
  derivative-level transfer problem flagged in Step 2 lives inside}
  \textbf{(R)}, not in Step 2 itself. If one weakens \textbf{(R)} to weak
  convergence, Lemma~\ref{lem:sigma-equality} degrades to
  Lemma~\ref{lem:sigma-semicontinuity} and the consequences are catalogued
  there.
\end{remark}

\subsection{Step 2: transfer audit --- what survives the limit}
\label{subsec:s39-step2}

This is the step at which an argument of this shape is most likely to be
wrong, so it is treated separately and pessimistically. There are two
logically distinct questions, and conflating them is the trap:

\begin{enumerate}[label=(Q\arabic*)]
\item Does a given piece of machinery \emph{hold for $U$}? Here the honest
  answer, for everything the program needs, is \emph{yes and for a trivial
  reason}: $U$ is a smooth solution satisfying scale-invariant Type-I bounds,
  so each lemma may be \emph{re-derived on $U$ from scratch}. Nothing is
  transferred as a limit of inequalities.
\item Does a \emph{quantity} converge, i.e.\ is $\sigma^*(U)$ related to
  $\lim_j\sigma^*(u_{\lambda_j})$? Here the answer is a one-directional
  inequality in general, and the reframe must be designed so that it needs
  only the available direction.
\end{enumerate}

\subsubsection*{(Q1): the machinery holds on $U$ --- and is log-free there}

\begin{proposition}[Re-derivation on the profile]\label{prop:transfer-audit}
  Let $U\in\mathcal{A}(c_I)$ be nontrivial. Then:
  \begin{enumerate}[label=(\alph*)]
  \item \emph{[Proved]} The magnitude and direction
    equations \eqref{eq:mag-eq}--\eqref{eq:dir-eq} hold for $U$ pointwise on
    $\{\omega_U\neq0\}$, and the orthogonal split
    Lemma~\ref{lem:grad-omega-split} holds pointwise. Reason: both are
    algebraic consequences of the vorticity equation for a smooth solution,
    proved in Lemma~\ref{lem:direction-eq} with no use of the domain,
    of finiteness of any global norm, or of a finite initial time.
  \item \emph{[Proved]} Proposition~\ref{prop:sigma-star-invariant} (exact
    scale invariance of $\sigma^*$) holds for $U$, and by
    Lemma~\ref{lem:typeI-scaleinv-bounds} the induced action on the profile
    class is $\sigma^*(U_\mu)(s)=\sigma^*(U)(\mu^2 s)$.
  \item \emph{[Proved]} Lemma~\ref{lem:gram-negativity} (Gram negativity)
    holds for $U$: it is the pointwise linear-algebra fact that
    $\nabla m\otimes\nabla m$ contracted against the positive semidefinite
    Gram tensor $G_{jk}=\partial_j\hat e\cdot\partial_k\hat e$ is
    non-negative, plus one integration by parts on a compactly supported
    cutoff. Neither step sees the limit procedure.
  \item \emph{[Proved-modulo-(R)]} On the profile the logarithmic factor
    $\mathcal{L}$ of \S\ref{sec:moser-s33}--\S\ref{sec:cf-s37} may be
    replaced by an absolute constant: for every $s<0$,
    \[
      \bigl\|\nabla U(\cdot,s)\bigr\|_{L^\infty(\R^3)}
      \;\leq\; C_1(c_I)\,|s|^{-1}
      \qquad (s<0),
    \]
    and in particular $\|\nabla U(\cdot,s)\|_{L^\infty}\leq C(c_I)\,M_U(s)$
    at every time with $M_U(s)\geq c_I^{-1}|s|^{-1}$ (i.e.\ whenever the
    Type-I upper bound is saturated up to a factor).
  \item \emph{[Proved-modulo-(R),(d)]} Theorem~\ref{thm:apriori-sigma} holds
    for $U$ in the log-free form
    \begin{equation}\label{eq:profile-apriori}
      \langle\sigma^*_U\rangle(s_0)
      \;:=\;(r^*)^{-3}\iint_{Q(r^*/2)}|\nabla\hat e_U|^2\,
        \mathbf{1}_{\{|\omega_U|\geq M_U/2\}}\,dy\,ds
      \;\leq\; C(c_I),
      \qquad r^* = M_U(s_0)^{-1/2},
    \end{equation}
    for every $s_0<0$ with $Q(r^*/2)\subset\R^3\times(-\infty,0)$.
  \end{enumerate}
\end{proposition}
\begin{proof}
  (a)--(c) are as stated: inspect the proofs of
  Lemma~\ref{lem:direction-eq}, Proposition~\ref{prop:sigma-star-invariant}
  and Lemma~\ref{lem:gram-negativity} and observe that each uses only
  pointwise identities for a smooth divergence-free field and integration by
  parts against a compactly supported cutoff.

  (d) By \textbf{(R)} applied to the limit (equivalently, by passing the
  \textbf{(R)} bounds through the $C^\infty_{\mathrm{loc}}$ convergence),
  $\|\nabla U\|_{L^\infty(B_R(y_0)\times[-2|s|,-|s|/2])}\leq
  C_1(c_I)(1+R)|s|^{-1}$; the point is that the exponent $-1$ and the
  constant are \emph{independent of $y_0$}, because
  \eqref{eq:profile-bounds} is uniform in $y$. Taking $R=1$ and varying
  $y_0$ over $\R^3$ gives the global-in-space bound. Dimensional check:
  $[\nabla U] = \mathcal{T}^{-1} = L^{-2}$ in units $\nu=1$, and
  $[|s|^{-1}]=L^{-2}$. \emph{Why the logarithm disappears.} In
  Remark~\ref{rem:marginality} the factor
  $\mathcal{L}=\log(e+\|u\|_{H^3}/M)$ measures the ratio between the resolved
  scale and the smallest active scale of $u$; a Calder\'on--Zygmund bound on
  $\nabla u$ in $L^\infty$ must pay it. On a Type-I \emph{profile} that ratio
  is $O(1)$ by scale invariance --- $U$ has, by construction, no structure
  below the scale $|s|^{1/2}$ that is not already visible at scale
  $|s|^{1/2}$ --- so the CZ step is replaced by the interior estimate (R),
  which is logarithm-free. This is a genuine (if modest) gain of the reframe:
  every marginal $\mathcal{L}$-degraded estimate in
  \S\ref{sec:moser-s33}--\S\ref{sec:verification-s36} becomes
  $\mathcal{L}$-free on the profile.

  (e) Re-run the proof of Theorem~\ref{thm:apriori-sigma} verbatim on $U$:
  test the magnitude equation \eqref{eq:mag-eq} (valid on $U$ by (a)) with
  $m\eta^2$, $\eta$ the composite cutoff of
  Definition~\ref{def:composite-cutoff} built from $M_U(s)$; all cutoffs are
  compactly supported in $Q(r^*)$, so no decay of $U$ at spatial infinity is
  needed. The right-hand terms are estimated by
  Lemma~\ref{lem:mag-rhs} with $\mathcal{L}$ replaced by $C(c_I)$ using (d)
  in place of the CZ-with-logarithm bound, and with
  Lemma~\ref{lem:local-enstrophy-typeI} re-derived on $U$ (its proof is the
  vorticity Caccioppoli inequality on a compactly supported cutoff, again
  local). The conclusion of Theorem~\ref{thm:apriori-sigma} then reads
  $\nu\langle\sigma^*_U\rangle\leq C(c_I)$, i.e.\ \eqref{eq:profile-apriori}
  in units $\nu=1$. Dimensional check in physical units:
  $[\nu\langle\sigma^*\rangle]=1$ by Remark~\ref{rem:s39-units}. Note that
  \eqref{eq:profile-apriori} is invariant under $U\mapsto U_\mu$, consistent
  with (b) and Lemma~\ref{lem:typeI-scaleinv-bounds}.
\end{proof}

\begin{remark}[Two things that do \emph{not} transfer, and must not be used]
  \label{rem:s39-nontransfer}
  \begin{enumerate}[label=(\roman*)]
  \item \textbf{Proposition~\ref{prop:localized-cf} (localized
    Constantin--Fefferman) does not transfer to the profile.} Its far-field
    step bounds
    $\int_{|y|\geq\rho'}|\omega(x^*-y)|\,|y|^{-3}dy$ by
    $C(\rho')^{-3/2}\Omega^{1/2}$ with
    $\Omega=\|\omega\|_{L^2}^2$; for $U\in\mathcal{A}(c_I)$ on $\R^3$ there is
    \emph{no} reason for $\|\omega_U(\cdot,s)\|_{L^2(\R^3)}$ to be finite, and
    the naive substitute $|\omega_U|\leq M_U$ makes the far-field integral
    logarithmically divergent at large $|y|$. Consequently \emph{the entire
    H2/Regime-B analysis of \S\ref{sec:cf-s37} is unavailable on the profile
    as written}, and Step 4 below may not invoke it. This is a real
    obstruction discovered in this session, not a bookkeeping remark; it is
    recorded in the ledger.
  \item \textbf{The Bridge Lemma may not be pulled back.} Conjecture
    \ref{conj:bridge} is an implication with hypothesis
    $\sigma^*\leq\varepsilon_0$. Even if one proved $\sigma^*(U)=0$, the
    semicontinuity direction of Lemma~\ref{lem:sigma-semicontinuity} does
    \emph{not} give $\sigma^*(u_{\lambda_j})\to0$, so H1 cannot be applied to
    the original sequence. Any argument that routes the profile conclusion
    back through H1 on $u_{\lambda_j}$ is invalid. The dichotomy of
    \S\ref{subsec:s39-step5} is designed to obtain its contradiction
    \emph{entirely on $U$}, for exactly this reason.
  \end{enumerate}
\end{remark}

\subsubsection*{(Q2): the $\sigma^*$ transfer, and its direction}

\begin{lemma}[Lower semicontinuity --- the generic case]
  \label{lem:sigma-semicontinuity}
  Suppose $u_{\lambda_j}\rightharpoonup U$ only weakly in
  $H^1_{\mathrm{loc}}$, and suppose the normalizations converge:
  $M_{u_{\lambda_j}}(s)\to M_U(s)$ and the base points
  $y^*_j\to y^*$. Then
  \begin{equation}\label{eq:sigma-lsc}
    \sigma^*(U)(s) \;\leq\; \liminf_{j\to\infty}\,\sigma^*(u_{\lambda_j})(s) .
  \end{equation}
  In general \emph{no} inequality holds in the opposite direction.
  \emph{[Proved, modulo the level-set null-set technicality noted in the
  proof]}
\end{lemma}
\begin{proof}
  On the indicator set $\{|\omega|\geq M/4\}$ the integrand is
  $|\nabla\hat e|^2 = (|\nabla\omega|^2-|\nabla|\omega||^2)/|\omega|^2$
  (Lemma~\ref{lem:grad-omega-split}) with denominator bounded below by
  $M^2/16>0$, so $|\nabla\hat e|^2\mathbf{1}$ is a nonnegative quadratic
  form in $\nabla\omega$ with locally uniformly continuous coefficients, and
  the map $\nabla\omega\mapsto\iint|\nabla\hat e|^2\mathbf{1}$ is convex and
  weakly lower semicontinuous on $L^2_{\mathrm{loc}}$; \eqref{eq:sigma-lsc}
  follows. (The technicality: the indicator's boundary
  $\{|\omega_U| = M_U/4\}$ must have measure zero for the coefficient
  convergence; this holds for a.e.\ choice of the threshold constant $1/4$ by
  Fubini, and the definition of $\sigma^*$ is stable under replacing $1/4$ by
  any nearby constant, so this is harmless but should be stated.)
  Failure of the reverse: weak convergence permits oscillation to be lost in
  the limit, i.e.\ $|\nabla\hat e_j|^2\rightharpoonup |\nabla\hat e_U|^2 +
  \mu$ with $\mu\geq0$ a defect measure; nothing forces $\mu=0$.
\end{proof}

\begin{lemma}[Equality under (R) and (N)]\label{lem:sigma-equality}
  Under \textbf{(R)} (so that the convergence is
  $C^\infty_{\mathrm{loc}}$) and the second half of \textbf{(N)}
  ($M_{u_{\lambda_j}}(s)\to M_U(s)$, with base points converging), one has
  for a.e.\ $s<0$
  \[
    \sigma^*(U)(s) \;=\; \lim_{j\to\infty}\sigma^*(u_{\lambda_j})(s),
  \]
  and likewise for $\langle\sigma^*\rangle$.
  \emph{[Proved-modulo-(R),(N)]}
\end{lemma}
\begin{proof}
  Since $M_U(s)>0$, the set $\{|\omega_U|\geq M_U/4\}$ sits inside
  $\{|\omega_U|>0\}$ where $\hat e_U$ is smooth; on it the integrand of
  $\sigma^*$ is bounded, $|\nabla\hat e|^2\leq 16|\nabla\omega|^2/M^2$, and
  $C^\infty_{\mathrm{loc}}$ convergence gives locally uniform convergence of
  $|\nabla\hat e_j|^2$ there. The radius $r^*_j = M_{u_{\lambda_j}}^{-1/2}\to
  M_U^{-1/2}=r^*_U$ and the threshold $M_{u_{\lambda_j}}/4\to M_U/4$ by
  hypothesis, so the domains of integration converge in measure (again up to
  the null level set of the previous proof). Dominated convergence
  concludes.
\end{proof}

\begin{remark}[The honest verdict on Step 2]\label{rem:s39-step2-verdict}
  \begin{enumerate}[label=(\roman*)]
  \item The derivative-level obstruction is \textbf{real but relocated}. It
    does not appear as an unrepairable defect in Step 2; it appears as the
    hypothesis \textbf{(R)}. With \textbf{(R)} one has equality
    (Lemma~\ref{lem:sigma-equality}); without it one has only
    \eqref{eq:sigma-lsc}.
  \item \textbf{The available direction is the one the argument needs.}
    \eqref{eq:sigma-lsc} says: \emph{smallness of $\sigma^*$ along a
    sequence descends to the limit}. That is precisely the implication used
    in Proposition~\ref{prop:two-step-reduction} to absorb the oscillatory
    middle case into the aligned horn. It is \emph{not} the implication
    needed to push a limit conclusion back onto the sequence, which is why
    Remark~\ref{rem:s39-nontransfer}(ii) forbids that route.
  \item \textbf{Which horn is hurt.} Neither horn of
    \S\ref{subsec:s39-step5} needs the reverse inequality, because the
    dichotomy is intrinsic to $U$: $\sigma^*(U)$ either vanishes
    identically on some past interval or it does not, and both cases are
    attacked on $U$. What semicontinuity \emph{kills} is the
    apparently-attractive hybrid strategy ``prove $\sigma^*(U)=0$, conclude
    $\sigma^*(u_{\lambda_j})\leq\varepsilon_0$, apply H1+H2 to the original
    solution''. That strategy is unavailable, full stop. Any future draft
    that appears to use it should be treated as containing an error.
  \item \textbf{Where a defect measure would actually bite.} If one runs the
    argument with only weak convergence, the defect measure $\mu$ of
    Lemma~\ref{lem:sigma-semicontinuity} is concentrated coherence deficit
    that the limit does not see. The disordered horn (Step 4) would then be
    attacking the wrong object: the disorder could be entirely in $\mu$. This
    is the precise sense in which Step 4 --- not Step 3 --- is the horn that
    depends on \textbf{(R)}.
  \end{enumerate}
\end{remark}

\subsection{Step 3: the aligned horn --- $\sigma^*(U)=0$ forces triviality}
\label{subsec:s39-step3}

\begin{lemma}[Local alignment from a vanishing deficit]\label{lem:alignment-local}
  Let $U\in\mathcal{A}(c_I)$ be nontrivial and let $s<0$ satisfy
  $\sigma^*(U)(s)=0$, with base point $y^*$ (so
  $|\omega_U(y^*,s)|\geq M_U(s)/2$) and $r^*=M_U(s)^{-1/2}$. Then there is
  $e_0\in S^2$ with
  \[
    \hat e_U(\cdot,s) \;\equiv\; e_0
    \qquad\text{on the connected component } \mathcal{O}\ni y^*
    \text{ of } \bigl\{|\omega_U(\cdot,s)|>M_U(s)/4\bigr\}\cap B_{r^*}(y^*).
  \]
  \emph{[Proved]}
\end{lemma}
\begin{proof}
  $\sigma^*(U)(s)=0$ and $|\nabla\hat e_U|^2\geq0$ force
  $|\nabla\hat e_U|^2=0$ a.e.\ on
  $B_{r^*}(y^*)\cap\{|\omega_U|\geq M_U/4\}$. On the open subset
  $\{|\omega_U|>M_U/4\}\cap B_{r^*}(y^*)$ the field $\hat e_U$ is smooth
  (denominator bounded below), so $\nabla\hat e_U\equiv0$ there by
  continuity, and $\hat e_U$ is constant on each connected component.
  $y^*$ lies in this open set since $|\omega_U(y^*,s)|\geq M_U/2>M_U/4$.
\end{proof}

\begin{remark}[The gap is spatial globality, and it is large]
  \label{rem:A-up-discussion}
  Lemma~\ref{lem:alignment-local} is genuinely local: one time, one ball of
  radius $r^*$, and only the connected component of the high-vorticity set
  containing the base point. The Liouville mechanism of
  Lemma~\ref{lem:harmonic-reduction} below, by contrast, needs alignment on
  \emph{all} of $\R^3$, because it runs a Liouville theorem for bounded
  harmonic functions on $\R^3$ and no ball version of that exists. We
  therefore record the upgrade as a named input:
  \begin{itemize}
  \item[\textbf{(A-up)}] \emph{Global alignment upgrade.} If
    $U\in\mathcal{A}(c_I)$ and $\hat e_U(\cdot,s_1)\equiv e_0$ on a
    nonempty open subset of $\{|\omega_U(\cdot,s_1)|>M_U(s_1)/4\}$ for some
    $s_1<0$, then $\omega_U(y,s)\parallel e_0$ for all
    $(y,s)\in\R^3\times(-\infty,0)$.
    \emph{Status: \textbf{Open}, and this is the largest single gap opened by
    this section. It should not be regarded as a technicality.}
  \end{itemize}
  Three candidate mechanisms, none executed here:
  \emph{($\alpha$) Unique continuation} for the differential inequality
  satisfied by $w=|\nabla\hat e|^2$: $w\geq0$ vanishes on an open set, and
  $w$ satisfies a drift-diffusion inequality (\S\ref{sec:moser-s33},
  eq.~\eqref{eq:weak-w}) whose drift is $U+2\nabla\log m$; strong-maximum /
  unique-continuation principles for such operators require the drift to lie
  in a suitable Kato class, which is precisely the $K$-family problem of
  \S\ref{sec:moser-s33}--\S\ref{sec:verification-s36}. So this route reduces
  (A-up) to a problem of the \emph{same} kind the program already faces, not
  to a new one --- a mild positive, since the $K$-family is now controlled
  modulo the annulus residual (Proposition~\ref{prop:K-amended}).
  \emph{($\beta$) Backward uniqueness} in the sense of
  Escauriaza--Seregin--\v{S}ver\'ak \cite{EscauriazaSereginSverak2003},
  applied to the difference between $\omega_U$ and its projection onto
  $e_0$; the parabolic operator is right but the required decay at spatial
  infinity is not available for a profile.
  \emph{($\gamma$) Scale exhaustion}, which is available only under a
  strengthened hypothesis and is the one concrete partial result we can
  state:
\end{remark}

\begin{remark}[Partial reduction of (A-up) by scale exhaustion]
  \label{rem:scale-exhaustion}
  Suppose, more strongly, that $\sigma^*(U)(s)=0$ for \emph{every} $s<0$, and
  that the base points are centred, $|y^*(s)|\leq C_0|s|^{1/2}$. By
  Lemma~\ref{lem:typeI-scaleinv-bounds}, $\sigma^*(U_\mu)(-1)=\sigma^*(U)(-\mu^2)=0$
  for every $\mu>0$, so Lemma~\ref{lem:alignment-local} applies at every
  scale, giving alignment on a ball of radius $r^*(s)\asymp|s|^{1/2}$
  centred within distance $C_0|s|^{1/2}$ of the origin. As $s\to-\infty$
  these balls exhaust $\R^3$. \emph{Residual gap:} the alignment is only on
  the connected component of $\{|\omega_U|>M_U/4\}$ inside each ball, and
  (i) the direction $e_0=e_0(s)$ could a priori depend on $s$, (ii) the
  high-vorticity set need not fill, or even be connected across, the balls.
  So this reduces (A-up), \emph{under the strengthened hypothesis}, to a
  statement about the geometry of the high-vorticity set across scales
  ---  weaker than (A-up) but still open. The dichotomy of
  \S\ref{subsec:s39-step5} does \emph{not} deliver the strengthened
  hypothesis (it produces vanishing at a single time), so this remark is a
  pointer, not a proof.
\end{remark}

\begin{lemma}[Aligned profiles are two-dimensional]\label{lem:harmonic-reduction}
  Assume \textbf{(R)}. Let $U\in\mathcal{A}(c_I)$ satisfy
  $\omega_U(y,s)\parallel e_0$ for all $(y,s)\in\R^3\times(-\infty,0)$, with
  $e_0\in S^2$ fixed. Choose coordinates with $e_0=e_3$ and write
  $\omega_U=\zeta e_3$. Then
  \[
    \partial_3 U \equiv 0,
    \qquad U_3 = U_3(s),
    \qquad U_h := (U_1,U_2) = U_h(y_1,y_2,s),
  \]
  and $U_h$ is a smooth solution of the two-dimensional Navier--Stokes
  equations on $\R^2\times(-\infty,0)$ with
  $\sup_{\R^2}|U_h(\cdot,s)|\leq c_I|s|^{-1/2}$.
  \emph{[Proved-modulo-(R)]}
\end{lemma}
\begin{proof}
  \emph{Step 1.} $0=\operatorname{div}\omega_U=\partial_3\zeta$.

  \emph{Step 2.} Put $f:=\partial_3 U$. Then
  $\nabla\times f=\partial_3\omega_U=(\partial_3\zeta)e_3=0$ and
  $\nabla\cdot f=\partial_3(\operatorname{div}U)=0$, so
  $\Delta f=\nabla(\nabla\cdot f)-\nabla\times(\nabla\times f)=0$: each
  component of $f(\cdot,s)$ is harmonic on $\R^3$.

  \emph{Step 3.} By Proposition~\ref{prop:transfer-audit}(d),
  $\|\nabla U(\cdot,s)\|_{L^\infty(\R^3)}\leq C_1(c_I)|s|^{-1}<\infty$, so
  $f(\cdot,s)$ is a bounded harmonic function on $\R^3$, hence constant:
  $f(\cdot,s)=f(s)$ (Liouville). Therefore
  $U(y,s)=f(s)\,y_3+\tilde U(y_1,y_2,s)$. Since
  $\sup_{y\in\R^3}|U(y,s)|\leq c_I|s|^{-1/2}<\infty$ is finite \emph{uniformly
  in $y_3$}, we must have $f(s)=0$. Thus $\partial_3U\equiv0$.

  \emph{Step 4.} With $\partial_3U=0$:
  $\omega_1=\partial_2U_3-\partial_3U_2=\partial_2U_3$ and
  $\omega_2=\partial_3U_1-\partial_1U_3=-\partial_1U_3$; both vanish because
  $\omega_U\parallel e_3$. Hence $\nabla_hU_3=0$ and, with
  $\partial_3U_3=0$, $U_3=U_3(s)$.

  \emph{Step 5.} The horizontal momentum equations read
  $\partial_sU_h+(U_h\!\cdot\!\nabla_h)U_h+U_3\partial_3U_h
  =\Delta_hU_h-\nabla_hp$, and the $U_3\partial_3U_h$ term vanishes by
  Step 3. The third equation reads $\partial_sU_3=-\partial_3p$, so
  $\partial_3p$ is a function of $s$ alone and
  $p=p_h(y_1,y_2,s)-\partial_sU_3(s)\,y_3$, whence $\nabla_hp=\nabla_hp_h$.
  Therefore $(U_h,p_h)$ solves 2D NS. The bound is inherited from
  \eqref{eq:profile-bounds}.
\end{proof}

\begin{theorem}[Liouville theorem for bounded ancient 2D solutions
  \cite{KNSS2009}]\label{thm:knss}
  Let $v$ be a bounded ancient mild solution of the two-dimensional
  Navier--Stokes equations on $\R^2\times(-\infty,0)$. Then $v(y,s)=b(s)$ for
  some bounded measurable $b:(-\infty,0)\to\R^2$; in particular
  $\nabla v\equiv0$.
  \emph{[Cited. Koch--Nadirashvili--Seregin--\v{S}ver\'ak, Acta Math.\ 203
  (2009). The 3D analogue for general bounded ancient solutions is open;
  \cite{KNSS2009} also treats the 3D axisymmetric-without-swirl case. Caveat
  on applicability: the theorem is stated for \emph{mild} solutions, so
  membership of $U_h$ in that class must be checked --- for a smooth solution
  with $\|U_h\|_{L^\infty}<\infty$ and locally bounded derivatives this is
  the standard Duhamel representation modulo the additive $b(s)$ ambiguity,
  which is exactly what the conclusion encodes; we treat this as verbatim
  applicable but flag it rather than assert it.}
\end{theorem}

\begin{theorem}[Aligned horn]\label{thm:aligned-horn}
  Assume \textbf{(R)} and \textbf{(A-up)}. Let $U\in\mathcal{A}(c_I)$ and
  suppose $\sigma^*(U)(s_1)=0$ for some $s_1<0$ at which $M_U(s_1)>0$. Then
  $U$ is trivial: $\omega_U\equiv0$ and $U(y,s)=b(s)$.
  \emph{[Proved-modulo-(R),(A-up), and modulo the class caveat of
  Theorem~\ref{thm:knss}]}
\end{theorem}
\begin{proof}
  Lemma~\ref{lem:alignment-local} gives $\hat e_U(\cdot,s_1)\equiv e_0$ on a
  nonempty open subset of $\{|\omega_U(\cdot,s_1)|>M_U(s_1)/4\}$;
  \textbf{(A-up)} upgrades this to $\omega_U\parallel e_0$ on
  $\R^3\times(-\infty,0)$; Lemma~\ref{lem:harmonic-reduction} reduces $U$ to a
  bounded 2D ancient solution $U_h$ on $\R^2\times(-\infty,-\delta]$ for each
  $\delta>0$ (bounded there by $c_I\delta^{-1/2}$);
  Theorem~\ref{thm:knss} gives $U_h(y,s)=b_h(s)$ on $s\leq-\delta$, hence on
  $s<0$ since $\delta>0$ was arbitrary. Then
  $\zeta=\omega_3=\partial_1U_2-\partial_2U_1=0$, so $\omega_U\equiv0$, and
  with $U_3=U_3(s)$ we get $U(y,s)=b(s)$.
\end{proof}

\begin{remark}[Where the reframe's leverage actually sits, in Step 3]
  The chain is: vanishing \emph{scalar} coherence deficit $\Rightarrow$
  vorticity parallel to a \emph{fixed} direction $\Rightarrow$ the flow is
  two-dimensional $\Rightarrow$ 2D Liouville. The essential point is that the
  2D reduction is what makes a Liouville theorem available at all: the
  general 3D bounded-ancient Liouville problem is open \cite{KNSS2009}, so no
  version of Step 3 can work without producing a genuine dimensional
  reduction first. This is why $\sigma^*$ --- a measure of the failure of
  \emph{directional} coherence --- is the right quantity to have built the
  program around, and it is the strongest structural vindication of
  \S\ref{sec:pincer-s32} obtained so far. Note also that
  Lemma~\ref{lem:harmonic-reduction} uses the Type-I \emph{spatial}
  boundedness of the profile twice (Steps 3 and 4) --- on the original
  solution, on $\T^3$, the harmonic-Liouville step is unavailable, so this
  argument genuinely requires the profile.
\end{remark}

\subsection{Step 4: the disordered horn --- persistent disorder over an infinite past}
\label{subsec:s39-step4}

The premise now is $\sigma^*(U)(s)\geq\delta>0$ for all $s\leq -S_0$. The hoped-for
mechanism is that persistent geometric disorder costs persistent excess
dissipation, which an infinite past cannot pay. The first task is to state
precisely why the naive form of that argument \emph{cannot} work.

\begin{proposition}[No summable budget exists: the naive infinite-past
  argument fails structurally]\label{prop:no-summable-budget}
  Let $U\in\mathcal{A}(c_I)$ and $s_k=-2^k$, $r^*_k=M_U(s_k)^{-1/2}$. Then:
  \begin{enumerate}[label=(\alph*)]
  \item \emph{[Proved]} The dyadic past shells $(-2^{k+1},-2^k)$ and the
    self-similar cylinders $Q_k=B_{r^*_k/2}(y^*_k)\times[s_k-(r^*_k/2)^2,s_k]$
    are commensurate: $(r^*_k)^2 = M_U(s_k)^{-1}\geq |s_k|/c_I$, so each
    dyadic shell carries $O(1)$ self-similar cylinders, and the infinite past
    carries infinitely many.
  \item \emph{[Proved]} Every such cylinder satisfies the \emph{same}
    dimensionless bound $\langle\sigma^*_U\rangle(s_k)\leq C(c_I)$
    (Proposition~\ref{prop:transfer-audit}(e)), with no decay in $k$; and if
    $U$ is exactly self-similar ($U_\mu=U$ for all $\mu>0$) then
    $\langle\sigma^*_U\rangle(s_k)$ is \emph{independent} of $k$, by
    Lemma~\ref{lem:typeI-scaleinv-bounds}.
  \item \emph{[Proved]} Consequently no additive scale-invariant dissipation
    functional can furnish a finite total budget over the infinite past:
    $\sum_k\langle\sigma^*_U\rangle(s_k)=+\infty$ is consistent with (indeed
    forced by) the class $\mathcal{A}(c_I)$.
  \end{enumerate}
  \emph{Interpretation: the trap named in the program brief is not merely a
  technical difficulty --- for scale-invariant quantities it is a theorem
  that ``infinite past $+$ finite budget'' is vacuous. Any working
  infinite-past argument must instead use a quantity that is (i) bounded in a
  scale-invariant way and (ii) \emph{monotone} in logarithmic time.}
\end{proposition}
\begin{proof}
  (a) is the identity $r^*=M^{-1/2}$ together with
  \eqref{eq:profile-bounds}. (b) is
  Proposition~\ref{prop:transfer-audit}(e) plus
  Lemma~\ref{lem:typeI-scaleinv-bounds} applied with $\mu=2^{k/2}$, under
  which $Q_k$ is the image of $Q_0$. (c) is immediate from (b).
\end{proof}

Proposition~\ref{prop:no-summable-budget} dictates the correct framework.

\begin{lemma}[Logarithmic time and the profile amplitude]\label{lem:logtime}
  For $U\in\mathcal{A}(c_I)$ nontrivial set
  \[
    \tau := \log|s| \in\R \quad (\tau\to+\infty \iff s\to-\infty),
    \qquad
    \Phi(\tau) := |s|\,M_U(s)\big|_{s=-e^\tau} .
  \]
  Then
  \begin{enumerate}[label=(\alph*)]
  \item \emph{[Proved]} $0<\Phi(\tau)\leq c_I$ for all $\tau$;
  \item \emph{[Proved]} $\Phi$ is scale-covariant by translation only:
    $\Phi_{U_\mu}(\tau)=\Phi_U(\tau+2\log\mu)$. Hence $\Phi$ is a
    \emph{bounded} function on a \emph{line of infinite length}, and the
    scaling group acts by shifts;
  \item \emph{[Proved]} $M_U$ is locally Lipschitz in $s$ and, for a.e.\ $s$,
    \begin{equation}\label{eq:logPhi-ode}
      \frac{d}{d\tau}\log\Phi \;\geq\; 1 + |s|\,w_{\max}(s) - |s|\,\alpha_{\max}(s),
    \end{equation}
    where $w_{\max}(s)=|\nabla\hat e_U|^2(y_{\max}(s),s)$,
    $\alpha_{\max}(s)=\alpha_U(y_{\max}(s),s)$ and $y_{\max}(s)$ is a point
    at which $|\omega_U(\cdot,s)|$ attains its supremum --- \emph{provided
    the supremum is attained}; see Remark~\ref{rem:attainment}.
  \end{enumerate}
\end{lemma}
\begin{proof}
  (a) Upper bound: \eqref{eq:profile-bounds}. Positivity: nontriviality via
  Proposition~\ref{prop:profile-extraction}. (b) By
  Lemma~\ref{lem:typeI-scaleinv-bounds},
  $|s|M_{U_\mu}(s)=|s|\mu^2M_U(\mu^2s)=|\mu^2s|M_U(\mu^2s)$, and
  $\log|\mu^2s|=\tau+2\log\mu$.
  (c) $M_U$ is locally Lipschitz because $\nabla\omega_U$ and
  $\partial_s\omega_U$ are locally bounded by \textbf{(R)}. At a point of
  spatial maximum of $|\omega_U(\cdot,s)|$ one has $\nabla|\omega_U|=0$ and
  $\Delta|\omega_U|\leq0$, so \eqref{eq:mag-eq} (valid on $U$ by
  Proposition~\ref{prop:transfer-audit}(a); note $M_U>0$ so
  $y_{\max}\in\{\omega_U\neq0\}$) and Hamilton's trick give, for a.e.\ $s$,
  \[
    \frac{d}{ds}\log M_U \;\leq\; \alpha_{\max} - w_{\max}
    \qquad(\nu=1).
  \]
  Since $d/d\tau=-|s|\,d/ds$ and $\log\Phi=\tau+\log M_U$,
  \[
    \frac{d}{d\tau}\log\Phi
    = 1 - |s|\frac{d}{ds}\log M_U
    \geq 1 - |s|\bigl(\alpha_{\max}-w_{\max}\bigr),
  \]
  which is \eqref{eq:logPhi-ode}. Dimensional check: $[\,|s|\alpha\,]$ and
  $[\,|s|w\,]$ are both $\mathcal{T}\cdot L^{-2}=1$ in units $\nu=1$; $\tau$
  and $\log\Phi$ are dimensionless; every term of
  \eqref{eq:logPhi-ode} is dimensionless. \emph{Note the good sign:} the
  coherence deficit $w_{\max}$ enters \eqref{eq:logPhi-ode} with a
  \textbf{plus}, i.e.\ disorder \emph{drives $\Phi$ up}, and $\Phi$ is
  bounded --- this is the whole mechanism.
\end{proof}

\begin{remark}[Attainment of the supremum]\label{rem:attainment}
  On $\R^3$ the supremum $M_U(s)$ need not be attained. Two honest
  repairs: (i) include attainment in the non-degeneracy package
  \textbf{(N)} (it is implied by ``no vorticity escape'', which (N) already
  asserts in the form $M_U=\lim_jM_{u_{\lambda_j}}$ together with
  compactness of the maximizing points); or (ii) run the argument with
  almost-maximizers, replacing $\Delta|\omega|\leq0$ by
  $\Delta|\omega|\leq\epsilon$ at an $\epsilon$-almost-maximum and letting
  $\epsilon\to0$ --- routine but with a bookkeeping cost we do not pay here.
  We adopt (i) and record it.
\end{remark}

\begin{proposition}[Averaged necessary condition over the infinite past]
  \label{prop:logtime-necessary}
  Let $U\in\mathcal{A}(c_I)$ be nontrivial with $M_U$ attaining its suprema.
  Then for every $\tau_0$,
  \begin{equation}\label{eq:necessary-average}
    \limsup_{\tau_1\to+\infty}\ \frac{1}{\tau_1-\tau_0}
    \int_{\tau_0}^{\tau_1}\Bigl[\,|s|\,\alpha_{\max} - |s|\,w_{\max}\Bigr]\,d\tau
    \;\geq\; 1 .
  \end{equation}
  \emph{[Proved-modulo-(R) and Remark~\ref{rem:attainment}]}
\end{proposition}
\begin{proof}
  Integrate \eqref{eq:logPhi-ode} from $\tau_0$ to $\tau_1$:
  \[
    \log\Phi(\tau_1)-\log\Phi(\tau_0)
    \;\geq\; \int_{\tau_0}^{\tau_1}\bigl[1+|s|w_{\max}-|s|\alpha_{\max}\bigr]d\tau .
  \]
  By Lemma~\ref{lem:logtime}(a) the left side is at most
  $\log c_I-\log\Phi(\tau_0)=:C_0<\infty$, a constant independent of
  $\tau_1$. Rearranging and dividing by $\tau_1-\tau_0$,
  \[
    \frac{1}{\tau_1-\tau_0}\int_{\tau_0}^{\tau_1}
      \bigl[|s|\alpha_{\max}-|s|w_{\max}\bigr]d\tau
    \;\geq\; 1 - \frac{C_0}{\tau_1-\tau_0},
  \]
  and letting $\tau_1\to\infty$ gives \eqref{eq:necessary-average}.
\end{proof}

\begin{remark}[What Proposition~\ref{prop:logtime-necessary} is and is not]
  It \emph{is} the correct form of the infinite-past leverage: the infinite
  past has infinite \emph{logarithmic} length, and a bounded quantity cannot
  grow at a positive average logarithmic rate over infinite length. It
  \emph{is not} yet a contradiction: it says only that, on average, the
  normalized stretching at the vorticity maximum must exceed the normalized
  coherence sink there by at least $1$. Note that this is a statement about
  \emph{pointwise} values at the maximum, whereas $\sigma^*$ is a
  \emph{ball average}; and note that the required margin is exactly $1$, so
  the target below is \emph{marginal}, not subcritical. Marginality is what
  defeated \S19--\S20, and honesty requires flagging it. The difference from
  \S20 is nonetheless real: this is a comparison of two averages by a fixed
  constant, not a scale recursion requiring bootstrapped a priori smallness,
  so it cannot fail in the specific manner of \S\ref{sec:audit-s31},
  Defect~1.
\end{remark}

\begin{conjecture}[Disorder--depletion target for profiles]
  \label{target:disorder-depletion}
  There is a function $c:(0,\infty)\to(0,1]$ such that: for every
  $c_I\geq1$, every nontrivial $U\in\mathcal{A}(c_I)$ whose suprema are
  attained, and every $\delta>0$, if $\sigma^*(U)(s)\geq\delta$ for all
  $s\leq-S_0$ then for some $\tau_0$
  \[
    \limsup_{\tau_1\to+\infty}\ \frac{1}{\tau_1-\tau_0}
    \int_{\tau_0}^{\tau_1}\Bigl[\,|s|\,\alpha_{\max} - |s|\,w_{\max}\Bigr]d\tau
    \;\leq\; 1 - c(\delta).
  \]
  \emph{Status: \textbf{Open target}, sharply posed. Together with
  Proposition~\ref{prop:logtime-necessary} it is a contradiction, hence it
  closes the disordered horn.}
\end{conjecture}

\begin{remark}[Decomposition of Conjecture~\ref{target:disorder-depletion}
  into its two hard halves]\label{rem:target-decomposition}
  \begin{enumerate}[label=(T4\alph*)]
  \item \emph{Average-to-point for $w$ [Open].} One needs the ball average
    $\sigma^*\geq\delta$ to force a log-time-averaged lower bound on
    $|s|w_{\max}$, i.e.\ on $w$ \emph{at the vorticity maximum}. Averages do
    not control point values in general. The natural mechanism is a
    weak-Harnack inequality for the parabolic operator governing $w$: the
    $\sigma^*$-ball is centred at the vorticity maximum and has radius
    exactly $r^*$, so a Harnack inequality on $Q(r^*)$ would transfer the
    average to the centre at a slightly later time with a universal constant.
    Applicability requires the drift $U+2\nabla\log m$ of
    \eqref{eq:weak-w} to lie in a Kato/Morrey class --- which is precisely
    the $K$-family controlled, modulo the annulus residual, by
    Proposition~\ref{prop:K-amended}. So (T4a) plugs into machinery the
    program already has; it is not a new front.
  \item \emph{Depletion of $|s|\alpha_{\max}$ on profiles [Open, and
    currently \emph{worse} than on the original solution].} One needs
    $|s|\alpha_{\max}\leq 1-c(\delta)+|s|w_{\max}$ on average --- a
    Constantin--Fefferman depletion statement in self-similar variables. By
    Remark~\ref{rem:s39-nontransfer}(i), Proposition~\ref{prop:localized-cf}
    is \emph{not available on the profile}: its far-field step needs
    $\|\omega\|_{L^2(\R^3)}<\infty$, which fails for
    $U\in\mathcal{A}(c_I)$, and the crude substitute $|\omega_U|\leq M_U$
    makes the far-field kernel integral logarithmically divergent. A profile
    version therefore requires a genuinely different far-field control ---
    e.g.\ exploiting $\sup|U|\leq c_I|s|^{-1/2}$ through the Biot--Savart
    \emph{velocity} representation rather than the vorticity one, or a
    local-energy Type-I bound $\sup_{s}\int_{B_R}|U|^2\leq CR^3|s|^{-1}$.
    Not attempted. \emph{This is the harder of the two halves and is the
    honest reason Step 4 does not close in this session.}
  \end{enumerate}
\end{remark}

\subsection{Step 5: the reframed dichotomy, assembled}
\label{subsec:s39-step5}

\begin{proposition}[Two-step reduction: the oscillatory case is aligned]
  \label{prop:two-step-reduction}
  Let $U\in\mathcal{A}(c_I)$ be nontrivial with
  $\liminf_{s\to-\infty}\sigma^*(U)(s)=0$; pick $s_k\to-\infty$ with
  $\sigma^*(U)(s_k)\to0$ and set $\mu_k=|s_k|^{1/2}$, so that
  $U_{\mu_k}\in\mathcal{A}(c_I)$ and $\sigma^*(U_{\mu_k})(-1)\to0$. Assume
  \textbf{(R)} for the family $\{U_{\mu_k}\}$ and
  \begin{itemize}
  \item[\textbf{(N$'$)}] $\liminf_k\Phi_U(\log|s_k|)>0$ and no vorticity
    escape in the second extraction (i.e.\ the analogue of \textbf{(N)} for
    $\{U_{\mu_k}\}$).
  \end{itemize}
  Then there exists a nontrivial $V\in\mathcal{A}(c_I)$ with
  $\sigma^*(V)(-1)=0$.
  \emph{[Proved-modulo-(R),(N$'$)]}
\end{proposition}
\begin{proof}
  Lemma~\ref{lem:typeI-scaleinv-bounds} places every $U_{\mu_k}$ in
  $\mathcal{A}(c_I)$ with the same constant and gives
  $\sigma^*(U_{\mu_k})(-1)=\sigma^*(U)(-\mu_k^2)=\sigma^*(U)(s_k)\to0$.
  Compactness exactly as in Proposition~\ref{prop:profile-extraction} (the
  hypotheses are identical, the class being scaling-closed) yields
  $V\in\mathcal{A}(c_I)$ with $U_{\mu_k}\to V$ in
  $C^\infty_{\mathrm{loc}}$. Nontriviality of $V$: by \textbf{(N$'$)},
  $M_{U_{\mu_k}}(-1)=|s_k|M_U(s_k)=\Phi_U(\log|s_k|)\geq\phi_0>0$ along the
  sequence, and no-escape gives $M_V(-1)\geq\phi_0>0$. Finally
  $\sigma^*(V)(-1)=0$: under \textbf{(R)} this is
  Lemma~\ref{lem:sigma-equality}; even with only weak convergence it follows
  from the lower semicontinuity \eqref{eq:sigma-lsc} of
  Lemma~\ref{lem:sigma-semicontinuity}, which is the available direction.
\end{proof}

\begin{theorem}[Profile-rigidity dichotomy --- conditional]
  \label{thm:rigidity-dichotomy}
  Assume \textbf{(R)}, \textbf{(N)}, \textbf{(N$'$)}, \textbf{(A-up)}, the
  attainment convention of Remark~\ref{rem:attainment}, and
  Conjecture~\ref{target:disorder-depletion}. Then there is no Type-I
  singularity at a point $x^*$ satisfying \textbf{(N)}.
  \emph{[Conditional on the five named inputs. Not a theorem of this
  document in any stronger sense.]}
\end{theorem}
\begin{proof}
  Suppose there were. Proposition~\ref{prop:profile-extraction} produces a
  nontrivial $U\in\mathcal{A}(c_I)$. Exactly one of the following holds.

  \emph{Case 1: $\liminf_{s\to-\infty}\sigma^*(U)(s)=0$.}
  Proposition~\ref{prop:two-step-reduction} produces a nontrivial
  $V\in\mathcal{A}(c_I)$ with $\sigma^*(V)(-1)=0$ and $M_V(-1)>0$;
  Theorem~\ref{thm:aligned-horn} forces $V$ trivial --- contradiction.

  \emph{Case 2: $\liminf_{s\to-\infty}\sigma^*(U)(s)=:2\delta>0$.} Then
  $\sigma^*(U)(s)\geq\delta$ for all $s\leq-S_0$, some $S_0$.
  Conjecture~\ref{target:disorder-depletion} and
  Proposition~\ref{prop:logtime-necessary} then assert
  $1\leq 1-c(\delta)<1$ --- contradiction.
\end{proof}

\begin{remark}[What the reframe bought, stated without inflation]
  \label{rem:s39-bought}
  The reframe does not prove $\sigma^*$-decay. What it does is replace one
  open item (``force $\sigma^*(t)\to0$ dynamically'', for which every
  candidate mechanism so far has been a critical or supercritical recursion)
  by a \emph{different} set of open items, of which the two substantive ones
  --- \textbf{(A-up)} and Conjecture~\ref{target:disorder-depletion} --- are
  attacked by unrelated techniques (unique continuation / Liouville rigidity;
  log-time monotonicity). The concrete gains recorded here are:
  \begin{enumerate}[label=(\arabic*)]
  \item \emph{The aligned horn is genuinely closed modulo one hypothesis.}
    Lemmas~\ref{lem:alignment-local} and \ref{lem:harmonic-reduction} plus
    the cited 2D Liouville theorem give a complete argument once (A-up) is
    granted, and the 2D reduction in
    Lemma~\ref{lem:harmonic-reduction} is a full proof, not a sketch.
  \item \emph{Profiles are logarithm-free}
    (Proposition~\ref{prop:transfer-audit}(d),(e)): every marginal
    $\mathcal{L}$-degraded estimate of
    \S\ref{sec:moser-s33}--\S\ref{sec:verification-s36} improves to an
    absolute constant on the profile.
  \item \emph{The naive infinite-past argument is proved impossible}
    (Proposition~\ref{prop:no-summable-budget}), and the correct framework
    --- bounded amplitude $\Phi$ monotone in log-time --- is identified,
    with one \emph{proved} averaged necessary condition
    \eqref{eq:necessary-average} extracted from it.
  \item \emph{A previously unnoticed obstruction is recorded}
    (Remark~\ref{rem:s39-nontransfer}(i)): the \S\ref{sec:cf-s37}
    Constantin--Fefferman machinery does \emph{not} survive the passage to
    the profile, because $\|\omega\|_{L^2(\R^3)}=\infty$ there. Any future
    session that uses Proposition~\ref{prop:localized-cf} on a profile is in
    error.
  \end{enumerate}
\end{remark}

\subsection{Ledger and distance to the goal (post-S39)}

\begin{center}
\begin{tabular}{lll}
\hline
Item & Status & Where \\
\hline
Type-I ancient class scaling-closed & Proved & Lem.~\ref{lem:typeI-scaleinv-bounds} \\
Rescaled bounds $\lambda$-uniform & Proved & \eqref{eq:rescaled-bounds} \\
Scale-uniform interior regularity (R) & \textbf{Conditional} (Serrin-type) & \S\ref{subsec:s39-step1} \\
Non-degeneracy / no escape (N) & \textbf{Assumed}; CKN gives it only on $Q_1(0)$ & Rem.~\ref{rem:N-discussion} \\
Profile extraction & Proved mod (R),(N) & Prop.~\ref{prop:profile-extraction} \\
Eqns \eqref{eq:mag-eq},\eqref{eq:dir-eq}, split, Gram on $U$ & \textbf{Proved} (re-derived, not transferred) & Prop.~\ref{prop:transfer-audit}(a)--(c) \\
Profiles are $\mathcal{L}$-free & Proved mod (R) & Prop.~\ref{prop:transfer-audit}(d) \\
$\nu\langle\sigma^*\rangle\leq C(c_I)$ on $U$ & Proved mod (R) & Prop.~\ref{prop:transfer-audit}(e) \\
$\sigma^*$ transfer: lower semicontinuity & \textbf{Proved} (available direction) & Lem.~\ref{lem:sigma-semicontinuity} \\
$\sigma^*$ transfer: equality & Proved mod (R),(N) & Lem.~\ref{lem:sigma-equality} \\
Pull-back of H1 to the sequence & \textbf{Forbidden} (no reverse inequality) & Rem.~\ref{rem:s39-nontransfer}(ii) \\
Prop.~\ref{prop:localized-cf} on profiles & \textbf{Fails} ($\Omega=\infty$) & Rem.~\ref{rem:s39-nontransfer}(i) \\
Local alignment from $\sigma^*=0$ & \textbf{Proved} & Lem.~\ref{lem:alignment-local} \\
Global alignment upgrade (A-up) & \textbf{Open} --- largest new gap & Rem.~\ref{rem:A-up-discussion} \\
Aligned $\Rightarrow$ 2D ($\partial_3U=0$) & \textbf{Proved} mod (R) & Lem.~\ref{lem:harmonic-reduction} \\
2D bounded ancient Liouville & Cited \cite{KNSS2009}; class caveat & Thm.~\ref{thm:knss} \\
Aligned horn & Proved mod (R),(A-up) & Thm.~\ref{thm:aligned-horn} \\
Naive finite-budget argument & \textbf{Proved impossible} & Prop.~\ref{prop:no-summable-budget} \\
Log-time amplitude ODE & Proved mod (R), attainment & Lem.~\ref{lem:logtime} \\
Averaged necessary condition & \textbf{Proved} mod (R), attainment & Prop.~\ref{prop:logtime-necessary} \\
Disorder--depletion target & \textbf{Open target}, marginal & Conj.~\ref{target:disorder-depletion} \\
\quad (T4a) average-to-point for $w$ & Open; plugs into $K$-family & Rem.~\ref{rem:target-decomposition} \\
\quad (T4b) profile depletion of $\alpha$ & Open; needs new far field & Rem.~\ref{rem:target-decomposition} \\
Two-step reduction & Proved mod (R),(N$'$) & Prop.~\ref{prop:two-step-reduction} \\
Reframed dichotomy & Conditional on 5 inputs & Thm.~\ref{thm:rigidity-dichotomy} \\
Type-II exclusion & Open (outside all present methods) & Rem.~\ref{rem:s39-distance} \\
\hline
\end{tabular}
\end{center}

\begin{remark}[Distance to the goal]\label{rem:s39-distance}
  Suppose both horns closed completely --- \textbf{(A-up)} proved and
  Conjecture~\ref{target:disorder-depletion} proved, with \textbf{(R)},
  \textbf{(N)}, \textbf{(N$'$)} discharged from the Type-I literature. The
  conclusion would be: \emph{no Type-I singularity}. That is strictly less
  than the Clay problem. What would remain:
  \begin{enumerate}[label=(\arabic*)]
  \item \textbf{Type-II exclusion.} Every estimate in
    \S\ref{sec:moser-s33}--\S\ref{sec:rigidity-s39} uses the Type-I rates,
    and the profile-extraction of \S\ref{subsec:s39-step1} uses them
    \emph{essentially}: without $M(t)\lesssim(T-t)^{-1}$ the rescaled family
    \eqref{eq:rescaling} has no uniform bound and there is no limit object at
    all. A Type-II singularity would rescale to nothing, and the entire
    reframe would be vacuous against it. This is outside these methods, and
    no part of this document should be read as progress on it.
  \item \textbf{The nontriviality input.} Even in the Type-I case,
    \textbf{(N)} is not a technicality: it is the assertion that the
    singularity's strength does not escape the rescaling window. Discharging
    it needs a Type-I-improved $\varepsilon$-regularity theorem
    \cite{Seregin2012}; until then, ``no Type-I singularity'' should be read
    as ``no Type-I singularity at a non-degenerate point''.
  \item \textbf{Marginality of the surviving target.}
    Conjecture~\ref{target:disorder-depletion} must beat the constant $1$
    exactly. The program's history (\S\ref{sec:audit-s31}) is that marginal
    targets are where errors hide; the mitigating structural difference is
    noted after Proposition~\ref{prop:logtime-necessary}, but it is a
    mitigation, not a proof.
  \end{enumerate}
  Accordingly the honest one-line summary of this section is: \emph{the
  single deep gap of \S\ref{sec:closure-s38} has been replaced by two
  structurally unrelated gaps, one of which (the aligned horn) is now
  complete except for a single clearly stated propagation hypothesis, and one
  of which (the disordered horn) has yielded a proved averaged necessary
  condition and a sharply-posed marginal target.} No claim is made beyond
  this.
\end{remark}

% ══════════════════════════════════════════════════════════════════════════════
\section{Execution of the Level Iteration for the Annulus Residual:
an Exact Inventory and a Structural Obstruction (Session S41)}
\label{sec:annulus-s41}

This section carries out the programme announced in
Remark~\ref{rem:level-iteration}: the nested De~Giorgi level iteration that
was to remove the annulus residual $\mathcal{A}_{\mathrm{ann}}$
(Definition~\ref{def:annulus-residual}) from
Proposition~\ref{prop:K-amended}, the last item standing between H1's
estimate layer and a proof. The outcome is \emph{mixed, and on the decisive
point negative}; it is reported as such.

\begin{enumerate}[label=(\alph*)]
  \item The nested levels and cutoffs are constructed and their derivative
    bounds computed; the constant $C\,4^{k}$ asserted in
    Remark~\ref{rem:level-iteration} is confirmed, and its origin
    identified: it is $|\chi_k'|^2$, produced when the second-order term
    $\nu\Delta m$ inside $D_t\chi_k$ is integrated by parts
    (\S\ref{subsec:s41-setup}).
  \item An \emph{exact} inventory of the $\chi$-derivative families is
    derived from the $K$-identity of Proposition~\ref{prop:K-self}
    (\S\ref{subsec:s41-inventory}). This is a genuine positive result: four
    of the eight families have a \emph{favourable sign} and may simply be
    discarded, and two more fall in the inhomogeneous class already allowed
    by Proposition~\ref{prop:K-amended}. $\mathcal{A}_{\mathrm{ann}}$
    shrinks to an irreducible core of two families
    (Proposition~\ref{prop:s41-core}).
  \item The recursion for that core is derived
    (\S\ref{subsec:s41-recursion}). Its coefficient in front of the
    level-$(k+1)$ quantity is \emph{never smaller than $1$}, for
    \emph{any} choice of the Young parameters $\delta_k$; the proposed
    $\delta_k=\delta_0 8^{-k}$ makes it strictly worse
    (Proposition~\ref{prop:s41-nocontraction}). The unrolled iteration is
    weaker than the trivial inequality and carries no information
    (Proposition~\ref{prop:s41-vacuous}).
  \item Neither classical iteration lemma applies: the linear
    (Ladyzhenskaya--Ural'tseva / Giaquinta--Giusti hole-filling) lemma needs
    a contraction factor $<1$; the nonlinear (Stampacchia--De~Giorgi) lemma
    needs a superlinear gain running from coarse levels to fine ones. Here
    the recursion runs the \emph{opposite} way --- the level sets
    \emph{grow} with $k$ --- and is homogeneous of degree one
    (Remark~\ref{rem:s41-whichlemma}).
  \item The $k\to\infty$ limit does not annihilate the residual: it
    converges to a \emph{flux through the level set} $\{|\omega|=M/8\}$, and
    a co-area identity shows that \emph{no} admissible choice of levels ---
    nested, geometric, uniform, or optimally selected --- reduces the
    residual below the level density of the underlying measure, whose only
    a priori bound is a fixed multiple of the band integral
    (\S\ref{subsec:s41-coarea}). This is the crux, and it is not
    bookkeeping.
  \item The core residual does admit the unconditional bound
    $\mathcal{A}_{\mathrm{ann}}^{\mathrm{core}}
    \leq C\,\mathcal{L}\,\|w\|_{L^\infty}M^{1/2}$
    (Proposition~\ref{prop:s41-crude}) --- which is \emph{exactly} the bound
    of Proposition~\ref{prop:crude-K} and therefore falls short of
    absorption by exactly one logarithm.
\end{enumerate}

\noindent Accordingly $\mathcal{A}_{\mathrm{ann}}$ is reclassified from
\emph{open-mechanical} to \emph{open}, of the same analytic type as the
$c_K$-decorrelation question (Conjecture~\ref{conj:K-refined}); the
consolidated table of \S\ref{sec:verification-s36} is amended accordingly.
Nothing else in the document is affected: this section touches only the
$\mathcal{A}_{\mathrm{ann}}$ row of H1's estimate layer, and in particular
makes no statement about H2, $\sigma^*$-decay, Type-II exclusion, or the
withdrawn \S\ref{sec:audit-s31} material.

\subsection{Admissible levels, the nested family, and the derivative bounds}
\label{subsec:s41-setup}

Throughout, $m=|\omega|$, $w=|\nabla\hat{e}|^2$,
$M=M(t)=\|\omega(\cdot,t)\|_{L^\infty}$, $r^*=M^{-1/2}$, and $\psi$ is the
fixed space--time cutoff of Definition~\ref{def:composite-cutoff}
($\psi\equiv1$ on $Q(\rho)$, $\operatorname{supp}\psi\subseteq Q(\rho')$,
$\tfrac{r^*}{2}\leq\rho<\rho'\leq r^*$). Only the $|\omega|$-level of the
composite cutoff is iterated below; the spatial nesting $\rho<\rho'$ is
untouched, and terms carrying $\nabla\psi$ or $\partial_t\psi$ remain in the
$(\rho'-\rho)^{-2}$ family of Proposition~\ref{prop:first-rung} and are not
annulus terms.

\begin{definition}[Admissible level cutoffs]\label{def:s41-admissible}
  For $\tfrac{1}{16}\leq\theta_-<\theta_+\leq\tfrac14$ let
  $\mathcal{C}(\theta_-,\theta_+)$ be the set of nondecreasing
  $\chi\in C^\infty([0,\infty);[0,1])$ with $\chi\equiv0$ on $[0,\theta_-]$
  and $\chi\equiv1$ on $[\theta_+,\infty)$. For $\chi\in\mathcal{C}$ write
  $\eta_\chi=\psi\cdot\chi(m/M)$.
\end{definition}

\begin{lemma}[Uniform admissibility of every level pair]
  \label{lem:s41-admissible}
  Let $\chi\in\mathcal{C}(\theta_-,\theta_+)$ with
  $\tfrac1{16}\leq\theta_-<\theta_+\leq\tfrac14$. Then
  \begin{enumerate}[label=(\roman*)]
    \item $0\leq\chi\leq1$ and $\chi\equiv1$ on $\{m\geq M/4\}$, hence a
      fortiori on $\{m\geq M/2\}$;
    \item $\operatorname{supp}\chi(m/M)\subseteq\{m\geq M/16\}$, so
      $M/16\leq m\leq M$ on $\operatorname{supp}\eta_\chi$;
    \item $\chi'\geq0$.
  \end{enumerate}
  These, together with the derivative bound
  $\|\chi'\|_\infty\leq2/(\theta_+-\theta_-)\leq32$, are the properties of
  the $\chi$-factor used in
  \S\ref{sec:magnitude-s34}--\S\ref{sec:verification-s36}
  (see the proof). Consequently
  Proposition~\ref{prop:K-amended} holds verbatim with $\eta$ replaced by
  $\eta_\chi$, with constants independent of the pair
  $(\theta_-,\theta_+)$ in that range; in particular the inhomogeneous term
  \begin{equation}\label{eq:s41-I}
    \mathcal{I} \;:=\; C\,\mathcal{L}^2\,M^{3/2}\,
      \bigl(\nu^{-1}+\nu^{-1/2}(\sigma^*)^{1/2}\bigr)
  \end{equation}
  is level-independent, as is the a priori enstrophy budget
  \begin{equation}\label{eq:s41-budget}
    \mathcal{E} \;:=\; \nu\iint_{Q(r^*)}|\nabla\omega|^2
    \;=\; \nu\iint_{Q(r^*)}\bigl(|\nabla m|^2+m^2w\bigr)
    \;\leq\; C\,\mathcal{L}\,M^{1/2}.
  \end{equation}
\end{lemma}
\begin{proof}
  (i)--(iii) are immediate from Definition~\ref{def:s41-admissible}. For the
  remaining assertions: Lemma~\ref{lem:w2-cancel} uses only
  $\hat{e}\cdot\partial_k\hat{e}=0$ and is cutoff-blind;
  Lemma~\ref{lem:gram-negativity} uses only $\eta\geq0$ and the Gram
  positivity; Lemma~\ref{lem:stretch-couple} and \S\ref{sec:verification-s36}
  (C2) use only $\eta\leq1$, the bound
  $|\alpha|\leq\|\nabla u\|_{L^\infty}\leq CM\mathcal{L}$ and
  Theorem~\ref{thm:apriori-sigma}, whose own proof needs only
  $\{m\geq M/2\}\subseteq\{\chi=1\}$, i.e.~(i). Wherever $m$ occurs in a
  denominator the level enters through the lower bound of (ii), replacing
  the numerical factor $m^{-1}\leq 8/M$ (levels $\tfrac18,\tfrac14$) by
  $m^{-1}\leq16/M$: a change of absolute constants only. Likewise
  Lemma~\ref{lem:mag-rhs}(c) of \S\ref{sec:magnitude-s34} uses the
  $\chi$-factor only through $\|\chi'\|_\infty$ and through the pointwise
  dominations $m^2w\leq|\nabla\omega|^2$, $|\nabla m|^2\leq|\nabla\omega|^2$,
  which are level-free.
  \eqref{eq:s41-budget} is Lemma~\ref{lem:grad-omega-split} combined with
  Lemma~\ref{lem:local-enstrophy-typeI} and $M^2(r^*)^3=M^{1/2}$; it
  involves no cutoff at all.
\end{proof}

\begin{remark}[Why the level pair is genuinely free]\label{rem:s41-free}
  Lemma~\ref{lem:s41-admissible} is what makes any level-selection argument
  legitimate downstream: a bound proved for \emph{one} admissible
  $\chi\in\mathcal{C}(\theta_-,\theta_+)$, $\theta_\pm$ in the stated range,
  feeds Proposition~\ref{prop:first-rung} and
  Theorem~\ref{thm:apriori-sigma} exactly as the original
  $(\tfrac18,\tfrac14)$ pair does. This positive observation is used in
  \S\ref{subsec:s41-coarea}; it is also the reason the failure established
  there is a failure of the \emph{method} and not of a particular
  normalisation.
\end{remark}

\begin{definition}[The nested family of Remark~\ref{rem:level-iteration}]
  \label{def:s41-nested}
  Put $\theta_k=\tfrac18+2^{-k-3}$ for $k=0,1,2,\dots$, so $\theta_0=\tfrac14$,
  $\theta_k\downarrow\tfrac18$, and
  \[
    \ell_k \;:=\; \theta_k-\theta_{k+1} \;=\; 2^{-k-3}-2^{-k-4}\;=\;2^{-k-4}.
  \]
  Fix once and for all a nondecreasing $\Theta\in C^\infty(\R;[0,1])$ with
  $\Theta\equiv0$ on $(-\infty,0]$, $\Theta\equiv1$ on $[1,\infty)$,
  $0\leq\Theta'\leq2$, $|\Theta''|\leq16$, and set
  \[
    \chi_k(s)\;=\;\Theta\Bigl(\frac{s-\theta_{k+1}}{\ell_k}\Bigr)
      \;\in\;\mathcal{C}(\theta_{k+1},\theta_k),
    \qquad
    \eta_k \;=\; \psi\,\chi_k(m/M),
  \]
  \[
    \mathrm{Ann}_k \;=\; \{\theta_{k+1}M\leq m\leq\theta_kM\}
      \cap\operatorname{supp}\psi \;\supseteq\;
      \operatorname{supp}\bigl[\chi_k'(m/M)\bigr]\cap\operatorname{supp}\psi .
  \]
\end{definition}

\begin{lemma}[Derivative bounds; the origin of $4^{k}$]\label{lem:s41-derivs}
  For every $k\geq0$,
  \begin{align}
    0\;\leq\;\chi_k' &\;\leq\; 2/\ell_k \;=\; 2^{k+5},
      \label{eq:s41-chip}\\
    |\chi_k''| &\;\leq\; 16/\ell_k^2 \;=\; 2^{2k+12},
      \label{eq:s41-chipp}\\
    \bigl|\nabla[\chi_k(m/M)]\bigr| &\;=\; \chi_k'\,|\nabla m|/M
      \;\leq\; 2^{k+5}\,|\nabla m|/M .
      \label{eq:s41-gradchi}
  \end{align}
  Hence the combinations occurring \emph{before} the integration by parts of
  Lemma~\ref{lem:three-ibp}(b), namely $\chi_k$ and $\chi_k\chi_k'$, obey
  $\chi_k\chi_k'\leq 2^{5}\cdot2^{k}$, whereas the combinations occurring
  \emph{after} it, namely $(\chi_k')^2$ and $\chi_k|\chi_k''|$, obey
  $(\chi_k')^2+\chi_k|\chi_k''|\leq 2^{13}\cdot4^{k}$. The constant
  $C\,4^{k}$ of Remark~\ref{rem:level-iteration} is therefore correct, and
  it is the \emph{square} of the cutoff derivative: it is created by
  $\nu\Delta m$ inside $D_t\chi_k$ when that term is integrated by parts, as
  it must be (Lemma~\ref{lem:three-ibp}(b)).
\end{lemma}
\begin{proof}
  Chain rule: $\chi_k'(s)=\Theta'((s-\theta_{k+1})/\ell_k)/\ell_k$ and
  $\chi_k''(s)=\Theta''(\cdot)/\ell_k^2$; insert
  $\ell_k^{-1}=2^{k+4}$ and $\ell_k^{-2}=2^{2k+8}$, and use
  $\|\Theta'\|_\infty\leq2$, $\|\Theta''\|_\infty\leq16$; then
  $2\cdot2^{k+4}=2^{k+5}$ and $16\cdot 2^{2k+8}=2^{2k+12}=2^{12}\cdot4^{k}$.
  \eqref{eq:s41-gradchi} is $\nabla[\chi_k(m/M)]=\chi_k'(m/M)\nabla m/M$
  ($M=M(t)$ is $x$-independent). Finally
  $(\chi_k')^2\leq2^{2k+10}$ and $\chi_k|\chi_k''|\leq2^{2k+12}$, whose sum
  is $\leq2^{2k+13}=2^{13}4^{k}$.
\end{proof}

\begin{lemma}[The level-$k$ annulus carries full level-$(k+1)$ weight]
  \label{lem:s41-fullweight}
  For every $k\geq0$: $\chi_{k+1}(m/M)\equiv1$ on $\mathrm{Ann}_k$, hence
  $\eta_{k+1}^2=\psi^2$ there. Moreover $\chi_{k+1}\geq\chi_k$ pointwise, the
  interiors of the $\mathrm{Ann}_k$ are pairwise disjoint, and
  $\bigcup_{k\geq0}\mathrm{Ann}_k
   =\{M/8<m\leq M/4\}\cap\operatorname{supp}\psi=:\mathrm{Ann}$,
  the annulus of Definition~\ref{def:annulus-residual} up to the level set
  $\{m=M/8\}$.
\end{lemma}
\begin{proof}
  $\chi_{k+1}\equiv1$ on $[\theta_{k+1},\infty)$ while
  $\mathrm{Ann}_k\subseteq\{m/M\geq\theta_{k+1}\}$. For monotonicity: if
  $s\geq\theta_{k+1}$ then $\chi_{k+1}(s)=1\geq\chi_k(s)$; if
  $s\leq\theta_{k+1}$ then $\chi_k(s)=0\leq\chi_{k+1}(s)$. Disjointness and
  the union statement follow from $\theta_{k+1}<\theta_k$ and
  $\theta_k\downarrow\tfrac18$.
\end{proof}

\subsection{Exact inventory of the $\chi$-derivative families}
\label{subsec:s41-inventory}

Definition~\ref{def:annulus-residual} specifies $\mathcal{A}_{\mathrm{ann}}$
by the \emph{shapes} of its terms, not by a term-by-term derivation. Since
the whole question is whether these terms can be removed, we first derive
them exactly. The following identity is the one behind
Proposition~\ref{prop:K-self}; it is exact (no inequalities), and we record
it in the form in which the cutoff derivatives are visible.

\begin{lemma}[The exact weighted magnitude identity]\label{lem:s41-identity}
  Let $\eta=\psi\chi(m/M)$ with $\chi\in\mathcal{C}$. Testing
  \eqref{eq:mag-eq} against $m\,w\,\eta^2$ and integrating by parts in space
  and in time (using $\divg u=0$, so that
  $\iint(D_tf)g=-\iint f\,D_tg$ up to time slices) gives
  \begin{align}
    \nu\iint|\nabla m|^2w\eta^2
    &+\nu\iint m^2w^2\eta^2
      \;+\;[\text{time slices}] \notag\\
    &=\;\underbrace{\iint\alpha\,m^2w\,\eta^2}_{R_1}
      \;\underbrace{-\;\nu\iint m\,\eta^2\,\nabla m\cdot\nabla w}_{R_2}
      \;\underbrace{-\;\nu\iint m\,w\,\nabla m\cdot\nabla(\eta^2)}_{R_3}
      \notag\\
    &\qquad
      \;+\;\underbrace{\tfrac12\iint m^2\,(D_tw)\,\eta^2}_{R_4}
      \;+\;\underbrace{\tfrac12\iint m^2\,w\,D_t(\eta^2)}_{R_5} .
      \label{eq:s41-identity}
  \end{align}
\end{lemma}
\begin{proof}
  Multiply \eqref{eq:mag-eq} by $mw\eta^2$ and integrate:
  $\iint(D_tm)mw\eta^2=\iint\alpha m^2w\eta^2
  +\nu\iint(\Delta m)mw\eta^2-\nu\iint m^2w^2\eta^2$.
  Expand $\nabla(mw\eta^2)=w\eta^2\nabla m+m\eta^2\nabla w+mw\nabla(\eta^2)$
  in $\nu\iint(\Delta m)mw\eta^2=-\nu\iint\nabla m\cdot\nabla(mw\eta^2)$,
  and $\iint(D_tm)mw\eta^2=\tfrac12\iint D_t(m^2)\,w\eta^2
  =-\tfrac12\iint m^2D_t(w\eta^2)+[\text{slices}]$. Rearranging gives
  \eqref{eq:s41-identity}. The middle viscous cross-term $R_2$ is the one
  absorbed by Cauchy--Schwarz in Proposition~\ref{prop:K-self}; $R_1$ is C2
  and $R_4$ is C1 of \S\ref{sec:verification-s36}.
\end{proof}

\begin{proposition}[Inventory: four families die by sign, two are
  inhomogeneous, two survive]\label{prop:s41-core}
  Write $\chi=\chi(m/M)$, $\chi'=\chi'(m/M)$, $\chi''=\chi''(m/M)$, and
  isolate in \eqref{eq:s41-identity} all terms carrying a
  $\chi$-derivative (all remaining terms carry either $\eta^2$ in full, or
  $\nabla\psi$/$\partial_t\psi$ and belong to the $(\rho'-\rho)$-family).
  Then, exactly:
  \begin{enumerate}[label=\textup{(G\arabic*)}]
    \item from $R_3$:
      $\displaystyle -\frac{2\nu}{M}\iint m\,w\,|\nabla m|^2\,\psi^2\,\chi\chi'
      \;\leq\;0$;
    \item from $R_5$ through the sink $-\nu mw$ of \eqref{eq:mag-eq}:
      $\displaystyle -\frac{\nu}{M}\iint m^3w^2\,\psi^2\,\chi\chi'\;\leq\;0$;
    \item from $R_5$ through $\nu\Delta m$, after one integration by parts,
      the $(\chi')^2$ piece:
      $\displaystyle -\frac{\nu}{M^2}\iint m^2w\,|\nabla m|^2\,\psi^2\,(\chi')^2
      \;\leq\;0$;
    \item from the same integration by parts, the $\nabla(m^2)$ piece:
      $\displaystyle -\frac{2\nu}{M}\iint m\,w\,|\nabla m|^2\,\psi^2\,\chi\chi'
      \;\leq\;0$.
  \end{enumerate}
  \begin{enumerate}[label=\textup{(I\arabic*)}]
    \item from $R_5$ through $\alpha m$:
      $\displaystyle \frac{1}{M}\iint \alpha\,m^3w\,\psi^2\,\chi\chi'$,
      bounded in absolute value by
      $\|\chi\chi'\|_\infty\,\|\nabla u\|_{L^\infty}\nu^{-1}\mathcal{E}
      \leq C\|\chi\chi'\|_\infty\mathcal{L}^2M^{3/2}\nu^{-1}$;
    \item from $R_5$ through $-m\dot M/M$: the same bound, using
      $|\dot M|/M\leq\|\nabla u\|_{L^\infty}\leq CM\mathcal{L}$.
  \end{enumerate}
  \begin{enumerate}[label=\textup{(S\arabic*)}]
    \item the $\chi\chi''$ piece of the integration by parts in (G3):
      $\displaystyle -\frac{\nu}{M^2}\iint m^2w\,|\nabla m|^2\,\psi^2\,\chi\chi''$,
      of \emph{indefinite} sign, with
      $|{\cdot}|\leq\tfrac{1}{16}\|\chi\chi''\|_\infty\,
      \nu\iint_{\{\theta_-M\leq m\leq\theta_+M\}}|\nabla m|^2w\,\psi^2$;
    \item the cross family, arising both from $R_5$ (through
      $\nu\Delta m$, the $\nabla w$ piece) and from $R_4$ (the
      $\nu\Delta\partial_k\hat{e}$ piece of the differentiated direction
      equation, cf.\ Lemma~\ref{lem:three-ibp}(a)):
      $\displaystyle \frac{c\,\nu}{M}\iint m^2\,\psi^2\,\chi\chi'\,
      \nabla m\cdot\nabla w$, $|c|\leq2$, of indefinite sign.
  \end{enumerate}
  Consequently, discarding (G1)--(G4) and absorbing (I1)--(I2) into
  $\mathcal{I}$ of \eqref{eq:s41-I},
  \begin{equation}\label{eq:s41-core}
    \mathcal{A}_{\mathrm{ann}}
    \;\leq\;
    \mathcal{A}^{\mathrm{core}}_{\mathrm{ann}}
    \;:=\;
    \Bigl|\text{\upshape(S1)}\Bigr| + \Bigl|\text{\upshape(S2)}\Bigr| .
  \end{equation}
\end{proposition}
\begin{proof}
  Write $\nabla(\eta^2)=\chi^2\nabla(\psi^2)+2\psi^2\chi\chi'\nabla m/M$ and
  $D_t(\eta^2)=\chi^2D_t(\psi^2)
  +2\psi^2\chi\chi'\bigl(D_tm/M-m\dot M/M^2\bigr)$; the $\nabla(\psi^2)$ and
  $D_t(\psi^2)$ parts are $(\rho'-\rho)$-family terms by
  Definition~\ref{def:composite-cutoff}, not annulus terms. Substituting the
  first display into $R_3$ gives (G1), whose sign is
  $\leq0$ because $m,w\geq0$ and $\chi,\chi'\geq0$
  (Lemma~\ref{lem:s41-admissible}(iii)). Substituting the second into $R_5$
  and then $D_tm=\alpha m+\nu\Delta m-\nu mw$ from \eqref{eq:mag-eq} gives
  four contributions: the $-\nu mw$ one is (G2), $\leq0$ for the same
  reason; the $\alpha m$ and $-m\dot M/M$ ones are (I1)--(I2), estimated by
  $|\alpha|\leq\|\nabla u\|_{L^\infty}$, $m/M\leq1$, $\psi^2\leq1$ and
  $\nu\iint_{Q(\rho')}m^2w\leq\mathcal{E}$ from
  \eqref{eq:s41-budget}, together with $M\mathcal{L}\cdot\mathcal{L}M^{1/2}
  =\mathcal{L}^2M^{3/2}$; the $\nu\Delta m$ one is
  $\frac{\nu}{M}\iint m^2w\psi^2\chi\chi'\Delta m
  =-\frac{\nu}{M}\iint\nabla m\cdot\nabla\bigl(m^2w\psi^2\chi\chi'\bigr)$
  (no boundary terms: the integrand vanishes off
  $\operatorname{supp}\psi\cap\{\chi'\neq0\}$), and expanding the gradient by
  the product rule yields exactly the four pieces (G3), (G4), (S1), (S2)
  --- using $\nabla[\chi\chi'(m/M)]
  =\bigl((\chi')^2+\chi\chi''\bigr)\nabla m/M$ --- plus one further piece
  carrying $\nabla(\psi^2)$, again a $(\rho'-\rho)$-family term. Signs: in
  (G3) the factor $(\chi')^2\geq0$ enters with an overall minus, in (G4)
  the factor $\chi\chi'\geq0$ likewise. The stated bound in (S1) uses
  $m^2/M^2\leq\tfrac{1}{16}$ on $\operatorname{supp}\chi'$.
  For (S2), the $R_4$ contribution has the same shape by
  Lemma~\ref{lem:three-ibp}(a) with $\eta^2$ replaced by $m^2\eta^2$, and
  $|\nabla w|\leq2w^{1/2}|\nabla^2\hat{e}|$ is not yet used.
\end{proof}

\begin{remark}[What is and is not established by
  Proposition~\ref{prop:s41-core}]\label{rem:s41-caveat}
  Two honest caveats.
  \begin{enumerate}[label=(\roman*)]
    \item Families (G1)--(G4), (I1)--(I2), (S1) are derived from the
      \emph{exact} identity \eqref{eq:s41-identity} and are therefore
      complete for $R_1,R_2,R_3,R_5$. Family $R_4$ (coupling C1) is
      \emph{not} displayed term-by-term in \S\ref{sec:verification-s36};
      what is asserted above for it is only that its
      $\chi$-derivative terms have the shape (S2), by the structural
      parallel with Lemma~\ref{lem:three-ibp}(a). If a full audit of C1
      produced $\chi$-derivative terms of a shape not listed here, the
      inventory would have to be extended --- but this could only
      \emph{enlarge} $\mathcal{A}^{\mathrm{core}}_{\mathrm{ann}}$, so the
      negative conclusion of \S\ref{subsec:s41-recursion}--\S\ref{subsec:s41-coarea}
      is unaffected; only the positive statement
      \eqref{eq:s41-core} is conditional on that shape.
    \item Definition~\ref{def:annulus-residual} writes the annulus weight as
      $\psi|\chi'|$. The derivation above shows the weight is in fact
      $\psi^2$ times a $\chi$-derivative factor: the $\chi$-derivatives are
      produced by $\nabla(\eta^2)$ and $D_t(\eta^2)$, both of which retain
      the full spatial factor $\psi^2$. We read
      Definition~\ref{def:annulus-residual} accordingly. Every estimate
      below uses only $\psi\leq1$, hence holds in either reading.
  \end{enumerate}
\end{remark}

\subsection{The recursion, and why no choice of $\delta_k$ contracts}
\label{subsec:s41-recursion}

\begin{definition}[Level quantities]\label{def:s41-levelquantities}
  With $\eta_k$ as in Definition~\ref{def:s41-nested}, set
  \[
    Y_k=\nu\iint|\nabla m|^2\,w\,\eta_k^2,\qquad
    D_k=\nu\iint|\nabla^2\hat{e}|^2\,\eta_k^2,
  \]
  and let $D_\flat=\nu\iint_{\{m>M/8\}}|\nabla^2\hat{e}|^2\psi^2$, so
  $D_k\leq D_\flat$ for all $k$ by Lemma~\ref{lem:s41-fullweight}. By that
  lemma $\chi_{k+1}\geq\chi_k$, hence
  \begin{equation}\label{eq:s41-monotone}
    Y_0\;\leq\;Y_1\;\leq\;Y_2\;\leq\;\cdots\;\uparrow\;
    Y_\infty=\nu\iint_{\{m>M/8\}}|\nabla m|^2w\,\psi^2 .
  \end{equation}
\end{definition}

\begin{proposition}[The level-$k$ recursion]\label{prop:s41-recursion}
  For every $k\geq0$ and every $\delta_k>0$,
  \begin{equation}\label{eq:s41-recursion}
    Y_k \;\leq\; G_k \;+\;\Lambda_k\,Y_{k+1},
    \qquad
    G_k \;=\; C_*\bigl(M^2D_\flat+2^{k}\mathcal{I}\bigr)+\delta_k M^2D_\flat,
    \qquad
    \Lambda_k \;=\; 2^{8}\,4^{k}\;+\;2^{8}\,\delta_k^{-1}\,4^{k}.
  \end{equation}
  (The factor $2^{k}$ on $\mathcal{I}$ records that families (I1)--(I2) of
  Proposition~\ref{prop:s41-core} carry $\|\chi_k\chi_k'\|_\infty\leq2^{k+5}$;
  it only strengthens the negative conclusions below.)
\end{proposition}
\begin{proof}
  Proposition~\ref{prop:K-amended} at level $k$
  (legitimate by Lemma~\ref{lem:s41-admissible}) gives
  $Y_k\leq C_*(M^2D_k+\mathcal{I})+\mathcal{A}_k
  \leq C_*(M^2D_\flat+2^{k}\mathcal{I})+\mathcal{A}^{\mathrm{core}}_k$, where
  $\mathcal{A}^{\mathrm{core}}_k$ is \eqref{eq:s41-core} formed with
  $\chi_k$. For (S1), by Proposition~\ref{prop:s41-core},
  Lemma~\ref{lem:s41-derivs} and Lemma~\ref{lem:s41-fullweight}
  ($\eta_{k+1}^2=\psi^2$ on $\mathrm{Ann}_k$, and
  $\operatorname{supp}\chi_k'\subseteq\mathrm{Ann}_k$),
  \[
    |\text{\upshape(S1)}_k|
    \;\leq\;\tfrac{1}{16}\,2^{2k+12}\;
      \nu\iint_{\mathrm{Ann}_k}|\nabla m|^2w\,\psi^2
    \;\leq\; 2^{8}\,4^{k}\,Y_{k+1}.
  \]
  For (S2), $|c|\leq2$ and $|\nabla w|\leq2w^{1/2}|\nabla^2\hat{e}|$ give
  $|\text{\upshape(S2)}_k|\leq\frac{4\nu}{M}\iint m^2\psi^2\chi_k\chi_k'
  |\nabla m|\,w^{1/2}|\nabla^2\hat{e}|$; then Young's inequality
  $4XZ\leq\delta_kX^2+4\delta_k^{-1}Z^2$ with
  $X=m|\nabla^2\hat{e}|\,\psi\mathbf 1_{\mathrm{Ann}_k}$ and
  $Z=\tfrac{m}{M}\chi_k\chi_k'|\nabla m|w^{1/2}\psi$ gives
  \[
    |\text{\upshape(S2)}_k| \;\leq\;
      \delta_k\,\nu\iint m^2|\nabla^2\hat{e}|^2\psi^2
        \mathbf 1_{\mathrm{Ann}_k}
      \;+\;4\delta_k^{-1}\,\frac{\nu}{M^2}
        \iint m^2(\chi_k\chi_k')^2|\nabla m|^2w\,\psi^2
    \;\leq\; \delta_k M^2D_\flat + 2^{8}\delta_k^{-1}4^{k}Y_{k+1},
  \]
  using $m\leq M$, $m^2/M^2\leq\tfrac1{16}$ and
  $(\chi_k\chi_k')^2\leq2^{2k+10}$, so
  $4\cdot\tfrac1{16}\cdot2^{2k+10}=2^{8}4^{k}$.
\end{proof}

\begin{proposition}[No choice of Young parameters produces a contraction]
  \label{prop:s41-nocontraction}
  In \eqref{eq:s41-recursion}, for every sequence $(\delta_k)$ of positive
  numbers and every $k$,
  \[
    \Lambda_k \;\geq\; 2^{8}\,4^{k} \;\geq\; 2^{8}\;>\;1,
    \qquad \Lambda_k\to\infty .
  \]
  In particular:
  \begin{enumerate}[label=(\roman*)]
    \item the choice proposed in Remark~\ref{rem:level-iteration},
      $\delta_k=\delta_0 8^{-k}$, gives
      $\Lambda_k=2^{8}4^{k}+2^{8}\delta_0^{-1}32^{k}$, which grows
      \emph{faster} than the unmodified $4^{k}$: the Young device makes the
      recursion strictly worse, not better;
    \item no other choice helps, because the term $2^{8}4^{k}$ --- coming
      from family (S1), which is not a product of a dissipative factor with
      anything else --- is present irrespective of $\delta_k$: Young's
      inequality has no purchase on it;
    \item if one additionally demands, as any absorption of the extracted
      dissipation into $M^2D_\flat$ requires, that $\sum_k\delta_k\leq\tfrac12$,
      then $\delta_k\leq\tfrac12$ and hence
      $2^{8}\delta_k^{-1}4^{k}\geq2^{9}4^{k}$: the two geometric factors act
      on the \emph{same} side of the inequality and reinforce each other.
      The phrase ``geometric beats geometric'' has no realisation here.
  \end{enumerate}
\end{proposition}
\begin{proof}
  Immediate from \eqref{eq:s41-recursion}: $\Lambda_k$ is a sum of two
  nonnegative terms, the first of which is $2^{8}4^{k}$ and does not involve
  $\delta_k$. For (i), $8^{-k}$ inverted is $8^{k}$ and
  $8^{k}4^{k}=32^{k}$. For (iii), $\sum_k\delta_k\leq\tfrac12$ forces
  $\delta_k\leq\tfrac12$, i.e.\ $\delta_k^{-1}\geq2$.
\end{proof}

\begin{proposition}[The unrolled iteration is weaker than trivial]
  \label{prop:s41-vacuous}
  With $G_k$ and $\Lambda_k$ as in \eqref{eq:s41-recursion}, unrolling
  \eqref{eq:s41-recursion} $N$ times gives
  \[
    Y_0 \;\leq\; \sum_{j=0}^{N-1}G_j\prod_{i<j}\Lambda_i
      \;+\;\Bigl(\prod_{i<N}\Lambda_i\Bigr)Y_N,
    \qquad
    \prod_{i<N}\Lambda_i \;\geq\; 2^{8N}\,4^{N(N-1)/2}.
  \]
  Since $Y_N\geq Y_0$ by \eqref{eq:s41-monotone} and
  $\prod_{i<N}\Lambda_i\geq1$, the remainder term alone already exceeds
  $Y_0$: the unrolled inequality is implied by the trivial statement
  $Y_0\leq Y_0$ and carries no information, for every $N$. Letting
  $N\to\infty$ does not help: both the sum and the remainder coefficient
  diverge, and $Y_N\uparrow Y_\infty$, which is not known to vanish (indeed
  $Y_\infty\geq Y_0$).
\end{proposition}
\begin{proof}
  Induction on $N$ using \eqref{eq:s41-recursion} and
  $\Lambda_i\geq2^{8}4^{i}$, so
  $\prod_{i<N}\Lambda_i\geq2^{8N}4^{0+1+\cdots+(N-1)}$. The monotonicity
  $Y_N\geq Y_0$ is \eqref{eq:s41-monotone}.
\end{proof}

\begin{remark}[Which iteration lemma would be needed, and why none applies]
  \label{rem:s41-whichlemma}
  The task of Remark~\ref{rem:level-iteration} is to decide whether the
  scheme is a linear telescoping (Ladyzhenskaya--Ural'tseva Ch.~II,
  Giaquinta--Giusti \cite{Giaquinta1982,Giusti2003}) or needs a nonlinear
  iteration lemma (Stampacchia--De~Giorgi type). Neither is available.
  \begin{enumerate}[label=(\roman*)]
    \item \textbf{Linear (hole-filling) lemma.} If $f\geq0$ is bounded on
      $[\rho_0,\rho_1]$ and $f(s)\leq\vartheta f(t)+A(t-s)^{-\gamma}+B$ for
      all $\rho_0\leq s<t\leq\rho_1$ with $\vartheta<1$, then
      $f(s)\leq C(\gamma,\vartheta)\bigl(A(t-s)^{-\gamma}+B\bigr)$. The
      hypothesis $\vartheta<1$ is what makes the geometric selection of
      radii work. Here $\vartheta$ is $\Lambda_k\geq2^{8}$
      (Proposition~\ref{prop:s41-nocontraction}), so the lemma does not
      apply; and by Proposition~\ref{prop:s41-nocontraction}(ii) no
      rearrangement of the estimate reduces $\Lambda_k$ below $1$.
    \item \textbf{Nonlinear (Stampacchia) lemma.} If
      $Y_{k+1}\leq C\,b^{k}Y_k^{1+\beta}$ with $b>1$, $\beta>0$ and
      $Y_0\leq C^{-1/\beta}b^{-1/\beta^2}$, then $Y_k\to0$. Three of its
      hypotheses fail simultaneously here: (a) the direction --- our
      inequality bounds the \emph{coarse} level by the \emph{fine} one
      ($Y_k$ by $Y_{k+1}$), because the truncation levels
      $\theta_k$ \emph{decrease}, so the level sets $\{m>\theta_{k+1}M\}$
      \emph{increase} with $k$, exactly the opposite of De~Giorgi
      truncation, where superlevel sets shrink to measure zero; (b)
      homogeneity --- \eqref{eq:s41-recursion} is homogeneous of degree one
      in $Y$, with no superlinear gain to beat $b^{k}$; (c) smallness ---
      $Y_0$ small is the conclusion sought, not an available hypothesis.
    \item \textbf{The structural reason.} De~Giorgi truncation gains because
      the sets on which the iterated quantity lives shrink. Here the
      truncation is in an \emph{auxiliary} variable, $|\omega|$, which on
      the region of interest is confined to the fixed band
      $[M/16,M]$; no set shrinks, no measure decays, and there is
      nothing for a nonlinear gain to feed on.
  \end{enumerate}
\end{remark}

\subsection{The $k\to\infty$ limit is a level-set flux; the co-area
obstruction}
\label{subsec:s41-coarea}

The failure above is not an artefact of the geometric choice
$\theta_k=\tfrac18+2^{-k-3}$. The following identifies the obstruction
exactly and shows it is common to \emph{all} level selections.

\begin{definition}[Level measure]\label{def:s41-levelmeasure}
  For $F\geq0$ measurable with $\iint_{Q(\rho')}F<\infty$, let $\mu_F$ be
  the push-forward of $F\,dx\,dt$ under
  $(x,t)\mapsto m(x,t)/M(t)$, so that
  $\iint F\,g(m/M)\,dx\,dt=\int_0^\infty g\,d\mu_F$ for every bounded Borel
  $g$. Write $\mu_F=\varrho_F\,ds+\mu_F^{\mathrm{sing}}$ for its Lebesgue
  decomposition.
\end{definition}

\begin{proposition}[Co-area obstruction: no level selection gains]
  \label{prop:s41-coarea}
  Let $\tfrac1{16}\leq\theta_-<\theta_+\leq\tfrac14$ and
  $F\geq0$ as above. Then:
  \begin{enumerate}[label=(\alph*)]
    \item for every $\chi\in\mathcal{C}(\theta_-,\theta_+)$,
      \[
        \iint F\,\chi'(m/M)\,dx\,dt \;=\; \int_{\theta_-}^{\theta_+}\chi'\,d\mu_F
        \;\geq\; \operatorname*{ess\,inf}_{s\in(\theta_-,\theta_+)}\varrho_F(s),
      \]
      because $\chi'\geq0$ and $\int_{\theta_-}^{\theta_+}\chi'\,ds=1$;
    \item conversely
      $\displaystyle\inf_{\chi\in\mathcal{C}(\theta_-,\theta_+)}
       \iint F\chi'(m/M)
       \;\leq\;\frac{\mu_F\bigl([\theta_-,\theta_+]\bigr)}{\theta_+-\theta_-}$,
      attained up to $\epsilon$ by $\chi'\equiv(\theta_+-\theta_-)^{-1}$
      mollified;
    \item hence the residual produced by \emph{any} admissible cutoff is a
      \emph{level density} of $\mu_F$, and the only a priori information
      available about it --- namely the total mass
      $\mu_F([\theta_-,\theta_+])\leq\iint_{Q(\rho')}F$ --- bounds it from
      above only by $(\theta_+-\theta_-)^{-1}\leq16$ times that total mass,
      and bounds it from below not at all;
    \item for the nested family of Definition~\ref{def:s41-nested},
      $\chi_k'\,ds\rightharpoonup\delta_{1/8}$ weakly-$*$, so if $\varrho_F$
      is continuous at $\tfrac18$ and $\mu_F^{\mathrm{sing}}$ vanishes near
      $\tfrac18$ then
      \[
        \iint F\,\chi_k'(m/M)\;\longrightarrow\;\varrho_F(\tfrac18)
        \;=\;\int \!M(t)\!\!\int_{\{m=M(t)/8\}}\frac{F}{|\nabla m|}
          \,d\mathcal{H}^2\,dt
      \]
      (the last equality by the co-area formula, valid for a.e.\ $t$ at
      which $M(t)/8$ is a regular value of $m(\cdot,t)$). The limit is a
      \emph{flux through the level set} $\{|\omega|=M/8\}$; it is zero only
      if $F$ vanishes there.
  \end{enumerate}
\end{proposition}
\begin{proof}
  (a) is the definition of $\mu_F$ together with
  $\int\chi'\,d\mu_F\geq\int\chi'\varrho_F\,ds
  \geq\operatorname{ess\,inf}\varrho_F\cdot\int\chi'\,ds$ and
  $\int_{\theta_-}^{\theta_+}\chi'=\chi(\theta_+)-\chi(\theta_-)=1$.
  (b) is the same identity for the stated $\chi$.
  (c) restates (a)--(b).
  (d): $\chi_k'\geq0$, $\int\chi_k'=1$ and
  $\operatorname{supp}\chi_k'\subseteq[\theta_{k+1},\theta_k]\downarrow\{\tfrac18\}$,
  which is weak-$*$ convergence to $\delta_{1/8}$; the co-area identity is
  the substitution $s=M(t)\sigma$ in
  $\iint F g(m/M)=\int\!dt\int_0^\infty g(s/M)
  \bigl(\int_{\{m=s\}}F/|\nabla m|\,d\mathcal{H}^2\bigr)ds$.
\end{proof}

\begin{remark}[Reading of Proposition~\ref{prop:s41-coarea}]
  \label{rem:s41-reading}
  Applied to $F=|\nabla m|^2w$ --- the integrand of family (S1) --- part (d)
  says that the level iteration does not annihilate the residual; it
  converts it into the trace
  $\int M\!\int_{\{m=M/8\}}|\nabla m|\,w\,d\mathcal{H}^2dt$. Controlling
  that trace requires a bound for $w$ \emph{on a level set of $|\omega|$},
  weighted by $|\nabla m|$: precisely the pointwise information that H1 is
  trying to produce. Part (c) says the same thing statically: the annulus
  residual is, for every admissible cutoff, a level density of $\mu_F$,
  and the only a priori bound is $16\times$ the band mass
  $\nu\iint_{\mathrm{Ann}}|\nabla m|^2w\,\psi^2$ --- which is the very
  quantity $Y$ restricted to the band where $\eta^2$ degenerates. Choosing a
  ``good'' level by pigeonhole gains nothing either: with $N$ equally spaced
  levels the derivative bound is $|\chi'|\leq16N$ while the best of $N$
  disjoint annuli carries at most $1/N$ of the band mass, and the two
  factors cancel \emph{exactly}. This is the sharp form of the obstruction:
  the co-area normalisation $\int|\chi'|\,ds=1$ is scale-free, so no
  redistribution of levels can beat it.
\end{remark}

\begin{remark}[Why the same terms are harmless at the quadratic level]
  \label{rem:s41-contrast}
  It is worth recording precisely why Lemma~\ref{lem:mag-rhs}(c) could
  dispose of the annulus terms of the \emph{quadratic} magnitude Caccioppoli
  \eqref{eq:mag-caccioppoli} ``with no absorption argument needed'', while
  the same device fails here. There, the annulus integrands were
  $|\nabla m|^2$ and $m^2w$, both dominated pointwise by $|\nabla\omega|^2$
  and hence by the unconditional budget \eqref{eq:s41-budget}. Here the
  integrands are $|\nabla m|^2w$ and $m^2w^2$: each is a budgeted quadratic
  quantity \emph{times $w$}, and $w$ is exactly the unknown of H1. The
  annulus residual is thus not a defect of the cutoff, but the quartic
  nature of the $K$-identity making itself felt at the boundary of the
  level set.
\end{remark}

\subsection{What is nevertheless proved, and the exact closure target}
\label{subsec:s41-target}

\begin{proposition}[Unconditional crude bound for the core residual]
  \label{prop:s41-crude}
  With $\mathcal{A}^{\mathrm{core}}_{\mathrm{ann}}$ as in
  \eqref{eq:s41-core} formed with any
  $\chi\in\mathcal{C}(\theta_-,\theta_+)$,
  $\tfrac1{16}\leq\theta_-<\theta_+\leq\tfrac14$, and any $\delta>0$,
  \[
    \mathcal{A}^{\mathrm{core}}_{\mathrm{ann}}
    \;\leq\; \delta\,M^2D_\flat
    \;+\; C\bigl(1+\delta^{-1}\bigr)\,
      \|w\|_{L^\infty(Q(\rho')\cap\{m\geq M/16\})}\;\mathcal{E}
    \;\leq\; C\bigl(1+\delta^{-1}\bigr)\,\mathcal{L}\,\|w\|_{L^\infty}\,M^{1/2}
      \;+\;\delta M^2D_\flat,
  \]
  with $C$ depending only on $\|\Theta'\|_\infty,\|\Theta''\|_\infty$ and
  $\theta_+-\theta_-$. Equivalently, in the $K$-units of
  Proposition~\ref{prop:crude-K} (divide by $M^2$ and use
  $M^{-3/2}=(r^*)^3$):
  \[
    M^{-2}\,\mathcal{A}^{\mathrm{core}}_{\mathrm{ann}}
    \;\leq\; C\bigl(1+\delta^{-1}\bigr)\,\mathcal{L}\,\|w\|_{L^\infty}\,(r^*)^3
    \;+\;\delta\,D_\flat .
  \]
\end{proposition}
\begin{proof}
  For (S1): $|\nabla m|^2\leq|\nabla\omega|^2$
  (Lemma~\ref{lem:grad-omega-split}), $\psi\leq1$, $w\leq\|w\|_\infty$ on the
  support, and $\|\chi\chi''\|_\infty\leq C$; apply \eqref{eq:s41-budget}.
  For (S2): Young as in the proof of
  Proposition~\ref{prop:s41-recursion} with parameter $\delta$, then
  $|\nabla m|^2w\leq\|w\|_\infty|\nabla\omega|^2$ in the second factor and
  \eqref{eq:s41-budget} again.
\end{proof}

\begin{remark}[The shortfall is exactly one logarithm]\label{rem:s41-shortfall}
  Proposition~\ref{prop:s41-crude} reproduces, term for term, the bound of
  Proposition~\ref{prop:crude-K}. The quantity that
  Proposition~\ref{prop:second-rung} is able to absorb is
  $\delta\,\nu\,(r^*)^3\|w\|_{L^\infty}$; what
  Proposition~\ref{prop:s41-crude} delivers is
  $C\mathcal{L}\|w\|_{L^\infty}(r^*)^3$. Absorption therefore requires
  $C\mathcal{L}\leq\delta\nu$, which fails for fixed $\nu$ as
  $\mathcal{L}=\log(e+\|u\|_{H^3}/M)\to\infty$. So the annulus residual is
  short of closure by exactly the same single logarithm that
  Proposition~\ref{prop:crude-K} identified for $K$ itself --- and, by
  Remark~\ref{rem:s41-reading}, no level scheme recovers it.
\end{remark}

\begin{remark}[The exact end target, for the record]\label{rem:s41-endtarget}
  For the ``modulo $\mathcal{A}_{\mathrm{ann}}$'' caveat of
  Proposition~\ref{prop:K-amended} to be removable, what must be proved is:
  \emph{there exists one admissible $\chi\in\mathcal{C}(\theta_-,\theta_+)$
  with $\tfrac1{16}\leq\theta_-<\theta_+\leq\tfrac14$ and constants
  $C$ such that, with $\eta=\eta_\chi$,}
  \begin{equation}\label{eq:s41-endtarget}
    \mathcal{A}_{\mathrm{ann}}
    \;\leq\; \tfrac12\,\nu\iint|\nabla m|^2\,w\,\eta^2
      \;+\; C\,M^2\,\nu\iint|\nabla^2\hat{e}|^2\eta^2
      \;+\; C\,\mathcal{L}^2M^{3/2}
        \bigl(\nu^{-1}+\nu^{-1/2}(\sigma^*)^{1/2}\bigr).
  \end{equation}
  Two points must be stressed. First, the absorbed first term must carry
  \emph{the same} $\eta$ as the left-hand side of
  Proposition~\ref{prop:K-amended}: an annulus bound in terms of a
  \emph{fuller} weight (e.g.\ $\psi^2$ on the band, or $\eta_{k+1}^2$) does
  not absorb, and this --- not the value of any constant --- is what defeats
  the nested scheme. Second, by Lemma~\ref{lem:s41-admissible} it suffices
  to prove \eqref{eq:s41-endtarget} for a \emph{single} admissible level
  pair, since all such pairs are interchangeable downstream; so the
  difficulty is genuinely analytic, not a matter of choosing levels
  cleverly.
\end{remark}

\begin{conjecture}[Band absorption --- the replacement for
  Remark~\ref{rem:level-iteration}]\label{target:band-absorption}
  There exist $\theta_-<\theta_+$ in $[\tfrac1{16},\tfrac14]$, a cutoff
  $\chi\in\mathcal{C}(\theta_-,\theta_+)$ and a constant $C$ such that, with
  $\eta=\eta_\chi$ and
  $\mathrm{Band}=\{\theta_-M\leq m\leq\theta_+M\}\cap\operatorname{supp}\psi$,
  \begin{align}
    \nu\iint_{\mathrm{Band}}|\nabla m|^2\,w\,\psi^2
    &\;\leq\; 2^{-11}\,\nu\iint|\nabla m|^2\,w\,\eta^2
      \;+\; C\,M^2\,\nu\iint|\nabla^2\hat{e}|^2\eta^2
      \;+\; C\,\mathcal{L}^2M^{3/2}\nu^{-1},
      \label{eq:s41-bandY}\\
    M^2\,\nu\iint_{\mathrm{Band}}|\nabla^2\hat{e}|^2\,\psi^2
    &\;\leq\; C\,M^2\,\nu\iint|\nabla^2\hat{e}|^2\eta^2
      \;+\; C\,\mathcal{L}^2M^{3/2}\nu^{-1}.
      \label{eq:s41-bandD}
  \end{align}
\end{conjecture}

\begin{remark}[Status of Conjecture~\ref{target:band-absorption}]
  \label{rem:s41-band-status}
  Conjecture~\ref{target:band-absorption} implies
  \eqref{eq:s41-endtarget}. Indeed, for a single level pair of width
  $\theta_+-\theta_-\geq\tfrac1{16}$ one has
  $\|\chi\chi'\|_\infty\leq32$ and $\|\chi\chi''\|_\infty\leq2^{12}$, so
  Proposition~\ref{prop:s41-core} and the Young step of
  Proposition~\ref{prop:s41-recursion} with $\delta=1$ give
  $|\text{\upshape(S1)}|+|\text{\upshape(S2)}|
  \leq 2^{9}\,\nu\iint_{\mathrm{Band}}|\nabla m|^2w\psi^2
  + M^2\nu\iint_{\mathrm{Band}}|\nabla^2\hat{e}|^2\psi^2$; inserting
  \eqref{eq:s41-bandY}--\eqref{eq:s41-bandD} and
  $2^{9}\cdot2^{-11}=\tfrac14<\tfrac12$ yields
  \eqref{eq:s41-endtarget}. We do not claim the converse. The conjecture
  asserts that the
  radial--angular coupling density on the transition band is not
  concentrated there relative to its value on $\{m\geq\theta_+M\}$. This is
  a statement of exactly the same analytic type as
  Conjecture~\ref{conj:K-refined} (a decorrelation/equipartition
  hypothesis), restricted to the band, and it is \emph{open}. In particular
  the claim of Remark~\ref{rem:if-verified} --- that verification of C1--C2
  makes Conjecture~\ref{conj:K-refined} unnecessary --- must be qualified:
  it holds on $\{m\geq M/4\}$, where the level cutoff is inactive, and not
  on the transition band, where a hypothesis of that type is still required.

  \medskip
  \noindent\textbf{Amended (S44).} The phrase ``exactly the same analytic
  type'' above is accurate as a statement about the \emph{form} of the two
  hypotheses and must not be read as an implication in either direction:
  \S\ref{sec:relation-s44} shows that neither conjecture implies the other,
  nor any nontrivial part of the other, from the a priori structure they are
  formulated against (Propositions~\ref{prop:s44-a-fails}
  and~\ref{prop:s44-b-fails}; the exact scope of that claim is
  Remark~\ref{rem:s44-scope}). The reason is that the two are supported on
  disjoint regions of the $|\omega|$-level decomposition
  (Remark~\ref{rem:s44-disjoint}). Three further amendments. (i) The
  description ``a decorrelation/equipartition hypothesis'' applies to
  \eqref{eq:s41-bandY} but \emph{not} to \eqref{eq:s41-bandD}, which is a bare
  level non-concentration statement of a different type
  (Remark~\ref{rem:s44-bandD}). (ii) The absolute constants of
  \eqref{eq:s41-bandY}--\eqref{eq:s41-bandD} do \emph{not} make this
  conjecture weaker than Conjecture~\ref{conj:K-refined}: reduced to a
  hypothesis on the band correlation constant by the route on which the
  latter was calibrated, \eqref{eq:s41-bandY} requires
  $c_K^{\mathrm{band}}\leq C/\mathcal{L}$, the same decaying form
  (Proposition~\ref{prop:s44-linfty-route},
  Remark~\ref{rem:s44-strength}). (iii) A single hypothesis covers the
  $Y$-halves of both targets, namely Conjecture~\ref{conj:lc-s44} --- the
  $c_K$ hypothesis itself, asserted on $\{m\geq M/16\}$ rather than on
  $\{m\geq M/4\}$ (Proposition~\ref{prop:s44-lc-implies}); the
  volume-normalised alternative Conjecture~\ref{conj:jle-s44} also suffices
  but is strictly stronger in the concentrated regime
  (Remark~\ref{rem:s44-jle-status}).
\end{remark}

\subsection{Amended status of $\mathcal{A}_{\mathrm{ann}}$}
\label{subsec:s41-status}

\begin{center}
\begin{tabular}{lll}
\hline
Sub-item of $\mathcal{A}_{\mathrm{ann}}$ & Status & Where \\
\hline
Families (G1)--(G4) & \textbf{Proved} good-signed; discarded & Prop.~\ref{prop:s41-core} \\
Families (I1)--(I2) & \textbf{Proved} in the allowed class $\mathcal{I}$ & Prop.~\ref{prop:s41-core} \\
Nested cutoffs, $4^{k}$ constant & \textbf{Proved} (constant confirmed) & Lem.~\ref{lem:s41-derivs} \\
Level pairs interchangeable & \textbf{Proved} & Lem.~\ref{lem:s41-admissible} \\
Recursion with $\delta_k$ & \textbf{Proved not a contraction} & Prop.~\ref{prop:s41-nocontraction} \\
$\delta_k=\delta_08^{-k}$ as proposed & \textbf{Proved counterproductive} ($32^{k}$) & Prop.~\ref{prop:s41-nocontraction}(i) \\
Unrolled iteration & \textbf{Proved vacuous} & Prop.~\ref{prop:s41-vacuous} \\
$k\to\infty$ limit & \textbf{Proved} to be a level-set flux & Prop.~\ref{prop:s41-coarea}(d) \\
Any other level selection & \textbf{Proved} no gain (co-area) & Prop.~\ref{prop:s41-coarea}(c) \\
Core residual (S1)+(S2), crude bound & \textbf{Proved}; one logarithm short & Prop.~\ref{prop:s41-crude} \\
Core residual, absorption & \textbf{Open} (was ``open-mechanical'') & Conj.~\ref{target:band-absorption} \\
Relation to Conj.~\ref{conj:K-refined} & \textbf{Proved independent} (S44); common hyp.\ isolated & \S\ref{sec:relation-s44}, Conj.~\ref{conj:lc-s44} \\
\hline
\end{tabular}
\end{center}

\begin{remark}[Honest summary of S41]\label{rem:s41-summary}
  The level iteration of Remark~\ref{rem:level-iteration} does not close, and
  the reason is structural rather than computational: the residue of a
  level-set truncation is, by the co-area normalisation
  $\int|\chi'|=1$, a level density of the underlying measure, and no
  redistribution of levels --- geometric, uniform, nested, or selected ---
  makes a level density smaller than a fixed multiple of the band mass.
  What was gained instead is a substantial reduction: six of the eight
  $\chi$-derivative families are removed outright, four by sign and two into
  the already-allowed inhomogeneous class, leaving a single explicit core
  \eqref{eq:s41-core} with a sharp, unconditional, one-logarithm-short
  bound and a sharply stated open target
  (Conjecture~\ref{target:band-absorption}). The correct classification of
  $\mathcal{A}_{\mathrm{ann}}$ is therefore \emph{open}, of the same type as
  Conjecture~\ref{conj:K-refined}, and \S\ref{sec:verification-s36} is
  amended accordingly. No other item of the program is affected by this
  section, and in particular nothing here bears on H2, on
  $\sigma^*$-decay, or on Type-II exclusion.
\end{remark}

% ══════════════════════════════════════════════════════════════════════════════
\section{The Relationship Between the $c_K$ and Band-Absorption Hypotheses
(Session S44)}
\label{sec:relation-s44}

Remark~\ref{rem:s41-summary} classifies $\mathcal{A}_{\mathrm{ann}}$ as
``open, of the same analytic type as Conjecture~\ref{conj:K-refined}'', and
Remark~\ref{rem:s41-band-status} repeats the phrase; both note that the crude
bounds fail by the same single logarithm
(Proposition~\ref{prop:crude-K}, Proposition~\ref{prop:s41-crude},
Remark~\ref{rem:s41-shortfall}). This section asks whether the resemblance is
substantive, i.e.\ what the \emph{logical} relationship between
Conjecture~\ref{conj:K-refined} and Conjecture~\ref{target:band-absorption}
actually is. The outcome is mixed, and is reported as such.

\begin{enumerate}[label=(\alph*)]
  \item Conjecture~\ref{conj:K-refined} does \emph{not} imply
    Conjecture~\ref{target:band-absorption}
    (Proposition~\ref{prop:s44-a-fails}); the reason is structural, not a
    matter of constants: the two hypotheses are supported on \emph{disjoint}
    regions of the $|\omega|$-level decomposition
    (Remark~\ref{rem:s44-disjoint}). The band-transported form of the $c_K$
    hypothesis does imply the $Y$-half \eqref{eq:s41-bandY}: through the
    $\|w\|_{L^\infty}$ route with no extra ingredient
    (Proposition~\ref{prop:s44-linfty-route}), or through the budget route
    with the volume hypothesis (NB) named explicitly
    (Proposition~\ref{prop:s44-transport}).
  \item Conjecture~\ref{target:band-absorption} does not imply
    Conjecture~\ref{conj:K-refined}, nor any upper bound on $c_K$ whatsoever
    (Proposition~\ref{prop:s44-b-fails}).
  \item There are common hypotheses from which the $Y$-halves of both targets
    follow, one for each of the two available routes
    (Remark~\ref{rem:s44-two-routes}). The weaker and preferred one is
    Conjecture~\ref{conj:lc-s44}: \emph{the $c_K$ hypothesis itself, asserted
    on $\{m\geq M/16\}$ instead of on $\{m\geq M/4\}$}
    (Proposition~\ref{prop:s44-lc-implies}). The alternative,
    Conjecture~\ref{conj:jle-s44}, replaces the decaying constant by an $O(1)$
    one at the cost of normalising by the ambient parabolic volume, and is a
    demanding hypothesis rather than a cheap one
    (Remark~\ref{rem:s44-jle-status}). The angular half \eqref{eq:s41-bandD}
    of Conjecture~\ref{target:band-absorption} is of a different type and is
    covered by neither (Remark~\ref{rem:s44-bandD}).
  \item The strength comparison suggested by the ``$C/\mathcal{L}$ versus
    $O(1)$'' discrepancy does \emph{not} survive: \eqref{eq:s41-bandY} is
    $O(1)$ in its bookkeeping only. Reduced to a hypothesis on a correlation
    constant by the route on which Conjecture~\ref{conj:K-refined} was
    calibrated, it requires $c_K^{\mathrm{band}}\leq C/\mathcal{L}$ --- the
    same decaying form (Proposition~\ref{prop:s44-linfty-route}). The
    volume-normalised alternative buys an $O(1)$ constant only at the price of
    a volume hypothesis, symmetrically on both sides
    (\eqref{eq:s44-NB} and \eqref{eq:s44-NBE}), and is \emph{strictly
    stronger} than the corresponding correlation bound
    (Lemma~\ref{lem:s44-xi-vs-cK}). See
    Proposition~\ref{prop:s44-calibration}, Remark~\ref{rem:s44-strength} and
    Remark~\ref{rem:s44-numerics}. The first draft of this section concluded
    the opposite from an inverted lemma; that conclusion is withdrawn and the
    correction is recorded in Remark~\ref{rem:s44-corrigendum}.
\end{enumerate}

\noindent This section changes no proved statement. It rearranges open
hypotheses; its only edits elsewhere are a cross-reference appended to
Remark~\ref{rem:s41-band-status} and one row added to the table of
\S\ref{subsec:s41-status}. Nothing here bears on H2, on $\sigma^*$-decay, on
the hypotheses of \S\ref{sec:rigidity-s39}, on Type-II exclusion, or on the
withdrawn \S\ref{sec:audit-s31} material.

\subsection{Normalisation; the two hypotheses side by side}
\label{subsec:s44-setup}

Throughout, the notation is that of \S\ref{sec:annulus-s41}: $m=|\omega|$,
$w=|\nabla\hat{e}|^2$, $M=M(t)$, $r^*=M^{-1/2}$, $\psi$ the fixed space--time
cutoff of Definition~\ref{def:composite-cutoff} with
$\tfrac{r^*}{2}\leq\rho<\rho'\leq r^*$, $\mathcal{E}$ the enstrophy budget
\eqref{eq:s41-budget}, and $\mathcal{I}$ the allowed inhomogeneous class
\eqref{eq:s41-I}. We abbreviate its dominant constituent
\[
  \mathcal{I}_0 \;:=\; \mathcal{L}^2M^{3/2}\nu^{-1}
  \qquad(\text{so }\mathcal{I}\geq C^{-1}\mathcal{I}_0
  \text{ and }\mathcal{I}_0\text{ is exactly the last term of }
  \eqref{eq:s41-bandY}\text{--}\eqref{eq:s41-bandD}).
\]
Constants $C$ depend only on the standing Type-I data $(C_I,\nu)$ and on
$\|\Theta'\|_\infty,\|\Theta''\|_\infty$, never on $M$ or $\mathcal{L}$.
We write $|Q(\rho)|=c_5\rho^5$ for the parabolic volume ($c_5$ absolute) and
\[
  Q^\sharp \;:=\; Q(\rho')\cap\{m\geq M/16\},
  \qquad
  E \;=\; Q(r^*)\cap\{m\geq M/4\},
  \qquad
  \mathrm{Band} \;=\; \{\theta_-M\leq m\leq\theta_+M\}\cap\operatorname{supp}\psi ,
\]
with $\tfrac1{16}\leq\theta_-<\theta_+\leq\tfrac14$ admissible in the sense of
Definition~\ref{def:s41-admissible}.

\begin{definition}[Volume-normalised coupling excess]\label{def:s44-excess}
  For a measurable weight $\varpi$ with $0\leq\varpi\leq1$ supported in
  $Q(r^*)$, and provided both factors in the denominator are positive and
  finite, set
  \[
    \Xi[\varpi]
    \;:=\;
    \frac{|Q(r^*)|\;\displaystyle\iint|\nabla m|^2\,w\,\varpi}
         {\Bigl(\displaystyle\iint|\nabla m|^2\,\varpi\Bigr)
          \Bigl(\displaystyle\iint w\,\varpi\Bigr)} .
  \]
  Thus $\Xi[\varpi]\leq\kappa_0$ says: \emph{the coupling density
  $|\nabla m|^2w$ is at most $\kappa_0$ times what it would be if
  $|\nabla m|^2$ and $w$ were uncorrelated and spread over the whole cylinder
  $Q(r^*)$.} It is the same ratio as $c_K$ except that the normalising volume
  is the ambient $|Q(r^*)|$ and not the measure of the region carrying
  $\varpi$. If either denominator vanishes then $|\nabla m|^2w\,\varpi=0$
  a.e., so $\Xi[\varpi]$ is left undefined and every statement below in which
  it occurs holds trivially; we do not repeat this convention.
\end{definition}

\begin{lemma}[$\Xi$ versus $c_K$: the volume ratio]\label{lem:s44-xi-vs-cK}
  With $c_K$ as in Conjecture~\ref{conj:K-refined},
  \begin{equation}\label{eq:s44-xi-ck}
    \Xi[\mathbf{1}_E] \;=\; c_K\,\frac{|Q(r^*)|}{|E|}\;\geq\;c_K ,
  \end{equation}
  and more generally, for any measurable $R\subseteq Q(r^*)$ of positive
  measure and any weight $\varpi$ supported in $R$,
  \[
    \Xi[\varpi] \;=\; c_K[\varpi]\,\frac{|Q(r^*)|}{|R|},
    \qquad
    c_K[\varpi]\;:=\;
    \frac{|R|\iint|\nabla m|^2w\,\varpi}
      {\bigl(\iint|\nabla m|^2\varpi\bigr)\bigl(\iint w\,\varpi\bigr)} ,
  \]
  $c_K[\mathbf{1}_E]$ being exactly $c_K$ and
  $c_K[\mathbf{1}_{\mathrm{Band}}\psi^2]$ exactly the $c_K^{\mathrm{band}}$ of
  \eqref{eq:s44-cKband}. Consequently:
  \begin{enumerate}[label=(\roman*)]
    \item a bound $\Xi[\varpi]\leq\kappa$ is \emph{strictly stronger} than the
      correlation bound $c_K[\varpi]\leq\kappa$, by exactly the volume ratio
      $|Q(r^*)|/|R|$;
    \item Conjecture~\ref{conj:K-refined} gives only
      $\Xi[\mathbf{1}_E]\leq (C/\mathcal{L})\,|Q(r^*)|/|E|$, which is $O(1)$
      \emph{only} under a volume lower bound on $E$; and
    \item $\Xi[\varpi]\geq1$ whenever $|\nabla m|^2$ and $w$ are supported in a
      common set of measure $\leq|Q(r^*)|$ and are uncorrelated on it --- for
      $|\nabla m|^2\equiv G$, $w\equiv W$ on $R$ and both vanishing off $R$ one
      has $c_K[\varpi]=1$ but $\Xi[\varpi]=|Q(r^*)|/|R|$.
  \end{enumerate}
\end{lemma}
\begin{proof}
  Writing out the two averages in Conjecture~\ref{conj:K-refined} gives
  $c_K = |E|\iint_E|\nabla m|^2w
  \big/\bigl(\iint_E|\nabla m|^2\cdot\iint_E w\bigr)$, and
  Definition~\ref{def:s44-excess} with $\varpi=\mathbf{1}_E$ gives the same
  expression with $|E|$ replaced by $|Q(r^*)|$; dividing yields
  \eqref{eq:s44-xi-ck}. The general case is identical, and (i)--(iii) follow
  from $|R|\leq|Q(r^*)|$ and direct substitution.
\end{proof}

\begin{remark}[Corrigendum: this lemma was stated inverted in the first draft
  of S44]\label{rem:s44-corrigendum}
  In the first draft of this section Lemma~\ref{lem:s44-xi-vs-cK} was recorded
  as $\Xi[\mathbf{1}_E]=c_K|E|/|Q(r^*)|\leq c_K$ --- the reciprocal of the
  correct volume factor, hence the wrong direction. The error was caught in
  the independent check of this section and is recorded here, in the style of
  the corrigendum of \S\ref{sec:magnitude-s34} and the retraction of
  \S\ref{sec:secondrung-s35}, because it \emph{reverses} the interpretation of
  the volume-normalised statistic $\Xi$ and therefore the answer to the
  strength question of \S\ref{subsec:s44-strength}. What the first draft
  concluded --- that $\Xi$-boundedness is weaker than $c_K$-boundedness, that
  Conjecture~\ref{conj:K-refined} may be weakened to an $O(1)$ statement
  without loss, and that the measured $c_K\approx1$ is direct evidence for
  $\Xi=O(1)$ --- is \textbf{withdrawn}; the corrected statements are
  Proposition~\ref{prop:s44-calibration}(iii),
  Remark~\ref{rem:s44-strength} and Remark~\ref{rem:s44-numerics} below.
  Nothing else in this section used the lemma: the two non-implications
  (Propositions~\ref{prop:s44-a-fails} and~\ref{prop:s44-b-fails}),
  Lemma~\ref{lem:s44-budget}, Proposition~\ref{prop:s44-transport},
  Proposition~\ref{prop:s44-calibration}(i)--(ii) and
  Proposition~\ref{prop:s44-jle-implies} are all independent of it and were
  re-checked line by line. Two internal consistency checks confirm the
  corrected direction: Remark~\ref{rem:s44-jle-status} computes
  $\Xi=|Q(r^*)|/|S|$ under co-concentration on a set $S$, and
  Proposition~\ref{prop:s44-transport} uses
  $\Xi[\mathbf{1}_{\mathrm{Band}}\psi^2]
  =c_K^{\mathrm{band}}|Q(r^*)|/|\mathrm{Band}|$; both agree with
  \eqref{eq:s44-xi-ck} and contradicted the first-draft lemma.
\end{remark}

\begin{remark}[The two hypotheses live on disjoint regions]
  \label{rem:s44-disjoint}
  Conjecture~\ref{conj:K-refined} is a statement about
  $E=Q(r^*)\cap\{m\geq M/4\}$. Conjecture~\ref{target:band-absorption} is a
  statement about $\mathrm{Band}\subseteq\{m\leq M/4\}$. The two regions meet
  only in the level set $\{m=M/4\}$, which carries no five-dimensional measure
  whenever $M(t)/4$ is a regular value of $m(\cdot,t)$ for a.e.\ $t$. Neither
  hypothesis constrains any integral supported in the other's region. This
  already makes it implausible that either implies the other, and
  \S\ref{subsec:s44-nonimplications} confirms it. It also locates precisely
  what Remark~\ref{rem:s41-band-status} was right to flag: ``the same analytic
  type'' is a statement about the \emph{form} of the two hypotheses, not about
  their content, and the qualification that remark attaches to
  Remark~\ref{rem:if-verified} is exactly the disjointness recorded here.
\end{remark}

\begin{lemma}[Budget lemma]\label{lem:s44-budget}
  Let $0\leq\varpi\leq1$ be supported in $Q^\sharp$. Then
  \[
    \nu\iint|\nabla m|^2\,\varpi \;\leq\;\mathcal{E},
    \qquad
    \iint w\,\varpi \;\leq\;\frac{256\,\mathcal{E}}{\nu M^2},
  \]
  and consequently, whenever $\Xi[\varpi]$ is defined,
  \begin{equation}\label{eq:s44-master}
    \nu\iint|\nabla m|^2\,w\,\varpi
    \;\leq\;
    \Xi[\varpi]\;\frac{256\,\mathcal{E}^2}{\nu\,M^2\,|Q(r^*)|}
    \;\leq\;
    C\,\Xi[\varpi]\;\mathcal{I}_0 .
  \end{equation}
\end{lemma}
\begin{proof}
  The first bound is $|\nabla m|^2\leq|\nabla m|^2+m^2w$ together with
  \eqref{eq:s41-budget} and $\varpi\leq1$. The second is the same estimate
  after using $m\geq M/16$ on $\operatorname{supp}\varpi$, so that
  $m^2w\geq M^2w/256$. Multiplying the two bounds, dividing by $|Q(r^*)|$ and
  using Definition~\ref{def:s44-excess} gives the first inequality of
  \eqref{eq:s44-master}. For the second, insert
  $|Q(r^*)|=c_5(r^*)^5=c_5M^{-5/2}$ and
  $\mathcal{E}\leq C_{\mathcal{E}}\mathcal{L}M^{1/2}$:
  \[
    \frac{\mathcal{E}^2}{M^2|Q(r^*)|}
    \;\leq\;\frac{C_{\mathcal{E}}^2\mathcal{L}^2M}{c_5\,M^{2}M^{-5/2}}
    \;=\;\frac{C_{\mathcal{E}}^2}{c_5}\,\mathcal{L}^2M^{3/2} . \qedhere
  \]
\end{proof}

Lemma~\ref{lem:s44-budget} is the whole mechanism of this section, and it is
worth saying at once what it replaces. Proposition~\ref{prop:crude-K} bounds
the coupling by $\|w\|_{L^\infty}$ times the enstrophy budget, thereby
\emph{discarding} the a priori $w$-budget of Theorem~\ref{thm:apriori-sigma};
\eqref{eq:s44-master} uses both budgets and pays for it with a correlation
hypothesis. The gain is not the logarithm per se: it is the elimination of
$\|w\|_{L^\infty}$, which is the unknown that H1 is trying to produce.

\subsection{Test configurations: what a non-implication can mean here}
\label{subsec:s44-config}

Neither conjecture can be refuted or verified by exhibiting a Navier--Stokes
solution; both are asserted only \emph{near a putative Type-I singularity},
whose existence is the open question. What can be settled --- and is settled
below --- is whether either implication follows from the a priori structure
that the two hypotheses are formulated against.

\begin{definition}[Admissible test configuration]\label{def:s44-config}
  Fix $\nu>0$, $M\geq1$, $r^*=M^{-1/2}$, $\mathcal{L}\geq1$, and $\psi$ as in
  Definition~\ref{def:composite-cutoff} with $\rho=\tfrac34r^*$, $\rho'=r^*$.
  An \emph{admissible test configuration} is a triple $(m,w,h)$ on $Q(r^*)$
  with $m$ Lipschitz and $w,h\geq0$ measurable, such that
  \begin{enumerate}[label=\textup{(A\arabic*)}]
    \item $0\leq m\leq M$ and $\sup_{Q(r^*)}m=M$;
    \item $\nu\iint_{Q(r^*)}\bigl(|\nabla m|^2+m^2w\bigr)
      \leq C_{\mathcal{E}}\,\mathcal{L}M^{1/2}$, i.e.\ \eqref{eq:s41-budget};
    \item $|\nabla w|\leq 2\,w^{1/2}h^{1/2}$ a.e.\ --- the pointwise link
      $|\nabla w|\leq2w^{1/2}|\nabla^2\hat{e}|$ used throughout
      \S\ref{sec:secondrung-s35}--\S\ref{sec:annulus-s41}, with $h$ standing
      for $|\nabla^2\hat{e}|^2$.
  \end{enumerate}
  In such a configuration $c_K$, $\Xi[\varpi]$,
  \eqref{eq:s41-bandY}--\eqref{eq:s41-bandD} and \eqref{eq:s44-master} are all
  well defined, with $|\nabla^2\hat{e}|^2$ read as $h$.
\end{definition}

\begin{remark}[Scope of the non-implications; stated as a limitation]
  \label{rem:s44-scope}
  An admissible test configuration need not come from a solution of the
  Navier--Stokes equations: $w$ need not equal $|\nabla\hat{e}|^2$ for a unit
  field $\hat{e}$, $h$ need not be its second-derivative density, and no
  evolution equation is imposed. Consequently
  Propositions~\ref{prop:s44-b-fails} and~\ref{prop:s44-a-fails} do
  \emph{not} say that the two conjectures are independent as statements about
  Navier--Stokes solutions; that question is not decidable by any argument
  available here. What they do say, and all they say, is: \emph{no implication
  between Conjecture~\ref{conj:K-refined} and
  Conjecture~\ref{target:band-absorption} can be derived from} (A1)--(A3),
  i.e.\ from the a priori budgets and the pointwise structure in terms of
  which the two statements are phrased. Any implication would have to invoke
  further PDE structure, and none is on offer in
  \S\ref{sec:magnitude-s34}--\S\ref{sec:annulus-s41}. That is the precise
  sense in which the ``same analytic type'' language of
  Remarks~\ref{rem:s41-summary} and~\ref{rem:s41-band-status} must not be read
  as ``one covers the other''.
\end{remark}

\subsection{The two non-implications}
\label{subsec:s44-nonimplications}

\begin{proposition}[Band absorption bounds no correlation]
  \label{prop:s44-b-fails}
  There is a family of admissible test configurations, indexed by
  $\ell\in(0,r^*/4)$, for which the left-hand sides of
  \eqref{eq:s41-bandY} and \eqref{eq:s41-bandD} \emph{vanish identically} for
  every admissible level pair --- so
  Conjecture~\ref{target:band-absorption} holds with $C=0$ --- while
  \[
    c_K \;\geq\;\frac{|E|}{8\,\ell^{5}}\;\longrightarrow\;\infty
    \qquad(\ell\to0),
  \]
  with $\mathcal{L}$ held fixed. Hence
  Conjecture~\ref{target:band-absorption} implies no upper bound on $c_K$ at
  all, a fortiori not Conjecture~\ref{conj:K-refined}.
\end{proposition}
\begin{proof}
  Let $\varphi(s)=\operatorname{dist}(s,2\mathbb{Z})\in[0,1]$ be the unit
  triangular wave, so $|\varphi'|=1$ a.e. Fix $\ell\in(0,r^*/4)$, let
  $K\subseteq B_\rho$ be an axis-aligned cube of side $\ell$ and $I\subseteq
  (t_0-\rho^2,t_0)$ an interval of length $\ell^2$, and put $N=K\times I$,
  $|N|=\ell^5$, and let $N_0=K_0\times I$ with $K_0$ the concentric cube of
  side $\ell/2$. Let $\xi$ be a Lipschitz cutoff, independent of $t$, with
  $\xi\equiv1$ on $K_0$, $\operatorname{supp}\xi\subseteq K$ and
  $|\nabla\xi|\leq8/\ell$, and set
  \[
    m(x,t)\;=\;\frac{M}{2}+\frac{M}{8}\,\xi(x)\,
      \varphi\!\Bigl(\frac{x_1}{\epsilon}\Bigr)
    \quad\text{on }N,
    \qquad\text{so}\qquad
    |\nabla m|^2 \;=\; G \;:=\;\frac{M^2}{64\,\epsilon^2}
    \ \text{ on }N_0,
  \]
  and $|\nabla m|^2\leq 4G$ on $N$ once $\epsilon\leq\ell/8$ (indeed
  $|\nabla m|\leq\tfrac{M}{8}(|\nabla\xi|+\epsilon^{-1})
  \leq\tfrac{M}{4\epsilon}=2G^{1/2}$). Here $\epsilon>0$ is chosen below, and
  $\epsilon\leq\ell/8$ holds once $\ell/r^*$ is small enough in terms of $\nu$,
  since the normalisation below forces
  $\epsilon^2\leq C\nu M^{3/2}\ell^{5}$, i.e.\
  $\epsilon^2/\ell^2\leq C\nu(\ell/r^*)^{3}$.
  Since $m/M\in[\tfrac38,\tfrac58]$ on $N$, we have
  $N\subseteq\{m\geq M/4\}$, so $N\subseteq E$ and $N$ is disjoint from every
  admissible band; and $m\equiv M/2$ on $\partial N$, so the definition
  extends Lipschitz across $\partial N$. Extend $m$ to the rest of $Q(r^*)$
  so that $\sup m=M$, $m$ is constant on a small ball $P_1\subseteq
  E\setminus\overline{N}$, $|E|\geq\tfrac12|Q(r^*)|$ (a normalisation
  independent of $\ell$), and $|\nabla m|\leq CM/r^*$ off $N$; the last
  contributes at most $C\nu M^2(r^*)^{-2}|Q(r^*)| = C\nu M^{1/2}$ to (A2).

  Let $\zeta$ be Lipschitz with $0\leq\zeta\leq1$,
  $\operatorname{supp}\zeta\subseteq N_0$ and $\zeta\not\equiv0$, and set
  $w=W\zeta^2$ on $N_0$ plus a further bump $W'\ge0$ supported in $P_1$, and
  $w=0$ elsewhere; set $h=W|\nabla\zeta|^2$ on $N_0$ and correspondingly on
  $P_1$, so that (A3) holds with
  $|\nabla w|=2W\zeta|\nabla\zeta|=2w^{1/2}h^{1/2}$.

  Choose $\epsilon$ and $W,W'$ so that
  $4\nu G|N| = \tfrac12 C_{\mathcal{E}}\mathcal{L}M^{1/2}$,
  $\nu M^2(W|N|+W'|P_1|)\leq\tfrac14C_{\mathcal{E}}\mathcal{L}M^{1/2}$, and
  $C\nu M^{1/2}\leq\tfrac14C_{\mathcal{E}}\mathcal{L}M^{1/2}$ (enlarging
  $C_{\mathcal{E}}$ if necessary); then (A1)--(A2) hold. Since $w=h=0$ on
  $\{m\leq M/4\}$, both left-hand sides of
  \eqref{eq:s41-bandY}--\eqref{eq:s41-bandD} vanish for every admissible pair.

  For $c_K$: on $E$ the product $|\nabla m|^2w$ is supported in $N_0$, where
  $|\nabla m|^2\equiv G$, so $\iint_E|\nabla m|^2w = G\iint_{N_0} w$; also
  $\iint_E w \ge \iint_{N_0} w>0$ and
  \[
    \iint_E|\nabla m|^2 \;\leq\; 4G|N| + C\frac{M^2}{(r^*)^2}|Q(r^*)|
    \;=\;4G|N| + C M^{1/2}
    \;\leq\; 6\,G|N|
  \]
  by the third normalisation above. Hence
  \[
    c_K \;=\;\frac{|E|\,G\iint_{N_0} w}
      {\bigl(\iint_E|\nabla m|^2\bigr)\bigl(\iint_E w\bigr)}
    \;\geq\;\frac{|E|\,G\iint_{N_0} w}
      {6 G|N|\cdot\bigl(\iint_{N_0}w+W'|P_1|\bigr)}
    \;\geq\;\frac{|E|}{6|N|}\cdot
      \frac{\iint_{N_0}w}{\iint_{N_0}w+W'|P_1|} ,
  \]
  and taking $W'|P_1|\leq\tfrac13\iint_{N_0} w$ gives
  $c_K\geq |E|/(8|N|)$, which is the claim, since
  $|E|\geq\tfrac12|Q(r^*)|$ is independent of $\ell$ while
  $|N|=\ell^5\to0$.
\end{proof}

\begin{remark}[The conclusions do not compare either]\label{rem:s44-conclusions}
  One might hope that even if the hypotheses do not compare, their
  \emph{conclusions} do: Conjecture~\ref{target:band-absorption} yields
  $K$-absorption (through Remark~\ref{rem:s41-band-status} and
  Proposition~\ref{prop:K-amended}), and $K$-absorption is what
  Conjecture~\ref{conj:K-refined} was introduced to supply. But $K$-absorption
  is an \emph{upper bound} on $Y=\nu\iint|\nabla m|^2w\eta^2$, whereas $c_K$ is
  a \emph{ratio}; a small $Y$ is perfectly compatible with a large $c_K$,
  since the denominator
  $\bigl(\iint_E|\nabla m|^2\bigr)\bigl(\iint_Ew\bigr)/|E|$ may be smaller
  still. That is exactly what the configuration of
  Proposition~\ref{prop:s44-b-fails} realises. So the conclusions do not
  compare in that direction either.
\end{remark}

\begin{proposition}[The $c_K$ hypothesis does not imply band absorption]
  \label{prop:s44-a-fails}
  There is a family of admissible test configurations, indexed by
  $\ell\in(0,r^*/4)$, with
  \[
    c_K \;=\; 0
  \]
  --- the strongest possible form of Conjecture~\ref{conj:K-refined} --- for
  which every admissible level pair of width
  $\theta_+-\theta_-\geq\tfrac1{16}$ violates \eqref{eq:s41-bandY}: writing
  $\Gamma$ for the quantity defined in the proof,
  \[
    \nu\iint_{\mathrm{Band}}|\nabla m|^2w\,\psi^2 \;\geq\;\tfrac14\,\Gamma,
    \qquad
    \nu\iint|\nabla m|^2w\,\eta_\chi^2 \;\leq\; 8\,\Gamma,
    \qquad
    M^2\nu\iint h\,\eta_\chi^2 \;+\;\mathcal{I}_0
    \;\leq\; C\Bigl[\Bigl(\frac{\ell}{r^*}\Bigr)^{\!3}
      +\Bigl(\frac{\ell}{r^*}\Bigr)^{\!5}\Bigr]\Gamma .
  \]
  Hence the constant $C$ required in \eqref{eq:s41-bandY} diverges as
  $\ell/r^*\to0$ with $\mathcal{L}$ held fixed, and
  Conjecture~\ref{target:band-absorption} fails along the family.
\end{proposition}
\begin{proof}
  Let $\varphi$ be as in the previous proof. Fix $\ell\in(0,r^*/4)$, let
  $F=K\times I$ with $K\subseteq B_\rho$ an axis-aligned cube of side $\ell$
  and $I$ an interval of length $\ell^2$ inside $\{\psi=1\}$, so
  $|F|=\ell^5$, and let $N=K'\times I$ where $K'$ is the concentric cube of
  side $\tfrac98\ell$, so that $|N|=(\tfrac98)^3|F|\leq2|F|$. Let $\xi$ be a
  Lipschitz cutoff, independent of $t$, with $\xi\equiv1$ on $K$,
  $\operatorname{supp}\xi\subseteq K'$ and $|\nabla\xi|\leq16/\ell$, and set
  \[
    m(x,t)\;=\;\frac{M}{16}+\frac{3M}{16}\,\xi(x)\,
      \varphi\!\Bigl(\frac{x_1}{\epsilon}\Bigr)
    \quad\text{on }N,
    \qquad
    |\nabla m|^2\;=\;G\;:=\;\frac{9M^2}{256\,\epsilon^2}\ \text{ on }F ,
  \]
  with $|\nabla m|^2\leq4G$ on $N$ whenever $\epsilon\leq\ell/16$ (since
  $|\nabla m|\leq\tfrac{3M}{16}(|\nabla\xi|+\epsilon^{-1})
  \leq\tfrac{3M}{8\epsilon}=2G^{1/2}$), and $m\equiv M/16$ on $\partial N$, so
  that the definition extends Lipschitz across $\partial N$.
  Then $m/M$ takes values in $[\tfrac1{16},\tfrac14]$ only. Moreover, provided
  $\epsilon\leq\ell/16$ (verified below), for every
  $[\theta_-,\theta_+]\subseteq[\tfrac1{16},\tfrac14]$ of length
  $\geq\tfrac1{16}$,
  \begin{equation}\label{eq:s44-thirdmass}
    \bigl|\{m/M\in[\theta_-,\theta_+]\}\cap F\bigr|
    \;\geq\;\frac34\cdot\frac{1/16}{3/16}\,|F| \;=\;\frac{|F|}{4}.
  \end{equation}
  Indeed, $F$ is a product of the cube $K$ with $I$, and the $x_1$-side of $K$,
  an interval of length $\ell$, contains at least
  $\lfloor \ell/2\epsilon\rfloor\geq \ell/2\epsilon-1$ full periods of
  $\varphi(\cdot/\epsilon)$; on each full period the set where
  $\varphi\in[a,b]$ has measure exactly $(b-a)\cdot2\epsilon$, so the total is
  at least $(b-a)(\ell-2\epsilon)\geq\tfrac34(b-a)\ell$ when
  $\epsilon\leq\ell/16$. Take $[a,b]$ to be the preimage of
  $[\theta_-,\theta_+]$, of length $(\theta_+-\theta_-)/(3/16)\geq\tfrac13$.

  Let $\zeta$ be Lipschitz and independent of $t$, with $\zeta\equiv1$ on $K$,
  $\operatorname{supp}\zeta\subseteq K'$ and $|\nabla\zeta|\leq16/\ell$, and put
  $w=W\zeta^2$, $h=W|\nabla\zeta|^2\leq 256W\ell^{-2}$ on $N$, so (A3) holds as
  before. Off $N$, extend $m$ Lipschitz with $|\nabla m|\leq CM/r^*$ so that
  $\sup m=M$, and set $w=h=0$ there except on two disjoint balls
  $P_1,P_2\subseteq E$ on which, respectively, $m$ is constant and a small
  bump of $w$ is placed, and $m$ is non-constant with $w=0$. Then
  $\iint_E|\nabla m|^2w=0$ while both denominators of $c_K$ are positive, so
  $c_K=0$.

  Fix $\epsilon,W$ by
  \[
    \nu G|F| \;=\;\gamma\, C_{\mathcal{E}}\mathcal{L}M^{1/2},
    \qquad
    \nu M^2W|F| \;=\;\gamma\, C_{\mathcal{E}}\mathcal{L}M^{1/2},
    \qquad \gamma:=\tfrac1{32} ,
  \]
  which is consistent with (A2): on $N$ one has $|\nabla m|^2\leq4G$ and
  $w\leq W$ with $|N|\leq2|F|$, so those two contributions are at most
  $8\gamma$ and $2\gamma$ times $C_{\mathcal{E}}\mathcal{L}M^{1/2}$, i.e.\
  at most $\tfrac14+\tfrac1{16}$ of the budget; the far field
  ($\leq C\nu M^{1/2}\leq\tfrac14C_{\mathcal{E}}\mathcal{L}M^{1/2}$ after
  enlarging $C_{\mathcal{E}}$) and the two small bumps on $P_1,P_2$ use the
  rest. Set
  \[
    \Gamma \;:=\;\nu\,G\,W\,|F|
    \;=\;\bigl(\nu G|F|\bigr)\,W
    \;=\;\gamma C_{\mathcal{E}}\mathcal{L}M^{1/2}\cdot
      \frac{\gamma C_{\mathcal{E}}\mathcal{L}M^{-3/2}}{\nu|F|}
    \;=\;\frac{\gamma^2\,C_{\mathcal{E}}^2\,\mathcal{L}^2}{\nu\,M\,|F|} .
  \]

  \emph{First display.} By \eqref{eq:s44-thirdmass}, $\zeta\equiv1$ and
  $\psi\equiv1$ on $F$, the left-hand side of \eqref{eq:s41-bandY} is at least
  $\nu GW|F|/4=\Gamma/4$.

  \emph{Second display.} $|\nabla m|^2w$ is supported in $N$, where
  $|\nabla m|^2\leq4G$ and $w\leq W$; since $\eta_\chi^2\leq\psi^2\leq1$ and
  $|N|\leq2|F|$, $\nu\iint|\nabla m|^2w\eta_\chi^2\leq4\nu GW|N|\leq8\Gamma$.

  \emph{Third display.} $h$ is supported in $N\setminus F$ and bounded by
  $256W\ell^{-2}$ there, and $|N\setminus F|\leq|N|\leq2|F|$, so
  \[
    \frac{M^2\nu\iint h\,\eta_\chi^2}{\Gamma}
    \;\leq\;\frac{\nu M^2\cdot 256W\ell^{-2}\cdot 2|F|}{\nu GW|F|}
    \;=\;\frac{512\,M^2}{G\,\ell^{2}}
    \;=\;\frac{512\cdot 256\,\epsilon^2}{9\,\ell^{2}} .
  \]
  From $\nu G|F|=\gamma C_{\mathcal{E}}\mathcal{L}M^{1/2}$, $|F|=\ell^5$ and
  $\mathcal{L}\geq1$ one gets
  $\epsilon^2 = \tfrac{9\gamma^{-1}}{256}\,\nu M^{3/2}\ell^{5}
  /(C_{\mathcal{E}}\mathcal{L}) \leq C\nu M^{3/2}\ell^{5}$, whence the ratio is
  at most $C\nu M^{3/2}\ell^{3}=C\nu(\ell/r^*)^{3}$; the same bound gives
  $\epsilon^2/\ell^2\leq C\nu(\ell/r^*)^{3}$, hence $\epsilon\leq\ell/16$ once
  $(\ell/r^*)^{3}\leq(256\,C\nu)^{-1}$, which is the standing smallness
  assumption on $\ell/r^*$ and validates both \eqref{eq:s44-thirdmass} and the
  bound $|\nabla m|^2\leq4G$ on $N$. Finally
  \[
    \frac{\mathcal{I}_0}{\Gamma}
    \;=\;\frac{\mathcal{L}^2M^{3/2}\nu^{-1}}
      {\gamma^2C_{\mathcal{E}}^2\mathcal{L}^2\,(\nu M|F|)^{-1}}
    \;=\;\frac{M^{5/2}|F|}{\gamma^2C_{\mathcal{E}}^2}
    \;=\;\frac{1}{\gamma^2C_{\mathcal{E}}^2}
      \Bigl(\frac{\ell}{r^*}\Bigr)^{5} ,
  \]
  using $M^{5/2}=(r^*)^{-5}$ and $|F|=\ell^5$.

  \emph{Conclusion.} If \eqref{eq:s41-bandY} held with some constant $C$, then
  $\tfrac14\Gamma\leq 2^{-11}\cdot8\Gamma
  + C'\bigl[(\ell/r^*)^{3}+(\ell/r^*)^{5}\bigr]\Gamma$, i.e.\
  $\tfrac14-2^{-8}\leq C'\bigl[(\ell/r^*)^{3}+(\ell/r^*)^{5}\bigr]$, which
  fails for $\ell/r^*$ small. Since $\mathcal{L}$ is held fixed while
  $\ell\to0$, no choice of $C$ works along the family.
\end{proof}

\begin{remark}[Reading of Proposition~\ref{prop:s44-a-fails}]
  \label{rem:s44-reading}
  The configuration is the natural adversary: $|\omega|$ oscillates rapidly
  \emph{inside} the transition band, so that $|\nabla m|^2$ and $w$ are both
  large on a common set whose $m$-values never leave $[\tfrac1{16},\tfrac14]$.
  Conjecture~\ref{conj:K-refined} is blind to it, being a statement about
  $\{m\geq M/4\}$, where in this configuration nothing happens. Two further
  observations. (i) The oscillation is what makes $|\nabla m|$ large while
  $m$ stays inside the band; there is no contradiction with $m\in[M/16,M/4]$,
  and this is exactly the mechanism identified in
  Remark~\ref{rem:s41-contrast} as the reason the band terms are not
  controlled by the enstrophy budget alone. (ii) The same configuration has
  $\Xi[\mathbf{1}_{\mathrm{Band}}\psi^2]\geq c(r^*/\ell)^{5}\to\infty$,
  consistently with Proposition~\ref{prop:s44-jle-implies} below: the
  hypothesis that \emph{does} imply \eqref{eq:s41-bandY} is violated here, as
  it must be.
\end{remark}

\subsection{What has to be added: the band transport}
\label{subsec:s44-transport}

Proposition~\ref{prop:s44-a-fails} shows that
Conjecture~\ref{conj:K-refined} must at least be \emph{transported} to the
band before any implication can be hoped for. It is then a genuine
implication, and the extra ingredient is explicit.

\begin{proposition}[Band-transported $c_K$ implies \eqref{eq:s41-bandY}]
  \label{prop:s44-transport}
  Let $\theta_\pm$ be admissible and suppose
  \begin{equation}\label{eq:s44-cKband}
    c_K^{\mathrm{band}}
    \;:=\;
    \frac{|\mathrm{Band}|\displaystyle\iint_{\mathrm{Band}}
        |\nabla m|^2w\,\psi^2}
         {\Bigl(\displaystyle\iint_{\mathrm{Band}}|\nabla m|^2\psi^2\Bigr)
          \Bigl(\displaystyle\iint_{\mathrm{Band}}w\,\psi^2\Bigr)}
    \;\leq\;\frac{C_1}{\mathcal{L}} .
  \end{equation}
  If in addition
  \begin{equation}\label{eq:s44-NB}
    \textup{(NB)}\qquad
    |\mathrm{Band}| \;\geq\;\frac{\kappa_1\,|Q(r^*)|}{\mathcal{L}}
    \qquad\text{for some }\kappa_1>0 ,
  \end{equation}
  then \eqref{eq:s41-bandY} holds --- indeed in the stronger form
  \[
    \nu\iint_{\mathrm{Band}}|\nabla m|^2w\,\psi^2
    \;\leq\;\frac{C\,C_1}{\kappa_1}\;\mathcal{I}_0 ,
  \]
  with neither the $Y$-term nor the $D$-term needed on the right.
\end{proposition}
\begin{proof}
  By Definition~\ref{def:s44-excess},
  $\Xi[\mathbf{1}_{\mathrm{Band}}\psi^2]
  = c_K^{\mathrm{band}}\,|Q(r^*)|/|\mathrm{Band}|
  \leq (C_1/\mathcal{L})(\mathcal{L}/\kappa_1)=C_1/\kappa_1$ by
  \eqref{eq:s44-cKband} and \eqref{eq:s44-NB}. Now apply
  \eqref{eq:s44-master} with $\varpi=\mathbf{1}_{\mathrm{Band}}\psi^2$, which
  is supported in $Q^\sharp$ because $\theta_-\geq\tfrac1{16}$.
\end{proof}

\begin{remark}[(NB) is not free, and is of the same family]
  \label{rem:s44-NB}
  Hypothesis \eqref{eq:s44-NB} says the transition band is not too thin. It is
  not automatic: by the co-area formula, as in
  Proposition~\ref{prop:s41-coarea}(d),
  \[
    |\mathrm{Band}|
    \;=\;\int\!\!\int_{\theta_-M}^{\theta_+M}
      \int_{\{m=s\}}\frac{d\mathcal{H}^2}{|\nabla m|}\,ds\,dt ,
  \]
  so a thin band is precisely a band on which $|\nabla m|$ is large --- the
  dangerous case, and the one realised in
  Proposition~\ref{prop:s44-a-fails}. Thus (NB) is itself a non-concentration
  statement about $|\nabla m|$ on level sets of $|\omega|$, and it should not
  be presented as a harmless technical addendum. Two consequences are worth
  recording. First, the honest answer to ``does
  Conjecture~\ref{conj:K-refined} imply
  Conjecture~\ref{target:band-absorption}?'' is: \emph{not as stated; its band
  transport implies only the $Y$-half \eqref{eq:s41-bandY}, and only when
  supplemented by} \eqref{eq:s44-NB}. Second --- and this is the point missed
  in the first draft of this section, see Remark~\ref{rem:s44-corrigendum} ---
  \emph{the bulk side needs a volume hypothesis of exactly the same kind}. By
  Lemma~\ref{lem:s44-xi-vs-cK}, passing from
  Conjecture~\ref{conj:K-refined} to the volume-normalised bound
  $\Xi[\mathbf{1}_E]=O(1)$ used in
  Proposition~\ref{prop:s44-calibration}(i) requires
  \begin{equation}\label{eq:s44-NBE}
    \textup{(NB-E)}\qquad
    |E| \;\geq\;\frac{\kappa_2\,|Q(r^*)|}{\mathcal{L}}
    \qquad\text{for some }\kappa_2>0 ,
  \end{equation}
  the exact analogue of \eqref{eq:s44-NB} with $\mathrm{Band}$ replaced by
  $E$. The symmetry is complete: neither the bulk route nor the band route is
  free of a volume condition, and in both cases the condition fails precisely
  when the region carrying the coupling is thin --- the filamentary regime.
\end{remark}

\begin{remark}[The two routes, and what each costs]\label{rem:s44-two-routes}
  It is worth separating, once and for all, the two ways a correlation
  hypothesis can be cashed in; the strength comparison of
  \S\ref{subsec:s44-strength} is entirely a matter of which route is used.
  \begin{enumerate}[label=\textup{(R\arabic*)}]
    \item \emph{The $\|w\|_{L^\infty}$ route} (that of
      Proposition~\ref{prop:crude-K} and Proposition~\ref{prop:s41-crude}):
      bound $\iint w\,\varpi\leq\|w\|_{L^\infty}|R|$, so that the volume
      factor in $c_K[\varpi]$ cancels and one lands on the Moser-absorbable
      quantity $\delta\nu(r^*)^3\|w\|_{L^\infty}$ of
      Proposition~\ref{prop:second-rung}. No volume hypothesis; the price is
      that a factor $\mathcal{L}$ must be found inside the correlation
      constant, and that the closure runs through $\|w\|_{L^\infty}$
      (Proposition~\ref{prop:s44-linfty-route}).
    \item \emph{The budget route} (Lemma~\ref{lem:s44-budget}): bound
      $\iint w\,\varpi$ by the a priori $w$-budget of
      Theorem~\ref{thm:apriori-sigma}, so that $\|w\|_{L^\infty}$ never
      appears and one lands directly in $\mathcal{I}_0$. No logarithm is
      needed in the correlation constant; the price is the volume factor
      $|Q(r^*)|/|R|$, i.e.\ a hypothesis of the form \eqref{eq:s44-NB} or
      \eqref{eq:s44-NBE}.
  \end{enumerate}
  Neither route dominates: (R1) is cheaper exactly when
  $|R|/|Q(r^*)|<1/\mathcal{L}$, (R2) exactly when
  $|R|/|Q(r^*)|>1/\mathcal{L}$. This trade-off, and not any difference between
  the two conjectures, is what the ``$C/\mathcal{L}$ versus $O(1)$''
  discrepancy in their statements records.
\end{remark}

\begin{proposition}[The $\|w\|_{L^\infty}$ route: both targets require a
  \emph{decaying} correlation constant]\label{prop:s44-linfty-route}
  Let $R\subseteq Q^\sharp$ be measurable, $\varpi=\mathbf{1}_R\psi^2$, and
  suppose $c_K[\varpi]\leq C_1/\mathcal{L}$, with $c_K[\varpi]$ as in
  Lemma~\ref{lem:s44-xi-vs-cK}. Then
  \[
    \frac{\nu}{M^2}\iint_R|\nabla m|^2\,w\,\psi^2
    \;\leq\; C\,C_1\,
      \|w\|_{L^\infty(Q(\rho')\cap\{m\geq M/16\})}\,(r^*)^3 ,
  \]
  which is precisely the quantity Proposition~\ref{prop:second-rung} absorbs
  with factor $\delta$, provided the absolute constant $CC_1$ satisfies
  $CC_1\leq\delta\nu$. Taken with $R=E$ this is the conclusion
  Conjecture~\ref{conj:K-refined} was designed to give
  (Proposition~\ref{prop:crude-K} with its logarithm removed); taken with
  $R=\mathrm{Band}$ it removes exactly the shortfall of
  Remark~\ref{rem:s41-shortfall} from the $Y$-half of the core residual. No
  volume hypothesis is used in either case.
\end{proposition}
\begin{proof}
  By the definition of $c_K[\varpi]$ in Lemma~\ref{lem:s44-xi-vs-cK} and
  $\iint w\,\varpi\leq\|w\|_{L^\infty}|R|$,
  \[
    \iint_R|\nabla m|^2w\,\psi^2
    \;=\;\frac{c_K[\varpi]}{|R|}
      \Bigl(\iint_R|\nabla m|^2\psi^2\Bigr)
      \Bigl(\iint_R w\,\psi^2\Bigr)
    \;\leq\; c_K[\varpi]\,\|w\|_{L^\infty}\iint_R|\nabla m|^2\psi^2 .
  \]
  Multiply by $\nu/M^2$, use $\nu\iint_R|\nabla m|^2\psi^2\leq\mathcal{E}
  \leq C_{\mathcal{E}}\mathcal{L}M^{1/2}$ from \eqref{eq:s41-budget}, and
  $M^{-3/2}=(r^*)^3$; the hypothesis $c_K[\varpi]\leq C_1/\mathcal{L}$ cancels
  the $\mathcal{L}$.
\end{proof}

\begin{remark}[A bookkeeping caveat, flagged and not resolved]
  \label{rem:s44-channel-caveat}
  Proposition~\ref{prop:s44-linfty-route} delivers the band term into the
  \emph{Moser channel} $\delta\nu(r^*)^3\|w\|_{L^\infty}$, which is the target
  Remark~\ref{rem:s41-shortfall} names, and not into the right-hand side of
  \eqref{eq:s41-endtarget}, which has no $\|w\|_{L^\infty}$ slot. These are
  two different closures of the same residual: the first leaves
  Proposition~\ref{prop:K-amended} in the form
  $Y\leq CM^2D_{\mathrm{ang}}+\mathcal{I}+C\|w\|_{L^\infty}M^{1/2}$, the last
  term being absorbed downstream by Proposition~\ref{prop:second-rung}; the
  second removes the residual outright. The document states both targets and
  does not reconcile them --- the discrepancy is pre-existing (compare
  Remark~\ref{rem:s41-shortfall} with \eqref{eq:s41-endtarget}) --- and it is
  recorded here rather than papered over. Separately, the constant condition
  $CC_1\leq\delta\nu$ is a genuine if non-divergent residue; it is present
  identically in Conjecture~\ref{conj:K-refined}, which this document does not
  track either, and nothing below claims to remove it.
\end{remark}

\subsection{The strength comparison: where the logarithm really comes from}
\label{subsec:s44-strength}

\begin{proposition}[Calibration]\label{prop:s44-calibration}
  Let $\kappa_0\geq1$. Then:
  \begin{enumerate}[label=(\roman*)]
    \item \textup{(bulk)} If $\Xi[\mathbf{1}_{Q^\sharp}\psi^2]\leq\kappa_0$
      then, for every admissible $\eta_\chi$,
      \[
        \nu\iint|\nabla m|^2\,w\,\eta_\chi^2 \;\leq\; C\kappa_0\,\mathcal{I}_0 ,
      \]
      i.e.\ the coupling term $K$ lands in the allowed inhomogeneous class
      \eqref{eq:s41-I} outright --- with no $D_{\mathrm{ang}}$ term, no
      $\|w\|_{L^\infty}$ and no annulus residual.
    \item \textup{(band)} If
      $\Xi[\mathbf{1}_{\mathrm{Band}}\psi^2]\leq\kappa_0$ then
      \eqref{eq:s41-bandY} holds, again in the stronger form with right-hand
      side $C\kappa_0\mathcal{I}_0$.
    \item \textup{(comparison --- corrected)} The hypotheses used in (i) and
      (ii) are \emph{stronger}, not weaker, than the corresponding
      correlation bounds. By Lemma~\ref{lem:s44-xi-vs-cK},
      $\Xi[\mathbf{1}_E]\leq\kappa_0$ implies $c_K\leq\kappa_0$ but not
      conversely, the gap being the factor $|Q(r^*)|/|E|$; and
      Conjecture~\ref{conj:K-refined} gives only
      $\Xi[\mathbf{1}_E]\leq(C/\mathcal{L})|Q(r^*)|/|E|$, which is $O(1)$
      exactly under \eqref{eq:s44-NBE}. Symmetrically on the band, with
      \eqref{eq:s44-NB}. Neither $\Xi$-bound is therefore available from the
      $c_K$-type hypotheses without a volume condition.
  \end{enumerate}
\end{proposition}
\begin{proof}
  (i) and (ii) are \eqref{eq:s44-master} applied with
  $\varpi=\mathbf{1}_{Q^\sharp}\psi^2$ and
  $\varpi=\mathbf{1}_{\mathrm{Band}}\psi^2$ respectively, using
  $\eta_\chi^2\leq\psi^2\mathbf{1}_{\{m\geq M/16\}}$
  (Lemma~\ref{lem:s41-admissible}(ii)) in (i) and
  $\mathrm{Band}\subseteq Q^\sharp$ in (ii). (iii) is
  Lemma~\ref{lem:s44-xi-vs-cK}.
\end{proof}

\begin{remark}[Answer to the strength question, corrected]
  \label{rem:s44-strength}
  The premise behind the question ``is
  Conjecture~\ref{target:band-absorption} strictly weaker than
  Conjecture~\ref{conj:K-refined}, needing only an $O(1)$ equipartition
  constant where $c_K$ needs a logarithmically decaying one?'' does not
  survive the correction recorded in Remark~\ref{rem:s44-corrigendum}. The
  answer is \textbf{no}, and the corrected picture is the following.
  \begin{enumerate}[label=(\arabic*)]
    \item \emph{\eqref{eq:s41-bandY} is $O(1)$ in its bookkeeping, not in its
      strength.} It is stated against $Y$, $M^2D_{\mathrm{ang}}$ and
      $\mathcal{I}_0$ rather than against a correlation constant, which is
      what makes its constants look absolute. Expressed as a hypothesis on the
      band correlation constant it is \emph{not} an $O(1)$ statement: by
      Proposition~\ref{prop:s44-linfty-route} with $R=\mathrm{Band}$, what the
      $\|w\|_{L^\infty}$ route requires is
      $c_K^{\mathrm{band}}\leq C/\mathcal{L}$ --- the exact analogue of
      Conjecture~\ref{conj:K-refined}, on the band. In this precise and
      limited sense the ``same analytic type'' language of
      Remark~\ref{rem:s41-band-status} is \emph{vindicated}: both targets
      require a logarithmically decaying correlation constant, on their
      respective regions.
    \item \emph{The alternative --- an $O(1)$ constant --- is bought with a
      volume hypothesis, symmetrically on both sides.} The budget route
      (Proposition~\ref{prop:s44-calibration}) does replace $C/\mathcal{L}$ by
      $O(1)$, but only in the volume-normalised statistic $\Xi$; and by
      Lemma~\ref{lem:s44-xi-vs-cK} an $O(1)$ bound on $\Xi$ is
      \emph{strictly stronger} than an $O(1)$ bound on the correlation
      constant, by the factor $|Q(r^*)|/|R|$. Recovering it from a $c_K$-type
      hypothesis costs \eqref{eq:s44-NB} on the band and \eqref{eq:s44-NBE} on
      the bulk. Remark~\ref{rem:s44-two-routes} states the trade-off: neither
      route dominates, the crossover being at $|R|/|Q(r^*)|=1/\mathcal{L}$.
    \item \emph{Consequently the logarithm is not an ``artefact''.} It is what
      the $\|w\|_{L^\infty}$ route costs, and that route is the one requiring
      no information about the measure of the region carrying the coupling.
      Removing the logarithm means acquiring that information instead. The
      first draft of this section asserted that
      Conjecture~\ref{conj:K-refined} could be weakened to an $O(1)$ statement
      without loss; that assertion is \textbf{withdrawn}
      (Remark~\ref{rem:s44-corrigendum}).
    \item \emph{Net verdict on the strength question.}
      Conjecture~\ref{target:band-absorption} is \textbf{not} shown to be
      weaker than Conjecture~\ref{conj:K-refined}. The two are incomparable as
      statements (Propositions~\ref{prop:s44-a-fails}
      and~\ref{prop:s44-b-fails}), and when each is reduced to a hypothesis on
      a correlation constant by the same route they require bounds of the
      \emph{same} form. What separates them is their support, not their
      strength.
  \end{enumerate}
  One caveat, recorded rather than concealed: landing $K$, or the band term,
  in $\mathcal{I}_0=\mathcal{L}^2M^{3/2}\nu^{-1}$ is ``allowed'' only in the
  sense of \S\ref{sec:verification-s36}~(C2), i.e.\ it is an
  $\mathcal{L}^2$-degraded inhomogeneous term of the same marginal class as
  every other such term in the scheme. It therefore inherits, and does not
  resolve, the standing marginality issue of Remark~\ref{rem:marginality}; it
  introduces no \emph{new} degradation, which is all that is claimed. The
  $\|w\|_{L^\infty}$ route carries its own caveat,
  Remark~\ref{rem:s44-channel-caveat}.
\end{remark}

\begin{remark}[What the existing numerics do and do not support]
  \label{rem:s44-numerics}
  The $c_K$ measurements of S34--S35 (\texttt{EXP-L4-KCORR-*}) reported
  $c_K\approx1$ and were explicitly caveated as unable to distinguish $O(1)$
  from the stronger $c_K\leq C/\mathcal{L}$. That caveat stands, and after the
  correction of Remark~\ref{rem:s44-corrigendum} a second must be added
  beside it.
  \begin{enumerate}[label=(\roman*)]
    \item \emph{A measured $c_K\approx1$ does not support
      $\Xi[\mathbf{1}_E]=O(1)$.} By Lemma~\ref{lem:s44-xi-vs-cK},
      $\Xi[\mathbf{1}_E]=c_K\,|Q(r^*)|/|E|\geq c_K$, so such a measurement
      bounds $\Xi$ only once the volume fraction $|E|/|Q(r^*)|$ is known and
      bounded below. The first draft of this section asserted the reverse
      inequality and drew the opposite conclusion; that is withdrawn.
    \item \emph{The gap is worst where it matters.} The correction factor
      $|Q(r^*)|/|E|$ is largest exactly when the near-maximal vorticity set is
      thin --- the filamentary regime expected near a putative singularity ---
      so an $O(1)$ correlation measurement in a benign configuration says
      correspondingly less about $\Xi$ in a concentrated one. Volume fractions
      quoted from the S42--S43 diagnostics should be treated as motivation for
      measuring $|E|/|Q(r^*)|$ directly and \emph{not} as evidence about it:
      those diagnostics take the near-maximal set over the whole periodic box
      rather than over $Q(r^*)$, which is not the quantity appearing in
      Lemma~\ref{lem:s44-xi-vs-cK}.
    \item \emph{What should be measured.} Three quantities, none of which the
      existing suite reports in the form needed: the volume fractions
      $|E|/|Q(r^*)|$ and $|\mathrm{Band}|/|Q(r^*)|$ at the working scale
      $r^*$, to test \eqref{eq:s44-NBE} and \eqref{eq:s44-NB}; the band
      correlation constant $c_K^{\mathrm{band}}$ of \eqref{eq:s44-cKband},
      which by Proposition~\ref{prop:s44-a-fails} cannot be inferred from the
      $E$-measurement; and the $\mathcal{L}$-dependence of both correlation
      constants, since Proposition~\ref{prop:s44-linfty-route} needs decay and
      not merely boundedness.
  \end{enumerate}
  The plain statement, replacing the first draft's: \emph{the existing $O(1)$
  measurements are consistent with, but do not establish, either
  Conjecture~\ref{conj:K-refined} or Conjecture~\ref{conj:lc-s44}; and they
  bear on Conjecture~\ref{conj:jle-s44} only in combination with a volume
  measurement that has not been made.}
\end{remark}

\subsection{Common hypotheses}
\label{subsec:s44-common}

Two hypotheses are recorded, one for each of the routes of
Remark~\ref{rem:s44-two-routes}. Each implies the $Y$-halves of \emph{both}
targets; neither implies the other; and the first is the weaker of the two in
the regime that matters. Both are open. Neither covers the angular half
\eqref{eq:s41-bandD}, for the reason given in Remark~\ref{rem:s44-bandD}.

\begin{conjecture}[Level-uniform decorrelation]\label{conj:lc-s44}
  There is $C_1$, depending only on the standing Type-I data, such that near a
  putative Type-I singularity, at the working scale $r^*=M(t_0)^{-1/2}$, with
  $Q^\sharp=Q(\rho')\cap\{m\geq M/16\}$ and $c_K[\,\cdot\,]$ as in
  Lemma~\ref{lem:s44-xi-vs-cK},
  \begin{equation}\label{eq:s44-lc}
    \textup{(LC)}\qquad
    c_K\bigl[\mathbf{1}_{Q^\sharp}\psi^2\bigr]
    \;\leq\;\frac{C_1}{\mathcal{L}} .
  \end{equation}
  This is \emph{verbatim} the hypothesis of Conjecture~\ref{conj:K-refined} ---
  same normalisation, same $1/\mathcal{L}$ decay --- asserted on
  $\{m\geq M/16\}$ instead of on $\{m\geq M/4\}$, i.e.\ extended downwards
  across the transition band.
\end{conjecture}

\begin{proposition}[(LC) covers both $Y$-halves]\label{prop:s44-lc-implies}
  Assume \eqref{eq:s44-lc}. Then, with
  $\|w\|_\infty=\|w\|_{L^\infty(Q(\rho')\cap\{m\geq M/16\})}$,
  \[
    \frac{\nu}{M^2}\iint|\nabla m|^2\,w\,\eta_\chi^2
    \;+\;\frac{\nu}{M^2}\iint_{\mathrm{Band}}|\nabla m|^2\,w\,\psi^2
    \;\leq\; C\,C_1\,\|w\|_\infty\,(r^*)^3
  \]
  for every admissible $\chi$ and every admissible level pair. Both terms are
  therefore of the class Proposition~\ref{prop:second-rung} absorbs with
  factor $\delta$, subject to the constant condition $CC_1\leq\delta\nu$ of
  Proposition~\ref{prop:s44-linfty-route}. In particular \eqref{eq:s44-lc}
  simultaneously supplies the conclusion sought from
  Conjecture~\ref{conj:K-refined} and removes the shortfall of
  Remark~\ref{rem:s41-shortfall} from the $Y$-half of the core residual, with
  \emph{no} volume hypothesis.
\end{proposition}
\begin{proof}
  Both $\{\eta_\chi>0\}$ and $\mathrm{Band}$ are contained in $Q^\sharp$ and
  the integrand is nonnegative, so both left-hand terms are bounded by
  $\tfrac{\nu}{M^2}\iint_{Q^\sharp}|\nabla m|^2w\,\psi^2$; apply
  Proposition~\ref{prop:s44-linfty-route} with $R=Q^\sharp$.
\end{proof}

\begin{remark}[Why (LC) is the right reading of ``same analytic type'']
  \label{rem:s44-lc-reading}
  Remark~\ref{rem:s41-summary} said that $\mathcal{A}_{\mathrm{ann}}$ is
  ``open, of the same analytic type as Conjecture~\ref{conj:K-refined}''.
  Sections~\ref{subsec:s44-nonimplications} show that this cannot mean one
  implies the other. Conjecture~\ref{conj:lc-s44} is what it \emph{can}
  legitimately mean: the two targets are two instances of a \emph{single}
  hypothesis, differing only in the region on which it is asserted, and the
  single hypothesis stated on the union of those regions covers both. This
  costs nothing in form: it is the $c_K$ hypothesis, unchanged in
  normalisation and in the decay rate of its constant, differing from
  Conjecture~\ref{conj:K-refined} only in the region on which it is asserted.

  Two things this does \emph{not} say, lest the sentence above be read too
  generously. First, \eqref{eq:s44-lc} does not imply
  Conjecture~\ref{conj:K-refined}: a bound on a single correlation ratio over
  $Q^\sharp$ does not restrict to the sub-region $E$, the ratio being
  non-monotone under restriction. What
  Proposition~\ref{prop:s44-lc-implies} supplies is the \emph{conclusion}
  sought from Conjecture~\ref{conj:K-refined}, not that conjecture itself.
  Second and conversely, Conjecture~\ref{conj:K-refined} does not imply
  \eqref{eq:s44-lc}: by Proposition~\ref{prop:s44-a-fails} the extension of
  the assertion across the band is exactly the part it does not supply. (LC)
  is therefore a genuinely new hypothesis of an old form, not a repackaging of
  an existing one.
\end{remark}

\begin{conjecture}[Joint level equipartition]\label{conj:jle-s44}
  There is $\kappa_0$, depending only on the standing Type-I data, such that
  near a putative Type-I singularity, at the working scale
  $r^*=M(t_0)^{-1/2}$, with $Q^\sharp=Q(\rho')\cap\{m\geq M/16\}$:
  \begin{align}
    \textup{(JLE-1)}\qquad
    \iint_{Q^\sharp}|\nabla m|^2\,w\,\psi^2
    &\;\leq\;\frac{\kappa_0}{|Q(r^*)|}
      \Bigl(\iint_{Q^\sharp}|\nabla m|^2\,\psi^2\Bigr)
      \Bigl(\iint_{Q^\sharp}w\,\psi^2\Bigr) ,
      \label{eq:s44-jle}\\
    \textup{(JLE-2)}\qquad
    \nu\iint_{\mathrm{Band}}|\nabla^2\hat{e}|^2\,\psi^2
    &\;\leq\;\kappa_0\,\nu\iint|\nabla^2\hat{e}|^2\,\eta_\chi^2
      \;+\;\kappa_0\,\mathcal{L}^2M^{-1/2}\nu^{-1} ,
      \label{eq:s44-jle2}
  \end{align}
  the latter for some admissible level pair $\theta_\pm$ of width
  $\theta_+-\theta_-\geq\tfrac1{16}$ and some
  $\chi\in\mathcal{C}(\theta_-,\theta_+)$.
\end{conjecture}

\begin{proposition}[Both targets follow from
  Conjecture~\ref{conj:jle-s44}]\label{prop:s44-jle-implies}
  Assume \textup{(JLE-1)}. Then
  \begin{enumerate}[label=(\roman*)]
    \item $\nu\iint|\nabla m|^2w\,\eta_\chi^2\leq C\kappa_0\,\mathcal{I}_0$
      for every admissible $\chi$; in particular the conclusion that
      Conjecture~\ref{conj:K-refined} was introduced to supply holds, and
      Conjecture~\ref{conj:K-refined} itself is not needed;
    \item \eqref{eq:s41-bandY} holds for every admissible level pair, in the
      stronger form with right-hand side $C\kappa_0\mathcal{I}_0$.
  \end{enumerate}
  If in addition \textup{(JLE-2)} holds for the pair $\theta_\pm$ and cutoff
  $\chi$ named there, then \eqref{eq:s41-bandD} holds for that pair; hence
  Conjecture~\ref{target:band-absorption} holds in full, hence
  \eqref{eq:s41-endtarget} holds, and the caveat ``modulo
  $\mathcal{A}_{\mathrm{ann}}$'' may be removed from
  Proposition~\ref{prop:K-amended}.
\end{proposition}
\begin{proof}
  \textup{(JLE-1)} is precisely $\Xi[\mathbf{1}_{Q^\sharp}\psi^2]\leq\kappa_0$,
  so (i) and (ii) are Proposition~\ref{prop:s44-calibration}(i)--(ii), using
  $\mathrm{Band}\subseteq Q^\sharp$,
  $\eta_\chi^2\leq\psi^2\mathbf{1}_{\{m\geq M/16\}}$, and the nonnegativity of
  the integrand for the restriction to subregions. Multiplying
  \eqref{eq:s44-jle2} by $M^2$ gives \eqref{eq:s41-bandD} with $C=\kappa_0$,
  since $M^2\cdot\mathcal{L}^2M^{-1/2}\nu^{-1}=\mathcal{L}^2M^{3/2}\nu^{-1}$.
  The last two implications are those of
  Remark~\ref{rem:s41-band-status}.
\end{proof}

\begin{remark}[Status, failure mode, and what (JLE-1) really asserts]
  \label{rem:s44-jle-status}
  (JLE-1) says: \emph{on the region where the level cutoff lives, the coupling
  density $|\nabla m|^2w$ is no larger than a fixed multiple of what
  uncorrelated densities spread over the ambient parabolic cylinder would
  give.} It is an equipartition/non-concentration hypothesis of exactly the
  kind described in the remark following
  Conjecture~\ref{conj:K-refined}, with two differences: the normalising
  volume is $|Q(r^*)|$ rather than the measure of the level region, and the
  constant is $O(1)$ rather than $O(1/\mathcal{L})$. It is \emph{sufficient}
  for both $Y$-halves and \emph{necessary} for neither
  (Proposition~\ref{prop:s44-b-fails} exhibits a configuration in which
  \eqref{eq:s41-bandY} holds while $\Xi$ is unbounded).

  \medskip
  \noindent\textbf{But it is a demanding hypothesis, not a cheap one.} This is
  the point the first draft of this section got backwards
  (Remark~\ref{rem:s44-corrigendum}). By Lemma~\ref{lem:s44-xi-vs-cK},
  (JLE-1) is equivalent to
  $c_K[\mathbf{1}_{Q^\sharp}\psi^2]\leq\kappa_0\,|Q^\sharp|/|Q(r^*)|$, i.e.\ it
  demands a correlation constant not merely bounded but \emph{small in
  proportion to the volume fraction of $Q^\sharp$}. Its exact failure mode is
  co-concentration: if $|\nabla m|^2$ and $w$ are both of sizes $G,W$ on a
  common set $S$ and negligible elsewhere, then $\Xi=|Q(r^*)|/|S|$, which is
  large precisely when the vorticity structure is filamentary --- the regime
  of interest. Neither (JLE-1) nor \eqref{eq:s44-lc} implies the other:
  (JLE-1) gives $c_K[\mathbf{1}_{Q^\sharp}\psi^2]\leq\kappa_0$, which is
  $O(1)$ and not $O(1/\mathcal{L})$, so it does not give \eqref{eq:s44-lc};
  and \eqref{eq:s44-lc} gives
  $\Xi[\mathbf{1}_{Q^\sharp}\psi^2]\leq(C_1/\mathcal{L})|Q(r^*)|/|Q^\sharp|$,
  which is $O(1)$ only under a volume hypothesis of the form
  \eqref{eq:s44-NBE}. They are two genuinely different sufficient hypotheses,
  matched to the two routes of Remark~\ref{rem:s44-two-routes}.

  \medskip
  \noindent\textbf{Which to target.} On present evidence
  Conjecture~\ref{conj:lc-s44} is the better target: it needs no information
  about the measure of any level region, it is the existing $c_K$ hypothesis
  with its region enlarged, and it degrades gracefully under concentration,
  whereas (JLE-1) degrades like $|Q(r^*)|/|S|$. Conjecture~\ref{conj:jle-s44}
  is retained because it is the natural hypothesis for the budget route, it
  yields the stronger conclusion (landing in $\mathcal{I}_0$ outright, with no
  $\|w\|_{L^\infty}$ and hence none of the caveats of
  Remark~\ref{rem:s44-channel-caveat}), and it becomes the more attractive of
  the two in any regime where the near-maximal vorticity set is \emph{not}
  thin.
\end{remark}

\begin{remark}[The angular half is \emph{not} of the $c_K$ type]
  \label{rem:s44-bandD}
  \eqref{eq:s41-bandD}, equivalently (JLE-2), contains no product of two
  densities and no correlation: it is a bare level non-concentration statement
  for the angular curvature density $\nu|\nabla^2\hat{e}|^2\psi^2$. Nothing in
  Conjecture~\ref{conj:K-refined}, in \eqref{eq:s44-jle}, or in any
  decorrelation hypothesis bears on it. It would follow from either of two
  unrelated inputs: an a priori bound
  $\nu\iint_{Q^\sharp}|\nabla^2\hat{e}|^2\psi^2
  \leq C\mathcal{L}^2M^{-1/2}\nu^{-1}$ --- none is available, since
  $\nu\iint|\nabla^2\hat{e}|^2\eta^2$ is a left-hand quantity of the first
  rung and is controlled only \emph{with} the cutoff $\eta$ --- or a level
  non-concentration hypothesis for that density. So the assertion of
  Remarks~\ref{rem:s41-summary} and~\ref{rem:s41-band-status} that
  $\mathcal{A}_{\mathrm{ann}}$ is ``of the same analytic type as
  Conjecture~\ref{conj:K-refined}'' is accurate for the $Y$-half of the
  closure target and \emph{inaccurate for the $D$-half}. This is the sharpest
  single correction this section makes to the surrounding text.
\end{remark}

\begin{remark}[Level refinement moves the two halves in opposite directions]
  \label{rem:s44-refinement}
  A small observation, recorded because it is the natural next thing to try
  and it does not work. Remark~\ref{rem:s41-reading} shows that pigeonholing
  over $N$ equally spaced levels gains nothing for a $\chi'$-weighted
  residual, because $|\chi'|\leq CN$ cancels the $1/N$ of the lightest
  annulus. The left-hand side of \eqref{eq:s41-bandD} carries \emph{no}
  $\chi'$ factor, so for it the pigeonhole is a genuine gain: at least half of
  $N$ equally spaced bands carry at most $4/N$ of
  $\nu\iint_{Q^\sharp}|\nabla^2\hat{e}|^2\psi^2$ each, and the constant
  multiplying that band in the Young step of
  Proposition~\ref{prop:s41-recursion} (with $\delta=1$) is $N$-independent.
  But the same refinement multiplies the $Y$-side constants by
  $\|\chi\chi'\|_\infty^2+\|\chi\chi''\|_\infty=O(N^2)$, so the coefficient
  $2^{-11}$ in \eqref{eq:s41-bandY} would have to be replaced by $cN^{-2}$
  while the band mass on that side decreases by at best $4/N$: a net loss of a
  factor $\sim N$. Since a single pair $\theta_\pm$ must serve both halves of
  Conjecture~\ref{target:band-absorption}, the refinement cannot be used.
  This refines, and does not contradict, Remark~\ref{rem:s41-reading}.
\end{remark}

\subsection{Status table (S44)}
\label{subsec:s44-status}

\begin{center}
\begin{tabular}{lll}
\hline
Question & Answer & Where \\
\hline
(a) Conj.~\ref{conj:K-refined} $\Rightarrow$ Conj.~\ref{target:band-absorption}?
  & \textbf{No} (disjoint supports) & Prop.~\ref{prop:s44-a-fails} \\
(a$'$) Band-transported $c_K$ $\Rightarrow$ \eqref{eq:s41-bandY}?
  & \textbf{Yes} --- via $\|w\|_{L^\infty}$, no extra ingredient
    & Prop.~\ref{prop:s44-linfty-route} \\
(a$''$) \dots{} or via the budget route?
  & \textbf{Yes}, but only given (NB) \eqref{eq:s44-NB}
    & Prop.~\ref{prop:s44-transport} \\
(b) Conj.~\ref{target:band-absorption} $\Rightarrow$ Conj.~\ref{conj:K-refined}?
  & \textbf{No}; no bound on $c_K$ at all & Prop.~\ref{prop:s44-b-fails} \\
(b$'$) Do the two \emph{conclusions} compare?
  & \textbf{No} & Rem.~\ref{rem:s44-conclusions} \\
(c) Common sufficient hypothesis (preferred)?
  & \textbf{Yes} --- (LC), the $c_K$ hypothesis on $\{m\geq M/16\}$
    & Conj.~\ref{conj:lc-s44}, Prop.~\ref{prop:s44-lc-implies} \\
(c$'$) Common hypothesis for the budget route?
  & \textbf{Yes} --- (JLE-1), but demanding
    & Conj.~\ref{conj:jle-s44}, Prop.~\ref{prop:s44-jle-implies} \\
(c$''$) Does either cover \eqref{eq:s41-bandD}?
  & \textbf{No} --- different analytic type & Rem.~\ref{rem:s44-bandD} \\
(d) Is \eqref{eq:s41-bandY} weaker than Conj.~\ref{conj:K-refined}?
  & \textbf{No}; $O(1)$ in bookkeeping only
    & Prop.~\ref{prop:s44-linfty-route}, Rem.~\ref{rem:s44-strength} \\
(d$'$) Is $\Xi=O(1)$ weaker than $c_K=O(1)$?
  & \textbf{No} --- \emph{stronger}, by $|Q(r^*)|/|R|$
    & Lem.~\ref{lem:s44-xi-vs-cK}, Prop.~\ref{prop:s44-calibration}(iii) \\
(d$''$) Numerics ($c_K\approx1$) support which?
  & Neither, on their own; a volume measurement is missing
    & Rem.~\ref{rem:s44-numerics} \\
(d$'''$) First-draft answer to (d)
  & \textbf{Withdrawn} (lemma was inverted) & Rem.~\ref{rem:s44-corrigendum} \\
(e) Independent as stated?
  & \textbf{Yes}, over the a priori structure (A1)--(A3)
    & Rem.~\ref{rem:s44-scope} \\
Volume conditions (NB), (NB-E)
  & \textbf{Symmetric}; needed by the budget route only
    & Rem.~\ref{rem:s44-NB}, Rem.~\ref{rem:s44-two-routes} \\
Level refinement as a route
  & \textbf{Net loss} & Rem.~\ref{rem:s44-refinement} \\
(LC) and (JLE-1) themselves
  & \textbf{Open}; mutually independent & Rem.~\ref{rem:s44-jle-status} \\
\hline
\end{tabular}
\end{center}

\begin{remark}[Honest summary of S44]\label{rem:s44-summary}
  The resemblance noted in Remarks~\ref{rem:s41-summary}
  and~\ref{rem:s41-band-status} is real at the level of \emph{form} and
  illusory at the level of \emph{content}: the two conjectures are supported
  on disjoint regions of the $|\omega|$-level decomposition, and neither
  implies the other, nor any nontrivial part of the other, from the a priori
  structure available --- Remark~\ref{rem:s44-scope} states the exact scope of
  that claim, which is a statement about admissible test configurations and
  not about Navier--Stokes solutions.

  The positive outcome is that a single hypothesis covers both $Y$-halves, and
  that it is a familiar one: Conjecture~\ref{conj:lc-s44} is the $c_K$
  hypothesis of Conjecture~\ref{conj:K-refined} with its region enlarged from
  $\{m\geq M/4\}$ to $\{m\geq M/16\}$, and nothing else. That is the
  legitimate content of ``same analytic type''
  (Remark~\ref{rem:s44-lc-reading}): not that either target subsumes the
  other, but that both are instances of one hypothesis stated on adjacent
  regions.

  The negative outcomes are three. First, the two conjectures are genuinely
  incomparable, so the loose language of
  Remarks~\ref{rem:s41-summary} and~\ref{rem:s41-band-status} must not be read
  as an implication. Second, the angular half \eqref{eq:s41-bandD} is of a
  different type entirely and is covered by no decorrelation hypothesis
  (Remark~\ref{rem:s44-bandD}). Third --- and this reverses what the first
  draft of this section claimed --- \eqref{eq:s41-bandY} is \emph{not} weaker
  than Conjecture~\ref{conj:K-refined}: its apparently absolute constants are
  an artefact of its bookkeeping, and reduced to a hypothesis on a correlation
  constant it needs the same $1/\mathcal{L}$ decay
  (Proposition~\ref{prop:s44-linfty-route}). The volume-normalised statistic
  $\Xi$ that would avoid the logarithm is \emph{stronger}, not weaker, than
  the correlation constant (Lemma~\ref{lem:s44-xi-vs-cK}), so the existing
  $O(1)$ numerics do not deliver it; what they would need to be supplemented
  by is listed in Remark~\ref{rem:s44-numerics}(iii). The withdrawal is
  recorded in Remark~\ref{rem:s44-corrigendum}.

  No proved statement of the document is altered, and nothing here bears on
  H2, on $\sigma^*$-decay, on the hypotheses of \S\ref{sec:rigidity-s39}, or
  on Type-II exclusion.
\end{remark}

% ════════════════════════════════════════════════════════════════════════════
% ══════════════════════════════════════════════════════════════════════════════
\begin{thebibliography}{99}

\bibitem{CKN1982}
  L.~Caffarelli, R.~Kohn, L.~Nirenberg,
  Partial regularity of suitable weak solutions of the Navier--Stokes equations,
  \emph{Comm. Pure Appl. Math.}, 35(6):771--831, 1982.

\bibitem{BenichouBrezis}
  Ph.~Bénilan, H.~Brezis, M.~G.~Crandall,
  A semilinear equation in $L^1(\mathbb{R}^N)$,
  \emph{Ann. Scuola Norm. Sup. Pisa}, 2(4):523--555, 1975.

\bibitem{Fefferman2000}
  C.~L.~Fefferman,
  \emph{Existence and Smoothness of the Navier--Stokes Equation},
  Clay Mathematics Institute, 2000.

\bibitem{Giaquinta1982}
  M.~Giaquinta, E.~Giusti,
  On the regularity of the minima of variational integrals,
  \emph{Acta Math.}, 148:31--46, 1982.

\bibitem{Lin1998}
  F.~Lin,
  A new proof of the Caffarelli--Kohn--Nirenberg theorem,
  \emph{Comm. Pure Appl. Math.}, 51(3):241--257, 1998.

\bibitem{Prodi1959}
  G.~Prodi,
  Un teorema di unicità per le equazioni di Navier--Stokes,
  \emph{Ann. Mat. Pura Appl.}, 48:173--182, 1959.

\bibitem{Serrin1962}
  J.~Serrin,
  On the interior regularity of weak solutions of the Navier--Stokes equations,
  \emph{Arch. Rat. Mech. Anal.}, 9:187--195, 1962.

\bibitem{Stein1970}
  E.~M.~Stein,
  \emph{Singular Integrals and Differentiability Properties of Functions},
  Princeton Univ. Press, 1970.

\bibitem{Stredulinsky1984}
  E.~W.~Stredulinsky,
  Higher integrability from reverse H\"older inequalities,
  \emph{Indiana Univ. Math. J.}, 29:407--413, 1980.

\bibitem{Krylov1996}
  N.~V.~Krylov,
  \emph{Lectures on Elliptic and Parabolic Equations in H\"older Spaces},
  Graduate Studies in Mathematics, Amer. Math. Soc., 1996.

\bibitem{BenilBrezis1975}
  Ph.~B\'enilan, H.~Br\'ezis, M.~G.~Crandall,
  A semilinear equation in $L^1(\mathbb{R}^N)$,
  \emph{Ann. Sc. Norm. Sup. Pisa}, 2(4):523--555, 1975.

\bibitem{ConstantinFefferman1993}
  P.~Constantin, C.~Fefferman,
  Direction of vorticity and the problem of global regularity
  for the Navier--Stokes equations,
  \emph{Indiana Univ. Math. J.}, 42(3):775--789, 1993.

\bibitem{Constantin1994}
  P.~Constantin,
  Geometric Statistics in Turbulence,
  \emph{SIAM Review}, 36(1):73--98, 1994.

\bibitem{KozonoTaniuchi2000}
  H.~Kozono, Y.~Taniuchi,
  Bilinear estimates in BMO and the Navier--Stokes equations,
  \emph{Math. Z.}, 235(1):173--194, 2000.

\bibitem{NecasRuzickaSverak1996}
  J.~Ne\v{c}as, M.~R\r{u}\v{z}i\v{c}ka, V.~\v{S}ver\'ak,
  On Leray's self-similar solutions of the Navier-Stokes equations,
  \emph{Acta Math.}, 176:283--294, 1996.

\bibitem{EscauriazaSereginSverak2003}
  L.~Escauriaza, G.~Seregin, V.~\v{S}ver\'ak,
  $L_{3,\infty}$-solutions of Navier-Stokes equations and backward uniqueness,
  \emph{Uspekhi Mat. Nauk}, 58:3--44, 2003.

\bibitem{KNSS2009}
  G.~Koch, N.~Nadirashvili, G.~Seregin, V.~\v{S}ver\'ak,
  Liouville theorems for the Navier--Stokes equations and applications,
  \emph{Acta Math.}, 203(1):83--105, 2009.

\bibitem{Seregin2012}
  G.~Seregin,
  A certain necessary condition of potential blow up for Navier--Stokes
  equations,
  \emph{Comm. Math. Phys.}, 312(3):833--845, 2012.

\bibitem{LadyzhenskayaSeregin1999}
  O.~A.~Ladyzhenskaya, G.~A.~Seregin,
  On partial regularity of suitable weak solutions to the three-dimensional
  Navier--Stokes equations,
  \emph{J. Math. Fluid Mech.}, 1(4):356--387, 1999.

\bibitem{Giusti2003}
  E.~Giusti,
  \emph{Direct Methods in the Calculus of Variations},
  World Scientific, 2003. (Iteration Lemma~6.1; hole-filling.)

\bibitem{LadyzhenskayaUraltseva1968}
  O.~A.~Ladyzhenskaya, N.~N.~Ural'tseva,
  \emph{Linear and Quasilinear Elliptic Equations},
  Academic Press, 1968. (Chapter~II: De~Giorgi level truncation.)

\end{thebibliography}

\end{document}
